\documentclass[11pt]{article}
% Shared preamble for all solution files.
\usepackage[margin=1.1in]{geometry}
\usepackage{enumitem}
\usepackage{tikz}
\usepackage[hidelinks]{hyperref}
\usepackage{amsmath,amssymb,amsthm}
\newcommand{\fa}{\mathfrak a}
\newcommand{\fb}{\mathfrak b}
\newcommand{\fc}{\mathfrak c}
\newcommand{\fp}{\mathfrak p}
\newcommand{\fq}{\mathfrak q}
\newcommand{\fm}{\mathfrak m}
\newcommand{\fn}{\mathfrak n}
\newcommand{\fr}{\mathfrak r}
\newcommand{\Nil}{\mathfrak N}
\newcommand{\Jac}{\mathfrak R}
\newcommand{\rad}{r}
\newcommand{\Spec}{\operatorname{Spec}}
\newcommand{\Max}{\operatorname{Max}}
\newcommand{\Ann}{\operatorname{Ann}}
\newcommand{\Ker}{\operatorname{Ker}}
\newcommand{\Z}{\mathbf Z}
\newcommand{\Q}{\mathbf Q}
\newcommand{\R}{\mathbf R}
\newcommand{\C}{\mathbf C}
\newcommand{\ci}{\mathrm{c}}
\newcommand{\ex}{\mathrm{e}}

\setlength{\parindent}{0pt}
\setlength{\parskip}{0.55em}
% \exercise{label}{statement}  -- followed by the solution text
\newcommand{\exercise}[2]{\par\bigskip\noindent{\large\textbf{Exercise #1}}\par\smallskip
  \begin{quote}\small\textbf{Problem.}\ #2\end{quote}\par}
\newcommand{\pf}{\noindent\textbf{Solution.}\ }

\title{Hartshorne, Chapter IV: Curves\\[2pt]\large Solutions to the exercises}
\author{}
\date{}
\begin{document}
\maketitle
\noindent\textbf{Conventions.} As in the book, $k$ is an algebraically closed field and ``curve'' means a complete nonsingular curve over $k$ unless an exercise says otherwise. Results of the text are cited by their numbers in Hartshorne (``Thm.~1.3'', ``Prop.~4.6'', \dots), and results of other chapters as ``II, \dots'' or ``III, \dots''. The numbering and statements of the exercises follow the edition scanned in this repository; each exercise is restated before its solution, and starred exercises are marked ``Starred in the book'' at the end of the statement.

Honest notes. Several printed statements need a hypothesis or correction, and the solutions say what is assumed or proved instead: characteristic assumptions (char $\neq2$ in much of Section 4; characteristic $0$ in Ex.~5.2, 5.3, 5.5, 5.7 and part of 6.7); the parenthetical remark in Ex.~2.3(h) about a tacnode is wrong (the dual branch has orders $(3,4)$); Ex.~4.16(d) as printed fails for supersingular curves; Ex.~5.5(b) is false as printed in the cuspidal case. \textbf{Ex.~6.7 is only partly proved:} the first half is proved in full, and the second half is proved for plane curves and curves on a quadric, while for the remaining nondegenerate curves it is proved conditionally on one lemma (the tangent line at a general point meets the curve nowhere else, in characteristic $0$) which I could not prove; the solution says so precisely. Some solutions use standard facts beyond the text (EGA smoothness criteria, miracle flatness, Weierstrass preparation), which are named where used.

\section*{1. The Riemann-Roch theorem}

Conventions. Throughout, $k$ is algebraically closed, a \emph{curve} is a complete nonsingular curve over $k$ (so projective), $K$ denotes a canonical divisor (and also, when no confusion is possible, a function field), and for a divisor $D$ on a curve we write $\mathcal L(D)$ for the associated invertible sheaf, $H^0(X,\mathcal L(D))=\{f\in K(X)^*:(f)+D\ge0\}\cup\{0\}$ and $l(D)=\dim_kH^0(X,\mathcal L(D))$. Results of the text are cited as ``Thm.~1.3'', ``Lemma~1.2'', \dots, those of other chapters as ``II, Ex.~6.11'' etc. Recall the following facts, used repeatedly: if $f:X\to\mathbf P^1$ is the morphism defined by a nonconstant $f\in K(X)$, then $f$ is finite, $f^*(\infty)=(f)_\infty$ is the divisor of poles of $f$, and $\deg f=[K(X):k(f)]=\deg (f)_\infty$ (II, 6.8, 6.9, and $e_P=-\operatorname{ord}_P(f)$ at a pole, because $1/f$ is a local parameter at $\infty$). For a point $P$ the sheaf sequence $0\to\mathcal L(D-P)\to\mathcal L(D)\to k(P)\to0$ gives $l(D-P)\le l(D)\le l(D-P)+1$.

\exercise{1.1}{Let $X$ be a curve, and let $P\in X$ be a point. Then there exists a nonconstant rational function $f\in K(X)$, which is regular everywhere except at $P$.}
\pf
Let $g$ be the genus and put $n=2g+1$. For $m\ge 2g-1$ we have $\deg(K-mP)=2g-2-m<0$, so $l(K-mP)=0$ by Lemma~1.2, and the Riemann--Roch theorem (Thm.~1.3) gives $l(mP)=m+1-g$. Hence $l(nP)-l((n-1)P)=1$ (both $n$ and $n-1=2g$ are $\ge 2g-1$), so $H^0(X,\mathcal L((n-1)P))$ is a proper subspace of $H^0(X,\mathcal L(nP))$. Pick $f$ in the difference. Then $(f)+nP\ge0$ but $(f)+(n-1)P\not\ge0$, which means $\operatorname{ord}_P(f)=-n$ and $\operatorname{ord}_Q(f)\ge0$ for all $Q\ne P$. Thus $f$ is regular on $X-\{P\}$ and has a pole of order $n=2g+1\ge1$ at $P$; in particular $f$ is nonconstant. \hfill$\square$

\exercise{1.2}{Again let $X$ be a curve, and let $P_1,\dots,P_r\in X$ be points. Then there is a rational function $f\in K(X)$ having poles (of some order) at each of the $P_i$, and regular elsewhere.}
\pf
For each $i$ choose, by Ex.~1.1, a rational function $f_i$ regular on $X-\{P_i\}$ with a pole at $P_i$, and put $f=f_1+\dots+f_r$ (if $r=0$ take $f=0$). At $P_i$, the function $f_i$ has a pole while $\sum_{j\ne i}f_j$ is regular, so $f$ has a pole at $P_i$ (a sum of a function with a pole and a regular function has a pole). At any $Q\notin\{P_1,\dots,P_r\}$ every $f_j$ is regular, hence so is $f$. \hfill$\square$

\exercise{1.3}{Let $X$ be an integral, separated, regular, one-dimensional scheme of finite type over $k$, which is \emph{not} proper over $k$. Then $X$ is affine. [\emph{Hint}: Embed $X$ in a (proper) curve $\overline X$ over $k$, and use (Ex.~1.2) to construct a morphism $f:\overline X\to\mathbf P^1$ such that $f^{-1}(\mathbf A^1)=X$.]}
\pf
\emph{Step 1: $X$ is an open subscheme of a curve $\overline X$.} Let $L=K(X)$; it has transcendence degree $1$ over $k$. Let $\overline X$ be the (unique) nonsingular projective curve with function field $L$ (I, 6.9 and its scheme-theoretic form II, 6.7), so the closed points of $\overline X$ are the discrete valuation rings of $L/k$. For a closed point $P\in X$, the local ring $\mathcal O_P$ is a regular local ring of dimension $1$, hence a DVR of $L/k$. Since $\overline X\to\operatorname{Spec}k$ is proper, the valuative criterion (II, 4.7) shows that the generic point $\operatorname{Spec}L\to\overline X$ extends to a morphism $\operatorname{Spec}\mathcal O_P\to\overline X$. This extension is defined on an open neighborhood $U_P$ of $P$ in $X$: if $\operatorname{Spec}B\subseteq\overline X$ is an affine open containing the image of the closed point and $\operatorname{Spec}A\subseteq X$ is an affine neighborhood of $P=\mathfrak p$, the morphism is a $k$-algebra map $B\to A_{\mathfrak p}$, the finitely many generators of $B$ go to elements of $A_{\mathfrak p}\subseteq L$ with a common denominator $s\in A-\mathfrak p$, so it comes from $B\to A_s$, which is compatible with the inclusions in $L$. Two such morphisms agree on the overlaps (they agree at the generic point, $X$ is reduced and $\overline X$ is separated, II, 4.1 and its proof), so they glue to a morphism $g:X\to\overline X$ which is the identity on function fields. For each closed point $P$ the map $\mathcal O_{\overline X,g(P)}\to\mathcal O_P$ is an inclusion of DVRs of $L$ with the second dominating the first; since a valuation ring is maximal for domination (I, Thm.~6.1A) it is an isomorphism. If $g(P)=g(Q)$, then $\mathcal O_P=\mathcal O_Q$ as subrings of $L$, and the two morphisms $\operatorname{Spec}\mathcal O_P\to X$ centered at $P$ and $Q$ agree on $\operatorname{Spec}L$; as $X$ is separated, the uniqueness part of the valuative criterion (II, 4.3) gives $P=Q$. So $g$ is injective on closed points (and sends the generic point to the generic point, so $g$ is injective). Next, $g$ is a local isomorphism near each closed point $P$: take affine opens $\operatorname{Spec}A\ni P$ and $\operatorname{Spec}B\ni g(P)$ with $g(\operatorname{Spec}A)\subseteq\operatorname{Spec}B$, and let $P,g(P)$ correspond to $\mathfrak p\subset A$, $\mathfrak q=\mathfrak p\cap B$. Then $B\subseteq A\subseteq L$ and $A$ is a finitely generated $B$-algebra, $A=B[a_1,\dots,a_m]$; as $A_{\mathfrak p}=B_{\mathfrak q}$ we can write $a_i=b_i/s_i$ with $b_i,s_i\in B$, $s_i\notin\mathfrak q$. With $s=\prod s_i\notin\mathfrak p$ we get $A_s=B_s$, i.e.\ $g$ maps the open set $D(s)\ni P$ of $\operatorname{Spec}A$ isomorphically onto the open set $D(s)$ of $\operatorname{Spec}B$. The set of points near which $g$ is an open immersion is open and contains all closed points, hence is all of $X$. An injective local open immersion is an open immersion. We regard $X\subseteq\overline X$ as an open subscheme. (This is the standard embedding of a one-dimensional separated scheme into its complete model.)

\emph{Step 2.} Since $\overline X$ is proper over $k$ and $X$ is not, $X\ne\overline X$. The complement $\overline X-X$ is a proper closed subset of an irreducible curve, hence a nonempty finite set $\{P_1,\dots,P_r\}$ of closed points. By Ex.~1.2 there is $f\in K(\overline X)$ with a pole at every $P_i$ and regular elsewhere. The corresponding morphism $f:\overline X\to\mathbf P^1$ is finite (nonconstant morphism of curves, II, 6.8) and $f^{-1}(\infty)=\{P_1,\dots,P_r\}$ (the points where $f$ has a pole). Hence $f^{-1}(\mathbf A^1)=\overline X-\{P_i\}=X$, and the restriction $X=f^{-1}(\mathbf A^1)\to\mathbf A^1$ of a finite morphism is finite, in particular affine. As $\mathbf A^1$ is affine, $X$ is affine. \hfill$\square$

\exercise{1.4}{Show that a separated, one-dimensional scheme of finite type over $k$, none of whose irreducible components is proper over $k$, is affine. [\emph{Hint}: Combine (Ex.~1.3) with (III, Ex.~3.1, Ex.~3.2, Ex.~4.2).]}
\pf
Let $X$ be such a scheme; it is noetherian. Recall (III, Ex.~3.1, 3.2) that a noetherian scheme is affine if and only if each of its irreducible components, with the reduced induced structure, is affine (first pass to $X_{\mathrm{red}}$, then to the components). So let $X_i$ be an irreducible component with its reduced structure. It is integral, separated, of finite type over $k$, and $\dim X_i\le1$; it cannot have dimension $0$, because a zero-dimensional integral scheme of finite type over $k$ is $\operatorname{Spec}k$, which is proper. So $\dim X_i=1$.

Let $\pi:\widetilde X_i\to X_i$ be the normalization (II, Ex.~3.8). It is a finite surjective morphism, $\widetilde X_i$ is integral, normal of dimension one, hence regular, and separated and of finite type over $k$ (being finite over $X_i$). It is not proper over $k$: otherwise $X_i$, the image under the surjective map $\pi$ of a proper scheme, would be proper ($X_i$ is separated and universally closed, since for every $Z$ the map $\widetilde X_i\times Z\to X_i\times Z$ is surjective and $\widetilde X_i\times Z\to Z$ is closed), contrary to the hypothesis. By Ex.~1.3, $\widetilde X_i$ is affine. By Chevalley's theorem (III, Ex.~4.2: if $f:X'\to Y'$ is finite surjective of noetherian separated schemes and $X'$ is affine, so is $Y'$), $X_i$ is affine. Hence all irreducible components of $X$ are affine, so $X$ is affine. \hfill$\square$

\exercise{1.5}{For an effective divisor $D$ on a curve $X$ of genus $g$, show that $\dim|D|\le\deg D$. Furthermore, equality holds if and only if $D=0$ or $g=0$.}
\pf
Write $D=P_1+\dots+P_n$ with $n=\deg D$ (points may repeat) and put $D_j=P_1+\dots+P_j$, $D_0=0$. Since $l(D_j)\le l(D_{j-1})+1$ (see the conventions above, with $D_{j-1}=D_j-P_j$) and $l(0)=\dim H^0(X,\mathcal O_X)=1$ ($X$ is a complete variety), we get $l(D)\le n+1$, i.e. $\dim|D|=l(D)-1\le\deg D$.

If $D=0$ then $\dim|D|=0=\deg D$. If $g=0$, then $\deg(K-D)=-2-\deg D<0$, so $l(K-D)=0$ (Lemma~1.2) and Riemann--Roch gives $l(D)=\deg D+1$, i.e. equality.

Conversely suppose equality holds, $D\ne0$ and $n\ge1$. Then $l(D)=n+1$, which forces $l(D_j)=l(D_{j-1})+1$ for all $j$; in particular $l(P_1)=2$. Take $f\in H^0(X,\mathcal L(P_1))$ nonconstant. Its divisor of poles is $\le P_1$ and nonzero, so $(f)_\infty=P_1$, and the finite morphism $f:X\to\mathbf P^1$ has degree $\deg(f)_\infty=1$. Thus $K(X)=k(f)\cong K(\mathbf P^1)$ and $X\cong\mathbf P^1$ (I, 6.12), so $g=0$. Hence equality holds only for $D=0$ or $g=0$. \hfill$\square$

\exercise{1.6}{Let $X$ be a curve of genus $g$. Show that there is a finite morphism $f:X\to\mathbf P^1$ of degree $\le g+1$. (Recall that the degree of a finite morphism of curves $f:X\to Y$ is defined as the degree of the field extension $[K(X):K(Y)]$ (II, Section~6).)}
\pf
Fix a point $P$ and let $D=(g+1)P$. By Riemann--Roch, $l(D)\ge\deg D+1-g=2$ (as $l(K-D)\ge0$). So $H^0(X,\mathcal L(D))$ contains a nonconstant function $f$. Then $(f)_\infty\le D$, and $(f)_\infty\ne0$ because $f$ is nonconstant. The associated morphism $f:X\to\mathbf P^1$ is finite of degree $\deg f=\deg(f)_\infty\le\deg D=g+1$. \hfill$\square$

\exercise{1.7}{A curve $X$ is called \emph{hyperelliptic} if $g\ge2$ and there exists a finite morphism $f:X\to\mathbf P^1$ of degree $2$.
\begin{enumerate}[label=(\alph*),leftmargin=*]
\item If $X$ is a curve of genus $g=2$, show that the canonical divisor defines a complete linear system $|K|$ of degree $2$ and dimension $1$, without base points. Use (II, 7.8.1) to conclude that $X$ is hyperelliptic.
\item Show that the curves constructed in (1.1.1) all admit a morphism of degree $2$ to $\mathbf P^1$. Thus there exist hyperelliptic curves of any genus $g\ge2$.
\end{enumerate}
\emph{Note.} We will see later (Ex.~3.2) that there exist nonhyperelliptic curves. See also (V, Ex.~2.10).}
\pf
\textbf{(a)} By Example~1.3.3, $\deg K=2g-2=2$ and $l(K)=g=2$, so $|K|$ is a complete linear system of degree $2$ and dimension $l(K)-1=1$. Suppose $P$ were a base point of $|K|$: then every section of $\mathcal L(K)$ vanishes at $P$, i.e. $l(K-P)=l(K)=2$. By Riemann--Roch applied to the divisor $P$, $l(P)-l(K-P)=1+1-2=0$, so $l(P)=2$. But $l(P)=2$ forces $X\cong\mathbf P^1$ (see the proof of Ex.~1.5), contradicting $g=2$. So $|K|$ has no base points. Hence $\mathcal L(K)$ is generated by its two-dimensional space of global sections, and by (II, 7.8.1) it defines a morphism $\varphi:X\to\mathbf P^1$ with $\varphi^*\mathcal O(1)\cong\mathcal L(K)$. As $\deg\mathcal L(K)=2>0$, $\varphi$ is not constant, hence finite, and $2=\deg\varphi^*\mathcal O(1)=\deg\varphi\cdot\deg\mathcal O(1)=\deg\varphi$ (II, 6.9). So $X$ is hyperelliptic.

\textbf{(b)} Let $Q\cong\mathbf P^1\times\mathbf P^1\subseteq\mathbf P^3$ be a nonsingular quadric surface and let $X\subseteq Q$ be an irreducible nonsingular divisor of type $(g+1,2)$, $g\ge2$, which has $p_a=g$ (Remark~1.1.1 and III, Ex.~5.6); so $X=Z(F)$ for a bihomogeneous form $F$ of bidegree $(g+1,2)$ in $((x_0,x_1),(y_0,y_1))$, and $F$ is irreducible. Consider the first projection $\mathrm{pr}_1:X\to\mathbf P^1$ (to the factor $(x_0:x_1)$, in which $F$ has degree $g+1$, while $F$ has degree $2$ in $(y_0,y_1)$). It is a morphism and is not constant on the integral curve $X$ unless $X$ lies in a fibre $\{a\}\times\mathbf P^1$, which would mean that the linear form defining $a$ divides $F$, impossible since $F$ is irreducible of bidegree $(g+1,2)$ with $g+1\ge3$. So it is finite. Its degree is $[K(X):k(x)]$ where $x=x_1/x_0$. In the chart $x_0y_0\ne0$ we have $K(X)=k(x)(y)$ with $y=y_1/y_0$ and $y$ satisfies the dehomogenization $f(x,y)=a(x)y^2+b(x)y+c(x)=0$ of $F$, where $a,c$ are not identically zero (otherwise $y_1$ or $y_0$ would divide $F$). This polynomial is irreducible in $k(x)[y]$, since by Gauss' lemma a factorization would give a factorization of $F$. Hence $[K(X):k(x)]=2$, and $\mathrm{pr}_1:X\to\mathbf P^1$ has degree $2$. Since $g\ge2$, $X$ is hyperelliptic of genus $g$ for every $g\ge2$. \hfill$\square$

\exercise{1.8}{\emph{$p_a$ of a Singular Curve.} Let $X$ be an integral projective scheme of dimension $1$ over $k$, and let $\widetilde X$ be its normalization (II, Ex.~3.8). Then there is an exact sequence of sheaves on $X$,
\[
0\to\mathcal O_X\to f_*\mathcal O_{\widetilde X}\to\sum_{P\in X}\widetilde{\mathcal O}_P/\mathcal O_P\to0,
\]
where $\widetilde{\mathcal O}_P$ is the integral closure of $\mathcal O_P$. For each $P\in X$, let $\delta_P=\operatorname{length}(\widetilde{\mathcal O}_P/\mathcal O_P)$.
\begin{enumerate}[label=(\alph*),leftmargin=*]
\item Show that $p_a(X)=p_a(\widetilde X)+\sum_{P\in X}\delta_P$. [\emph{Hint}: Use (III, Ex.~4.1) and (III, Ex.~5.3).]
\item If $p_a(X)=0$, show that $X$ is already nonsingular and in fact isomorphic to $\mathbf P^1$. This strengthens (1.3.5).
\item If $P$ is a node or an ordinary cusp (I, Ex.~5.6, Ex.~5.14), show that $\delta_P=1$. [\emph{Hint}: Show first that $\delta_P$ depends only on the analytic isomorphism class of the singularity at $P$. Then compute $\delta_P$ for the node and cusp of suitable plane cubic curves. See (V, 3.9.3) for another method.]
\end{enumerate}
(Part (c) is starred in the book.)}
\pf
Here $f:\widetilde X\to X$ is the normalization, a finite birational morphism; $\widetilde X$ is an integral normal, i.e.\ nonsingular, projective curve. The quotient $\mathcal T=\sum_P\widetilde{\mathcal O}_P/\mathcal O_P$ is a coherent sheaf supported on the finite set of singular points, and $\dim_k H^0(X,\mathcal T)=\sum_P\operatorname{length}_{\mathcal O_P}(\widetilde{\mathcal O}_P/\mathcal O_P)=\sum_P\delta_P$, because every composition factor of an $\mathcal O_P$-module of finite length is $k(P)=k$. Also $H^i(X,\mathcal T)=0$ for $i>0$ (the sheaf is a finite sum of skyscraper sheaves, hence flasque).

\textbf{(a)} For an integral projective scheme of dimension $1$ one has $p_a=1-\chi(\mathcal O)$ (III, Ex.~5.3). The Euler characteristic is additive in exact sequences (III, Ex.~5.1), so $\chi(f_*\mathcal O_{\widetilde X})=\chi(\mathcal O_X)+\chi(\mathcal T)=\chi(\mathcal O_X)+\sum_P\delta_P$. Since $f$ is finite, hence affine, $H^i(X,f_*\mathcal O_{\widetilde X})\cong H^i(\widetilde X,\mathcal O_{\widetilde X})$ (III, Ex.~4.1), so $\chi(f_*\mathcal O_{\widetilde X})=\chi(\mathcal O_{\widetilde X})$. Therefore $\chi(\mathcal O_{\widetilde X})=\chi(\mathcal O_X)+\sum\delta_P$, and applying $1-(\cdot)$ gives $p_a(X)=p_a(\widetilde X)+\sum_P\delta_P$.

\textbf{(b)} By Prop.~1.1, $p_a(\widetilde X)=g(\widetilde X)\ge0$, and each $\delta_P\ge0$. If $p_a(X)=0$, part (a) forces $\delta_P=0$ for all $P$ and $g(\widetilde X)=0$. Now $\delta_P=0$ means $\mathcal O_P=\widetilde{\mathcal O}_P$, so $\mathcal O_X\to f_*\mathcal O_{\widetilde X}$ is an isomorphism. As $f$ is finite, $\widetilde X=\operatorname{Spec}_X f_*\mathcal O_{\widetilde X}=\operatorname{Spec}_X\mathcal O_X=X$, i.e.\ $f$ is an isomorphism. So $X\cong\widetilde X$ is nonsingular of genus $0$, hence $X\cong\mathbf P^1$ by Example~1.3.5.

\textbf{(c)} \emph{Reading of the terms.} In I, Ex.~5.6 and 5.14 a node is a double point of a plane curve with two distinct tangent directions, and an ordinary cusp is a double point whose tangent cone is a double line and whose cubic term is nonzero on that line, i.e.\ the singularity of $y^2=x^3$. Concretely: in suitable coordinates $x,y$ at $P$ (so that $\widehat{\mathcal O}_P=k[[x,y]]/(h)$ for the local equation $h$ of the curve), $h=xy+(\text{terms of order}\ge3)$ for a node, and $h=y^2+(\text{cubic form with nonzero coefficient }a\text{ of }x^3)+(\text{terms of order}\ge4)$ for an ordinary cusp. We use only these descriptions, in any characteristic. (In characteristic $\ne2,3$ they say exactly that $\widehat{\mathcal O}_P\cong k[[x,y]]/(xy)$, resp.\ $k[[x,y]]/(y^2-x^3)$.)

\emph{Step 1: $\delta_P$ depends only on $R=\widehat{\mathcal O}_P$.} Put $\mathcal O=\mathcal O_P$, $\widetilde{\mathcal O}=\widetilde{\mathcal O}_P$ (a semilocal Dedekind domain, finite over $\mathcal O$ with the same fraction field). The quotient $\widetilde{\mathcal O}/\mathcal O$ has finite length, so it is unchanged by the faithfully flat completion: $\widetilde{\mathcal O}/\mathcal O\cong\widehat{\widetilde{\mathcal O}}/R$, where $\widehat{\widetilde{\mathcal O}}=\widetilde{\mathcal O}\otimes_{\mathcal O}R=\prod_{Q\mapsto P}\widehat{\mathcal O}_{\widetilde X,Q}$ is a finite product of complete DVRs (the $\mathfrak m_P$-adic topology on $\widetilde{\mathcal O}$ coincides with the topology of its Jacobson radical). Choose $0\ne c\in\mathcal O$ with $c\widetilde{\mathcal O}\subseteq\mathcal O$; $c$ is a nonzerodivisor on $R$ by flatness, and $R\subseteq\widehat{\widetilde{\mathcal O}}\subseteq c^{-1}R\subseteq Q(R)$ (total ring of fractions), with $Q(\widehat{\widetilde{\mathcal O}})=Q(R)$. A finite product of DVRs is integrally closed in its total ring of fractions, and it is integral over $R$; therefore $\widehat{\widetilde{\mathcal O}}$ is exactly the integral closure $\overline R$ of $R$ in $Q(R)$. Hence
\[
\delta_P=\dim_k\overline R/R,
\]
which depends only on the isomorphism class of the complete local $k$-algebra $R$, i.e.\ on the analytic type of the singularity.

\emph{Step 2: node.} Let $h=xy+h_3+h_4+\dots\in k[[x,y]]$ with $h_i$ homogeneous of degree $i$. Write $h=xy+g$ with $g\in(x,y)^3=x(x,y)^2+y(x,y)^2$, say $g=xa+yb$ with $a,b\in(x,y)^2$. Then $h=(x+b)(y+a)-ab$ with $ab\in(x,y)^4$. Putting $x_1=x+b$, $y_1=y+a$ (new coordinates, congruent to $x,y$ modulo $(x,y)^2$) we get $h=x_1y_1+g_1$ with $g_1\in(x,y)^4$. Repeating, we obtain coordinates $x_n,y_n$ with $x_{n+1}-x_n,\ y_{n+1}-y_n\in(x,y)^{n+1}$ and $h=x_ny_n+g_n$, $g_n\in(x,y)^{n+2}$. The sequences converge in $k[[x,y]]$ (this works in every characteristic), and $h=x_\infty y_\infty$ with $(x_\infty,y_\infty)$ a coordinate system. So $R\cong k[[x,y]]/(xy)$. Since $(x)\cap(y)=(xy)$ the map $R\to k[[x]]\times k[[y]]$ is injective; the target $k[[x]]\times k[[y]]$ is a product of DVRs, finite over $R$ (generated as an $R$-module by $(1,0),(0,1)$), with total ring of fractions $k((x))\times k((y))=Q(R)$; hence it is $\overline R$. The image of $R$ is $\{(a(x),b(y)):a(0)=b(0)\}$ (it contains $x\mapsto(x,0)$, $y\mapsto(0,y)$, $1\mapsto(1,1)$). So $\dim_k\overline R/R=1$, i.e.\ $\delta_P=1$.

\emph{Step 3: ordinary cusp.} Now $h=y^2+(\text{cubic with }x^3\text{-coefficient }a\ne0)+(\text{order}\ge4)$. Then $h(0,y)=y^2\cdot(\text{unit})$, so by the Weierstrass preparation theorem $h=U\cdot(y^2+\alpha(x)y+\beta(x))$ with $U$ a unit and $\alpha,\beta\in xk[[x]]$. Comparing the terms of degree $2$ and $3$ of $h$ shows $\operatorname{ord}\alpha\ge2$, $\operatorname{ord}\beta\ge3$, and the coefficient of $x^3$ in $\beta$ equals $a/U(0)\ne0$; so $\operatorname{ord}\beta=3$. Thus $R\cong k[[x]][y]/(y^2+\alpha y+\beta)=k[[x]]\oplus k[[x]]y$. Write $\alpha=x\alpha'$, $\beta=x^2\beta'$ with $\alpha'\in xk[[x]]$, $\operatorname{ord}\beta'=1$, and set $w=y/x$. Then $w^2+\alpha'w+\beta'=0$ is an Eisenstein equation over $k[[x]]$, so $A=k[[x]][w]=k[[x]]\oplus k[[x]]w$ is a complete DVR (with uniformizer $w$), finite over $R$, with $Q(A)=A[1/x]=Q(R)$; hence $A=\overline R$. Since $y=xw$, $R=k[[x]]\oplus k[[x]]xw$ and $A/R\cong k[[x]]w/xk[[x]]w\cong k$. Thus $\delta_P=1$.

(As a check, the plane cubics $y^2z=x^3+x^2z$ and $y^2z=x^3$ have $p_a=1$, normalization $\mathbf P^1$, and exactly one singular point, a node resp.\ cusp, so (a) gives $\delta=1$ for them, in agreement.) \hfill$\square$

\exercise{1.9}{\emph{Riemann-Roch for Singular Curves.} Let $X$ be an integral projective scheme of dimension $1$ over $k$. Let $X_{\mathrm{reg}}$ be the set of regular points of $X$.
\begin{enumerate}[label=(\alph*),leftmargin=*]
\item Let $D=\sum n_iP_i$ be a divisor with support in $X_{\mathrm{reg}}$, i.e., all $P_i\in X_{\mathrm{reg}}$. Then define $\deg D=\sum n_i$. Let $\mathcal L(D)$ be the associated invertible sheaf on $X$, and show that
\[
\chi(\mathcal L(D))=\deg D+1-p_a.
\]
\item Show that any Cartier divisor on $X$ is the difference of two very ample Cartier divisors. (Use (II, Ex.~7.5).)
\item Conclude that every invertible sheaf $\mathcal L$ on $X$ is isomorphic to $\mathcal L(D)$ for some divisor $D$ with support in $X_{\mathrm{reg}}$.
\item Assume furthermore that $X$ is a locally complete intersection in some projective space. Then by (III, 7.11) the dualizing sheaf $\omega_X$ is an invertible sheaf on $X$, so we can define the \emph{canonical divisor} $K$ to be a divisor with support in $X_{\mathrm{reg}}$ corresponding to $\omega_X$. Then the formula of (a) becomes
\[
l(D)-l(K-D)=\deg D+1-p_a.
\]
\end{enumerate}
(Starred in the book.)}
\pf
Let $\mathcal K$ be the sheaf of rational functions on $X$. The set $X-X_{\mathrm{reg}}$ is a proper closed subset of the integral one-dimensional scheme $X$, hence a finite set of closed points. For $P\in X_{\mathrm{reg}}$ the local ring $\mathcal O_P$ is a DVR with uniformizer $t_P$; a divisor $D=\sum n_PP$ supported in $X_{\mathrm{reg}}$ is a Cartier divisor, and $\mathcal L(D)\subseteq\mathcal K$ is the subsheaf with $\mathcal L(D)_P=t_P^{-n_P}\mathcal O_P$ for $P\in X_{\mathrm{reg}}$ and $=\mathcal O_P$ at all other points (on a neighborhood of a regular point it is the usual $\mathcal L(D)$).

\textbf{(a)} Let $P\in X_{\mathrm{reg}}$. Then $\mathcal L(D)\subseteq\mathcal L(D+P)$ and the quotient is a skyscraper sheaf at $P$ with stalk $t_P^{-n-1}\mathcal O_P/t_P^{-n}\mathcal O_P\cong k(P)$ (where $n=n_P$); so there is an exact sequence $0\to\mathcal L(D)\to\mathcal L(D+P)\to k(P)\to0$. By additivity of $\chi$ (III, Ex.~5.1) and $\chi(k(P))=1$, we get $\chi(\mathcal L(D+P))=\chi(\mathcal L(D))+1$, while $\deg(D+P)=\deg D+1$. So the formula holds for $D$ iff it holds for $D+P$ (and likewise for $D-P$). Any divisor supported in $X_{\mathrm{reg}}$ is reached from $0$ by finitely many steps $\pm P$ with $P\in X_{\mathrm{reg}}$, so it suffices to treat $D=0$, where the formula says $\chi(\mathcal O_X)=1-p_a$; this is (III, Ex.~5.3) (for an integral projective scheme of dimension $r$, $p_a=(-1)^r(\chi(\mathcal O_X)-1)$).

\textbf{(b)} We first record: (i) if $\mathcal M$ is very ample and $\mathcal N$ is generated by global sections, then $\mathcal M\otimes\mathcal N$ is very ample (II, Ex.~7.5). Indeed, $\mathcal M$ gives a closed immersion $i:X\to\mathbf P^N$, $\mathcal N$ gives a morphism $j:X\to\mathbf P^M$; then $(i,j):X\to\mathbf P^N\times\mathbf P^M$ is a closed immersion (it is the graph of $j$ followed by $i\times\mathrm{id}$), and composing with the Segre embedding gives an embedding by $\mathcal M\otimes\mathcal N$. In particular $\mathcal A^m$ is very ample for $m\ge1$ if $\mathcal A=\mathcal O_X(1)$ is.

Let $D$ be a Cartier divisor and $\mathcal A=\mathcal O_X(1)$ a very ample sheaf. By Serre (II, 5.17) $\mathcal L(D)\otimes\mathcal A^n$ is generated by global sections for $n\gg0$. Take such an $n$. Then $\mathcal M_1=\mathcal L(D)\otimes\mathcal A^{n+1}$ and $\mathcal M_2=\mathcal A^{n+1}$ are very ample by (i), and $\mathcal L(D)\cong\mathcal M_1\otimes\mathcal M_2^{-1}$. As $X$ is integral, every invertible sheaf is $\mathcal L(E)$ for a Cartier divisor $E$ (II, 6.15), so $\mathcal M_i\cong\mathcal L(D_i)$ with Cartier divisors $D_i$. Then $\mathcal L(D)\cong\mathcal L(D_1-D_2)$, so $D=D_1-D_2+(\varphi)$ for a rational function $\varphi$; put $D_1'=D_1+(\varphi)$, which is linearly equivalent to $D_1$, so $\mathcal L(D_1')\cong\mathcal M_1$ is very ample. Then $D=D_1'-D_2$ with both very ample.

\textbf{(c)} Let $\mathcal L$ be invertible. As above $\mathcal L\cong\mathcal L(D')$ for a Cartier divisor $D'$, and by (b) $D'=E_1-E_2$ with $\mathcal L(E_i)$ very ample, giving closed immersions $X\hookrightarrow\mathbf P^{N_i}$ with $\mathcal O(1)|_X\cong\mathcal L(E_i)$. Choose hyperplanes $H_i\subseteq\mathbf P^{N_i}$ containing none of the finitely many points of $X-X_{\mathrm{reg}}$ (possible: the hyperplanes through one given point form a proper linear subspace of the dual projective space, and $k$ is infinite). Since $X$ is integral and $X\not\subseteq H_i$, $E_i'=H_i\cap X$ is an effective Cartier divisor, with $\mathcal L(E_i')\cong\mathcal O(1)|_X\cong\mathcal L(E_i)$ and support in $X_{\mathrm{reg}}$. A Cartier divisor with support in $X_{\mathrm{reg}}$ is a divisor $\sum n_PP$ as in (a). Hence $\mathcal L\cong\mathcal L(E_1'-E_2')$ with $E_1'-E_2'$ supported in $X_{\mathrm{reg}}$.

\textbf{(d)} A locally complete intersection is Cohen--Macaulay (II, 8.21A, 8.23), so Cor.~7.7 of Ch.~III applies to the projective Cohen--Macaulay scheme $X$ of dimension $1$ and to the locally free sheaf $\mathcal F=\mathcal L(D)$: $H^1(X,\mathcal L(D))\cong H^0(X,\mathcal L(D)^\vee\otimes\omega_X)'$ (dual vector space). Here $\omega_X\cong\mathcal L(K)$ and $\mathcal L(D)^\vee\otimes\omega_X\cong\mathcal L(K-D)$. Therefore $\dim H^1(X,\mathcal L(D))=l(K-D)$, and $\chi(\mathcal L(D))=l(D)-l(K-D)$; combined with (a) this is the asserted formula. \hfill$\square$

\exercise{1.10}{Let $X$ be an integral projective scheme of dimension $1$ over $k$, which is locally complete intersection, and has $p_a=1$. Fix a point $P_0\in X_{\mathrm{reg}}$. Imitate (1.3.7) to show that the map $P\to\mathcal L(P-P_0)$ gives a one-to-one correspondence between the points of $X_{\mathrm{reg}}$ and the elements of the group $\operatorname{Pic}^0X$. This generalizes (II, 6.11.4) and (II, Ex.~6.7).}
\pf
\emph{Reading.} The superscript $0$ is lost in the printed statement (``$\operatorname{Pic}\ X$''); as in Example~1.3.7 the group is $\operatorname{Pic}^0X$, the subgroup of invertible sheaves of degree $0$ (the full group $\operatorname{Pic}X$ is much bigger, e.g.\ $\mathbf Z\oplus\operatorname{Pic}^0X$ for a plane cubic). Here the degree of an invertible sheaf $\mathcal M$ is defined by $\deg\mathcal M=\chi(\mathcal M)-\chi(\mathcal O_X)$; by Ex.~1.9(a),(c) every $\mathcal M$ is some $\mathcal L(D)$ with $D$ supported in $X_{\mathrm{reg}}$ and then $\deg\mathcal M=\deg D$, so $\deg$ is additive on $\mathcal L(D)\otimes\mathcal L(D')=\mathcal L(D+D')$ and $\operatorname{Pic}^0X=\{\mathcal M:\deg\mathcal M=0\}$ is a subgroup. Clearly $\mathcal L(P-P_0)\in\operatorname{Pic}^0X$ for $P\in X_{\mathrm{reg}}$.

\emph{Lemma.} Let $\mathcal M$ be an invertible sheaf on $X$ with $H^0(X,\mathcal M)\ne0$. Then $\deg\mathcal M\ge0$, and if $\deg\mathcal M=0$ then $\mathcal M\cong\mathcal O_X$. Indeed, a nonzero section $s$ gives an injection $s:\mathcal O_X\to\mathcal M$ ($X$ is integral and $\mathcal M$ is locally free of rank $1$), whose cokernel $\mathcal T$ vanishes at the generic point, hence has finite support, so $\chi(\mathcal T)=h^0(\mathcal T)=\sum\operatorname{length}\mathcal T_x\ge0$. Then $\deg\mathcal M=\chi(\mathcal T)\ge0$, with equality iff $\mathcal T=0$ iff $s$ is an isomorphism.

\emph{The canonical divisor has degree $0$.} By Ex.~1.9(d) and $p_a=1$: for $D=0$, $l(0)-l(K)=0$, and $l(0)=\dim H^0(\mathcal O_X)=1$, so $l(K)=1$; for $D=K$, $l(K)-l(0)=\deg K$, so $\deg K=0$. (By the Lemma, even $\omega_X\cong\mathcal O_X$.)

\emph{Surjectivity.} Let $\mathcal M\in\operatorname{Pic}^0X$; by Ex.~1.9(c), $\mathcal M=\mathcal L(D)$ with $D$ supported in $X_{\mathrm{reg}}$ and $\deg D=0$. Put $D'=D+P_0$, of degree $1$. By Ex.~1.9(d), $l(D')-l(K-D')=1+1-1=1$. The invertible sheaf $\mathcal L(K-D')$ has degree $-1<0$, so it has no nonzero sections by the Lemma, i.e.\ $l(K-D')=0$ and $l(D')=1$. Let $s$ be a nonzero section of $\mathcal L(D')$ (unique up to scalar) and $\mathcal T$ the cokernel of $s:\mathcal O_X\to\mathcal L(D')$. Then $\chi(\mathcal T)=\deg D'=1$, so $\mathcal T$ is a skyscraper sheaf $k(Q)$ at a closed point $Q$. Tensoring with $\mathcal L(D')^{-1}$ gives $0\to\mathcal L(D')^{-1}\to\mathcal O_X\to k(Q)\to0$, i.e.\ the ideal sheaf $\mathfrak m_Q$ of $Q$ is isomorphic to the invertible sheaf $\mathcal L(D')^{-1}$. A local ring of dimension $1$ whose maximal ideal is principal is a DVR, so $Q\in X_{\mathrm{reg}}$ and $\mathfrak m_Q=\mathcal L(-Q)$. Hence $\mathcal L(D')\cong\mathcal L(Q)$, i.e.\ $\mathcal L(D)\cong\mathcal L(Q-P_0)$.

\emph{Injectivity.} Suppose $\mathcal L(Q-P_0)\cong\mathcal L(Q'-P_0)$ with $Q,Q'\in X_{\mathrm{reg}}$. Then $\mathcal M=\mathcal L(Q)\cong\mathcal L(Q')$, and $h^0(\mathcal M)=1$ by the computation above (degree $1$). The canonical inclusion $1:\mathcal O_X\to\mathcal L(Q)$ has cokernel supported exactly at $Q$, and similarly for $Q'$; under the isomorphism these are two nonzero sections of $\mathcal M$, hence proportional, with cokernels that are isomorphic, so $Q=Q'$.

Hence $P\mapsto\mathcal L(P-P_0)$ is a bijection $X_{\mathrm{reg}}\to\operatorname{Pic}^0X$. \hfill$\square$

\section*{2. Hurwitz's theorem}

Conventions. $k$ is algebraically closed, ``curve'' means complete nonsingular curve over $k$, and $\mathcal L(D)$, $l(D)$, $K$ are as in Section~1. For a finite morphism $f:X\to Y$ of curves and $Q\in Y$ we use constantly $\sum_{P\mapsto Q}e_P=\deg f$ (II, 6.9). A finite morphism of nonsingular curves is flat, so $f_*\mathcal M$ is locally free of rank $\deg f$ for every invertible $\mathcal M$ on $X$. Hurwitz's theorem (Cor.~2.4) is used in the form $2g(X)-2=n(2g(Y)-2)+\deg R$ for $f$ separable of degree $n$, and $\deg R=\sum(e_P-1)$ if all ramification is tame, in particular always in characteristic $0$.

\exercise{2.1}{Use (2.5.3) to show that $\mathbf P^n$ is simply connected.}
\pf
``Simply connected'' means that every etale covering $f:X\to\mathbf P^n$ (a finite etale morphism) is trivial, i.e.\ $X$ is a finite disjoint union of copies of $\mathbf P^n$ (Example~2.5.3). We use induction on $n$; $n=0$ is trivial and $n=1$ is Example~2.5.3.

Let $n\ge2$ and let $f:X\to\mathbf P^n$ be etale and finite. Since $f$ is open and closed and $\mathbf P^n$ is connected, each connected component $X'$ of $X$ maps onto $\mathbf P^n$ with some degree $d'$, and $X'\to\mathbf P^n$ is again finite etale. It suffices to show $d'=1$ for each $X'$, because a finite etale morphism of degree $1$ is an isomorphism. So assume $X'$ is connected. It is smooth over $k$ (etale over $\mathbf P^n$), hence regular, connected, so an integral normal projective variety of dimension $n\ge2$ (finite over $\mathbf P^n$).

The sheaf $f^*\mathcal O(1)$ is ample on $X'$: for coherent $\mathcal F$ on $X'$, since $f$ is finite (hence affine), (III, Ex.~4.1) and the projection formula give $H^i(X',\mathcal F\otimes f^*\mathcal O(m))\cong H^i(\mathbf P^n,f_*\mathcal F\otimes\mathcal O(m))=0$ for $i>0$, $m\gg0$ (III, 5.2), and then Serre's criterion (III, 5.3) applies. Let $H\subseteq\mathbf P^n$ be a hyperplane. Then $f^*H=f^{-1}(H)$ is an effective Cartier divisor on $X'$ with $\mathcal L(f^*H)=f^*\mathcal O(1)$ ample, so by (III, 7.9) its support $f^{-1}(H)$ is connected. On the other hand $f^{-1}(H)\to H\cong\mathbf P^{n-1}$ is finite etale of degree $d'$ (base change), so by induction it is a trivial covering: $f^{-1}(H)\cong\bigsqcup_{d'}H$. Being connected, $d'=1$. Hence $X'\cong\mathbf P^n$, and $\mathbf P^n$ has no nontrivial etale covering. \hfill$\square$

\exercise{2.2}{\emph{Classification of Curves of Genus 2.} Fix an algebraically closed field $k$ of characteristic $\ne2$.
\begin{enumerate}[label=(\alph*),leftmargin=*]
\item If $X$ is a curve of genus $2$ over $k$, the canonical linear system $|K|$ determines a finite morphism $f:X\to\mathbf P^1$ of degree $2$ (Ex.~1.7). Show that it is ramified at exactly $6$ points, with ramification index $2$ at each one. Note that $f$ is uniquely determined, up to an automorphism of $\mathbf P^1$, so $X$ determines an (unordered) set of $6$ points of $\mathbf P^1$, up to an automorphism of $\mathbf P^1$.
\item Conversely, given six distinct elements $\alpha_1,\dots,\alpha_6\in k$, let $K$ be the extension of $k(x)$ determined by the equation $z^2=(x-\alpha_1)\cdots(x-\alpha_6)$. Let $f:X\to\mathbf P^1$ be the corresponding morphism of curves. Show that $g(X)=2$, the map $f$ is the same as the one determined by the canonical linear system, and $f$ is ramified over the six points $x=\alpha_i$ of $\mathbf P^1$, and nowhere else. (Cf.\ (II, Ex.~6.4).)
\item Using (I, Ex.~6.6), show that if $P_1,P_2,P_3$ are three distinct points of $\mathbf P^1$, then there exists a unique $\varphi\in\operatorname{Aut}\mathbf P^1$ such that $\varphi(P_1)=0$, $\varphi(P_2)=1$, $\varphi(P_3)=\infty$. Thus in (a), if we order the six points of $\mathbf P^1$, and then normalize by sending the first three to $0,1,\infty$, respectively, we may assume that $X$ is ramified over $0,1,\infty,\beta_1,\beta_2,\beta_3$, where $\beta_1,\beta_2,\beta_3$ are three distinct elements of $k$, $\ne0,1$.
\item Let $\Sigma_6$ be the symmetric group on $6$ letters. Define an action of $\Sigma_6$ on sets of three distinct elements $\beta_1,\beta_2,\beta_3$ of $k$, $\ne0,1$, as follows: reorder the set $0,1,\infty,\beta_1,\beta_2,\beta_3$ according to a given element $\sigma\in\Sigma_6$, then renormalize as in (c) so that the first three become $0,1,\infty$ again. Then the last three are the new $\beta_1',\beta_2',\beta_3'$.
\item Summing up, conclude that there is a one-to-one correspondence between the set of isomorphism classes of curves of genus $2$ over $k$, and triples of distinct elements $\beta_1,\beta_2,\beta_3$ of $k$, $\ne0,1$, modulo the action of $\Sigma_6$ described in (d). In particular, there are many non-isomorphic curves of genus $2$. We say that curves of genus $2$ depend on three parameters, since they correspond to the points of an open subset of $\mathbf A^3_k$ modulo a finite group.
\end{enumerate}}
\pf
\textbf{(a)} By Ex.~1.7(a), $|K|$ is base point free of degree $2$ and dimension $1$ and defines a finite morphism $f:X\to\mathbf P^1$ of degree $2$. Since $\operatorname{char}k\ne2$, the degree-$2$ extension $K(X)/k(f)$ is separable. For each $Q\in\mathbf P^1$, $\sum_{P\mapsto Q}e_P=2$, so $e_P\in\{1,2\}$; all ramification is tame since $\operatorname{char}k\ne2$. Hurwitz's theorem (2.4) gives $2\cdot2-2=2\cdot(0-2)+\deg R$, so $\deg R=6=\sum_P(e_P-1)=\#\{P:e_P=2\}$. Thus $f$ is ramified at exactly $6$ points, each with $e_P=2$; each lies alone in its fibre ($e_P=2=\deg f$), so the six ramification points map to six distinct branch points $B\subseteq\mathbf P^1$.

A different choice of basis of $H^0(X,\mathcal L(K))$ changes $f$ by an element of $\operatorname{PGL}_2=\operatorname{Aut}\mathbf P^1$ (I, Ex.~6.6): $f$ is the morphism defined by a basis of $H^0(X,\mathcal L(K))$, and changing the basis composes it with a projective linear transformation. Hence $f$ is unique up to $\operatorname{Aut}\mathbf P^1$, and $X$ determines the $6$-element set $B\subseteq\mathbf P^1$ up to $\operatorname{Aut}\mathbf P^1$.

\textbf{(b)} Put $h=\prod_{i=1}^6(x-\alpha_i)$. It is squarefree of degree $6$, hence not a square in $k(x)$, so $K=k(x)(z)$, $z^2=h$, is a field of degree $2$ over $k(x)$. Let $X$ be the nonsingular projective curve with function field $K$ and $f:X\to\mathbf P^1$ the morphism given by $k(x)\subseteq K$; it is finite of degree $2$ and separable.

\emph{The affine part.} Let $A=k[x,z]/(z^2-h)$. The Jacobian criterion: the partial derivatives are $2z$ and $-h'(x)$, which vanish simultaneously with $z^2=h$ only if $z=0$, $h(x)=0=h'(x)$, impossible since $h$ has simple roots. So $\operatorname{Spec}A$ is a nonsingular affine curve, $A$ is integrally closed, finite and free of rank $2$ over $k[x]$ (basis $1,z$); thus $A$ is the integral closure of $k[x]$ in $K$ and $\operatorname{Spec}A=f^{-1}(\mathbf A^1)$. For $a\in k$ we have $\dim_kA/(x-a)A=2=\sum_{P\mapsto a}e_P$ (since $A/(x-a)$ has this dimension and $\mathfrak m_P^{e_P}$-primary parts add up). The fibre ring is $k[z]/(z^2-h(a))$: if $h(a)\ne0$ it has two maximal ideals, so two points with $e=1$; if $h(a)=0$ it is $k[z]/(z^2)$, one point $P_a=(a,0)$ with $e=2$. At $P_a$ the maximal ideal is $(x-a,z)$ and $x-a=z^2\big/\big(h/(x-a)\big)$ with $h/(x-a)$ a unit at $a$, so $z$ is a local parameter at $P_a$, with $v_{P_a}(z)=1$.

\emph{The point at infinity.} Put $x'=1/x$, $w=z/x^3$; then $w^2=\prod(1-\alpha_ix')=:h^*(x')$, with $h^*(0)=1$ and $h^*$ squarefree (its roots are the $1/\alpha_i$, $\alpha_i\ne0$). The same argument applied to $B=k[x',w]/(w^2-h^*)$ (nonsingular, integral closure of $k[x']$) shows: over $x'=0$ the fibre ring is $k[w]/(w^2-1)$, so there are two points $\infty_\pm$ ($w=\pm1$), unramified.

Hence $f$ is ramified exactly at the six points $P_{\alpha_i}$ (with $e=2$), over the six points $x=\alpha_i$, and nowhere else. Hurwitz: $2g(X)-2=2(-2)+6=2$, so $g(X)=2$.

\emph{Comparison with $|K|$.} Let $F=f^*(\infty)=\infty_++\infty_-$; any fibre of $f$ is linearly equivalent to $F$. The function $z$ has $v_{P_{\alpha_i}}(z)=1$, is a unit at the other finite points, and at $\infty_\pm$: $z=w/x'^3$ with $w$ a unit and $x'$ a local parameter, so $v_{\infty_\pm}(z)=-3$. Hence $(z)=R-3F$, where $R=\sum P_{\alpha_i}$ is the ramification divisor, i.e.\ $R\sim3F$. By Prop.~2.3, $K_X\sim f^*K_{\mathbf P^1}+R\sim-2F+3F=F$. Then $\mathcal L(K_X)\cong\mathcal L(F)=f^*\mathcal O(1)$, and $f^*:H^0(\mathbf P^1,\mathcal O(1))\to H^0(X,f^*\mathcal O(1))$ is injective with $2$-dimensional source, while $\dim H^0(X,\mathcal L(K_X))=g=2$. So it is an isomorphism: $|K_X|=f^*|\mathcal O(1)|$, and the morphism defined by the complete linear system $|K_X|$ is exactly $f$ (for the basis $f^*s_0,f^*s_1$).

\textbf{(c)} By (I, Ex.~6.6), $\operatorname{Aut}\mathbf P^1=\operatorname{PGL}_2(k)$ acting by $x\mapsto(ax+b)/(cx+d)$. Existence: for $P_1,P_2,P_3\in k$ take $\varphi(x)=\dfrac{(x-P_1)(P_2-P_3)}{(x-P_3)(P_2-P_1)}$; if one of the $P_i$ is $\infty$ omit the two factors containing it (e.g.\ $P_3=\infty$: $\varphi=(x-P_1)/(P_2-P_1)$). One checks directly $\varphi(P_1)=0,\varphi(P_2)=1,\varphi(P_3)=\infty$. Uniqueness: if $\psi\in\operatorname{PGL}_2$ fixes $0,1,\infty$, then $\psi(\infty)=\infty$ gives $c=0$, $\psi(0)=0$ gives $b=0$, $\psi(1)=1$ gives $a=d$, so $\psi=\mathrm{id}$.

Now order the six points $B$ of (a) as $(p_1,\dots,p_6)$ and let $\varphi$ send $p_1,p_2,p_3$ to $0,1,\infty$. Replacing $f$ by $\varphi\circ f$ (again a degree-$2$ finite morphism, ramified exactly over $\varphi(B)$, by the uniqueness in (a)) we get a curve ramified over $0,1,\infty,\beta_1,\beta_2,\beta_3$ with $\beta_i=\varphi(p_{3+i})$ distinct elements of $\mathbf P^1-\{0,1,\infty\}=k-\{0,1\}$.

\textbf{(d)} Let $S$ be the set of ordered $6$-tuples $(p_1,\dots,p_6)$ of distinct points of $\mathbf P^1$, and $T$ the set of ordered triples $(\beta_1,\beta_2,\beta_3)$ of distinct elements of $k-\{0,1\}$. By (c), $\operatorname{PGL}_2$ acts on $S$ freely (a transformation fixing three points is the identity), and every orbit contains exactly one normalized tuple $(0,1,\infty,\beta_1,\beta_2,\beta_3)$; so $T\cong S/\operatorname{PGL}_2$. The group $\Sigma_6$ acts on $S$ by permuting the entries, and this commutes with the $\operatorname{PGL}_2$-action; hence it induces an action on $S/\operatorname{PGL}_2\cong T$. Concretely, $\sigma\in\Sigma_6$ sends the normalized tuple to the permuted tuple, which is then renormalized by the unique $\varphi$ of (c): this is exactly the procedure of (d), so it is a genuine action. Its orbit set is $T/\Sigma_6=\{\text{6-element subsets of }\mathbf P^1\}/\operatorname{PGL}_2$.

\textbf{(e)} Define $\Theta$ from the set of isomorphism classes of genus-$2$ curves to $T/\Sigma_6$ by $X\mapsto$ class of the branch set $B(X)$ (well defined by (a), (c), (d); isomorphic curves have the same canonical map up to $\operatorname{Aut}\mathbf P^1$).

\emph{Surjective.} Given a $6$-element set $B\subseteq\mathbf P^1$, pick $c\in k-B$ and apply the automorphism $x\mapsto1/(x-c)$ of $\mathbf P^1$, which moves $B$ into $k$ (the only point sent to $\infty$ is $c$). So we may assume $B=\{\alpha_1,\dots,\alpha_6\}\subseteq k$, and the curve of (b) has genus $2$ and branch set $B$.

\emph{Injective.} Let $X,X'$ be genus-$2$ curves with branch sets $B,B'=\sigma(B)$, $\sigma\in\operatorname{Aut}\mathbf P^1$. Replace $f'$ by $\sigma^{-1}f'$ so that both have branch set $B$, and (as above) move $B$ into $k$. Then $K(X)\supseteq k(x)$ is a quadratic extension, and since $\operatorname{char}k\ne2$ it is $k(x)(\sqrt u)$ with $u\in k(x)^*$; as $k$ is algebraically closed we may take $u=\prod_j(x-\gamma_j)$ with distinct $\gamma_j\in k$ (modify $u$ by squares and constants). The extension is ramified over a finite $a\in k$ iff $v_a(u)$ is odd, i.e.\ iff $a\in\{\gamma_j\}$, and over $\infty$ iff $\deg u$ is odd. (If $v_a(u)$ is odd and $P\mapsto a$ had $e_P=1$, then $v_P(u)=v_a(u)$ would be odd, but $u$ is a square in $K(X)$, so $v_P(u)$ is even. If $v_a(u)=2m$, then $y=\sqrt u/(x-a)^m$ satisfies $y^2=w$ with $w=u/(x-a)^{2m}$ a unit of $\mathcal O_a$; the fibre of $\mathcal O_a[y]=\mathcal O_a[Y]/(Y^2-w)$ over $a$ is $k[Y]/(Y^2-w(a))\cong k\times k$ since $\operatorname{char}k\ne2$, so the integral closure has two points over $a$, each with $e=1$. The same applies at $\infty$ with $x'=1/x$ and $u\,x'^{\,\deg u}$.) As $B\subseteq k$ is the branch set, $\{\gamma_j\}=B=\{\alpha_1,\dots,\alpha_6\}$ and $\deg u=6$ is even. Thus $K(X)\cong k(x)(z)$ with $z^2=\prod(x-\alpha_i)$, the same for $X'$; since a curve is determined by its function field (I, 6.12), $X\cong X'$. So $\Theta$ is a bijection.

Finally $T$ is infinite ($k$ is infinite) and $\Sigma_6$ is finite, so $T/\Sigma_6$ is infinite: there are infinitely many non-isomorphic curves of genus $2$, depending on $3$ parameters ($T$ is an open subset of $\mathbf A^3_k$). \hfill$\square$

\exercise{2.3}{\emph{Plane Curves.} Let $X$ be a curve of degree $d$ in $\mathbf P^2$. For each point $P\in X$, let $T_P(X)$ be the tangent line to $X$ at $P$ (I, Ex.~7.3). Considering $T_P(X)$ as a point of the dual projective plane $(\mathbf P^2)^*$, the map $P\to T_P(X)$ gives a morphism of $X$ to its dual curve $X^*$ in $(\mathbf P^2)^*$ (I, Ex.~7.3). Note that even though $X$ is nonsingular, $X^*$ in general will have singularities. We assume $\operatorname{char}k=0$ below.
\begin{enumerate}[label=(\alph*),leftmargin=*]
\item Fix a line $L\subseteq\mathbf P^2$ which is not tangent to $X$. Define a morphism $\varphi:X\to L$ by $\varphi(P)=T_P(X)\cap L$, for each point $P\in X$. Show that $\varphi$ is ramified at $P$ if and only if either (1) $P\in L$, or (2) $P$ is an \emph{inflection point} of $X$, which means that the intersection multiplicity (I, Ex.~5.4) of $T_P(X)$ with $X$ at $P$ is $\ge3$. Conclude that $X$ has only finitely many inflection points.
\item A line of $\mathbf P^2$ is a \emph{multiple tangent} of $X$ if it is tangent to $X$ at more than one point. It is a \emph{bitangent} if it is tangent to $X$ at exactly two points. If $L$ is a multiple tangent of $X$, tangent to $X$ at the points $P_1,\dots,P_r$, and if none of the $P_i$ is an inflection point, show that the corresponding point of the dual curve $X^*$ is an \emph{ordinary $r$-fold point}, which means a point of multiplicity $r$ with distinct tangent directions (I, Ex.~5.3). Conclude that $X$ has only finitely many multiple tangents.
\item Let $O\in\mathbf P^2$ be a point which is not on $X$, nor on any inflectional or multiple tangent of $X$. Let $L$ be a line not containing $O$. Let $\psi:X\to L$ be the morphism defined by projection from $O$. Show that $\psi$ is ramified at a point $P\in X$ if and only if the line $OP$ is tangent to $X$ at $P$, and in that case the ramification index is $2$. Use Hurwitz's theorem and (I, Ex.~7.2) to conclude that there are exactly $d(d-1)$ tangents of $X$ passing through $O$. Hence the degree of the dual curve (sometimes called the \emph{class} of $X$) is $d(d-1)$.
\item Show that for all but a finite number of points of $X$, a point $O$ of $X$ lies on exactly $(d+1)(d-2)$ tangents of $X$, not counting the tangent at $O$.
\item Show that the degree of the morphism $\varphi$ of (a) is $d(d-1)$. Conclude that if $d\ge2$, then $X$ has $3d(d-2)$ inflection points, properly counted. (If $T_P(X)$ has intersection multiplicity $r$ with $X$ at $P$, then $P$ should be counted $r-2$ times as an inflection point.) If $r=3$ we call it an \emph{ordinary inflection point}. Show that an ordinary inflection point of $X$ corresponds to an ordinary cusp of the dual curve $X^*$.
\item Now let $X$ be a plane curve of degree $d\ge2$, and assume that the dual curve $X^*$ has only nodes and ordinary cusps as singularities (which should be true for sufficiently general $X$). Then show that $X$ has exactly $\frac12d(d-2)(d-3)(d+3)$ bitangents. [\emph{Hint}: Show that $X$ is the normalization of $X^*$. Then calculate $p_a(X^*)$ two ways: once as a plane curve of degree $d(d-1)$, and once using (Ex.~1.8).]
\item For example, a plane cubic curve has exactly $9$ inflection points, all ordinary. The line joining any two of them intersects the curve in a third one.
\item A plane quartic curve has exactly $28$ bitangents. (This holds even if the curve has a tangent with four-fold contact, in which case the dual curve $X^*$ has a tacnode.)
\end{enumerate}}
\pf
\emph{Standing assumptions and notation.} $\operatorname{char}k=0$, and $X=Z(F)\subseteq\mathbf P^2$ is a nonsingular (hence irreducible) curve of degree $d$. For $d=1$ everything is degenerate ($X$ is a line, equal to all its tangents, the map $\varphi$ of (a) is constant and $X^*$ is a point), so we assume $d\ge2$ throughout; then no tangent line is contained in $X$, so every tangent meets $X$ at $P$ with a finite multiplicity. Genus: $g=(d-1)(d-2)/2$ (I, Ex.~7.2). We identify points of $\mathbf P^2$ with vectors of $k^3$ (up to scalars), and lines with covectors, written as vectors via $\ell\cdot p=0$; the dual plane is $\mathbf P(k^3)$ again, a point $\ell$ being the line $\{p:\ell\cdot p=0\}$ and a point $p$ of $\mathbf P^2$ being the line $p^*=\{\ell:\ell\cdot p=0\}$ of $(\mathbf P^2)^*$. We write $[a,b,c]=\det(a,b,c)=a\cdot(b\times c)$. The tangent map is $\gamma:X\to X^*\subseteq(\mathbf P^2)^*$, $\gamma(P)=[\nabla F(P)]$ (I, Ex.~7.3). For a line $L=\{\lambda\cdot p=0\}$ and a point $\ell$ of the dual plane (a line) the point $\ell\cap L$ is $\ell\times\lambda$, and the line through two points $p,q$ is $p\times q$.

\textbf{Local computations (L1)--(L3).} Let $P\in X$, let $t$ be a local parameter at $P$, identify $\widehat{\mathcal O}_{X,P}=k[[t]]$ and let $r(t)\in k[[t]]^3$ be a lift of the arc $\operatorname{Spec}k[[t]]\to X\subseteq\mathbf P^2$, with $r(0)=P$. Primes are $d/dt$.

(L1) $r(0),r'(0)$ are linearly independent. Indeed in an affine chart $r=(x(t),y(t),1)$; $k[[x,y]]\to k[[t]]$ is onto, so $x$ or $y$ has order exactly $1$, i.e.\ $(x'(0),y'(0))\neq0$, and $r'(0)=(x'(0),y'(0),0)$.

Put $\ell(t)=r\times r'$ and $D(t)=[r,r',r'']$. Then $\ell(0)\ne0$, and

(L2) $[\ell(t)]=\gamma(P(t))$ (the tangent line at $P(t)$) and $\ell\cdot r=\ell\cdot r'=\ell'\cdot r=0$. Indeed $F(r(t))=0$; differentiating, $\nabla F(r)\cdot r'=0$, and Euler's relation gives $\nabla F(r)\cdot r=dF(r)=0$; so $\nabla F(r(t))$ is orthogonal to $r,r'$ and over $k((t))$ it is a multiple $c(t)\ell(t)$ of $\ell=r\times r'$; since $\nabla F(r(0))\ne0$ (nonsingularity) and $\ell(0)\ne0$, $c$ is a unit of $k[[t]]$. The relations $\ell\cdot r=\ell\cdot r'=0$ are clear, and differentiating $\ell\cdot r=0$ gives $\ell'\cdot r=-\ell\cdot r'=0$. In particular $\ell(0)=T_P(X)$.

(L3) $\ell\times\ell'=D\,r$ (use $\ell'=r\times r''$ and $(a\times b)\times(a\times c)=[a,b,c]a$), and $\operatorname{ord}_tD=m_P-2$, where $m_P=\operatorname{ord}_t(\ell(0)\cdot r(t))=(T_P(X)\cdot X)_P\ge2$ is the intersection multiplicity of the tangent line with $X$ at $P$ (I, Ex.~5.4; for a line $\{\lambda\cdot p=0\}$ the multiplicity at $P\in X$ is $\operatorname{ord}_t(\lambda\cdot r(t))$). To compute $\operatorname{ord}D$ note that $D$ is changed by a unit factor if $r$ is replaced by $ur$ ($D\mapsto u^3D$), by a change of parameter $t=\tau(s)$ ($D\mapsto\tau'^3D$), or by a linear change of coordinates ($D\mapsto\det\cdot D$). So we may take affine coordinates with $P=(0,0)$ and $T_P(X)=\{y=0\}$; then $y$ has order $m_P\ge2$ on $X$ at $P$, so $x$ has order $1$ and we may take $t=x$, $r=(t,g(t),1)$ with $\operatorname{ord}g=m_P$. Then $D=\det\begin{pmatrix}t&g&1\\1&g'&0\\0&g''&0\end{pmatrix}=g''$, which has order exactly $m_P-2$ because $\operatorname{char}k=0$.

\textbf{(a)} The morphism: with $L=\{\lambda\cdot p=0\}$, $\varphi(P)=\nabla F(P)\times\lambda$, given by forms of degree $d-1$ with no common zero on $X$ (as $\nabla F(P)\times\lambda=0$ iff $T_P=L$). Let $v(t)=\ell(t)\times\lambda\in k[[t]]^3$, so $[v(t)]=\varphi(P(t))$, with $v(0)\ne0$ since $\ell(0)\ne\lambda$. As $v\cdot\lambda=0$, $v$ lies in the plane $\lambda^\perp$ representing $L$; choose a basis $e_1,e_2$ of $\lambda^\perp$ so that $v=v_0e_1+v_1e_2$ with $v_0(0)\ne0$. Then $u=v_1/v_0$ is the pull-back of an affine coordinate on $L$ at $\varphi(P)$, so
\[
e_P(\varphi)=\operatorname{ord}_t\bigl(u-u(0)\bigr)=1+\operatorname{ord}_tu'=1+\operatorname{ord}_t(v_0v_1'-v_1v_0')=1+\operatorname{ord}_t(v\times v'),
\]
using $\operatorname{char}k=0$ and $v\times v'=(v_0v_1'-v_1v_0')\,e_1\times e_2$ with $e_1\times e_2\ne0$ constant. By the identity $(a\times b)\times(c\times e)=[a,b,e]c-[a,b,c]e$,
\[
v\times v'=(\ell\times\lambda)\times(\ell'\times\lambda)=-[\ell,\lambda,\ell']\lambda=\pm(\ell\times\ell')\cdot\lambda\;\lambda=\pm D\,(r\cdot\lambda)\,\lambda
\]
by (L3). Hence, with $(L\cdot X)_P=\operatorname{ord}_t(\lambda\cdot r)\ge0$ ($>0$ iff $P\in L$):
\[
e_P(\varphi)=1+(m_P-2)+(L\cdot X)_P.\tag{$\ast$}
\]
(In particular $u'\neq0$, so $\varphi$ is nonconstant.) Both summands $m_P-2$ and $(L\cdot X)_P$ are $\ge0$, so $e_P(\varphi)>1$ iff $m_P\ge3$ (i.e.\ $P$ is an inflection point) or $P\in L$. Since $\varphi$ is a finite separable morphism of curves, it has only finitely many ramification points (Prop.~2.2(a)), and every inflection point is one of them. So $X$ has only finitely many inflection points.

\textbf{Lemma B ($\gamma:X\to X^*$ is birational; tangent to the dual branch).} Let $d\ge2$. The map $\gamma$ is not constant (otherwise all tangent lines coincide with one line $L_0\ni P$ for all $P$, so $X\subseteq L_0$). So $\gamma:X\to X^*$ is a finite surjective morphism onto an integral plane curve. By (L2), $\ell(0)\cdot r(0)=0=\ell'(0)\cdot r(0)$, so $\operatorname{span}(\ell(0),\ell'(0))\subseteq r(0)^\perp$, and by (L3) $\ell(0),\ell'(0)$ are linearly independent iff $\ell\times\ell'(0)=D(0)r(0)\ne0$ iff $m_P=2$. So at a non-inflection point $P$, the arc $\gamma(P(t))=[\ell(t)]$ is smooth at $t=0$ (an immersion), with tangent line equal to $\operatorname{span}(\ell(0),\ell'(0))=r(0)^\perp$, i.e.\ the line $P^*\subseteq(\mathbf P^2)^*$ of lines through $P$. To see $\deg\gamma=1$: choose $Q\in X^*$ with $X^*$ smooth at $Q$ and with $\gamma^{-1}(Q)$ consisting of $\deg\gamma$ points (char $0$), none of which is an inflection point (finitely many exceptions). If $P\ne P'$ both map to $Q$, the tangent line of $X^*$ at $Q$ equals both $P^*$ and $P'^*$, so $P=P'$. Hence $\deg\gamma=1$, and $\gamma$ is birational: since $X$ is nonsingular, $\gamma$ is the normalization of $X^*$.

\textbf{(b)} Let $L$ be a multiple tangent, tangent at exactly $P_1,\dots,P_r$ ($r\ge2$) and none an inflection point; let $Q=[L]\in X^*$. The points of $X$ over $Q$ are exactly $P_1,\dots,P_r$, and by Lemma B the branches of $X^*$ at $Q$ (the images of the formal neighborhoods of the $P_i$) are $r$ smooth branches with distinct tangent lines $P_i^*$. For a line $\Lambda\subseteq(\mathbf P^2)^*$ through $Q$, with local equation $\mu$, the intersection multiplicity is $(\Lambda\cdot X^*)_Q=\sum_i\operatorname{ord}_{P_i}(\gamma^*\mu)$: since $\mathcal O_{X^*,Q}\subseteq\widetilde{\mathcal O}$ is finite birational with finite-length quotient, $\dim\mathcal O/\mu\mathcal O=\dim\widetilde{\mathcal O}/\mu\widetilde{\mathcal O}$ (multiplicities $e(\mu;\cdot)$ are additive and vanish on finite-length modules), and the right side is the sum above. Each summand is $\ge1$, and $=1$ iff $\Lambda$ is not the tangent line $P_i^*$ of the $i$-th branch. Thus the multiplicity of $X^*$ at $Q$ (the minimum of $(\Lambda\cdot X^*)_Q$ over $\Lambda\ni Q$, I, Ex.~5.4(b)) is $r$, and the tangent cone consists of the $r$ distinct lines $P_1^*,\dots,P_r^*$: $Q$ is an ordinary $r$-fold point.

Consequence: $r\ge2$ means $X^*$ is singular at $Q$, and a plane curve has finitely many singular points; so there are finitely many multiple tangents without an inflection point of tangency. A multiple tangent tangent at an inflection point is the tangent line at one of the finitely many inflection points (part (a)). So there are finitely many multiple tangents.

\textbf{(c)} Choose coordinates with $O=(0:0:1)$ and $L=\{z=0\}$, so $\psi(x:y:z)=(x:y)$; $\psi$ is finite and separable (char $0$). Let $Q\in L$, say $Q=(1:b)$ (the other chart is symmetric); the fibre $\psi^{-1}(Q)$ is the set of points of $X$ on the line $OQ=\{y=bx\}$, and $x\ne0$ at those points (else the point is $O\notin X$). The local parameter $s=y/x-b$ at $Q$ pulls back to $(y-bx)/x$, and $x$ is a unit, so for $P\in\psi^{-1}(Q)$:
\[
e_P(\psi)=\operatorname{ord}_P(y-bx)=(OP\cdot X)_P,
\]
the intersection multiplicity of the line $OP$ with $X$ at $P$. In particular $\deg\psi=\sum_{P\mapsto Q}e_P=d$. A line through $P$ has multiplicity $\ge2$ iff it is the tangent $T_P(X)$, so $\psi$ is ramified at $P$ iff $OP$ is tangent to $X$ at $P$. In that case $O\in T_P(X)$ and, as $O$ is on no inflectional tangent, $P$ is not an inflection point, so $e_P=m_P=2$.

By Hurwitz (2.4, tame since char $0$): $2g-2=d(-2)+\deg R$ with $\deg R=\#\{P:e_P=2\}$, and $2g-2=d(d-3)$ by (I, Ex.~7.2). So $\deg R=d(d-3)+2d=d(d-1)$. Distinct ramification points have distinct tangent lines through $O$ (otherwise that line would be a multiple tangent through $O$), and every tangent line through $O$ is tangent at some ramification point. So exactly $d(d-1)$ tangent lines of $X$ pass through $O$.

Class: the line $O^*\subseteq(\mathbf P^2)^*$ meets $X^*$ in the points $\gamma(P)$ with $O\in T_P(X)$, i.e.\ $\ell(0)\cdot O=0$. The pull-back to $X$ of the equation of $O^*$ is $\ell(t)\cdot O$; at such a $P$ it vanishes to order exactly $1$, since $\ell(0)\cdot O=\ell'(0)\cdot O=0$ would force $O\perp\operatorname{span}(\ell(0),\ell'(0))=r(0)^\perp$ (Lemma B), i.e.\ $O=P\in X$. So $\gamma^*\mathcal O(1)$ has a divisor of $d(d-1)$ distinct simple points; thus $\deg\gamma^*\mathcal O(1)=d(d-1)$. As $\gamma$ is birational onto $X^*$ (Lemma B), $\deg X^*=d(d-1)$.

\textbf{(d)} Let $O\in X$ be such that $O$ is not an inflection point and does not lie on any inflectional tangent or multiple tangent (these are finitely many lines by (a),(b), each meeting $X$ in finitely many points since $X$ contains no line, and finitely many inflection points): this excludes only finitely many $O$. Let $L\not\ni O$ be a line and $\psi_O:X\to L$ the projection from $O$, extended over $O$ by $\psi_O(O)=T_O(X)\cap L$ (a morphism since $X$ is a nonsingular curve). A general line through $O$ meets $X$ in $O$ (simply) and $d-1$ other points, so $\deg\psi_O=d-1$. For $P\ne O$ the computation of (c) gives $e_P(\psi_O)=(OP\cdot X)_P$. At $P=O$ take affine coordinates with $O=(0,0)$, $X$ parametrized by $r=(t,g(t),1)$, $\operatorname{ord}g=m_O=2$: then $\psi_O=(x:y)=(1:t^{m_O-1}u(t))$ with $u$ a unit, so $e_O=m_O-1=1$: $\psi_O$ is unramified at $O$. For $P\ne O$, $\psi_O$ is ramified iff $OP=T_P(X)$, and then $P$ is not an inflection point (else $O$ would lie on an inflectional tangent) and $e_P=m_P=2$; moreover distinct such $P$ have distinct tangents through $O$ (no multiple tangent through $O$), and $T_P(X)\ne T_O(X)$ for $P\ne O$ since $T_O(X)$ is not a multiple tangent. By Hurwitz,
\[
d(d-3)=2g-2=(d-1)(-2)+\deg R,\qquad \deg R=d^2-d-2=(d+1)(d-2),
\]
and $\deg R$ is the number of points $P\ne O$ with $O\in T_P(X)$, i.e.\ the number of tangents through $O$ other than $T_O(X)$. This is $(d+1)(d-2)$, for all but finitely many $O\in X$.

\textbf{(e)} Take $L$ as in (a). The fibre $\varphi^{-1}(Q)$ over $Q\in L$ is $\{P:Q\in T_P(X)\}$. For all but finitely many $Q\in L$ the point $Q$ is not on $X$ and not on any inflectional or multiple tangent (these are finitely many lines, different from $L$ because $L$ is not tangent, so each meets $L$ in one point), so by (c) there are exactly $d(d-1)$ tangents through $Q$, tangent at $d(d-1)$ distinct points. Thus $\deg\varphi=d(d-1)$ ($\varphi$ is separable).

Hurwitz for $\varphi:X\to L\cong\mathbf P^1$: $d(d-3)=2g-2=d(d-1)(-2)+\deg R_\varphi$, so $\deg R_\varphi=d(d-3)+2d(d-1)=3d^2-5d$. By $(\ast)$, $\deg R_\varphi=\sum_P(e_P-1)=\sum_P(m_P-2)+\sum_P(L\cdot X)_P=\sum_P(m_P-2)+d$ (Bezout for a line: $L\cdot X$ has degree $d$; $L$ is not contained in $X$). Hence
\[
\sum_{P\in X}(m_P-2)=3d^2-5d-d=3d(d-2),
\]
i.e.\ $X$ has $3d(d-2)$ inflection points counted with weight $m_P-2$.

Ordinary inflection point ($m_P=3$): take coordinates as in (L3), $r=(t,g,1)$ with $g=at^3+\dots$, $a\ne0$. Then $\ell=r\times r'=(-g',\,1,\,tg'-g)$, so in the chart of $(\mathbf P^2)^*$ where the middle coordinate is $1$ the branch of $X^*$ at $\gamma(P)$ is $(u,v)=(-g',tg'-g)=(-3at^2+\cdots,\,2at^3+\cdots)$: orders $(2,3)$. Reparametrize by $s$ with $u=s^2$ (a square root of a unit exists in char $0$); then $v=\alpha(s^2)+s^3\beta(s^2)$ with $\beta(0)\ne0$ and $\operatorname{ord}\alpha\ge2$, so with $v'=v-\alpha(u)$: $v'^2=u^3\beta(u)^2$. So the branch has equation $v'^2-u^3\beta(u)^2=0$: multiplicity $2$, tangent cone the double line $v'^2$, nonzero $x^3$-coefficient: an ordinary cusp (in the sense of Ex.~1.8(c)). (If $T_P(X)$ happens to be tangent to $X$ at other points too, the point of $X^*$ has further branches; the assertion concerns the branch from $P$, and the point is an ordinary cusp when $T_P$ is tangent only at $P$.)

\textbf{(f)} By Lemma B, $X$ is the normalization of $X^*$, so $p_a(\widetilde{X^*})=g=(d-1)(d-2)/2$. By (c) $X^*$ has degree $d(d-1)$, so $p_a(X^*)=\frac{(D-1)(D-2)}2$ with $D=d(d-1)$. By Ex.~1.8(a),(c) (nodes and ordinary cusps have $\delta=1$, smooth points $\delta=0$):
\[
p_a(X^*)=g+\nu+\kappa,
\]
where $\nu$ is the number of nodes and $\kappa$ the number of cusps of $X^*$.

Determination of $\nu,\kappa$. Let $Q=[L]\in X^*$, with $P_1,\dots,P_r$ the points of tangency of $L$ ($r\ge1$). By Lemma B the branch at $P_i$ is smooth iff $P_i$ is not an inflection point; for an inflection point with contact $m$ the branch has orders $(m-1,m)$ (computed as in (e)), multiplicity $m-1\ge2$. So the multiplicity of $X^*$ at $Q$ is $\sum_i\mu_i$ with $\mu_i=1$ or $m_{P_i}-1$. Singular points have multiplicity $2$ by hypothesis, so either $r=2$ and both branches are smooth, meeting with distinct tangents $P_1^*\ne P_2^*$ (a node), or $r=1$ and $m_P=3$ (a cusp, as in (e)). Conversely every bitangent gives a point with $r=2$, hence a singular point, hence (by the hypothesis) a node with both tangency points non-inflectional; every inflection point gives a singular point, hence a cusp with $m_P=3$ and $r=1$. Therefore: nodes of $X^*$ $\leftrightarrow$ bitangents of $X$ ($\nu$ = number of bitangents), and cusps $\leftrightarrow$ inflection points, all ordinary, so by (e) $\kappa=3d(d-2)$. Then
\[
\nu=\frac{(D-1)(D-2)}2-\frac{(d-1)(d-2)}2-3d(d-2).
\]
With $(D-1)(D-2)=d^2(d-1)^2-3d(d-1)+2$ and $(d-1)(d-2)=d^2-3d+2$ the difference of the first two terms is $\frac12(d^4-2d^3-3d^2+6d)$; subtracting $\frac12(6d^2-12d)$ gives $\nu=\frac12(d^4-2d^3-9d^2+18d)=\frac12d(d-2)(d^2-9)=\tfrac12d(d-2)(d-3)(d+3)$.

\textbf{(g)} Let $d=3$. A line meets $X$ in a divisor of degree $3$, so $m_P\le3$ for every $P$ (and $X$ contains no line). By (e), $\sum(m_P-2)=3\cdot3\cdot1=9$, with weights $1$ at the inflection points ($m_P=3$): so there are exactly $9$ inflection points, all ordinary.

Let $P\ne Q$ be inflection points, and $H$ the hyperplane class. The line $PQ$ is not $T_P(X)$ (which meets $X$ only at $P$, with multiplicity $3$), nor $T_Q(X)$, so it meets $P$ and $Q$ with multiplicity $1$ and $PQ\cdot X=P+Q+R$ with a third point $R\ne P,Q$. We show that $R$ is an inflection point. Take an inflection point $O$ as origin of the group law on $X$ (genus $1$, Example~1.3.7), so that $\mathcal L(P-O)$ etc.\ are the elements of $\operatorname{Pic}^0X$. A point $S$ is an inflection point iff $3S\sim H$ (if $3S=T_S\cdot X\sim H$; conversely, $l(H)=3=h^0(\mathbf P^2,\mathcal O(1))$ by Riemann--Roch, so every divisor $\sim H$ is a line section, and $3S\sim H$ gives a line meeting $X$ only at $S$). Since $H\sim3O$, $S$ is an inflection point iff $3(S-O)\sim0$. Now $P+Q+R\sim H\sim3O$ gives $(R-O)\sim-(P-O)-(Q-O)$, so $3(R-O)\sim0$, and $R$ is an inflection point. \hfill(Here $d=3$ and char $0$ are used.)

\textbf{(h)} For a general quartic $d=4$ and (f) gives $\frac12\cdot4\cdot2\cdot1\cdot7=28$ bitangents. For \emph{every} nonsingular quartic one can argue directly as follows (characteristic $0$). A tangent line $T_P(X)$ meets $X$ in a divisor of degree $4$ which is of one of the forms $2P+Q+R$ ($Q\ne R$), $2P+2Q$ (bitangent), $3P+Q$ (ordinary flex), $4P$ (hyperflex, $m_P=4$), with $P,Q,R$ distinct. By Lemma B and (b), (e), the singular points of $X^*$ are: a node at each bitangent, a cusp of orders $(2,3)$ at each ordinary flex (its tangent line meets $X$ only at $P$ and simply at $Q$, so there is a single branch), and a unibranch point with orders $(3,4)$ at each hyperflex. For orders $(3,4)$ the value semigroup of $k[[u,v]]\subseteq k[[t]]$ is $\{0,3,4,6,7,8,\dots\}$ (no cancellation can create the values $1,2,5$), so $\delta=\#\text{gaps}=3$ (Step~1 of Ex.~1.8(c): $\delta=\dim k[[t]]/\widehat{\mathcal O}$). Let $\nu$ be the number of (honest) bitangents, $\kappa$ the number of ordinary flexes, $h$ the number of hyperflexes. By (e), $\kappa+2h=3\cdot4\cdot2=24$. By Ex.~1.8(a): $p_a(X^*)=\frac{11\cdot10}2=55=g+\nu+\kappa+3h=3+\nu+\kappa+3h$. Hence $\nu=52-\kappa-3h=28-h$. So the number of bitangents, counting a hyperflex line (contact $2P+2Q$ with $P=Q$) as a degenerate bitangent, is always $\nu+h=28$. With the literal definition (``tangent at exactly two points'') the number is $28-h$, which is $28$ exactly when there is no hyperflex. (Remark: the singularity of $X^*$ at a hyperflex line is not a tacnode as stated in the book, but a unibranch singularity with orders $(3,4)$; the count $28$ in the sense above is unaffected.) \hfill$\square$

\exercise{2.4}{\emph{A Funny Curve in Characteristic $p$.} Let $X$ be the plane quartic curve $x^3y+y^3z+z^3x=0$ over a field of characteristic $3$. Show that $X$ is nonsingular, every point of $X$ is an inflection point, the dual curve $X^*$ is isomorphic to $X$, but the natural map $X\to X^*$ is purely inseparable.}
\pf
Let $F=x^3y+y^3z+z^3x$ over the algebraically closed field $k$ of characteristic $3$. Then $F_x=3x^2y+z^3=z^3$, $F_y=x^3+3y^2z=x^3$, $F_z=y^3+3z^2x=y^3$. These vanish simultaneously only at $x=y=z=0$, so $X=Z(F)$ is nonsingular (and irreducible, as a nonsingular plane curve is irreducible: two components would meet by Bezout).

\emph{The tangent line.} At $P=(a:b:c)\in X$ the tangent line is $T_P=\{c^3x+a^3y+b^3z=0\}$, so the dual map is $\gamma(P)=(c^3:a^3:b^3)$.

\emph{Every point is an inflection point.} Let $P=(a,b,c)\in X$, and let $R=(x',y',z')\in T_P$ be any point of the tangent line different from $P$; the points of the line are $P+sR$. Since the characteristic is $3$, $(x+sx')^3=x^3+s^3x'^3$, hence
\[
F(P+sR)=F(P)+s\,(a^3y'+b^3z'+c^3x')+s^3\,(bx'^3+cy'^3+az'^3)+s^4F(R).
\]
Here $F(P)=0$ and $a^3y'+b^3z'+c^3x'=0$ because $R\in T_P$. So $F(P+sR)=s^3(\cdots)$ is divisible by $s^3$, i.e.\ the intersection multiplicity of $T_P$ with $X$ at $P$ is $\ge3$. Thus every point of $X$ is an inflection point (in the sense of Ex.~2.3: contact $\ge3$).

\emph{The dual curve.} Let $(u:v:w)=(c^3:a^3:b^3)=\gamma(P)$ in $(\mathbf P^2)^*$ with dual coordinates. Since $F$ has coefficients in $\mathbf F_3$, $F(a,b,c)^3=F(a^3,b^3,c^3)$, so $0=F(a^3,b^3,c^3)=F(v,w,u)=v^3w+w^3u+u^3v=F(u,v,w)$ (note $F(v,w,u)=F(u,v,w)$). So $\gamma(X)\subseteq Z(F)=:X'\subseteq(\mathbf P^2)^*$, an irreducible curve; $\gamma$ is not constant, so $X^*=\gamma(X)=X'$, which is the same quartic $F=0$ written in the dual coordinates: $X^*\cong X$ (via the identity of the coordinates $(u:v:w)\leftrightarrow(x:y:z)$).

\emph{Pure inseparability.} Under this identification $\gamma$ is $(a:b:c)\mapsto(c^3:a^3:b^3)$, that is, the Frobenius morphism $(a:b:c)\mapsto(a^3:b^3:c^3)$ of $X$ composed with the cyclic permutation of coordinates, an automorphism of $X$ (as $F(z,x,y)=F$). Equivalently, on function fields $\gamma^*$ sends $u/w\mapsto(c/b)^3$, so $\gamma^*K(X^*)=K(X)^3$, and $K(X)/K(X)^3$ is purely inseparable of degree $3$ (II, Example~2.4.3 / Prop.~2.5). So $X\to X^*$ is purely inseparable (of degree $3$). \hfill$\square$

\exercise{2.5}{\emph{Automorphisms of a Curve of Genus $\ge2$.} Prove the theorem of Hurwitz [1] that a curve $X$ of genus $g\ge2$ over a field of characteristic $0$ has at most $84(g-1)$ automorphisms. We will see later (Ex.~5.2) or (V, Ex.~1.11) that the group $G=\operatorname{Aut}X$ is finite. So let $G$ have order $n$. Then $G$ acts on the function field $K(X)$. Let $L$ be the fixed field. Then the field extension $L\subseteq K(X)$ corresponds to a finite morphism of curves $f:X\to Y$ of degree $n$.
\begin{enumerate}[label=(\alph*),leftmargin=*]
\item If $P\in X$ is a ramification point, and $e_P=r$, show that $f^{-1}f(P)$ consists of exactly $n/r$ points, each having ramification index $r$. Let $P_1,\dots,P_s$ be a maximal set of ramification points of $X$ lying over distinct points of $Y$, and let $e_{P_i}=r_i$. Then show that Hurwitz's theorem implies that
\[
(2g-2)/n=2g(Y)-2+\sum_{i=1}^s(1-1/r_i).
\]
\item Since $g\ge2$, the left hand side of the equation is $>0$. Show that if $g(Y)\ge0$, $s\ge0$, $r_i\ge2$, $i=1,\dots,s$ are integers such that
\[
2g(Y)-2+\sum_{i=1}^s(1-1/r_i)>0,
\]
then the minimum value of this expression is $1/42$. Conclude that $n\le84(g-1)$. See (Ex.~5.7) for an example where this maximum is achieved.
\end{enumerate}
\emph{Note.} It is known that this maximum is achieved for infinitely many values of $g$ (Macbeath [1]). Over a field of characteristic $p>0$, the same bound holds, provided $p>g+1$, with one exception, namely the hyperelliptic curve $y^2=x^p-x$, which has $p=2g+1$ and $2p(p^2-1)$ automorphisms (Roquette [1]). For other bounds on the order of the group of automorphisms in characteristic $p$, see Singh [1] and Stichtenoth [1].}
\pf
We work over $k$ algebraically closed of characteristic $0$ and take as given (as the statement does) that $G=\operatorname{Aut}X$ is finite, of order $n$. By Artin's theorem $K(X)/L$ is Galois with group $G$ and $[K(X):L]=n$; $f:X\to Y$ is the corresponding finite morphism of degree $n$ ($Y$ the nonsingular projective curve with function field $L$), separable since char $0$. Moreover $f\circ\sigma=f$ for every $\sigma\in G$.

\textbf{(a)} Let $Q\in Y$. The points of $f^{-1}(Q)$ correspond to the extensions to $K(X)$ of the valuation ring $\mathcal O_Q\subseteq L$. In a Galois extension, the Galois group acts transitively on the set of such valuation rings (Atiyah--Macdonald, Ex.~5.13), i.e.\ $G$ acts transitively on $f^{-1}(Q)$. For $\sigma\in G$ and $P\in X$, $\sigma$ induces an isomorphism $\mathcal O_{X,\sigma P}\to\mathcal O_{X,P}$ compatible with $f^\#$ (as $f\sigma=f$), so $e_{\sigma P}=e_P$. Hence all points of a fibre have the same ramification index $r$, and $n=\sum_{P\mapsto Q}e_P=r\cdot\#f^{-1}(Q)$, so $\#f^{-1}(Q)=n/r$.

Let $Q_i=f(P_i)$, $i=1,\dots,s$ be the branch points (a maximal set of ramification points over distinct points of $Y$ has one point $P_i$ over each branch point $Q_i$). By Hurwitz (2.4; char $0$ so the ramification is tame), $2g-2=n(2g(Y)-2)+\sum_P(e_P-1)$. The ramification points lie in the fibres over $Q_1,\dots,Q_s$, and the fibre over $Q_i$ has $n/r_i$ points, each with $e-1=r_i-1$, contributing $(n/r_i)(r_i-1)=n(1-1/r_i)$. Dividing by $n$:
\[
\frac{2g-2}n=2g(Y)-2+\sum_{i=1}^s\Bigl(1-\frac1{r_i}\Bigr).
\]

\textbf{(b)} Let $\gamma=g(Y)$ and $E=2\gamma-2+\sum_{i=1}^s(1-1/r_i)$, assumed $>0$, with each $r_i\ge2$ so each term $1-1/r_i\ge\frac12$.
\begin{itemize}
\item $\gamma\ge2$: $E\ge2$.
\item $\gamma=1$: then $E=\sum(1-1/r_i)>0$ forces $s\ge1$, and $E\ge\frac12$.
\item $\gamma=0$, $E=-2+\sum(1-1/r_i)$. If $s\le2$ then $E<0$. If $s\ge5$ then $E\ge-2+\frac52=\frac12$. If $s=4$: $E=2-\sum1/r_i>0$ and $\sum1/r_i\le2$ with equality only for $(2,2,2,2)$, so some $r_i\ge3$ and $\sum1/r_i\le\frac32+\frac13=\frac{11}6$, whence $E\ge\frac16$.
\item $\gamma=0$, $s=3$: $E=1-(\frac1{r_1}+\frac1{r_2}+\frac1{r_3})>0$ with $r_1\le r_2\le r_3$. We maximize $S=\sum1/r_i<1$. If $r_1\ge3$: $S\le1$ with equality only for $(3,3,3)$, so otherwise $S\le\frac13+\frac13+\frac14=\frac{11}{12}$. If $r_1=2$: $\frac1{r_2}+\frac1{r_3}<\frac12$; $r_2=2$ is impossible; $r_2=3$ forces $r_3\ge7$, so $S\le\frac12+\frac13+\frac17=\frac{41}{42}$ with equality for $(2,3,7)$; $r_2=4$ forces $r_3\ge5$, $S\le\frac12+\frac14+\frac15=\frac{19}{20}<\frac{41}{42}$; $r_2\ge5$ gives $S\le\frac12+\frac25=\frac9{10}<\frac{41}{42}$.
\end{itemize}
So in every case $E\ge\frac1{42}$, and the minimum value $\frac1{42}$ is attained for $\gamma=0$, $s=3$, $(r_1,r_2,r_3)=(2,3,7)$ (and only there). By (a), $\frac{2g-2}n=E\ge\frac1{42}$, i.e.\ $n\le42(2g-2)=84(g-1)$. (For $g=3$ the Klein quartic of Ex.~5.7 has $168=84\cdot2$ automorphisms, so the bound is sharp.) \hfill$\square$

\exercise{2.6}{\emph{$f_*$ for Divisors.} Let $f:X\to Y$ be a finite morphism of curves of degree $n$. We define a homomorphism $f_*:\operatorname{Div}X\to\operatorname{Div}Y$ by $f_*(\sum n_iP_i)=\sum n_if(P_i)$ for any divisor $D=\sum n_iP_i$ on $X$.
\begin{enumerate}[label=(\alph*),leftmargin=*]
\item For any locally free sheaf $\mathcal E$ on $Y$, of rank $r$, we define $\det\mathcal E=\bigwedge^r\mathcal E\in\operatorname{Pic}Y$ (II, Ex.~6.11). In particular, for any invertible sheaf $\mathcal M$ on $X$, $f_*\mathcal M$ is locally free of rank $n$ on $Y$, so we can consider $\det f_*\mathcal M\in\operatorname{Pic}Y$. Show that for any divisor $D$ on $X$,
\[
\det(f_*\mathcal L(D))\cong(\det f_*\mathcal O_X)\otimes\mathcal L(f_*D).
\]
Note in particular that $\det(f_*\mathcal L(D))\ne\mathcal L(f_*D)$ in general! [\emph{Hint}: First consider an effective divisor $D$, apply $f_*$ to the exact sequence $0\to\mathcal L(-D)\to\mathcal O_X\to\mathcal O_D\to0$, and use (II, Ex.~6.11).]
\item Conclude that $f_*D$ depends only on the linear equivalence class of $D$, so there is an induced homomorphism $f_*:\operatorname{Pic}X\to\operatorname{Pic}Y$. Show that $f_*\cdot f^*:\operatorname{Pic}Y\to\operatorname{Pic}Y$ is just multiplication by $n$.
\item Use duality for a finite flat morphism (III, Ex.~6.10) and (III, Ex.~7.2) to show that
\[
\det f_*\Omega_X\cong(\det f_*\mathcal O_X)^{-1}\otimes\Omega_Y^{\otimes n}.
\]
\item Now assume that $f$ is separable, so we have the ramification divisor $R$. We define the \emph{branch divisor} $B$ to be the divisor $f_*R$ on $Y$. Show that
\[
(\det f_*\mathcal O_X)^2\cong\mathcal L(-B).
\]
\end{enumerate}}
\pf
All sheaves $f_*\mathcal M$ below are locally free of rank $n$ ($f$ is finite and flat).

\emph{A determinant lemma.} Let $0\to\mathcal E'\to\mathcal E\to\mathcal T\to0$ be exact on the curve $Y$, with $\mathcal E',\mathcal E$ locally free of the same rank $n$ and $\mathcal T$ of finite support. Then $\det\mathcal E'\cong\det\mathcal E\otimes\mathcal L(-B_{\mathcal T})$, where $B_{\mathcal T}=\sum_Q\operatorname{length}(\mathcal T_Q)\,Q$. (This is a special case of the determinant calculus of II, Ex.~6.11(b),(d): $\det\mathcal T=\mathcal L(B_{\mathcal T})$ and $\det$ is multiplicative.) Direct proof: $\wedge^n\mathcal E'\to\wedge^n\mathcal E$ is an injection of invertible sheaves, so $\wedge^n\mathcal E'=\wedge^n\mathcal E\otimes\mathcal L(-B)$ for an effective divisor $B$. At a point $Q$ with uniformizer $t$ the elementary-divisor theorem over the DVR $\mathcal O_Q$ gives bases $e_i$ of $\mathcal E_Q$ and $t^{a_i}e_i$ of $\mathcal E'_Q$; then $\mathcal T_Q\cong\bigoplus\mathcal O_Q/t^{a_i}$ has length $\sum a_i$, which is the order of vanishing of $\wedge^n\mathcal E'_Q\to\wedge^n\mathcal E_Q$. So the coefficient of $Q$ in $B$ is $\operatorname{length}\mathcal T_Q$.

\textbf{(a)} Let $\mathcal A$ be an invertible sheaf on $X$ and $P\in X$. Apply $f_*$ (exact, $f$ finite) to $0\to\mathcal A(-P)\to\mathcal A\to\mathcal A\otimes k(P)\to0$: we get $0\to f_*\mathcal A(-P)\to f_*\mathcal A\to k(f(P))\to0$, where the cokernel is a skyscraper sheaf at $f(P)$ of length $1$. The lemma gives
\[
\det f_*(\mathcal A(-P))\cong\det f_*\mathcal A\otimes\mathcal L(-f(P)).\tag{1}
\]
(For an effective divisor $E$ one gets the same way $\det f_*\mathcal A(-E)\cong\det f_*\mathcal A\otimes\mathcal L(-f_*E)$: the cokernel $f_*(\mathcal A\otimes\mathcal O_E)$ has length at $Q$ equal to $\sum_{P\mapsto Q}\operatorname{mult}_P(E)$, as residue fields are $k$; this is the case suggested by the hint.)

Claim: $\det f_*\mathcal L(D)\cong\det f_*\mathcal O_X\otimes\mathcal L(f_*D)$ for every $D$. For $D=0$ it is trivial. If the claim holds for $D+P$, apply (1) with $\mathcal A=\mathcal L(D+P)$: $\det f_*\mathcal L(D)\cong\det f_*\mathcal L(D+P)\otimes\mathcal L(-f(P))\cong\det f_*\mathcal O_X\otimes\mathcal L(f_*(D+P)-f(P))=\det f_*\mathcal O_X\otimes\mathcal L(f_*D)$; conversely if it holds for $D$, (1) shows it for $D+P$. Every divisor is reached from $0$ by adding and subtracting points, which proves the claim.

In general $\det f_*\mathcal O_X\not\cong\mathcal O_Y$ (see (d): e.g.\ $f:\mathbf P^1\to\mathbf P^1$ of degree $2$ has $\deg B=2$, so $\deg\det f_*\mathcal O_X=-1$), so $\det f_*\mathcal L(D)\ne\mathcal L(f_*D)$ already for $D=0$.

\textbf{(b)} If $D\sim D'$ then $\mathcal L(D)\cong\mathcal L(D')$, so $\det f_*\mathcal L(D)\cong\det f_*\mathcal L(D')$, and by (a) $\mathcal L(f_*D)\cong\det f_*\mathcal L(D)\otimes(\det f_*\mathcal O_X)^{-1}\cong\mathcal L(f_*D')$; since $\operatorname{Pic}Y\cong\operatorname{Cl}Y$ (II, 6.16), $f_*D\sim f_*D'$. So $f_*$ descends to $\operatorname{Pic}X=\operatorname{Cl}X\to\operatorname{Cl}Y=\operatorname{Pic}Y$. For $Q\in Y$, $f_*f^*Q=f_*\sum_{P\mapsto Q}e_PP=(\sum_{P\mapsto Q}e_P)Q=nQ$ (II, 6.9). So $f_*f^*=n$ on $\operatorname{Div}Y$, hence on $\operatorname{Pic}Y$.

\textbf{(c)} For the finite flat morphism $f$, by (III, Ex.~7.2(a)) the sheaf $f^!\omega_Y$ is a dualizing sheaf for $X$, hence $\cong\omega_X$ ($\omega_X=\Omega_X$ is the dualizing sheaf, III, 7.12), where $f^!\mathcal G$ is the $\mathcal O_X$-module corresponding to $\mathcal Hom_Y(f_*\mathcal O_X,\mathcal G)$ (III, Ex.~6.10(a)). Taking $\mathcal F=\mathcal O_X$ in (III, Ex.~6.10(b)): $f_*\Omega_X\cong f_*f^!\Omega_Y\cong\mathcal Hom_Y(f_*\mathcal O_X,\Omega_Y)=(f_*\mathcal O_X)^\vee\otimes\Omega_Y$, as $f_*\mathcal O_X$ is locally free. Since $\det(\mathcal E\otimes\mathcal N)=\det\mathcal E\otimes\mathcal N^{\otimes\operatorname{rk}\mathcal E}$ for $\mathcal N$ invertible, and $\det(\mathcal E^\vee)=(\det\mathcal E)^{-1}$,
\[
\det f_*\Omega_X\cong\det\bigl((f_*\mathcal O_X)^\vee\bigr)\otimes\Omega_Y^{\otimes n}\cong(\det f_*\mathcal O_X)^{-1}\otimes\Omega_Y^{\otimes n}.
\]

\textbf{(d)} By Prop.~2.3, $\Omega_X\cong f^*\Omega_Y\otimes\mathcal L(R)$. By the projection formula $f_*\Omega_X\cong\Omega_Y\otimes f_*\mathcal L(R)$, so using (a),
\[
\det f_*\Omega_X\cong\Omega_Y^{\otimes n}\otimes\det f_*\mathcal L(R)\cong\Omega_Y^{\otimes n}\otimes\det f_*\mathcal O_X\otimes\mathcal L(B).
\]
Comparing with (c): $(\det f_*\mathcal O_X)^{-1}\otimes\Omega_Y^{\otimes n}\cong\Omega_Y^{\otimes n}\otimes\det f_*\mathcal O_X\otimes\mathcal L(B)$. Cancelling $\Omega_Y^{\otimes n}$: $(\det f_*\mathcal O_X)^{-2}\cong\mathcal L(B)$, i.e.\ $(\det f_*\mathcal O_X)^2\cong\mathcal L(-B)$. \hfill$\square$

\exercise{2.7}{\emph{Etale Covers of Degree $2$.} Let $Y$ be a curve over a field $k$ of characteristic $\ne2$. We show there is a one-to-one correspondence between finite etale morphisms $f:X\to Y$ of degree $2$, and $2$-torsion elements of $\operatorname{Pic}Y$, i.e., invertible sheaves $\mathcal L$ on $Y$ with $\mathcal L^2\cong\mathcal O_Y$.
\begin{enumerate}[label=(\alph*),leftmargin=*]
\item Given an etale morphism $f:X\to Y$ of degree $2$, there is a natural map $\mathcal O_Y\to f_*\mathcal O_X$. Let $\mathcal L$ be the cokernel. Then $\mathcal L$ is an invertible sheaf on $Y$, $\mathcal L\cong\det f_*\mathcal O_X$, and so $\mathcal L^2\cong\mathcal O_Y$ by (Ex.~2.6). Thus an etale cover of degree $2$ determines a $2$-torsion element in $\operatorname{Pic}Y$.
\item Conversely, given a $2$-torsion element $\mathcal L$ in $\operatorname{Pic}Y$, define an $\mathcal O_Y$-algebra structure on $\mathcal O_Y\oplus\mathcal L$ by $\langle a,b\rangle\cdot\langle a',b'\rangle=\langle aa'+\varphi(b\otimes b'),ab'+a'b\rangle$, where $\varphi$ is an isomorphism of $\mathcal L\otimes\mathcal L\to\mathcal O_Y$. Then take $X=\operatorname{Spec}(\mathcal O_Y\oplus\mathcal L)$ (II, Ex.~5.17). Show that $X$ is an etale cover of $Y$.
\item Show that these two processes are inverse to each other. [\emph{Hint}: Let $\tau:X\to X$ be the involution which interchanges the points of each fibre of $f$. Use the trace map $a\mapsto a+\tau(a)$ from $f_*\mathcal O_X\to\mathcal O_Y$ to show that the sequence of $\mathcal O_Y$-modules in (a) is split exact.]
\end{enumerate}
\emph{Note.} This is a special case of the more general fact that for $(n,\operatorname{char}k)=1$, the etale Galois covers of $Y$ with group $\mathbf Z/n\mathbf Z$ are classified by the etale cohomology group $H^1_{\mathrm{et}}(Y,\mathbf Z/n\mathbf Z)$, which is equal to the group of $n$-torsion points of $\operatorname{Pic}Y$. See Serre [6].}
\pf
Throughout, ``$2$-torsion'' correspondence means: isomorphism classes of finite etale $f:X\to Y$ of degree $2$ (isomorphism over $Y$) $\leftrightarrow$ isomorphism classes of invertible $\mathcal L$ with $\mathcal L^2\cong\mathcal O_Y$. (The trivial cover $Y\sqcup Y$ corresponds to $\mathcal L=\mathcal O_Y$.) We use residue fields $k(Q)=k$ at closed points, $k$ algebraically closed (the argument works for any $k$ with $\operatorname{char}\ne2$ using geometric fibres).

\textbf{(a)} Let $\mathcal A=f_*\mathcal O_X$, locally free of rank $2$ with the structure map $\mathcal O_Y\to\mathcal A$, $1\mapsto1$. For a closed point $Q$ the fibre $\mathcal A\otimes k(Q)$ is the ring of the scheme-theoretic fibre $f^{-1}(Q)$, which is nonzero, so $1\ne0$ in it. Therefore $1$ is part of a local basis of $\mathcal A$, the map $\mathcal O_Y\to\mathcal A$ is a locally split injection, and the cokernel $\mathcal L=\mathcal A/\mathcal O_Y$ is invertible. By the multiplicativity of $\det$ in short exact sequences of locally free sheaves (II, Ex.~5.16(d) / Ex.~6.11), $\det\mathcal A\cong\det\mathcal O_Y\otimes\det\mathcal L=\mathcal L$. Since $f$ is etale, it is unramified and separable, $R=0$ and $B=0$, and Ex.~2.6(d) gives $(\det\mathcal A)^2\cong\mathcal L(-B)=\mathcal O_Y$. So $\mathcal L^2\cong\mathcal O_Y$. (If $X$ is disconnected, then each component maps onto $Y$ because $f$ is flat and finite, hence open and closed, so each has degree $1$ and $X=Y\sqcup Y$, giving $\mathcal A=\mathcal O_Y^2$ and $\mathcal L\cong\mathcal O_Y$ directly. If $X$ is connected it is regular, hence integral, i.e.\ a curve, and Ex.~2.6 applies as above.) Isomorphic covers give isomorphic $\mathcal L$, so we get a well-defined map.

\textbf{(b)} Let $\mathcal L$ be $2$-torsion and $\varphi:\mathcal L\otimes\mathcal L\to\mathcal O_Y$ an isomorphism. Note $\varphi$ is symmetric, since $\mathcal L$ is locally free of rank $1$ and the swap on $\mathcal L\otimes\mathcal L$ is the identity (locally $b=\beta e$, $b'=\beta'e$). Locally $\mathcal L=\mathcal O_Ye$ and $u:=\varphi(e\otimes e)$ is a unit, and the algebra $\mathcal A=\mathcal O_Y\oplus\mathcal L$ is locally $\mathcal O_Y[e]/(e^2-u)$, which is commutative and associative; so $\mathcal A$ is a commutative associative $\mathcal O_Y$-algebra (these are local identities), finite locally free of rank $2$. Hence $f:X=\operatorname{Spec}\mathcal A\to Y$ is finite, flat, of degree $2$. At a closed point $Q$ the fibre is $k[e]/(e^2-u(Q))$ with $u(Q)\ne0$; as $\operatorname{char}k\ne2$, $e^2-u(Q)$ has two distinct roots, so the fibre is $k\times k$, reduced, i.e.\ $f$ is unramified. A flat unramified morphism is etale. So $X\to Y$ is an etale cover of degree $2$.

\textbf{(c)} \emph{(b) then (a).} Starting from $(\mathcal L,\varphi)$ and $\mathcal A=\mathcal O_Y\oplus\mathcal L$, the cokernel of $\mathcal O_Y\to\mathcal A$ is $\mathcal A/\mathcal O_Y=\mathcal L$. Also the isomorphism class of $X$ over $Y$ does not depend on $\varphi$: another choice is $\varphi'=c\varphi$ with $c\in H^0(\mathcal O_Y^*)=k^*$ (isomorphisms $\mathcal L\otimes\mathcal L\to\mathcal O_Y$ differ by units of $\mathcal O_Y$, which are constants since $Y$ is a complete variety); take $\lambda\in k$ with $\lambda^2c=1$; then $\langle a,b\rangle\mapsto\langle a,\lambda b\rangle$ is an isomorphism of algebras $(\mathcal O\oplus\mathcal L,\varphi)\to(\mathcal O\oplus\mathcal L,\varphi')$. So the composite $\mathcal L\mapsto X\mapsto\mathcal L$ is the identity.

\emph{(a) then (b).} Let $f:X\to Y$ be etale of degree $2$ and $\mathcal A=f_*\mathcal O_X$. Since $\operatorname{char}\ne2$, the trace $\operatorname{tr}:\mathcal A\to\mathcal O_Y$ (trace of multiplication, an $\mathcal O_Y$-linear map) satisfies $\operatorname{tr}(1)=2$, a unit; so $\frac12\operatorname{tr}$ splits $\mathcal O_Y\to\mathcal A$ (this is the trace $a\mapsto a+\tau(a)$ of the hint, divided by $2$), and $\mathcal A=\mathcal O_Y\oplus\mathcal A_0$ with $\mathcal A_0=\ker\operatorname{tr}\cong\mathcal A/\mathcal O_Y=\mathcal L$ (split exactness of the sequence of (a)). Locally $\mathcal A_0=\mathcal O_Ye$ with $\operatorname{tr}e=0$; by Cayley--Hamilton, $e^2-\operatorname{tr}(e)e+N(e)=0$, so $e^2=-N(e)=:u\in\mathcal O_Y$. Hence $\mathcal A_0\cdot\mathcal A_0\subseteq\mathcal O_Y$, and the multiplication of $\mathcal A=\mathcal O_Y\oplus\mathcal L$ is $\langle a,b\rangle\langle a',b'\rangle=\langle aa'+\varphi(b\otimes b'),ab'+a'b\rangle$ with $\varphi:\mathcal L\otimes\mathcal L\to\mathcal O_Y$ the multiplication. The map $\varphi$ is an isomorphism: locally it is $e\otimes e\mapsto u$, and $u$ is a unit, because otherwise at a point $Q$ with $u(Q)=0$ the fibre $k[e]/(e^2)$ would be nonreduced, contradicting that $f$ is etale. Thus $\mathcal A\cong\mathcal O_Y\oplus\mathcal L$ as in (b), and $X\cong\operatorname{Spec}(\mathcal O_Y\oplus\mathcal L)$ over $Y$. So the composite $X\mapsto\mathcal L\mapsto X'$ gives $X'\cong X$.

Hence the two processes are mutually inverse bijections. \hfill$\square$

\section*{3. Embeddings in projective space}

Conventions. $k$ is algebraically closed, ``curve'' means complete nonsingular curve over $k$ (as throughout Chapter~IV), $\ell(D)=\dim_k H^0(X,\mathcal L(D))$ and $\dim|D|=\ell(D)-1$. The degree of a curve $X\subseteq\mathbf P^n$ is the degree of the divisor $X\cdot H$ cut out by a hyperplane $H\not\supseteq X$ (Example~3.3.2). Put $\delta(d)=\frac12(d-1)(d-2)$; a nonsingular plane curve of degree $d$ has genus $\delta(d)$ (I, Ex.~7.2). Besides the results of Section~3 I use the following standard facts.

\textbf{(F1)} (Ex.~1.5) If $E$ is an effective divisor then $\dim|E|\le\deg E$, with equality only for $E=0$ or $g=0$. In particular, \emph{if $E$ is a divisor of degree $1$ with $\ell(E)\ge2$, then $g=0$} (then $E\sim E'$ with $E'$ effective, and $\dim|E'|\ge1=\deg E'$).

\textbf{(F2)} (Ex.~1.8) If $\nu:X\to\Gamma$ is the normalization of an integral projective curve $\Gamma$, then $p_a(\Gamma)=g(X)+\sum_{P}\delta_P$ with $\delta_P=\operatorname{length}(\tilde{\mathcal O}_P/\mathcal O_P)$, and $\delta_P>0$ exactly at the singular points; a node has $\delta_P=1$.

\textbf{(F3)} Let $\nu:X\to\Gamma\subseteq\mathbf P^n$ be a finite birational morphism from a curve $X$ onto an integral curve $\Gamma$ (so $\nu$ is the normalization). Then $\deg\Gamma=\deg\nu^*\mathcal O_\Gamma(1)$. Indeed, tensoring $0\to\mathcal O_\Gamma\to\nu_*\mathcal O_X\to\mathcal Q\to0$ ($\mathcal Q$ of finite length $\sum\delta_P$) with $\mathcal O(m)$ and using the projection formula, $\chi(\mathcal O_\Gamma(m))=m\deg\nu^*\mathcal O_\Gamma(1)+1-g(X)-\sum\delta_P$, and the left side is $m\deg\Gamma+1-p_a(\Gamma)$.

\textbf{Lemma A.} Let $X\subseteq\mathbf P^n$ be an integral curve of degree $d$ not contained in any hyperplane. Then $d\ge n$.

\emph{Proof.} Choose $n$ distinct points $P_1,\dots,P_n\in X$. There is a hyperplane $H$ through them (they impose at most $n$ linear conditions on the $n+1$ coefficients), and $H\not\supseteq X$. The effective divisor $X\cdot H$ has degree $d$ and multiplicity $\ge1$ at each $P_i\in H$, so $d\ge n$. \hfill$\square$

\exercise{3.1}{If $X$ is a curve of genus $2$, show that a divisor $D$ is very ample $\Leftrightarrow\deg D\ge5$. This strengthens (3.3.4).}
\pf
($\Leftarrow$) $2g+1=5$, so this is (3.2)(b).

($\Rightarrow$) Let $D$ be very ample, $d=\deg D$. Then $d>0$ (3.3) and $D\sim X\cdot H$ for a hyperplane section of an embedding, so $D$ may be taken effective. By the proof of (3.1) the complete system $|D|$ is base point free and defines a closed immersion $\varphi:X\to\mathbf P^n$, $n=\dim|D|$; its image is not contained in a hyperplane (the coordinates form a basis of $H^0(X,\mathcal L(D))$, so are linearly independent), and by Example~3.3.2 it has degree $d$.

If $n=0$, $\varphi$ is not an immersion of a curve. If $n=1$, $\varphi(X)$ is a closed irreducible $1$-dimensional subset of $\mathbf P^1$, so $X\cong\mathbf P^1$, contradicting $g=2$. Hence $n\ge2$, and Lemma~A gives $n\le d$.

\emph{$d\le2$ is impossible:} then $n=d=2$, so $\ell(D)=3$. For a point $P$, $\ell(D-P)\ge2$ while $\deg(D-P)=1$, so $g=0$ by (F1), a contradiction.

\emph{$d=3,4$ are impossible:} since $d>2g-2=2$, $D$ is nonspecial and Riemann--Roch gives $\ell(D)=d+1-g=d-1$, so $n=d-2$. For $d=3$ this is $n=1$, excluded above. For $d=4$ we get $n=2$: $X$ is isomorphic to a plane curve of degree $4$, whose genus is $\delta(4)=3\ne2$.

Hence $d\ge5$. \hfill$\square$

\exercise{3.2}{Let $X$ be a plane curve of degree $4$.
\begin{enumerate}[label=(\alph*),leftmargin=*]
\item Show that the effective canonical divisors on $X$ are exactly the divisors $X.L$, where $L$ is a line in $\mathbf P^2$.
\item If $D$ is any effective divisor of degree $2$ on $X$, show that $\dim|D|=0$.
\item Conclude that $X$ is not hyperelliptic (Ex.~1.7).
\end{enumerate}}
\pf
$X$ is a hypersurface, so a complete intersection, and by (II, Ex.~8.4(e)) $\omega_X\cong\mathcal O_X(4-2-1)=\mathcal O_X(1)$; also $g=\delta(4)=3$. Restriction of linear forms gives a map $H^0(\mathbf P^2,\mathcal O(1))\to H^0(X,\mathcal O_X(1))$, injective because $X$ lies on no line. Both sides have dimension $3$ ($\dim H^0(X,\omega_X)=g=3$), so it is an isomorphism: every section of $\omega_X\cong\mathcal O_X(1)$ is the restriction of a linear form.

(a) An effective canonical divisor is the divisor of zeros of a nonzero section of $\omega_X\cong\mathcal O_X(1)$, hence of the restriction of a linear form, i.e.\ it is $X.L$ for a line $L$ (and $L\not\supseteq X$). Conversely every $X.L$ is the divisor of a section of $\mathcal O_X(1)\cong\omega_X$, hence canonical, effective.

(b) Let $D\ge0$ have degree $2$. By Riemann--Roch, $\ell(D)=2+1-3+\ell(K-D)=\ell(K-D)$. Fix $K=X.L_0$ ($L_0=\{\lambda_0=0\}$); then $\mathcal L(K-D)\cong\{$sections $s$ of $\mathcal O_X(1)$ with $\operatorname{div}(s)\ge D\}$ via $f\mapsto f\lambda_0$, and by the first paragraph such sections are restrictions of linear forms $\lambda$ with $\operatorname{div}(\lambda|_X)\ge D$. If $D=P+Q$ with $P\ne Q$ this means $\lambda(P)=\lambda(Q)=0$, two independent conditions on the $3$-dimensional space of linear forms, so the space is $1$-dimensional. If $D=2P$ it means $\lambda(P)=0$ and $\operatorname{ord}_P(\lambda|_X)\ge2$, i.e.\ the line $\{\lambda=0\}$ is the tangent line of $X$ at $P$ ($X$ is nonsingular at $P$); again a $1$-dimensional space. So $\ell(K-D)=1=\ell(D)$ and $\dim|D|=0$.

(c) A hyperelliptic curve has a finite morphism $f:X\to\mathbf P^1$ of degree $2$. For distinct $a,b\in\mathbf P^1$ the effective divisors $f^*(a)\ne f^*(b)$ of degree $2$ are linearly equivalent (pull-back of $a\sim b$, II, Ex.~6.8), so $\dim|f^*(a)|\ge1$, contradicting (b). \hfill$\square$

\exercise{3.3}{If $X$ is a curve of genus $\ge2$ which is a complete intersection (II, Ex.~8.4) in some $\mathbf P^n$, show that the canonical divisor $K$ is very ample. Conclude that a curve of genus $2$ can never be a complete intersection in any $\mathbf P^n$. Cf.\ (Ex.~5.1).}
\pf
Let $X$ be the complete intersection of hypersurfaces of degrees $d_1,\dots,d_{n-1}$. By (II, Ex.~8.4(e)), $\omega_X\cong\mathcal O_X(m)$ with $m=\sum d_i-n-1$. Since $g\ge2$, $\deg K=2g-2>0$, while $\deg\mathcal O_X(m)=m\deg X$ with $\deg X>0$; so $m\ge1$. Compose the inclusion $X\hookrightarrow\mathbf P^n$ with the $m$-uple embedding $\mathbf P^n\hookrightarrow\mathbf P^N$ (II, Ex.~5.13): this is a closed immersion of $X$ whose pull-back of $\mathcal O(1)$ is $\mathcal O_X(m)\cong\omega_X$. Hence $\omega_X=\mathcal L(K)$ is very ample.

If $g=2$, this would make $K$, a divisor of degree $2g-2=2$, very ample, contradicting (3.1) (very ample needs degree $\ge5$). So a curve of genus $2$ is not a complete intersection in any $\mathbf P^n$. \hfill$\square$

\exercise{3.4}{Let $X$ be the $d$-uple embedding of $\mathbf P^1$ in $\mathbf P^d$, for any $d\ge1$ (I, Ex.~2.12). We call $X$ the \emph{rational normal curve of degree $d$} in $\mathbf P^d$.
\begin{enumerate}[label=(\alph*),leftmargin=*]
\item Show that $X$ is projectively normal, and that its homogeneous ideal can be generated by forms of degree $2$.
\item If $X$ is any curve of degree $d$ in $\mathbf P^n$, with $d\le n$, which is not contained in any $\mathbf P^{n-1}$, show that in fact $d=n$, $g(X)=0$, and $X$ differs from the rational normal curve of degree $d$ only by an automorphism of $\mathbf P^d$. Cf.\ (II, 7.8.5).
\item In particular, any curve of degree $2$ in any $\mathbf P^n$ is a conic in some $\mathbf P^2$.
\item A curve of degree $3$ in any $\mathbf P^n$ must be either a plane cubic curve, or the twisted cubic curve in $\mathbf P^3$.
\end{enumerate}}
\pf
Write $X=\{(s^d:s^{d-1}t:\dots:t^d)\}\subseteq\mathbf P^d$, $x_i\mapsto s^{d-i}t^i$.

(a) \emph{Projective normality.} $X\cong\mathbf P^1$ is normal, so by (II, Ex.~5.14(d)) it suffices to show that $H^0(\mathbf P^d,\mathcal O(l))\to H^0(X,\mathcal O_X(l))$ is surjective for all $l\ge0$. Now $\mathcal O_X(l)\cong\mathcal O_{\mathbf P^1}(dl)$, with $H^0$ spanned by the monomials $s^{dl-j}t^j$, $0\le j\le dl$. Given $j$, write $j=j_1+\dots+j_l$ with $0\le j_i\le d$ (possible as $j\le dl$); then $s^{dl-j}t^j=\prod_i s^{d-j_i}t^{j_i}$ is the image of the monomial $x_{j_1}\cdots x_{j_l}$. So the map is onto.

\emph{Ideal.} The ideal $I$ of $X$ is the kernel of $\theta:k[x_0,\dots,x_d]\to k[s,t]$, $x_i\mapsto s^{d-i}t^i$ (a form vanishes on $X$ iff its image vanishes, $k$ being infinite). Every monomial $m=x_{i_1}\cdots x_{i_l}$ goes to a monomial $s^{dl-w}t^w$, $w=\sum i_j$ (the \emph{weight}). Hence $I$ is spanned by the binomials $m-m'$ with $\theta(m)=\theta(m')$, i.e.\ with the same degree and weight (group the monomials of an element of $\ker\theta$ by image; the coefficients in each group sum to $0$). Let $J\subseteq I$ be the ideal generated by the quadrics $x_ix_j-x_{i+1}x_{j-1}$ for $0\le i$, $i+2\le j\le d$. Take a monomial $x_{i_1}\cdots x_{i_l}$, and suppose two of its indices satisfy $i\le j-2$. Replacing $x_ix_j$ by $x_{i+1}x_{j-1}$ changes it modulo $J$ and decreases $\sum i_r^2$ by $2(j-i)-2\ge2$; so after finitely many steps we reach a monomial congruent to it mod $J$ whose indices all lie in $\{a,a+1\}$ for some $a$. Given $(l,w)$, the multiset of indices in $\{a,a+1\}$ is unique: if $q\in\{0,\dots,l-1\}$ of the indices equal $a+1$ (and the others $a$), then $w=la+q$, which determines $a$ and $q$. Hence any two monomials of the same degree and weight are congruent mod $J$, so $I=J$ is generated by quadrics. (For $d=1$, $I=0$.)

(b) By Lemma~A, $d\ge n$, so $d=n$. Let $D=X\cdot H$ (degree $d$). The restriction $H^0(\mathbf P^n,\mathcal O(1))\to H^0(X,\mathcal L(D))$ is injective, as $X$ lies in no hyperplane, so $\ell(D)\ge n+1=d+1$. Choose distinct points $P_1,\dots,P_{d-1}\in X$; since $\ell(E-P)\ge\ell(E)-1$ always, $\ell(D-P_1-\dots-P_{d-1})\ge2$, while this divisor has degree $1$. By (F1), $g(X)=0$, so $X\cong\mathbf P^1$.

Then $\mathcal L(D)\cong\mathcal O_{\mathbf P^1}(d)$ and the embedding is given by linearly independent sections $\sigma_0,\dots,\sigma_d\in H^0(\mathbf P^1,\mathcal O(d))$. This space has dimension $d+1$, so the $\sigma_i$ form a basis, and $\sigma_i=\sum_ja_{ij}s^{d-j}t^j$ with $A=(a_{ij})$ invertible. Thus the embedding is the $d$-uple embedding followed by the automorphism $A$ of $\mathbf P^d$.

(c) Let $X$ have degree $2$ and let $\mathbf P^m$ be the linear span of $X$ (smallest linear subspace containing it). Then $X\subseteq\mathbf P^m$ is not in a hyperplane, so $m\le2$ by Lemma~A. $m=0$ is absurd and $m=1$ would make $X=\mathbf P^1$ of degree $1$. So $m=2$: $X$ is a plane curve of degree $2$, a conic.

(d) Similarly for degree $3$ with span $\mathbf P^m$: $m\le3$, $m=1$ is impossible, $m=2$ gives a plane cubic, and $m=3=d$ gives, by (b), the rational normal curve of degree $3$ up to an automorphism of $\mathbf P^3$, i.e.\ the twisted cubic. \hfill$\square$

\exercise{3.5}{Let $X$ be a curve in $\mathbf P^3$, which is not contained in any plane.
\begin{enumerate}[label=(\alph*),leftmargin=*]
\item If $O\notin X$ is a point, such that the projection from $O$ induces a birational morphism $\varphi$ from $X$ to its image in $\mathbf P^2$, show that $\varphi(X)$ must be singular. [\emph{Hint:} Calculate $\dim H^0(X,\mathcal O_X(1))$ two ways.]
\item If $X$ has degree $d$ and genus $g$, conclude that $g<\frac12(d-1)(d-2)$. (Use (Ex.~1.8).)
\item Now let $\{X_t\}$ be the flat family of curves induced by the projection (III, 9.8.3) whose fibre over $t=1$ is $X$, and whose fibre $X_0$ over $t=0$ is a scheme with support $\varphi(X)$. Show that $X_0$ always has nilpotent elements. Thus the example (III, 9.8.4) is typical.
\end{enumerate}}
\pf
Throughout, $\varphi$ is given by the three linear forms vanishing at $O$; these are sections of $\mathcal O_{\mathbf P^3}(1)$ with no common zero on $X$, so $\varphi^*\mathcal O_{\mathbf P^2}(1)\cong\mathcal O_X(1)$. Put $\Gamma=\varphi(X)$, an integral plane curve of degree $e$ (it is not a point, since $X$ lies on no plane through $O$).

(a) Suppose $\Gamma$ is nonsingular. Then $\varphi:X\to\Gamma$ is a finite birational morphism onto a normal curve, hence an isomorphism ($\mathcal O_\Gamma\to\varphi_*\mathcal O_X$ is a finite birational extension of integrally closed rings). Computation 1: since $X$ lies in no plane, restriction $H^0(\mathbf P^3,\mathcal O(1))\to H^0(X,\mathcal O_X(1))$ is injective, so $\dim H^0(X,\mathcal O_X(1))\ge4$. Computation 2: $\mathcal O_X(1)\cong\varphi^*\mathcal O_{\mathbf P^2}(1)$ and $\varphi$ is an isomorphism, so $H^0(X,\mathcal O_X(1))\cong H^0(\Gamma,\mathcal O_\Gamma(1))$. From $0\to\mathcal O_{\mathbf P^2}(1-e)\to\mathcal O_{\mathbf P^2}(1)\to\mathcal O_\Gamma(1)\to0$ and $H^1(\mathbf P^2,\mathcal O(m))=0$ (III, 5.1), $H^0(\mathbf P^2,\mathcal O(1))\to H^0(\Gamma,\mathcal O_\Gamma(1))$ is surjective, so $\dim H^0(\Gamma,\mathcal O_\Gamma(1))\le3$. Contradiction. So $\Gamma$ is singular.

(b) Here $X$ is in no plane (for a plane curve $g=\delta(d)$, so the inequality needs this hypothesis). By (3.10) there is $O\notin X$ such that $\varphi$ is birational onto its image $\Gamma$. Then $X$ is the normalization of $\Gamma$, so by (F3) $e=\deg\varphi^*\mathcal O(1)=\deg\mathcal O_X(1)=d$. By (F2), $\delta(d)=p_a(\Gamma)=g+\sum_P\delta_P$, and $\sum\delta_P\ge1$ because $\Gamma$ is singular by (a). Hence $g<\delta(d)=\frac12(d-1)(d-2)$.

(c) In coordinates with $O=(0:0:0:1)$ the family is $X_a=\sigma_a(X)$, $\sigma_a(x_0:x_1:x_2:x_3)=(x_0:x_1:x_2:ax_3)$ for $a\ne0$, extended to a flat family over $\mathbf A^1$ (III, 9.8.3). All fibres have the same Hilbert polynomial (III, 9.9; $\mathbf A^1$ is connected). For $a=1$ this is $P_X(m)=\chi(\mathcal O_X(m))=dm+1-g$ (Riemann--Roch). (As in (a), (b), $\varphi$ is taken birational onto $\Gamma$, which is possible by (3.10).) Suppose $X_0$ had no nilpotents. Then $X_0$ would be the reduced scheme on $\Gamma$ (a plane curve in the plane $x_3=0$), whose Hilbert polynomial is that of a plane curve of degree $e$: from $0\to\mathcal O_{\mathbf P^2}(m-e)\to\mathcal O_{\mathbf P^2}(m)\to\mathcal O_\Gamma(m)\to0$ one gets $P_\Gamma(m)=em+1-\delta(e)$. Comparing with $dm+1-g$: $e=d$ and $g=\delta(d)$, contradicting (b). So $X_0$ has nilpotent elements. (For the twisted cubic: $d=3$, $g=0$, and the nodal cubic has $p_a=1$, so the missing $p_a-g=1$ is carried by the embedded point, as computed in III, 9.8.4.) \hfill$\square$

\exercise{3.6}{\emph{Curves of Degree 4.}
\begin{enumerate}[label=(\alph*),leftmargin=*]
\item If $X$ is a curve of degree $4$ in some $\mathbf P^n$, show that either (1) $g=0$, in which case $X$ is either the rational normal quartic in $\mathbf P^4$ (Ex.~3.4) or the rational quartic curve in $\mathbf P^3$ (II, 7.8.6), or (2) $X\subseteq\mathbf P^2$, in which case $g=3$, or (3) $X\subseteq\mathbf P^3$ and $g=1$.
\item In the case $g=1$, show that $X$ is a complete intersection of two irreducible quadric surfaces in $\mathbf P^3$ (I, Ex.~5.11). [\emph{Hint:} Use the exact sequence $0\to\mathcal I_X\to\mathcal O_{\mathbf P^3}\to\mathcal O_X\to0$ to compute $\dim H^0(\mathbf P^3,\mathcal I_X(2))$, and thus conclude that $X$ is contained in at least two irreducible quadric surfaces.]
\end{enumerate}}
\pf
(a) Let $\mathbf P^m$ be the linear span of $X$; $X$ lies in no hyperplane of $\mathbf P^m$, so $m\le4$ by Lemma~A. $m=1$ is impossible ($X$ would be $\mathbf P^1$, of degree $1$).

$m=4$: by Ex.~3.4(b), $g=0$ and $X$ is the rational normal quartic.

$m=2$: $X$ is a nonsingular plane quartic, so $g=\delta(4)=3$.

$m=3$: $X\subseteq\mathbf P^3$ is not in a plane. By (3.10) there is a projection $O\notin X$ birational onto its image, so by Ex.~3.5(b), $g<\delta(4)=3$. If $g=2$: $D=X\cdot H$ has degree $4>2g-2$, so $\ell(D)=4+1-2=3$, but $\ell(D)\ge4$ as the four coordinates are independent sections of $\mathcal L(D)$ ($X$ in no plane). So $g\in\{0,1\}$. If $g=0$, $X$ is a nonsingular curve $\cong\mathbf P^1$ of degree $4$ in $\mathbf P^3$ not in a plane, so $\mathcal L(D)\cong\mathcal O_{\mathbf P^1}(4)$ and the embedding is given by a base-point-free $4$-dimensional subspace of $H^0(\mathbf P^1,\mathcal O(4))$ (of dimension $5$), i.e.\ $X$ is the projection of the rational normal quartic from a point of $\mathbf P^4$: a rational quartic curve in $\mathbf P^3$ of the type of (II, 7.8.6) (the one there is the projection of the normal quartic from $(0:0:1:0:0)$, the omitted monomial $s^2t^2$). This is case (1). If $g=1$, case (3).

(b) Let $X\subseteq\mathbf P^3$, $g=1$, $d=4$; $X$ is in no plane (a plane quartic has $g=3$). From $0\to\mathcal I_X(2)\to\mathcal O_{\mathbf P^3}(2)\to\mathcal O_X(2)\to0$ we get the exact sequence $0\to H^0(\mathcal I_X(2))\to H^0(\mathcal O_{\mathbf P^3}(2))\to H^0(\mathcal O_X(2))$. Here $\dim H^0(\mathcal O_{\mathbf P^3}(2))=10$ and $\deg\mathcal O_X(2)=8>2g-2=0$, so $\ell=8$ by Riemann--Roch. Hence $\dim H^0(\mathcal I_X(2))\ge2$: there are two linearly independent quadratic forms $Q_1,Q_2$ vanishing on $X$.

Every nonzero quadratic form $Q$ vanishing on $X$ is irreducible: otherwise $Q=\lambda\lambda'$ and $X\subseteq\{\lambda=0\}\cup\{\lambda'=0\}$ lies in a plane ($X$ irreducible). Let $Y=V(Q_1)$, an irreducible quadric surface of degree $2$; $Q_2\notin(Q_1)$ (non-proportional of equal degree), so $Y\not\subseteq V(Q_2)$. All components $Z_j$ of $Y\cap V(Q_2)$ have dimension $1$ (I, 7.2), and Bezout (I, 7.7) gives $\sum_j i(Y,V(Q_2);Z_j)\deg Z_j=2\cdot2=4$. The irreducible curve $X$ lies in $Y\cap V(Q_2)$, hence is one of the $Z_j$, say $Z_1$, with $\deg Z_1=4$. Thus $i_1=1$ and there is no other component: $Z:=\operatorname{Proj}k[x_0,\dots,x_3]/(Q_1,Q_2)$ has $X$ as its only irreducible component and is reduced at the generic point of $X$ ($i_1=1$ is the length of the local ring there).

Now $Q_1,Q_2$ is a regular sequence in the UFD $S=k[x_0,\dots,x_3]$, so $S/(Q_1,Q_2)$ is Cohen--Macaulay of dimension $2$; hence $Z$ has no embedded points, and being generically reduced with the single generic point of $X$, $Z$ is reduced, so $Z=X$ as schemes. Moreover depth $2\ge1$ means $(Q_1,Q_2)$ is saturated, so $I(X)=(Q_1,Q_2)$. Thus $X$ is the complete intersection of the two irreducible quadrics $V(Q_1),V(Q_2)$. \hfill$\square$

\exercise{3.7}{In view of (3.10), one might ask conversely, is every plane curve with nodes a projection of a nonsingular curve in $\mathbf P^3$? Show that the curve $xy+x^4+y^4=0$ (assume $\operatorname{char}k\ne2$) gives a counterexample.}
\pf
Reading: ``projection'' means a projection $\varphi:X\to\mathbf P^2$ from a point $O\notin X$ of a nonsingular curve $X\subseteq\mathbf P^3$ which is birational onto the plane curve in question (as in (3.10)). The affine curve $xy+x^4+y^4=0$ has projective closure $\Gamma=V(F)$, $F=xyz^2+x^4+y^4$, in $\mathbf P^2$ with coordinates $(x:y:z)$.

\emph{$\Gamma$ has exactly one singular point, a node.} In the chart $z=1$: $f=xy+x^4+y^4$, $f_x=y+4x^3$, $f_y=x+4y^3$. At a singular point, $y=-4x^3$ and $x=-4y^3=256x^9$. If $x\ne0$ then $256x^8=1$, and then $f=-4x^4+x^4+256x^{12}=-4x^4+x^4+x^4=-2x^4\ne0$ (char $\ne2$), a contradiction. So $x=0$, hence $y=0$: the only singular point in the chart is the origin. There $f=xy+(\text{terms of degree 4})$, so the multiplicity is $2$ with tangent cone $xy=0$, two distinct lines: a node. At infinity, $z=0$, the points satisfy $x^4+y^4=0$, so $x,y\ne0$, and $F_x=yz^2+4x^3=4x^3\ne0$: nonsingular.

\emph{$\Gamma$ is irreducible.} If $F=GH$ with $\deg G,\deg H\ge1$, then $G,H$ meet in $\deg G\deg H\ge3$ points (counted with multiplicity), all singular points of $\Gamma$ (and no component is repeated, since the singular set is finite). So $G,H$ meet only at the origin, which has multiplicity $2$ on $\Gamma$; so both pass through it with multiplicity $1$ and the tangent cones multiply to $xy$, so their tangents are the two different lines $x=0,y=0$, and the local intersection multiplicity at the origin is $1$. Then $\deg G\deg H=1$, a contradiction.

Suppose now $\Gamma=\varphi(X)$ with $X\subseteq\mathbf P^3$ a nonsingular curve and $\varphi$ birational. Then $X$ is the normalization of $\Gamma$ (finite birational from a nonsingular curve), so by (F2) $g(X)=p_a(\Gamma)-\delta_O=3-1=2$, and by (F3) $\deg X=\deg\varphi^*\mathcal O(1)=\deg\Gamma=4$ (the divisor $\varphi^*\mathcal O(1)=\mathcal O_X(1)$ has degree $\deg X$). But a curve of degree $4$ in $\mathbf P^3$ has $g\in\{0,1,3\}$ by Ex.~3.6(a) (if $X$ lies in a plane, $g=3$). So no such $X$ exists. \hfill$\square$

\exercise{3.8}{We say a (singular) integral curve in $\mathbf P^n$ is \emph{strange} if there is a point which lies on all the tangent lines at nonsingular points of the curve.
\begin{enumerate}[label=(\alph*),leftmargin=*]
\item There are many singular strange curves, e.g., the curve given parametrically by $x=t$, $y=t^p$, $z=t^{2p}$ over a field of characteristic $p>0$.
\item Show, however, that if $\operatorname{char}k=0$, there aren't even any singular strange curves besides $\mathbf P^1$.
\end{enumerate}}
\pf
(a) (The statement asks us to check that this curve is a singular strange curve.) With homogeneous coordinates $(w:x:y:z)$ and affine coordinates $x/w,y/w,z/w$, let $C$ be the image of $\mathbf P^1\to\mathbf P^3$, $(s:t)\mapsto(s^{2p}:s^{2p-1}t:s^pt^p:t^{2p})$ (the entries $s^{2p},t^{2p}$ have no common zero, so this is a morphism). $C$ is an integral closed curve, and $x/w=t/s$ is a coordinate on $\mathbf P^1$, so $k(C)=k(t)$, $\mathbf P^1\to C$ is birational. The only point of $C$ with $w=0$ is $Q=(0:0:0:1)$ ($s=0$), and the affine part is $\{(t,t^p,t^{2p})\}=\{y=x^p,z=x^{2p}\}\cong\mathbf A^1$, nonsingular.

\emph{$C$ is singular at $Q$.} In the chart $z\ne0$ with parameter $s=1/t$, the coordinates are $(w,x,y)/z=(s^{2p},s^{2p-1},s^p)$, so $\mathcal O_{C,Q}$ is the localization of $k[s^{2p},s^{2p-1},s^p]$ at the maximal ideal $(s^{2p},s^{2p-1},s^p)$. Every element of $\mathcal O_{C,Q}$ has a power series expansion in $s$ with no term $s^j$ for $0<j<p$ (true for the polynomials, and for inverses of units $c+h$, $h\in s^pk[[s]]$). As $p\ge2$, $s\notin\mathcal O_{C,Q}$. But the normalization of $C$ is $\mathbf P^1$, whose local ring at $s=0$ is $k[s]_{(s)}\ni s$. So $\mathcal O_{C,Q}$ is not integrally closed, hence not regular: $Q$ is a singular point.

\emph{Strangeness.} At a nonsingular point $P=(a,a^p,a^{2p})$ the curve is cut out in $\mathbf A^3$ by $y-x^p$, $z-x^{2p}$, whose differentials are $dy-px^{p-1}dx=dy$ and $dz$ (char $p$). So the tangent space is the $x$-direction, and the tangent line is $\{y=a^p,z=a^{2p}\}$, which contains the point at infinity $A=(0:1:0:0)$ of the $x$-axis. Thus $A$ lies on all tangent lines at nonsingular points: $C$ is a singular strange curve.

(b) Let $\operatorname{char}k=0$. I show more generally:

\textbf{Lemma D.} Let $\nu:Y\to\mathbf P^n$ be a morphism from a curve, birational onto its image $C$. Assume that there is a point $A$ and a cofinite set of $P\in Y$ such that the tangent line $L_P$ to $\nu$ at $P$ (the line through $\nu(P)$ in the direction $d\nu_P\ne0$) passes through $A$. Then $C$ is a line.

Applying it to the normalization of a strange integral curve $X$ (at a nonsingular point of $X$, $\nu$ is a local isomorphism, so $L_P$ is the tangent line of $X$) shows that a strange curve in characteristic $0$ is a line; lines are nonsingular, so no singular curve is strange, and $\mathbf P^1$ itself is strange (every point is on its only tangent line).

\emph{Proof of Lemma D.} Choose coordinates with $A=(0:\dots:0:1)$ and write $\nu=(f_0:\dots:f_n)$. If $f_0=\dots=f_{n-1}=0$, $C=A$ is a point; so choose $j<n$ with $f_j\ne0$ and put $u_i=f_i/f_j\in K(Y)$, $i<n$. For all but finitely many $P$ the $u_i$ are regular at $P$, $\nu(P)\ne A$ and $f_j(P)\ne0$, and the chart $x_j\ne0$ has affine coordinates $v_i=x_i/x_j$, with $v_i\circ\nu=u_i$ for $i\neq j,n$ and $v_n\circ\nu=f_n/f_j$. The tangent line at $\nu(P)$ is $\nu(P)+k\cdot(\frac{du_i}{dt}(P))_{i\ne j}$ ($t$ a uniformizer at $P$), and its point at infinity in the hyperplane $x_j=0$ is the direction vector $(\frac{du_i}{dt}(P))$. $A$ is the point at infinity in the direction $e_n$, so $A\in L_P$ means $\frac{du_i}{dt}(P)=0$ for all $i<n$, $i\ne j$, i.e.\ $du_i\in\mathfrak m_P\Omega_{Y,P}$. Suppose some $u_i$ is nonconstant. Since $\operatorname{char}k=0$, $K(Y)$ is separably generated over $k$ and $du_i\ne0$ in $\Omega_{K(Y)/k}$ (II, 8.6A); thus $du_i$ is a nonzero rational section of the invertible sheaf $\Omega_Y$, which vanishes at only finitely many points. This contradicts the vanishing at infinitely many $P$. Hence all $u_i$ ($i<n$) are constants $c_i$, and $\nu(P)=(c_0f_j:\dots:f_j:\dots:c_{n-1}f_j:f_n)$ lies on the line through $(c_0:\dots:c_{n-1}:0)$ and $A$. So $C$ is contained in a line, hence is a line. \hfill$\square$

\exercise{3.9}{Prove the following lemma of Bertini: if $X$ is a curve of degree $d$ in $\mathbf P^3$, not contained in any plane, then for almost all planes $H\subseteq\mathbf P^3$ (meaning a Zariski open subset of the dual projective space $(\mathbf P^3)^*$), the intersection $X\cap H$ consists of exactly $d$ distinct points, no three of which are collinear.}
\pf
Recall that $X$ is a nonsingular curve (Chapter IV convention). I first prove two lemmas, valid in all characteristics; they are also used in Ex.~3.10.

\textbf{Setting.} Let $n\ge3$, $Y$ a curve and $\nu:Y\to\mathbf P^n$ a morphism, birational onto its image $C$, with $C$ not contained in a hyperplane. There is a cofinite open set $Y_0\subseteq Y$ on which $\nu$ is an immersion ($d\nu_P\ne0$; this holds generically because $\nu$ is separable) and $\nu^{-1}(\nu(P))=\{P\}$. For $P\in Y_0$ let $L_P\subseteq\mathbf P^n$ be the tangent line. The map $P\mapsto L_P$ is a morphism $Y_0\to G(1,n)$ and extends to a morphism $\tau:Y\to G(1,n)$ ($Y$ nonsingular, $G$ projective; I, 6.8); write $L_P=\tau(P)$ for every $P$. Let $K=\{(P,Q)\in Y\times Y:L_P\cap L_Q\ne\emptyset\}$, a closed subset (incidence of lines is closed). Let $W\subseteq Y^3$ be the locally closed set of triples $(P,Q,R)$ of points with $\nu(P),\nu(Q),\nu(R)$ pairwise distinct and collinear. If $\nu(P)\ne\nu(Q)$, the line $\nu(P)\nu(Q)$ is not contained in $C$ ($C$ is not a line), so the fibres of $p_{12}:W\to Y\times Y$ are finite and $\dim W\le2$.

\textbf{Lemma B.} If $\dim W=2$, then $L_P\cap L_Q\ne\emptyset$ for all $P,Q\in Y$, i.e.\ $K=Y\times Y$.

\emph{Proof.} Let $W_1$ be an irreducible component of $W$ of dimension $2$. I claim $p_{12}(W_1)\subseteq K$. Granting this: $p_{12}(W_1)$ is constructible of dimension $2$ (finite fibres) in the irreducible surface $Y\times Y$, hence dense, and $K$ is closed, so $K=Y\times Y$.

For the claim, let $V=W_1\cap p_{12}^{-1}(Y\times Y\setminus K)$, open in $W_1$; suppose $V\neq\emptyset$, so $\dim V=2$. Look at $p_3:V\to Y$. I show that every fibre of $p_3$ is finite, giving $\dim V\le1$, a contradiction. Fix $R\in Y$ and let $\psi:Y\to\mathbf P^{n-1}$ be the projection from $\nu(R)$ (a morphism, extended at $R$); it is nonconstant because $C$ is not a line. The fibre $V_R$ consists of pairs $(P,Q)$ with $P\ne Q$, $\psi(P)=\psi(Q)$ and $(P,Q)\notin K$. I show that all but finitely many pairs $(P,Q)$ with $\psi(P)=\psi(Q)$ are in $K$. Write $\psi=\pi_{\nu(R)}\circ\nu$ and let $\psi$ factor as $Y\xrightarrow{f}\tilde Y\to C_R=\psi(Y)$, $\tilde Y$ the normalization of $C_R$.

\emph{Case 1: $K(Y)/K(\tilde Y)$ is inseparable.} If $u\in K(\tilde Y)$ had $du\ne0$ in $\Omega_{K(Y)/k}$, then $K(Y)$ would be separable over $k(u)$ and hence over $K(\tilde Y)$. So $du=0$ for all $u\in K(\tilde Y)$, in particular for the affine coordinates of $\psi$; so $d\psi_P=0$ for all $P$. As $d\psi_P=d\pi\circ d\nu_P$ and $\ker d\pi$ is the direction towards $\nu(R)$, each $L_P$ ($P\in Y_0$, $\nu(P)\ne\nu(R)$) passes through $\nu(R)$, and two such lines meet. The remaining pairs $(P,Q)$ have $P$ or $Q$ in the finite set $(Y\setminus Y_0)\cup\{R\}$, and each point has a finite fibre under $\psi$: finitely many pairs.

\emph{Case 2: $f$ separable.} Then $f$ is ramified at only finitely many points (2.2), and $C_R$ has finitely many singular points. Let $E\subseteq C_R$ be the finite set consisting of the singular points of $C_R$, the images of the ramification points of $f$, of $Y\setminus Y_0$ and of $R$. For $(P,Q)$ with $\psi(P)=\psi(Q)=T\notin E$: $T$ is a nonsingular point of $C_R$ (so $\tilde Y\to C_R$ is an isomorphism near $T$), $f$ is unramified at $P$ and $Q$, so $\psi^*$ of a local parameter at $T$ is a uniformizer at $P$ and $d\psi_P:T_PY\to T_TC_R$ is an isomorphism. Since $d\psi_P=d\pi\circ d\nu_P$, $d\nu_P\ne0$, $L_P\not\ni\nu(R)$, and $\pi(L_P)$ is the tangent line $\ell_T$ of $C_R$ at $T$. Hence $L_P$, and likewise $L_Q$, lies in the plane $\Pi_T=\pi^{-1}(\ell_T)\cup\{\nu(R)\}$. Two lines in a plane meet: $(P,Q)\in K$. The pairs with $\psi(P)=\psi(Q)\in E$ are finite in number (finite fibres). \hfill$\square$

\textbf{Lemma C.} A set $\mathcal L$ of lines in $\mathbf P^n$, any two of which meet, is either concurrent (all through one point) or contained in one plane.

\emph{Proof.} Suppose the lines are not concurrent. Then $\mathcal L$ contains two distinct lines $L_1,L_2$, meeting in a point $a$, and a line $L_3\not\ni a$. $L_3$ meets $L_1$ and $L_2$ in points $b\ne c$ (if $b=c$ then $b\in L_1\cap L_2=\{a\}$), both in the plane $\Pi=\langle L_1,L_2\rangle$, so $L_3\subseteq\Pi$, and $L_1\cap L_2\cap L_3=\emptyset$. If $M\in\mathcal L$ were not contained in $\Pi$, then $M\cap\Pi$ would be a single point, lying on each of $L_1,L_2,L_3$ (as $M$ meets each and each lies in $\Pi$), which is impossible. So $\mathcal L\subseteq\Pi$. \hfill$\square$

\emph{Step 1: $\dim W\le1$ for our nonplanar nonsingular $X\subseteq\mathbf P^3$.} Apply the setting with $Y=X$, $\nu$ the inclusion. If $\dim W=2$, then by Lemma B the tangent lines $T_PX$ pairwise meet, and by Lemma C they are either all in a plane $\Pi$ --- then $X\subseteq\Pi$ since $P\in T_PX$, contradiction --- or all pass through a point $A$, i.e.\ $X$ is strange. By Theorem~3.9 a strange curve is a line, or a conic in characteristic $2$, hence planar. Contradiction. So $\dim W\le1$.

\emph{Step 2: the dimension count.} Let $\mathcal B=\{(H,P,Q,R):H\in(\mathbf P^3)^*,(P,Q,R)\in W,P,Q,R\in H\}$. A plane contains $P,Q,R$ iff it contains the line $PQ$, so $\mathcal B\to W$ is a $\mathbf P^1$-bundle and $\dim\mathcal B\le\dim W+1\le2$. The image $\Sigma$ of $\mathcal B$ in $(\mathbf P^3)^*$ is constructible of dimension $\le2$, hence contained in a proper closed subset: for $H$ in a dense open set, $X\cap H$ contains no three distinct collinear points.

\emph{Step 3: distinctness.} By Bertini (II, 8.18) there is a dense open set of $H\in(\mathbf P^3)^*$ with $H\not\supseteq X$ and $H\cap X$ regular at every point, i.e.\ the zero-dimensional scheme $H\cap X$ is reduced. It has length $\deg(X\cdot H)=d$, so $X\cap H$ consists of exactly $d$ distinct points.

The intersection of the two dense open sets of Steps 2 and 3 is dense open, and for $H$ in it $X\cap H$ consists of $d$ distinct points, no three collinear. \hfill$\square$

\exercise{3.10}{Generalize the statement that ``not every secant is a multisecant'' as follows. If $X$ is a curve in $\mathbf P^n$, not contained in any $\mathbf P^{n-1}$, and if $\operatorname{char}k=0$, show that for almost all choices of $n-1$ points $P_1,\dots,P_{n-1}$ on $X$, the linear space $L^{n-2}$ spanned by the $P_i$ does not contain any further points of $X$.}
\pf
For $n=2$ there is one point and nothing to prove, so let $n\ge3$. ``Almost all'' means: for $(P_1,\dots,P_{n-1})$ in a dense open subset of $X^{n-1}$. I will prove this for the larger class of data $(Y,\nu)$: $Y$ a curve and $\nu:Y\to\mathbf P^n$ a morphism, birational onto its image and with image in no hyperplane (call this class $\mathcal C_n$); $X\subseteq\mathbf P^n$ is the case $\nu=$ inclusion. For $(Y,\nu)\in\mathcal C_n$ and $1\le k\le n-1$ let $B_k(Y,\nu)\subseteq Y^k$ be the set of $(P_1,\dots,P_k)$ such that either the points $\nu(P_1),\dots,\nu(P_k)$ do not span a linear space of dimension $k-1$, or there is $R\in Y\setminus\{P_1,\dots,P_k\}$ with $\nu(R)$ in their span. It is constructible.

\textbf{Claim $P(n,k)$:} $\dim B_k(Y,\nu)\le k-1$ for every $(Y,\nu)\in\mathcal C_n$ and $1\le k\le n-1$.

For $k=n-1$ and $\nu$ the inclusion this gives the exercise: the complement of $B_{n-1}$ is constructible of dimension $n-1$ in the irreducible variety $X^{n-1}$, hence contains a dense open set, and for $(P_i)$ in it, the $P_i$ span a $\mathbf P^{n-2}=L^{n-2}$ with $L\cap X=\{P_1,\dots,P_{n-1}\}$.

\emph{$k=1$.} $B_1=\{P:\nu^{-1}\nu(P)\ne\{P\}\}$ is finite since $\nu$ is birational.

\emph{$k=2$, $n\ge3$.} $B_2$ consists of pairs with $\nu(P_1)=\nu(P_2)$ (the diagonal plus finitely many pairs: dimension $\le1$), and pairs $(P_1,P_2)$ such that $\nu(R)\in$ line $\nu(P_1)\nu(P_2)$ for some $R\notin\{P_1,P_2\}$. If moreover $\nu(R)\in\{\nu(P_1),\nu(P_2)\}$, then $P_1$ or $P_2$ lies in the finite set of points over which $\nu$ is not injective, giving dimension $\le1$. Otherwise $(P_1,P_2,R)\in W$ (notation of Ex.~3.9), and by finiteness of the fibres of $p_{12}$ this part has dimension $\le\dim W$. I claim $\dim W\le1$: if $\dim W=2$, Lemmas B and C show that the tangent lines $L_P$ are all in a plane (then $C$ lies in a plane, contradicting non-degeneracy, $n\ge3$) or pass through a common point $A$; in the latter case Lemma D (char $0$, Ex.~3.8(b)) says $C$ is a line, again contradicting non-degeneracy. So $\dim B_2\le1$.

\emph{Induction step $P(n-1,k)\Rightarrow P(n,k+1)$, for $2\le k$, $k+1\le n-1$.} Let $(Y,\nu)\in\mathcal C_n$ and suppose $\dim B_{k+1}=k+1$. Then $B_{k+1}$ contains a nonempty open $V\subseteq Y^{k+1}$. For $O\in Y$ let $V_O=\{(P_1,\dots,P_k):(P_1,\dots,P_k,O)\in V\}$, open in $Y^k$; it is nonempty for $O$ in a dense open subset of $Y$. Since $\dim B_2\le1$, for all but finitely many $O$ the set $\{P:(P,O)\in B_2\}$ is finite; also $\nu^{-1}\nu(O)=\{O\}$ for all but finitely many $O$. Fix $O$ satisfying all three conditions. Let $\psi=\psi_O:Y\to\mathbf P^{n-1}$ be the projection from $\nu(O)$. For $P,R\ne O$ with $(P,O)\notin B_2$, $\psi(R)=\psi(P)$ forces $R=P$; so $\psi$ is injective on a cofinite set, and as $\operatorname{char}k=0$ it is birational onto its image. The image is in no hyperplane of $\mathbf P^{n-1}$ (else $C$ would lie in a hyperplane through $\nu(O)$). So $(Y,\psi)\in\mathcal C_{n-1}$ and, as $k\le(n-1)-1$, $P(n-1,k)$ gives $\dim B_k(Y,\psi)\le k-1$: its complement contains a nonempty open set $G\subseteq Y^k$. Choose $(P_1,\dots,P_k)\in V_O\cap G$ with all $P_i\ne O$. I show $(P_1,\dots,P_k,O)\notin B_{k+1}$, contradicting $V\subseteq B_{k+1}$.

\emph{Independence.} Note $\nu(P_i)\ne\nu(O)$. If $c_0\nu(O)+\sum c_i\nu(P_i)=0$ for vector representatives, applying the linear projection with kernel $\langle\nu(O)\rangle$ gives $\sum c_i\psi(P_i)=0$, so $c_i=0$ by goodness of $(P_i)$ for $\psi$, then $c_0=0$.

\emph{Span.} Let $\Lambda=\operatorname{span}(\nu(P_1),\dots,\nu(P_k),\nu(O))$ and let $R\in Y\setminus\{P_1,\dots,P_k,O\}$ with $\nu(R)\in\Lambda$. Then $\nu(R)\ne\nu(O)$ and $\psi(R)\in\operatorname{span}(\psi(P_1),\dots,\psi(P_k))$, with $R\notin\{P_i\}$, contradicting $(P_i)\in G$.

This proves $P(n,k+1)$. Since $P(m,2)$ holds for all $m\ge3$ and $P(m,1)$ for all $m$, induction on $k$ gives $P(n,k)$ for all $1\le k\le n-1$. \hfill$\square$

\exercise{3.11}{\begin{enumerate}[label=(\alph*),leftmargin=*]
\item If $X$ is a nonsingular variety of dimension $r$ in $\mathbf P^n$, and if $n>2r+1$, show that there is a point $O\notin X$, such that the projection from $O$ induces a closed immersion of $X$ into $\mathbf P^{n-1}$.
\item If $X$ is the Veronese surface in $\mathbf P^5$, which is the $2$-uple embedding of $\mathbf P^2$ (I, Ex.~2.13), show that each point of every secant line of $X$ lies on infinitely many secant lines. Therefore, the secant variety of $X$ has dimension $4$, and so in this case there is a projection which gives a closed immersion of $X$ into $\mathbf P^4$ (II, Ex.~7.7). (A theorem of Severi [1] states that the Veronese surface is the only surface in $\mathbf P^5$ for which there is a projection giving a closed immersion into $\mathbf P^4$. Usually one obtains a finite number of double points with transversal tangent planes.)
\end{enumerate}}
\pf
(a) \emph{Criterion.} For $O\notin X$ let $\pi:X\to\mathbf P^{n-1}$ be the projection; it is the morphism defined by the base-point-free linear system of hyperplanes through $O$. As in the proof of (3.4), by (II, 7.3) and (II, 7.8.2), $\pi$ is a closed immersion iff this system separates points and tangent vectors. It separates points $P\ne Q$ iff $\pi(P)\ne\pi(Q)$ iff $O$ is not on the line $PQ$. It separates tangent vectors at $P$ iff $d\pi_P$ is injective on $T_PX$; since $\ker d\pi_P$ is the direction of the line $OP$, this holds iff the line $OP$ is not contained in the embedded tangent space $\mathbf T_P(X)\subseteq\mathbf P^n$ ($r$-dimensional linear space), i.e.\ iff $O\notin\mathbf T_P(X)$. So $\pi$ is a closed immersion iff $O$ lies on no secant line $PQ$ ($P\ne Q$) and on no $\mathbf T_P(X)$.

\emph{Dimension count.} The set $I=\{(P,Q,x):P\ne Q,\ x\in PQ\}\subseteq(X\times X\setminus\Delta)\times\mathbf P^n$ is a $\mathbf P^1$-bundle over the irreducible variety $X\times X\setminus\Delta$ of dimension $2r$, so $\dim I=2r+1$, and its image $S_1\subseteq\mathbf P^n$ (the points on secant lines) is constructible of dimension $\le2r+1$. The set $T=\{(P,x):P\in X,x\in\mathbf T_P(X)\}\subseteq X\times\mathbf P^n$ is closed with all fibres over $X$ of dimension $r$ ($X$ nonsingular), so $\dim T\le2r$ and its image $S_2$ (the union of the embedded tangent spaces) has dimension $\le2r$. Since $n>2r+1$, $S_1\cup S_2\ne\mathbf P^n$. Take $O\notin S_1\cup S_2$; then $O\notin X$ (as $X\subseteq S_2$), and by the criterion the projection from $O$ is a closed immersion into $\mathbf P^{n-1}$.

(b) $X=v(\mathbf P^2)\subseteq\mathbf P^5$, with $v(a)=(a_ia_j)_{0\le i\le j\le2}$ for $a\in k^3\setminus0$ (coordinates $z_{ij}$, $z_{ji}:=z_{ij}$). Valid in all characteristics.

\emph{A secant meets $X$ only in its two defining points.} Let $p,q\in k^3$ be independent and suppose $\alpha v(p)+\beta v(q)=\lambda v(r)$ with $\alpha,\beta\ne0$, $\lambda\ne0$ (vectors in $k^6$). For $\xi,\eta\in(k^3)^*$ the form $z(\xi,\eta)=\sum\xi_i\eta_jz_{ij}$ is linear in $z$, and $z(\xi,\eta)(v(a))=\xi(a)\eta(a)$. Taking $\xi=\eta$ vanishing on $\langle p,q\rangle$ gives $\lambda\xi(r)^2=0$, so $r\in\langle p,q\rangle$, $r=xp+yq$. Choose $\xi,\eta$ with $(\xi(p),\xi(q))=(1,0)$ and $(\eta(p),\eta(q))=(0,1)$: $\lambda xy=0$. So $r\propto p$ or $q$, which forces $\beta=0$ or $\alpha=0$. Hence the line through $v(p),v(q)$ meets $X$ in exactly these two points, and different secant lines through $Z$ correspond to different pairs of points.

\emph{Points of $X$.} Each $v(p)\in X$ lies on infinitely many secants (join it to any other point of $X$).

\emph{Other points.} Let $Z=\alpha v(p)+\beta v(q)$ with $\alpha,\beta\ne0$ be a point of a secant, not in $X$. As $k$ is algebraically closed choose $a=\sqrt\alpha p$, $b=\sqrt\beta q$ (independent vectors), so $Z=v(a)+v(b)$ as vectors. For any $c,d\in k$ with $c^2+d^2=1$ put $a'=ca+db$, $b'=-da+cb$. A direct computation of the $(i,j)$ entry gives $a'_ia'_j+b'_ib'_j=(c^2+d^2)(a_ia_j+b_ib_j)=a_ia_j+b_ib_j$ (the cross terms $\pm cd(a_ib_j+b_ia_j)$ cancel), so $v(a')+v(b')=Z$. Also $\det\begin{pmatrix}c&d\\-d&c\end{pmatrix}=1$, so $a',b'$ are independent and $Z$ lies on the secant through $v(a'),v(b')$. The conic $c^2+d^2=1$ has infinitely many points ($k$ algebraically closed), and the map $(c,d)\mapsto(c:d)$ is at most $2$-to-$1$ on it, so the points $[a']=[ca+db]$ of the line $\mathbf P(\langle a,b\rangle)$ take infinitely many values. By the first part these give infinitely many distinct secant lines through $Z$.

\emph{Dimension of the secant variety.} Every point $Z=sv(a)+tv(b)$ of a secant line, written as the symmetric matrix $(sa_ia_j+tb_ib_j)$, has rank $\le2$, so $\det(z_{ij})=0$ on $S_1$ (the union of the secant lines): the secant variety $\overline{S_1}$ lies in the cubic hypersurface $\det(z_{ij})=0$, so $\dim\overline{S_1}\le4$. This is also what the first part gives: $I'=\{(P,Q,x)\in I:x\notin X\}$ is irreducible of dimension $5$, and every nonempty fibre of $I'\to\mathbf P^5$ is infinite, hence of dimension $\ge1$, so the image has dimension $\le4$. In fact the fibres are exactly $1$-dimensional, so $\dim\overline{S_1}=4$: if $Z=v(p)+v(q)=v(a)+v(b)$ with independent $p,q$ and $a,b$, then the bilinear forms $B(\xi,\eta)=\xi(p)\eta(p)+\xi(q)\eta(q)=\xi(a)\eta(a)+\xi(b)\eta(b)$ on $(k^3)^*$ have radical $\langle p,q\rangle^\perp=\langle a,b\rangle^\perp$, so $\langle p,q\rangle=\langle a,b\rangle$, and comparing coefficients, $(p,q)=(a,b)A$ with $A$ invertible and $AA^{t}=1$, a $1$-dimensional family (a conic) of possibilities.

\emph{Projection into $\mathbf P^4$.} By the criterion in (a), the projection from $O$ is a closed immersion into $\mathbf P^4$ whenever $O\notin S_1\cup S_2$. Here $\dim\overline{S_1}=4$ and $\dim S_2\le2r=4$ (from (a)), so $S_1\cup S_2$ lies in a proper closed subset of $\mathbf P^5$ and such $O$ exists. (An explicit projection is given in II, Ex.~7.7(b).) \hfill$\square$

\exercise{3.12}{For each value of $d=2,3,4,5$ and $r$ satisfying $0\le r\le\frac12(d-1)(d-2)$, show that there exists an irreducible plane curve of degree $d$ with $r$ nodes and no other singularities.}
\pf
A \emph{node} is a double point with two distinct tangent directions. The cases ($\delta(d)$ is the maximal $r$): $d=2$, $r=0$; $d=3$, $r=0,1$; $d=4$, $r=0,\dots,3$; $d=5$, $r=0,\dots,6$. Three constructions cover everything.

\textbf{$r=0$.} A nonsingular plane curve of degree $d$ exists for every $d$ (Bertini, II, 8.20.2); it is irreducible. For $d=2$ this finishes the proof.

\textbf{Lemma 1 (projection of curves of genus $g$).} Let $X$ be a curve of genus $g$, and $d\ge2g+1$ with $g<\delta(d)$. Then there is an irreducible plane curve of degree $d$ with exactly $\delta(d)-g$ nodes and no other singularities.

\emph{Proof.} By (3.2)(b) a divisor $D$ of degree $d$ is very ample, giving $X\hookrightarrow\mathbf P^N$; projecting from points off $X$ repeatedly (3.5, as in Cor.~3.6) we get a closed immersion $X\subseteq\mathbf P^3$ with $\mathcal O_X(1)$ of degree $d$ (each projection pulls $\mathcal O(1)$ back to $\mathcal O_X(1)$). $X$ is in no plane: a nonsingular plane curve of degree $d$ has genus $\delta(d)\ne g$. By (3.10) there is $O\notin X$ such that the projection $\varphi:X\to\Gamma\subseteq\mathbf P^2$ is birational onto $\Gamma$ and $\Gamma$ has at most nodes. $\Gamma$ is irreducible, $\deg\Gamma=d$ by (F3), so $p_a(\Gamma)=\delta(d)$, and by (F2), $\delta(d)=g+\#\{\text{nodes}\}$ since each node has $\delta_P=1$. \hfill$\square$

Thus: $g=0$ (take $X=\mathbf P^1$, any $d$) gives $r=\delta(d)$, i.e.\ $(d,r)=(3,1),(4,3),(5,6)$; $g=1$ (a smooth plane cubic) gives $(4,2)$ and $(5,5)$; $g=2$ ($d=5$) gives $(5,4)$, once a curve of genus $2$ is available, see below.

\textbf{Lemma 2 ($r=1,2$ and $d\ge2r+1$).} For $r\in\{1,2\}$ and $d\ge2r+1$ there is an irreducible plane curve of degree $d$ with exactly $r$ nodes and no other singularities. Moreover its normalization has genus $\delta(d)-r$.

\emph{Proof.} Let $p_1=(0:0:1)$, $p_2=(1:0:0)$ and use $(x_0:x_1:x_2)$. Let $S\subseteq\mathbf P^2\times(\mathbf P^1)^r$ be the closed subscheme defined by $x_0t_1=x_1s_1$ (and for $r=2$ also $x_1t_2=x_2s_2$), $(s_i:t_i)$ coordinates on the $i$-th $\mathbf P^1$ (the pencil of lines through $p_i$). Let $\pi:S\to\mathbf P^2$ and $f_i:S\to\mathbf P^1$ be the projections.

\emph{Structure of $S$.} Over $U=\mathbf P^2\setminus\{p_1,p_2\}$, $(s_1:t_1)=(x_0:x_1)$ and $(s_2:t_2)=(x_1:x_2)$ are determined, so $\pi^{-1}(U)\cong U$. Over $p_1$ the equations leave $(s_1:t_1)$ free and force $(s_2:t_2)=(0:1)$, so $E_1=\pi^{-1}(p_1)\cong\mathbf P^1$, and $f_1|_{E_1}$ is an isomorphism; similarly $E_2=\pi^{-1}(p_2)\cong\mathbf P^1$ with $f_2|_{E_2}$ an isomorphism. Near $E_1$ ($x_2=1$, $u=x_0$, $v=x_1$) the second equation reads $s_2=vt_2$ and just expresses $s_2$ as a function of $v$ (graph), while the first is $ut_1=vs_1$: the standard blow-up of the origin. So $S$ is a nonsingular surface (similarly near $E_2$), irreducible because every component has dimension $\ge2$ and so meets $\pi^{-1}(U)$. $S$ is projective: $S=\operatorname{Bl}_{\{p_1,\dots,p_r\}}\mathbf P^2$. (For $r=1$ only $p_1$, $E_1$ are used; $\pi$ is an isomorphism near $p_2$, and everything below about $p_2,E_2$ is void.)

\emph{Embedding.} Put $a_0=d-2r\ge1$. The composite of the $a_0$-uple embedding of $\mathbf P^2$, the $2$-uple embeddings of the $\mathbf P^1$'s and the Segre embedding is a closed immersion $\mathbf P^2\times(\mathbf P^1)^r\hookrightarrow\mathbf P^N$ pulling $\mathcal O(1)$ back to $\mathcal O(a_0,2,\dots,2)$ (II, Ex.~5.13, 5.11). So $S\subseteq\mathbf P^N$ is a nonsingular projective surface, and each $E_i$ is embedded by $\mathcal O_{\mathbf P^1}(2)$: a conic.

\emph{Bertini.} By (II, 8.18) applied to $S$ and to each conic $E_i$, there is a hyperplane $H'\subseteq\mathbf P^N$ such that $D=H'\cap S$ is regular at every point and $H'\cap E_i$ is regular, $i=1,\dots,r$ (finite intersection of dense open sets of hyperplanes). $D$ is a nonsingular, connected (indeed irreducible) curve (III, 7.9.1). $H'\cap E_i$ is a reduced scheme of length $\deg E_i=2$: two distinct points $a_i\ne b_i$, and $D\cap E_i=H'\cap E_i$, so $D$ meets $E_i$ transversally in $a_i,b_i$.

\emph{The plane curve.} $D\ne E_i$, so $D\not\subseteq\pi^{-1}\{p_1,p_2\}$; $C=\pi(D)$ is an integral plane curve, $\pi|_D:D\to C$ is finite and birational ($\pi$ is an isomorphism over $U$), so $D$ is the normalization of $C$. \emph{Degree:} take a line $\ell\subseteq\mathbf P^2$ missing $p_1,p_2$ and $\tilde\ell=\pi^{-1}(\ell)\cong\ell$. The section $\pi^*(\text{equation of }\ell)$ of $\pi^*\mathcal O(1)$ restricted to $D$ vanishes on $D\cap\tilde\ell$ (and $D\not\supseteq\tilde\ell$ as $D$ meets $E_1$ and $\tilde\ell$ does not), so $\deg(\pi^*\mathcal O(1)|_D)=\operatorname{length}(D\cap\tilde\ell)=\operatorname{length}(H'\cap\tilde\ell)=\deg(\mathcal O(1)|_{\tilde\ell})$ for the embedding. On $\tilde\ell\cong\ell$ the maps $\pi$ and $f_i$ are isomorphisms onto $\ell$, $\mathbf P^1$ (projection of $\ell$ from $p_i\notin\ell$), so $\mathcal O(a_0,2,\dots,2)|_{\tilde\ell}=\mathcal O_{\mathbf P^1}(a_0+2r)=\mathcal O(d)$. By (F3), $\deg C=d$.

\emph{Singularities.} $C\setminus\{p_1,p_2\}\cong D\setminus(E_1\cup E_2)$ is nonsingular. Over $p_i$ the normalization $D\to C$ has the two preimages $a_i\ne b_i$; at these points $\pi|_D$ is an immersion: locally $\pi$ is the blow-up map $(u,t)\mapsto(u,ut)$, whose differential on $\{u=0\}$ has kernel exactly the tangent line to $E_i$, and $D$ is transversal to $E_i$. So $C$ has at $p_i$ two smooth branches. The tangent line of the branch through $a_i$ is the line through $p_i$ of slope given by the limit of the pencil parameter: $f_i(a_i)$; similarly $f_i(b_i)$; these are distinct because $f_i|_{E_i}$ is injective. So $p_i$ is a node, and $C$ has exactly $r$ nodes. The genus of $D$ is $p_a(C)-r=\delta(d)-r$ by (F2). \hfill$\square$

In particular $(3,1),(4,1),(5,1),(5,2)$ are done, and for $(d,r)=(4,1)$ the curve $D$ is a curve of genus $2$.

\textbf{Genus $2$.} Taking $X=D$ in Lemma 1 with $d=5=2g+1$: $(5,4)$ is done.

\textbf{Lemma 3 ($(d,r)=(5,3)$).} Let $\Gamma\subseteq\mathbf P^2$ be a nonsingular quartic (genus $3$). Not being a line, or a conic in characteristic $2$, $\Gamma$ is not strange (Theorem~3.9); so for every point $P$ of $\mathbf P^2$ the set $\{A\in\Gamma:P\in T_A\Gamma\}$ is finite. Choose successively $p_1,p_2,p_3\in\Gamma$ such that (i) they are not collinear; (ii) $T_{p_i}\Gamma\not\ni p_j$ for $j\ne i$; (iii) for each pair $i\ne j$, the line $L_{ij}=p_ip_j$ is not tangent to $\Gamma$ at any point other than $p_i,p_j$. Each condition excludes only finitely many points at the stage where a point is chosen: for (ii), the points of $\Gamma\cap T_{p_i}\Gamma$ and the finite set $\{A:p_i\in T_A\Gamma\}$; for (iii), if $L_{ij}=T_A\Gamma$ for some $A\ne p_i,p_j$ then $A\in\{A:p_i\in T_A\Gamma\}$ (finite) and $p_j\in T_A\Gamma\cap\Gamma$ (finite). By (ii), $L_{ij}$ is transversal to $\Gamma$ at $p_i$ and at $p_j$, so $L_{ij}\cdot\Gamma=p_i+p_j+a_{ij}+b_{ij}$ with $a_{ij}\ne b_{ij}$, both different from $p_i,p_j$ (and from $p_k$, by (i)).

Choose coordinates $p_1=(1:0:0),p_2=(0:1:0),p_3=(0:0:1)$ and let $T(x)=(x_2x_3:x_1x_3:x_1x_2)$. On $U=\{x_1x_2x_3\ne0\}$, $T$ is invertible with $T\circ T=\operatorname{id}$ ($T(T(x))=x_1x_2x_3\cdot x$). $\Gamma\not\subseteq$ the coordinate triangle, so $T|_{\Gamma\cap U}$ extends to a morphism $\varphi:\Gamma\to\mathbf P^2$, birational onto the integral curve $\Gamma'=\overline{T(\Gamma\cap U)}$ (the rational map $T$ restricted to $\Gamma'$ is the inverse).

\emph{Degree.} $\varphi$ is given by the sections $s_1=x_2x_3,s_2=x_1x_3,s_3=x_1x_2$ of $\mathcal O_\Gamma(2)$, whose only common zeros are the points $p_1,p_2,p_3$. At $p_1$, with a local parameter, $\operatorname{ord}(x_2/x_1)=\operatorname{ord}(x_3/x_1)=1$ because $T_{p_1}\Gamma$ is neither $\{x_2=0\}=L_{13}$ nor $\{x_3=0\}=L_{12}$; so $s_1,s_2,s_3$ have orders $2,1,1$ at $p_1$, and after removing the common factor the divisor is lowered by exactly $p_1$. The same holds at $p_2,p_3$. Hence $\varphi^*\mathcal O(1)\cong\mathcal O_\Gamma(2)(-p_1-p_2-p_3)$ has degree $8-3=5$ and $\deg\Gamma'=5$ by (F3). Since $\Gamma$ is the normalization of $\Gamma'$, (F2) gives $\sum\delta_P=p_a(\Gamma')-3=6-3=3$.

\emph{Three nodes.} The points of $\Gamma$ with $x_3=0$ are $p_1,p_2,a=a_{12},b=b_{12}$. At $a$, $s_3=x_1x_2\ne0$ and $s_1=s_2=0$, so $\varphi(a)=\varphi(b)=q_3:=(0:0:1)$; and $\varphi(p_1),\varphi(p_2)\ne q_3$ (in the chart $x_1=1$, with $x_2=\alpha t+\cdots$, $x_3=\beta t+\cdots$, $\alpha\beta\ne0$, the sections are $(\alpha\beta t^2:\beta t:\alpha t)$, so $\varphi(p_1)=(0:\beta:\alpha)$; likewise $\varphi(p_2)$ has first and third coordinates nonzero and second coordinate $0$, and $\varphi(p_3)$ has third coordinate $0$, by the same transversality, so neither equals $q_3$). Since $s_1=s_2=0\ne s_3$ forces $x_3=0$, $x_1x_2\ne0$, no other point maps to $q_3$. So $q_3$ has exactly the two preimages $a,b$. Near $q_3$ the affine coordinates are $X=s_1/s_3=x_3/x_1$, $Y=s_2/s_3=x_3/x_2$. As $L_{12}=\{x_3=0\}$ meets $\Gamma$ transversally at $a$ and $b$, $x_3$ has a simple zero there, so $\varphi$ is an immersion at $a$ and $b$ and the two branches at $q_3$ are smooth, with tangent slopes $Y/X=x_1/x_2$ evaluated at $a$ and at $b$. These differ, because $a\ne b$ are points of the line $x_3=0$ and are determined by $(x_1:x_2)$. So $q_3$ is a node. By symmetry $q_1,q_2$ are nodes as well. Each node contributes $\delta=1$, so these three nodes account for $\sum\delta_P=3$: there are no other singular points. Thus $\Gamma'$ is an irreducible plane quintic with exactly $3$ nodes. \hfill$\square$

\textbf{Summary.}
\begin{itemize}
\item $d=2$: $r=0$ (smooth).
\item $d=3$: $r=0$ (smooth), $r=1$ (Lemma~2, or Lemma~1 with $g=0$).
\item $d=4$: $r=0$ (smooth), $r=1$ (Lemma~2), $r=2$ (Lemma~1, $g=1$), $r=3$ (Lemma~1, $g=0$).
\item $d=5$: $r=0$ (smooth), $r=1,2$ (Lemma~2), $r=3$ (Lemma~3), $r=4$ (Lemma~1, $g=2$), $r=5$ (Lemma~1, $g=1$), $r=6$ (Lemma~1, $g=0$).
\end{itemize}
All hypotheses of the lemmas are met: Lemma~1 needs $d\ge2g+1$ and $g<\delta(d)$, Lemma~2 needs $d\ge2r+1$ ($3\ge3$, $4\ge3$, $5\ge3$, $5\ge5$). The characteristic is arbitrary. \hfill$\square$

\section*{4. Elliptic curves}

Conventions. Throughout, $k$ is algebraically closed and curves are complete nonsingular curves over $k$. As in the text, characteristic $2$ is excluded wherever the cited results need it (Prop.~4.6, Cor.~4.7, Prop.~4.10 all assume $\operatorname{char}k\ne2$); Ex.~4.6 assumes $\operatorname{char}k=0$. For an elliptic curve $(X,P_0)$ the group law is the one transported from $\operatorname{Pic}^0X$ by $P\mapsto\mathcal L(P-P_0)$ (the discussion before Prop.~4.8); we write $\tau_a$ for translation by $a$, $\rho=-1_X$ for the inverse, and $n_X$ for multiplication by $n$. A divisor $D=\sum n_iP_i$ of degree $0$ corresponds to the point $\sum n_iP_i$ computed in the group law, and a divisor $D$ of degree $e$ to the point $\sum n_iP_i$ of the class $D-eP_0$. We write $j(\lambda)=2^8(\lambda^2-\lambda+1)^3/(\lambda^2(\lambda-1)^2)$ for the function of Thm.~4.1; it is constant on the $\Sigma_3$-orbits of Lemma~4.5.

\exercise{4.1}{Let $X$ be an elliptic curve over $k$, with $\operatorname{char}k\ne 2$, let $P\in X$ be a point, and let $R$ be the graded ring $R=\bigoplus_{n\ge0}H^0(X,\mathcal O_X(nP))$. Show that for suitable choice of $t,x,y$,
\[ R\cong k[t,x,y]/(y^2-x(x-t^2)(x-\lambda t^2)) \]
as a graded ring, where $k[t,x,y]$ is graded by setting $\deg t=1$, $\deg x=2$, $\deg y=3$.}
\pf
Write $R_n=H^0(X,\mathcal O_X(nP))\subseteq K(X)$; multiplication $R_n\times R_m\to R_{n+m}$ is the product of functions. Let $t=1\in R_1$ (multiplication by $t$ is the inclusion $R_n\subseteq R_{n+1}$). By the construction in the proof of Prop.~4.6 (with $P_0=P$) there are $x\in R_2$ and $y\in R_3$, with $1,x$ a basis of $R_2$ and $1,x,y$ a basis of $R_3$, such that in $K(X)$
\[ y^2=x(x-1)(x-\lambda),\qquad \lambda\ne0,1 . \]
(In that proof $x,y$ are changed only by $y\mapsto y+\tfrac12(a_1x+a_3)$, $x\mapsto (x-a)/(b-a)$, which stay in $R_3$, $R_2$ and keep the pole orders.) Since $x\in R_2\setminus R_1$ and $y\in R_3\setminus R_2$, and $x,y$ are regular off $P$, $x$ has a pole of order exactly $2$ at $P$ and $y$ of order exactly $3$. Homogenizing, we get in $R_6$ the relation $y^2=x(x-t^2)(x-\lambda t^2)$ (here $t^2x$ is just the function $x$ viewed in degree $4$, etc.). So there is a graded homomorphism
\[ \varphi: S:=k[t,x,y]/(F)\to R,\qquad F=y^2-x(x-t^2)(x-\lambda t^2). \]

\emph{$\varphi_n$ is surjective.} For $n\ge1$, $\dim R_n=n$ by Riemann--Roch (as $\deg nP=n>0=2g-2$). For $c\in\{0,1\}$ and $b\ge0$ with $2b+3c\le n$ the function $x^by^c$ lies in $R_n$ and has pole order exactly $2b+3c$ at $P$. Every integer $m\ge2$ is uniquely of the form $2b+3c$ with $c\in\{0,1\}$ (take $c\equiv m \bmod 2$), while $1=2b+3c$ is impossible; so the pole orders occurring are exactly $0,2,3,\dots,n$, all distinct. Functions with distinct pole orders at $P$ are linearly independent, so these are $n$ (for $n\ge2$; for $n=1$ just $\{1\}$) linearly independent elements of $R_n$, hence a basis. Thus $R_n$ is spanned by the images of the monomials $t^{n-2b-3c}x^by^c$, $c\le1$, and $\varphi_n$ is onto (for $n=0$ it is trivial).

\emph{$\varphi_n$ is injective.} In $S$ we may use $F$ to replace $y^2$ by $x(x-t^2)(x-\lambda t^2)$, which is homogeneous of the same degree $6$; so $S_n$ is spanned by the monomials $t^ax^by^c$ with $c\le 1$ and $a+2b+3c=n$. By the count just made there are $n$ such monomials if $n\ge1$ (and one if $n=0$). Hence $\dim S_n\le n=\dim R_n$, and since $\varphi_n$ is onto it is an isomorphism. Thus $\varphi$ is an isomorphism of graded rings, i.e.\ $\ker(k[t,x,y]\to R)=(F)$. \hfill$\square$

\exercise{4.2}{If $D$ is any divisor of degree $\ge3$ on the elliptic curve $X$, and if we embed $X$ in $\mathbf P^n$ by the complete linear system $|D|$, show that the image of $X$ in $\mathbf P^n$ is projectively normal. \emph{Note.} It is true more generally that if $D$ is a divisor of degree $\ge 2g+1$ on a curve of genus $g$, then the embedding of $X$ by $|D|$ is projectively normal (Mumford [4, p.~55]).}
\pf
Let $d=\deg D\ge3$, so $n=d-1$ ($h^0(D)=d$ by Riemann--Roch) and $D$ is very ample (Cor.~3.2); $\mathcal O_X(1)=\mathcal L(D)$. Let $S=k[x_0,\dots,x_n]/I(X)$ be the homogeneous coordinate ring and $S_m$ its degree $m$ part. Since $X$ is normal, by II, Ex.~5.14 ($\Gamma_*(\mathcal O_X)=\bigoplus_m H^0(X,\mathcal O_X(m))$ is the integral closure of $S$) $X$ is projectively normal if and only if the maps $S_m\to H^0(X,\mathcal L(mD))$ are surjective for all $m\ge0$. Since $S_m$ is the image of $\operatorname{Sym}^mH^0(\mathcal L(D))$, it suffices to prove:
\[ (\ast)\qquad H^0(\mathcal L(D))\otimes H^0(\mathcal L(mD))\to H^0(\mathcal L((m+1)D))\ \text{ is surjective for all }m\ge1 ;\]
then by induction every element of $H^0(\mathcal L(mD))$ is a polynomial in $H^0(\mathcal L(D))$ (the cases $m=0,1$ being trivial: $m=1$ since $|D|$ is complete).

Recall that on $X$ (genus $1$, $K=0$) a divisor $E$ of positive degree has $H^1(\mathcal L(E))\cong H^0(\mathcal L(-E))^\vee=0$, $h^0(E)=\deg E$.

\emph{Case $m\ge2$.} Since $|D|$ is base point free and $k$ is infinite we can choose $s_1,s_2\in H^0(\mathcal L(D))$ without common zero (take $s_1\ne0$, then $s_2$ outside the finitely many hyperplanes $\{s: s(p)=0\}$, $p\in\operatorname{Supp}\operatorname{div}s_1$). The Koszul sequence
\[ 0\to\mathcal L(-D)\xrightarrow{(s_2,-s_1)}\mathcal O_X^2\xrightarrow{(s_1,s_2)}\mathcal L(D)\to0 \]
is exact; tensoring with $\mathcal L(mD)$ and taking cohomology we get that $H^0(\mathcal L(mD))^2\to H^0(\mathcal L((m+1)D))$ is onto as soon as $H^1(\mathcal L((m-1)D))=0$, true because $m\ge2$. A fortiori $(\ast)$ holds.

\emph{Case $m=1$.} By Bertini (II, 8.18) there is a hyperplane $H\not\supseteq X$ meeting $X$ in a reduced divisor $E=p_1+\dots+p_d$ ($p_i$ distinct), $E\in|D|$, cut out by a section $s\in H^0(\mathcal L(D))$. Let $W\subseteq H^0(\mathcal L(2D))$ be the image of $H^0(\mathcal L(D))\otimes H^0(\mathcal L(D))$. From $0\to\mathcal L(D)\xrightarrow{s}\mathcal L(2D)\to\mathcal L(2D)|_E\to0$ and $H^1(\mathcal L(D))=0$ we get that $H^0(\mathcal L(2D))\to\mathcal L(2D)|_E$ is onto with kernel $sH^0(\mathcal L(D))\subseteq W$. So it suffices to show that $W$ maps onto $\mathcal L(2D)|_E$. Fix a generator $\varepsilon_i$ of the fibre of $\mathcal L(D)$ at $p_i$; this identifies $\mathcal L(D)|_E\cong k^d$ and $\mathcal L(2D)|_E\cong k^d$ (via $\varepsilon_i^2$) so that the multiplication becomes the coordinatewise product. Let $U\subseteq k^d$ be the image of $H^0(\mathcal L(D))\to\mathcal L(D)|_E$; the kernel is $ks$ (sections of $\mathcal L(D-E)\cong\mathcal O_X$), so $\dim U=d-1$: $U=\{a:\sum c_ia_i=0\}$ is a hyperplane. The image of $W$ is the span of $U\cdot U$.

Claim: all $c_i\ne0$, i.e.\ $e_i\notin U$. If $e_i\in U$ there is $u\in H^0(\mathcal L(D))$ vanishing at $p_j$ for $j\ne i$ but not at $p_i$. Then $\operatorname{div}u=\sum_{j\ne i}p_j+q$ for a point $q$, and $\operatorname{div}u\sim E$ gives $q\sim p_i$, so $q=p_i$ (distinct points of a curve of genus $1$ are not linearly equivalent, as $X\not\cong\mathbf P^1$), so $\operatorname{div}u=E$ and $u(p_i)=0$, a contradiction.

Now, as $d\ge3$, for each $i$ pick $j\ne k$ both different from $i$. Then $c_je_i-c_ie_j\in U$ and $c_ke_i-c_ie_k\in U$, and their product is $c_jc_ke_i$. Hence all $e_i\in\operatorname{span}(U\cdot U)$, so $W$ maps onto $k^d$ and $W=H^0(\mathcal L(2D))$.

This proves $(\ast)$, hence projective normality. \hfill$\square$

\exercise{4.3}{Let the elliptic curve $X$ be embedded in $\mathbf P^2$ so as to have the equation $y^2=x(x-1)(x-\lambda)$. Show that any automorphism of $X$ leaving $P_0=(0,1,0)$ fixed is induced by an automorphism of $\mathbf P^2$ coming from the affine $(x,y)$-plane given by
\[ \begin{cases} x'=ax+b\\ y'=cy.\end{cases}\]
In each of the four cases of (4.7), describe these automorphisms of $\mathbf P^2$ explicitly, and hence determine the structure of the group $G=\operatorname{Aut}(X,P_0)$.}
\pf
Throughout $\operatorname{char}k\ne2$. Let $x,y\in K(X)$ be the affine coordinates; $\{1,x\}$ is a basis of $H^0(\mathcal O(2P_0))$, $\{1,x,y\}$ of $H^0(\mathcal O(3P_0))$, $X\subset\mathbf P^2$ is embedded by $|3P_0|$, and $x,y$ have poles of exact orders $2,3$ at $P_0$ and nowhere else (proof of Prop.~4.6).

\emph{Reduction to affine maps.} Let $\sigma\in\operatorname{Aut}(X)$ with $\sigma(P_0)=P_0$. Then $\sigma^*x=x\circ\sigma$ has a pole of order exactly $2$ at $\sigma^{-1}(P_0)=P_0$ and no other pole, so $\sigma^*x\in H^0(\mathcal O(2P_0))$, $\sigma^*x=ax+b$ with $a\ne0$. Similarly $\sigma^*y\in H^0(\mathcal O(3P_0))$ with exact pole order $3$: $\sigma^*y=cy+dx+e$, $c\ne0$. Since $\sigma$ is an automorphism, $(\sigma^*y)^2=(\sigma^*x)(\sigma^*x-1)(\sigma^*x-\lambda)$, i.e.
\[ c^2y^2+2cy(dx+e)+(dx+e)^2=(ax+b)(ax+b-1)(ax+b-\lambda). \]
Now $y^2=x(x-1)(x-\lambda)\in k[x]$ and $K(X)$ has basis $1,y$ over $k(x)$; comparing the coefficients of $y$ gives $2c(dx+e)=0$, so $d=e=0$ (char $\ne2$). Thus $\sigma^*x=ax+b$, $\sigma^*y=cy$, i.e.\ $\sigma(x_0,y_0)=(ax_0+b,cy_0)$: $\sigma$ is induced by the automorphism $[x:y:z]\mapsto[ax+bz:cy:z]$ of $\mathbf P^2$ (which fixes $(0,1,0)$).

\emph{Which $(a,b,c)$ occur.} The map $(x,y)\mapsto(ax+b,cy)$ gives an automorphism of $X$ iff $c^2x(x-1)(x-\lambda)=\prod_{u\in\{0,1,\lambda\}}(ax+b-u)=a^3\prod_u\bigl(x-\tfrac{u-b}a\bigr)$ as polynomials, i.e.\ iff $c^2=a^3$ and $\ell(x)=ax+b$ permutes $\{0,1,\lambda\}$ (the roots $\frac{u-b}{a}=\ell^{-1}(u)$ must be $\{0,1,\lambda\}$). So $G=\{(a,b,c):c^2=a^3,\ \ell_{a,b}\{0,1,\lambda\}=\{0,1,\lambda\}\}$, with composition $(a,b,c)\circ(a',b',c')=(aa',ab'+b,cc')$. Sending $(a,b,c)$ to $\ell_{a,b}$ is a homomorphism of $G$ onto the group $H$ of affine maps $x\mapsto ax+b$ permuting $\{0,1,\lambda\}$ (each $a$ has the two square roots of $a^3$ available as $c$); its kernel is $\{(1,0,\pm1)\}=\{1,\rho\}$, where $\rho:y\mapsto -y$ is the inversion. So $|G|=2|H|$ and $G/\{1,\rho\}\cong H\subseteq\Sigma_3$. An affine map is determined by the images $p,q$ of $0,1$ and then sends $\lambda$ to $p+\lambda(q-p)$, which must be the remaining element. Going through the six permutations:
\begin{itemize}
\item identity: always.
\item $0\leftrightarrow1$ ($x\mapsto1-x$): need $1-\lambda=\lambda$, i.e.\ $\lambda=\tfrac12$.
\item $1\leftrightarrow\lambda$ ($x\mapsto\lambda x$): need $\lambda^2=1$, i.e.\ $\lambda=-1$.
\item $0\leftrightarrow\lambda$ ($x\mapsto\lambda+(1-\lambda)x$): need $2\lambda-\lambda^2=0$, i.e.\ $\lambda=2$.
\item $0\to1\to\lambda\to0$ ($x\mapsto1+(\lambda-1)x$) and its inverse $0\to\lambda\to1\to0$ ($x\mapsto\lambda-\lambda x$): need $\lambda^2-\lambda+1=0$.
\end{itemize}
This gives exactly the four cases of Cor.~4.7:

(i) \emph{$\lambda\notin\{-1,2,\tfrac12\}$ and $\lambda^2-\lambda+1\ne0$} (generic $j$): $H=1$, $G=\{(x,y)\mapsto(x,\pm y)\}\cong\mathbf Z/2$.

(ii) \emph{$\lambda\in\{-1,2,\tfrac12\}$} (these are three distinct values if $\operatorname{char}k\ne3$ and give $j=2^6 3^3=1728$; then $\lambda^2-\lambda+1\ne0$): $H=\{1,\ell\}$ with $\ell=-x,\ 2-x,\ 1-x$ respectively, i.e.\ $a=-1$ and $c^2=a^3=-1$. The automorphisms are $(x,y)\mapsto(\ell(x),\pm iy)$, $i^2=-1$ (together with $(x,\pm y)$). The element $g=(-1,b,i)$ satisfies $g^2=(1,0,-1)=\rho\ne1$, so $g$ has order $4$ and $G=\langle g\rangle\cong\mathbf Z/4$. For $\lambda=-1$, $g:(x,y)\mapsto(-x,iy)$.

(iii) \emph{$\lambda^2-\lambda+1=0$, $\operatorname{char}k\ne3$} ($\lambda=-\omega,-\omega^2$, $j=0$): $H=\langle\ell\rangle\cong\mathbf Z/3$, $\ell=1+(\lambda-1)x$. Here $a=\lambda-1=\lambda^2$, $a^3=\lambda^6=1$ (as $\lambda^3=\lambda\cdot\lambda^2=\lambda^2-\lambda=-1$), so $c=\pm1$ and the automorphisms are $(x,y)\mapsto(\ell^r(x),\pm y)$. For $g:(x,y)\mapsto(\ell(x),-y)$, $g^3=\rho$, so $g$ has order $6$ and $G=\langle g\rangle\cong\mathbf Z/6$.

(iv) \emph{$\operatorname{char}k=3$, $\lambda=-1$} ($j=0=1728$): then $-1=2=\tfrac12$ and $\lambda^2-\lambda+1=(\lambda+1)^2=0$, so all six conditions hold and $H$ is the full group of affine permutations $x\mapsto\pm x+b$, $b\in\mathbf F_3$, of $\{0,1,-1\}=\mathbf F_3$, $|G|=12$. Put $t=(1,1,1):(x,y)\mapsto(x+1,y)$ (order $3$) and $s=(-1,0,i):(x,y)\mapsto(-x,iy)$, where $i\in\mathbf F_9$, $i^2=-1$. Then $s^2=\rho$, so $s$ has order $4$, and a direct computation gives $sts^{-1}=t^{-1}$. Since $\langle t\rangle\cap\langle s\rangle=1$ and $3\cdot4=12$, $G=\langle t\rangle\rtimes\langle s\rangle\cong\mathbf Z/3\rtimes\mathbf Z/4$ with $s$ acting by inversion (the nonabelian ``dicyclic'' group of order $12$; it is a central extension of $\Sigma_3$ by $\{1,\rho\}$).

Here $\{1,\rho\}$ is central in all cases, which is why in (i)--(iii) $G$ is abelian ($G/\{1,\rho\}=H$ is cyclic). \hfill$\square$

\exercise{4.4}{Let $X$ be an elliptic curve in $\mathbf P^2$ given by an equation of the form
\[ y^2+a_1xy+a_3y=x^3+a_2x^2+a_4x+a_6. \]
Show that the $j$-invariant is a rational function of the $a_i$, with coefficients in $\mathbf Q$. In particular, if the $a_i$ are all in some field $k_0\subseteq k$, then $j\in k_0$ also. Furthermore, for every $\alpha\in k_0$, there exists an elliptic curve defined over $k_0$ with $j$-invariant equal to $\alpha$.}
\pf
As throughout this section of the text, $\operatorname{char}k\ne2$ (the definition of $j$ requires it). (The statement as printed omits the characteristic hypothesis; I use this standing assumption.)

\emph{Step 1: reduction to $y^2=$ cubic.} Put $\bar y=y+\tfrac12(a_1x+a_3)$. Then $\bar y^2=f(x)$ with
\[ f(x)=x^3+Ax^2+Bx+C,\quad A=a_2+\tfrac14a_1^2,\ B=a_4+\tfrac12a_1a_3,\ C=a_6+\tfrac14a_3^2. \]
Let $P_0=[0:1:0]$ be the point at infinity (the only point on $z=0$). It is nonsingular, and in the chart $y=1$ with coordinates $u=x/y$, $v=z/y$ the curve is $v+a_1uv+a_3v^2=u^3+\cdots$, whose tangent at the origin is $v=0$; so $u$ is a local parameter, $\operatorname{ord}v=3$ and $x/z=u/v$ has a pole of order exactly $2$ at $P_0$ and no other pole. Thus $x\in H^0(\mathcal O(2P_0))$ is nonconstant and defines the morphism of degree $2$, $x:X\to\mathbf P^1$, i.e.\ the morphism of $|2P_0|$, with $x(P_0)=\infty$. Its branch points are $\infty$ and the points with $\bar y=0$, i.e.\ the roots $e_1,e_2,e_3$ of $f$; since $X$ is nonsingular they are distinct (a multiple root of $f$ gives a singular point $(e,0)$). Applying the affine map $x\mapsto(x-e_1)/(e_2-e_1)$ we get the $\lambda$ of the text, $\lambda=(e_3-e_1)/(e_2-e_1)$, and by definition (and Thm.~4.1(a)) $j(X)=j(\lambda)$.

\emph{Step 2: $j$ as a symmetric function of the roots.} Substituting $\lambda=\frac{e_3-e_1}{e_2-e_1}$, $\lambda-1=\frac{e_3-e_2}{e_2-e_1}$ gives
\[ j(X)=2^8\frac{N^3}{\Delta_f},\qquad N=(e_2-e_1)^2-(e_2-e_1)(e_3-e_1)+(e_3-e_1)^2,\quad \Delta_f=\prod_{i<j}(e_i-e_j)^2 . \]
Now $N=\sum e_i^2-\sum_{i<j}e_ie_j=s_1^2-3s_2=A^2-3B$ and $\Delta_f$ is the discriminant of $f$: $\Delta_f=A^2B^2-4B^3-4A^3C+18ABC-27C^2$ (both are symmetric polynomials with integer coefficients in the $e_i$, hence polynomials in $s_1=-A,s_2=B,s_3=-C$). Therefore
\[ j(X)=\frac{2^8\,(A^2-3B)^3}{A^2B^2-4B^3-4A^3C+18ABC-27C^2}, \]
a rational function of $a_1,a_2,a_3,a_4,a_6$ with rational coefficients (only $\tfrac12,\tfrac14$ occur), whose denominator $\Delta_f\ne0$ because $X$ is nonsingular. (In the usual notation $b_2=a_1^2+4a_2$, $b_4=2a_4+a_1a_3$, $b_6=a_3^2+4a_6$, $c_4=b_2^2-24b_4$ one has $A^2-3B=c_4/16$, $\Delta_f=\Delta/16$ and this is $j=c_4^3/\Delta$; I checked this with a computer algebra system.) If all $a_i\in k_0$ then $A,B,C\in k_0$ and $j\in k_0$.

\emph{Step 3: every $\alpha\in k_0$ occurs.} We exhibit curves of the given form with coefficients in $k_0$ (these are then nonsingular cubics, i.e.\ elliptic curves, defined over $k_0$); the nonsingularity is the condition $\Delta\ne0$, where $\Delta=-b_2^2b_8-8b_4^3-27b_6^2+9b_2b_4b_6$, $b_8=a_1^2a_6+4a_2a_6-a_1a_3a_4+a_2a_3^2-a_4^2$ (so $\Delta=16\Delta_f$).
\begin{itemize}
\item $\alpha\ne0,1728$: $E_\alpha:\ y^2+xy=x^3-\dfrac{36}{\alpha-1728}x-\dfrac{1}{\alpha-1728}$. Here $a_1=1$, $a_2=a_3=0$, and a direct computation gives $c_4=\dfrac{\alpha}{\alpha-1728}$, $\Delta=\dfrac{\alpha^2}{(\alpha-1728)^3}\ne0$, so $j=c_4^3/\Delta=\alpha$.
\item $\alpha=1728$: $y^2=x^3-x$, with $c_4=48$, $\Delta=64\ne0$, $j=48^3/64=1728$ (this is Example 4.6.1; $\lambda=-1$).
\item $\alpha=0\ne1728$ (i.e.\ $\operatorname{char}k\ne3$): $y^2+y=x^3$, with $c_4=0$, $\Delta=-27\ne0$, $j=0$.
\end{itemize}
If $\operatorname{char}k=3$ then $1728=0$, so the case $\alpha=0$ is covered by the second line and the others by the first. \hfill$\square$

\exercise{4.5}{Let $X,P_0$ be an elliptic curve having an endomorphism $f:X\to X$ of degree $2$.
\begin{enumerate}[label=(\alph*),leftmargin=*]
\item If we represent $X$ as a $2$-$1$ covering of $\mathbf P^1$ by a morphism $\pi:X\to\mathbf P^1$ ramified at $P_0$, then as in (4.4), show that there is another morphism $\pi':X\to\mathbf P^1$ and a morphism $g:\mathbf P^1\to\mathbf P^1$, also of degree $2$, such that $\pi\circ f=g\circ\pi'$.
\item For suitable choices of coordinates in the two copies of $\mathbf P^1$, show that $g$ can be taken to be the morphism $x\mapsto x^2$.
\item Now show that $g$ is branched over two of the branch points of $\pi$, and that $g^{-1}$ of the other two branch points of $\pi$ consists of the four branch points of $\pi'$. Deduce a relation involving the invariant $\lambda$ of $X$.
\item Solving the above, show that there are just three values of $j$ corresponding to elliptic curves with an endomorphism of degree $2$, and find the corresponding values of $\lambda$ and $j$. [\emph{Answers}: $j=2^6\cdot3^3$; $j=2^6\cdot5^3$; $j=-3^3\cdot5^3$.]
\end{enumerate}}
\pf
Throughout $\operatorname{char}k\ne2$; ``endomorphism'' means $f(P_0)=P_0$ (any morphism $X\to X$ of degree $2$ becomes one after composing with a translation, which does not change degrees), so $f$ is a homomorphism (Lemma~4.9). Here $\pi$ is the morphism defined by $|2P_0|$, $\pi(P_0)=\infty$, with branch points $\pi(X[2])$ (the ramification points of $\pi$ are the $4$ points $R$ with $2R\sim2P_0$, i.e.\ $R=-R$).

\emph{Preliminaries.} Let $\rho=-1_X$, so $\pi\rho=\pi$ (the fibres of $\pi$ are the divisors $R+\rho R\in|2P_0|$, Lemma~4.2) and $K(X)^{\rho}=\pi^*K(\mathbf P^1)$ (both have index $2$ in $K=K(X)$).

\emph{Cross ratios.} For $4$ distinct points $z_1,\dots,z_4\in\mathbf P^1$ let $c(z_1,z_2,z_3,z_4)=\frac{(z_4-z_1)(z_2-z_3)}{(z_4-z_3)(z_2-z_1)}$, the image of $z_4$ under the Moebius map sending $z_1,z_2,z_3$ to $0,1,\infty$. When the labelling of the four points is permuted, $c$ runs through a $\Sigma_3$-orbit as in Lemma~4.5 (classical), and $j(c)$ is constant on it. \emph{Fact (R):} if $\varphi:X\to\mathbf P^1$ is any morphism of degree $2$, then it has exactly $4$ branch points (Hurwitz) and $j(X)=j(c(z_1,\dots,z_4))$ for these points in any order. Indeed, if $P_1$ is a ramification point of $\varphi$ then $\varphi^*\varphi(P_1)=2P_1$, so $\varphi$ is (up to an automorphism of $\mathbf P^1$) the morphism defined by $|2P_1|$ (II, 7.8.1); apply the construction of $j$ with $P_1$ as base point and Thm.~4.1(a) (the proof of (4.4)). Sending $\varphi(P_1)$ to $\infty$ and then two other branch points to $0,1$ by an affine map gives $c(a,b,\infty,c_3)=\frac{c_3-a}{b-a}=\lambda$.

\emph{(a)} Since $\deg f=2$ and $\operatorname{char}k\ne2$, $f$ is separable, so $\ker f=\{0,T\}$ with $T$ a point of order $2$, $T=-T$. Let $K=K(X)$ and $h=\pi\circ f:X\to\mathbf P^1$, $\deg h=4$, $L=h^*K(\mathbf P^1)\subseteq K$. Consider the automorphisms $\rho$, $\tau_T$, $\sigma=\rho\tau_T:x\mapsto -x-T=T-x$ of $X$ and $G'=\{1,\rho,\tau_T,\sigma\}\cong(\mathbf Z/2)^2$ (these are four distinct maps: $\rho,\sigma$ have fixed points, $\tau_T$ none, and $\rho\ne\sigma$). We have $h\rho=\pi f\rho=\pi\rho f=\pi f=h$ (as $f$ is a homomorphism, $f\rho=\rho f$) and $h\tau_T=\pi\circ(f(x)+f(T))=h$. So $L\subseteq K^{G'}$, and $[K:K^{G'}]=4=[K:L]$, hence $L=K^{G'}$ and $K/L$ is Galois with group $G'$; $h$ is the quotient map $X\to X/G'\cong\mathbf P^1$.

Let $Y=X/\langle\sigma\rangle$ (the curve with function field $K^\sigma$) and $\pi':X\to Y$ the quotient. $\sigma$ has fixed points $\{y:2y=T\}$, four of them ($2_X$ is a surjective homomorphism with kernel $X[2]$, which consists of the $4$ ramification points of $\pi$, so $\{y:2y=T\}$ is a coset of $X[2]$), so Hurwitz for the tame degree $2$ cover $\pi'$ gives $0=2(2g(Y)-2)+4$, $g(Y)=0$, $Y\cong\mathbf P^1$. Since $L\subseteq K^\sigma\subseteq K$, $h=g\circ\pi'$ for a morphism $g:Y\to X/G'=\mathbf P^1$ of degree $[K^\sigma:L]=2$. So $\pi\circ f=g\circ\pi'$ with $\pi'$ of degree $2$ and $g$ of degree $2$. (Also $K^{\rho}\supseteq L$ gives $\pi f=g_1\pi$ for the map $g_1$ induced by $f$; the new factorization through $\pi'$ is the one wanted.)

\emph{(b)} A degree-$2$ morphism $g:\mathbf P^1\to\mathbf P^1$ has two ramification points $r_1,r_2$ (Hurwitz: $-2=2(-2)+\deg R$, tame), with distinct images $b_1\ne b_2$ (else the fibre over $b_1$ would have degree $\ge4$). Choose coordinates $s$ on the source with $r_1=0,r_2=\infty$, and $t$ on the target with $b_1=0,b_2=\infty$. Then $g^*(0)=2\cdot0$, $g^*(\infty)=2\cdot\infty$, so $t\circ g=\kappa s^2$ and replacing $s$ by $\kappa^{-1/2}s$ we get $g(s)=s^2$.

\emph{(c)} Now work with the fixed points of the three involutions: $F_\rho=X[2]$, $F_\sigma=\{2y=T\}$, $F_{\tau}=\varnothing$. The ramification points of $h$ are the points with nontrivial stabilizer in $G'$, i.e.\ $F_\rho\cup F_\sigma$ (disjoint: $2x=0\ne T$), $8$ points. The $G'$-orbit of $x\in F_\rho$ is $\{x,x+T\}$, similarly for $F_\sigma$, and $h$ identifies exactly the $G'$-orbits. So $B_\rho=h(F_\rho)$ and $B_\sigma=h(F_\sigma)$ are disjoint sets of $2$ points each, and $B=B_\rho\cup B_\sigma$ is the branch locus of $h$ ($4$ points). Since $f(F_\rho\cup F_\sigma)\subseteq X[2]$ ($2f(x)=f(2x)\in\{0,f(T)\}=0$), we have $B=\pi f(F_\rho\cup F_\sigma)\subseteq\pi(X[2])$, and both sets have $4$ points; so $B$ is exactly the set of branch points of $\pi$.

For $x\in X$ with image $y=\pi'(x)\in Y$ we have $e_x(h)=|\mathrm{Stab}_{G'}(x)|$ and $e_x(\pi')=|\mathrm{Stab}(x)\cap\langle\sigma\rangle|$, hence $e_y(g)=e_x(h)/e_x(\pi')$. This is $2$ if $x\in F_\rho$, and $1$ otherwise. So $g$ is branched exactly over $B_\rho$, two of the four branch points of $\pi$. The branch points of $\pi'$ are the $4$ points $\pi'(y)$, $y\in F_\sigma$ (distinct, since $\sigma$ fixes them); they lie over $h(F_\sigma)=B_\sigma$, and $g^{-1}(B_\sigma)$ has at most $2\cdot2=4$ points. Hence $g^{-1}(B_\sigma)$ consists exactly of the four branch points of $\pi'$: \emph{$g$ is branched over two branch points of $\pi$, and the preimage of the other two is the set of the four branch points of $\pi'$}.

Normalize as in (b): $g(s)=s^2$ with $B_\rho=\{0,\infty\}$; scaling $t$ by a constant (and $s$ by its square root) we may assume $B_\sigma=\{1,\nu\}$ with $\nu\ne0,1,\infty$. So the branch points of $\pi$ are $0,1,\nu,\infty$, those of $\pi'$ are $\pm1,\pm\mu$ with $\mu^2=\nu$. Fact (R) applied to $\pi$ and to $\pi'$ gives $j(X)=j(\nu)$ (as $c(0,1,\infty,\nu)=\nu$) and
\[ j(X)=j(c(-1,1,-\mu,\mu)),\qquad c(-1,1,-\mu,\mu)=\frac{(\mu+1)(1+\mu)}{(2\mu)\,2}=\frac{(1+\mu)^2}{4\mu}. \]
\emph{Relation.} The invariant $\lambda$ of $X$ (the class of $\nu$ under $\Sigma_3$) satisfies
\[ j(\nu)=j\Bigl(\frac{(1+\mu)^2}{4\mu}\Bigr),\qquad \mu^2=\nu. \]

\emph{(d)} We compute: $c=\frac{(1+\mu)^2}{4\mu}$ gives $c-1=\frac{(1-\mu)^2}{4\mu}$, $c^2-c+1=\frac{\mu^4+14\mu^2+1}{16\mu^2}$, so with $\nu=\mu^2$
\[ j(c)=2^8\frac{(\mu^4+14\mu^2+1)^3/(16\mu^2)^3}{(1-\mu^2)^4/(256\mu^4)}=\frac{2^4(\nu^2+14\nu+1)^3}{\nu(\nu-1)^4}. \]
Thus the relation reads $\dfrac{2^8(\nu^2-\nu+1)^3}{\nu^2(\nu-1)^2}=\dfrac{2^4(\nu^2+14\nu+1)^3}{\nu(\nu-1)^4}$, i.e.\ $P(\nu)=0$ with
\[ P(\nu)=16(\nu^2-\nu+1)^3(\nu-1)^2-\nu(\nu^2+14\nu+1)^3=(\nu+1)^2(\nu^2-6\nu+1)(\nu^2-\nu+16)(16\nu^2-\nu+1) \]
(the factorization is checked by expanding; $P(0)\ne0\ne P(1)$). The values of $j(\nu)$ at the roots:
\begin{itemize}
\item $\nu=-1$: $j=2^8\cdot27/4=1728=2^6\cdot3^3$ (the $\Sigma_3$-orbit is $\{-1,2,\tfrac12\}$).
\item $\nu^2-6\nu+1=0$, $\nu=3\pm2\sqrt2$: here $\nu^2-\nu+1=5\nu$ and $(\nu-1)^2=4\nu$, so $j=2^8\cdot125\nu^3/(4\nu^3)=8000=2^6\cdot5^3$.
\item $\nu^2-\nu+16=0$, $\nu=\frac{1\pm3\sqrt{-7}}2$: $\nu^2-\nu+1=-15$ and $\nu^2(\nu-1)^2=(\nu^2-\nu)^2=256$, so $j=2^8(-15)^3/256=-3375=-3^3\cdot5^3$; the roots of $16\nu^2-\nu+1$ are the reciprocals of these, in the same orbit, with the same $j$.
\end{itemize}
So necessarily $j(X)\in\{1728,8000,-3375\}$ (in characteristic $0$; in characteristic $p>2$ the same computation gives the reductions of these values).

\emph{Converse.} Conversely let $\nu$ be a root of $P$ and $\mu^2=\nu$. Put $X_\nu:w^2=t(t-1)(t-\nu)$ (so $j(X_\nu)=j(\nu)$) and let $X'$ be the curve of genus $1$ with $K(X')=k(s,v)$, $v^2=(s^2-1)(s^2-\nu)$, a double cover of $\mathbf P^1_s$ branched over $\pm1,\pm\mu$; by (R) and the computation above $j(X')=j((1+\mu)^2/4\mu)=j(\nu)=j(X_\nu)$, so $X'\cong X_\nu$ (Thm.~4.1(b)). The inclusion $k(t,w)\subseteq k(s,v)$, $t\mapsto s^2$, $w\mapsto sv$ (note $s^2v^2=s^2(s^2-1)(s^2-\nu)$) has degree $2$, so it gives a morphism $X'\to X_\nu$ of degree $2$, hence after composing with the isomorphism and a translation an endomorphism of degree $2$ of $X_\nu$. So exactly the three values $j=2^6\cdot3^3,\ 2^6\cdot5^3,\ -3^3\cdot5^3$ occur (in characteristic $0$), with $\lambda\in\{-1,2,\tfrac12\}$, $\{3\pm2\sqrt2,\dots\}$ and $\{\tfrac{1\pm3\sqrt{-7}}2,\dots\}$ respectively. \hfill$\square$

\exercise{4.6}{(a) Let $X$ be a curve of genus $g$ embedded birationally in $\mathbf P^2$ as a curve of degree $d$ with $r$ nodes. Generalize the method of (Ex.~2.3) to show that $X$ has $6(g-1)+3d$ inflection points. A node does not count as an inflection point. Assume $\operatorname{char}k=0$.

(b) Now let $X$ be a curve of genus $g$ embedded as a curve of degree $d$ in $\mathbf P^n$, $n\ge3$, not contained in any $\mathbf P^{n-1}$. For each point $P\in X$, there is a hyperplane $H$ containing $P$, such that $P$ counts at least $n$ times in the intersection $H\cap X$. This is called an \emph{osculating hyperplane} at $P$, a generalization of the tangent line for curves in $\mathbf P^2$. If $P$ counts $n+1$ times in $H\cap X$, we say $H$ is a \emph{hyperosculating hyperplane} and $P$ is a \emph{hyperosculation point}. Use Hurwitz's theorem and induction on $n$ to show that $X$ has $n(n+1)(g-1)+(n+1)d$ hyperosculation points.

(c) If $X$ is an elliptic curve, for any $d\ge3$, embed $X$ as a curve of degree $d$ in $\mathbf P^{d-1}$, and conclude that $X$ has exactly $d^2$ points of order $d$ in its group law.}
\pf
Throughout $\operatorname{char}k=0$ (Hurwitz's formula then reads $2g-2=m(2g'-2)+\sum(e_P-1)$ for a finite morphism of degree $m$). The counts are made with multiplicities, as in Ex.~2.3(e); the precise multiplicities are given below.

\emph{(a)} Let $\nu:X\to C\subset\mathbf P^2$ be the birational morphism onto the plane curve $C$ of degree $d$ with $r$ nodes as only singularities; $\nu$ is an immersion (each branch at a node is smooth), and $\nu^*\mathcal O(1)$ has degree $d$. We assume $d\ge2$ (for a line every tangent line is the line itself, the Gauss map is constant and the formula fails). For $P\in X$ let $T_P$ be the tangent line to the branch of $C$ at $P$ and $r_P=\operatorname{ord}_P(\nu^*T_P)\ge2$ its contact order with the branch. We say $P$ is an \emph{inflection point} if $r_P\ge3$ and count it $r_P-2$ times (nodes themselves do not count; only the branch through $P$ enters). There are only finitely many inflection points (as in Ex.~2.3(a); this also follows from the Wronskian argument in (b), or from the local formula below, which shows that they are among the ramification points of the nonconstant morphism $\phi$ defined below). The map $\gamma:P\mapsto T_P$ is a rational map $X\to(\mathbf P^2)^*$, hence a morphism (II, 6.8), and by the local computation below its value at every $P$ is $T_P$.

\emph{Local computation.} Fix $P$ and a line $L\ne T_P$, and let $Q_0=L\cap T_P$. Choose affine coordinates $(X_1,Y_1)$ on $\mathbf P^2$ (line at infinity not through $P$ or $Q_0$) with $P=(0,0)$, $T_P=\{Y_1=0\}$, so $Q_0=(q_1,0)$. The branch at $P$ is parametrized by $t\mapsto(t,y(t))$ with $t$ a local parameter at $P$ and $y=ct^{r}+\cdots$, $r=r_P$, $c\ne0$ (formal derivative $y'=dy/dt$). The tangent line at the point with parameter $t$ is $Y_1-y(t)=y'(t)(X_1-t)$. Write points of $L$ as $Q_0+uv$ with $v=(v_X,v_Y)$, $v_Y\ne0$ (as $L\ne T_P$). Then $T_t\cap L$ has parameter
\[ u(t)=\frac{y+y'(q_1-t)}{v_Y-y'v_X}. \]
The denominator is a unit. If $q_1\ne0$ ($P\notin L$) the numerator has order exactly $r-1$ (because $y'q_1=rcq_1t^{r-1}+\cdots$ and $y$ has order $r$; here $\operatorname{char}k=0$ is used). If $q_1=0$ ($P\in L$) the numerator is $y-ty'=c(1-r)t^r+\cdots$ of order exactly $r$.

\emph{The Gauss map to a line.} Choose a line $L$ which meets $C$ transversally in $d$ distinct smooth points, passes through no node and no inflection point, and is not tangent to any branch; define $\phi:X\to L$, $\phi(P)=T_P\cap L$ (a morphism, being the composite of $\gamma$ with the projection of $(\mathbf P^2)^*$ from the point $[L]$). By the local computation $\phi(P)=Q_0$ corresponds to $u(0)=0$ and the ramification index of $\phi$ at $P$ is $e_P=r_P-1$ if $P\notin L$ and $e_P=r_P=2$ if $P\in L$ (an inflection point is not on $L$, and for $P\in L$ the contact is $2$). In particular $\phi$ is ramified exactly at the $d$ points of $\nu^{-1}(L)$ (with $e=2$) and at the inflection points.

\emph{Degree of $\phi$.} Let $Q\in L$ be general and $\psi_Q:X\to\mathbf P^1$ the projection from $Q$ (a morphism, since $Q\notin C$): $P\mapsto$ the line $QP$. Its fibre over a line $\ell\ni Q$ is $\nu^*(\ell)$, so $\deg\psi_Q=d$, and the ramification index at $P$ is the contact order of $\ell=QP$ with the branch at $P$; so $\psi_Q$ is ramified at $P$ iff $QP$ is tangent to the branch, i.e.\ $Q\in T_P$, and then $e_P=r_P$. For $Q$ outside the finitely many tangent lines at inflection points, $e_P=2$ at all these points, and Hurwitz gives
\[ 2g-2=-2d+\#\{P:Q\in T_P\},\quad\text{so}\quad \#\{P:Q\in T_P\}=2g-2+2d. \]
On the other hand $\phi^{-1}(Q)=\{P:Q\in T_P\}$ (as $T_P\ne L$), and for $Q$ outside the finite branch set of $\phi$ this fibre is reduced; so $m:=\deg\phi=2g-2+2d$. (Check: for $C$ smooth, $g=\frac{(d-1)(d-2)}2$ and $m=d(d-1)$, as in Ex.~2.3(e).)

\emph{Hurwitz for $\phi$.} With $\Phi=\sum_P(r_P-2)$, the number of inflection points counted with multiplicity,
\[ 2g-2=-2m+d+\Phi,\quad\text{hence}\quad \Phi=2g-2+2(2g-2+2d)-d=6(g-1)+3d. \]
So $X$ has $6(g-1)+3d$ inflection points. Since $g=\frac{(d-1)(d-2)}2-r$, this is $3d(d-2)-6r$ (for a smooth curve $3d(d-2)$, as in Ex.~2.3(e); for a nodal cubic, $3$). \hfill (The same number is obtained from (b) with $n=2$.)

\emph{(b)} Let $V\subseteq H^0(X,\mathcal L)$, $\mathcal L=\mathcal O_X(1)$ of degree $d$, be the $(n+1)$-dimensional image of $H^0(\mathbf P^n,\mathcal O(1))$ (injective because $X$ lies in no hyperplane); $V$ is base point free. For $P\in X$ the orders $\operatorname{ord}_P(s)$, $s\in V\setminus0$, take exactly $n+1$ values $0=a_0(P)<a_1(P)<\dots<a_n(P)$ (as $\dim\{s:\operatorname{ord}_Ps\ge j\}$ drops by $0$ or $1$ when $j\to j+1$); $a_i\ge i$. A hyperplane $H=\{s=0\}$ meets $X$ at $P$ with multiplicity $\operatorname{ord}_Ps$. The sections of order $\ge n$ at $P$ form a space of dimension $\ge1$ ($n$ linear conditions), so there is a hyperplane in which $P$ counts at least $n$ times, and the osculating hyperplane (max contact $a_n(P)$) is unique. $P$ is a hyperosculation point iff $a_n(P)\ge n+1$ iff $w(P):=\sum_{i=0}^n(a_i(P)-i)>0$ (if $a_n=n$ then $a_i=i$ for all $i$ since the $a_i$ are strictly increasing). We count $P$ with multiplicity $w(P)$ (for $n=2$ this is the $r_P-2$ of (a); for a point with $(a_i)=(0,1,\dots,n-1,n+1)$ it is $1$). We prove:
\[ \sum_{P\in X}w(P)=(n+1)d+n(n+1)(g-1). \]
(This replaces the inductive argument by the Wronskian, which proves the whole formula at once; for $n=1$ it is exactly Hurwitz's theorem $\sum(e_P-1)=2d+2g-2$.)

Let $s_0,\dots,s_n$ be a basis of $V$. At a point $P$ choose a local parameter $z$ and a generator $e$ of $\mathcal L$, write $s_i=f_ie$ with $f_i\in\mathcal O_{X,P}$, and let $D=d/dz$ be the derivation of $K(X)$ with $Dz=1$ (exists since $\operatorname{char}k=0$). Define the Wronskian $W=\det\bigl(D^jf_i\bigr)_{0\le i,j\le n}\in\mathcal O_{X,P}$.

(1) \emph{$W\,e^{\otimes(n+1)}(dz)^{\otimes n(n+1)/2}$ is independent of the choices.} If $e'=ue$ ($u$ a unit) then $f_i'=f_i/u$, and by Leibniz $D^j(f_i/u)=\sum_{l\le j}\binom jlD^{j-l}(u^{-1})D^lf_i$, so $(D^jf_i')=(D^lf_i)\cdot M$ with $M$ triangular with diagonal entries $u^{-1}$; thus $W'=u^{-(n+1)}W$ and $W'e'^{n+1}=We^{n+1}$. If $z'$ is another parameter, $d/dz'=(dz/dz')D$, and $(d/dz')^jf=(dz/dz')^jD^jf+\sum_{l<j}(\ast)D^lf$, so $W'=(dz/dz')^{n(n+1)/2}W$, which compensates $(dz')^{n(n+1)/2}=(dz'/dz)^{n(n+1)/2}(dz)^{n(n+1)/2}$. Hence the local expressions glue to a global section $\mathcal W$ of $\mathcal L^{\otimes(n+1)}\otimes\Omega_X^{\otimes n(n+1)/2}$ (regular, as $W\in\mathcal O_{X,P}$), of degree $(n+1)d+\frac{n(n+1)}2(2g-2)=(n+1)d+n(n+1)(g-1)$.

(2) \emph{$\mathcal W\ne0$.} The constants of $D$ in $K(X)$ are $k$ (if $Df=0$ for some $f\notin k$, then $z$ is algebraic over $k(f)$, hence separable over it, and $D$ vanishing on $k(f)$ would force $Dz=0$, contradicting $Dz=1$; here $\operatorname{char}k=0$ is used), and the classical Wronskian criterion says $W=0$ iff $f_0,\dots,f_n$ are linearly dependent over the constants. Since $s_0,\dots,s_n$ are linearly independent over $k$, $W\ne0$.

(3) \emph{$\operatorname{ord}_P(W)=w(P)$.} $W$ is multiplied by a nonzero constant if the basis of $V$ is changed (multilinearity), so we may take a basis with $\operatorname{ord}_Pf_i=a_i$, $f_i=c_iz^{a_i}+\cdots$, $c_i\ne0$. Then $D^jf_i=z^{a_i-j}\bigl(c_i\,a_i(a_i-1)\cdots(a_i-j+1)+O(z)\bigr)$, hence
\[ W=z^{\sum a_i-n(n+1)/2}\Bigl(\prod_ic_i\cdot\det\bigl(a_i^{\underline j}\bigr)+O(z)\Bigr),\quad a^{\underline j}=a(a-1)\cdots(a-j+1). \]
Since $a^{\underline j}$ is monic of degree $j$ in $a$, $\det(a_i^{\underline j})=\prod_{i<l}(a_l-a_i)\ne0$ (Vandermonde, characteristic $0$). So $\operatorname{ord}_PW=\sum a_i-\sum_{j=0}^nj=w(P)$.

Therefore $\sum_Pw(P)=\deg\mathcal W=(n+1)d+n(n+1)(g-1)$, which is the asserted number of hyperosculation points, counted with multiplicity $w(P)$.

\emph{(c)} Let $X$ be elliptic, $d\ge3$, and embed $X$ by the complete linear system $|dP_0|$ into $\mathbf P^{d-1}$ (Cor.~3.2; $h^0(dP_0)=d$ so the image lies in no hyperplane). Take $n=d-1$, $g=1$: $\sum w(P)=(n+1)d=d^2$.

For $P\in X$ and $0\le k\le d-1$, $\deg(dP_0-kP)=d-k\ge1$ so $h^0(dP_0-kP)=d-k$ (Riemann--Roch); for $k=d$, $h^0(dP_0-dP)=1$ if $dP\sim dP_0$ and $0$ otherwise. If $dP\not\sim dP_0$ the dimension of $\{s:\operatorname{ord}_Ps\ge k\}$ drops by $1$ at each step $k=0,\dots,d-1$, so the orders are $0,1,\dots,d-1$ and $w(P)=0$. If $dP\sim dP_0$ the dimensions are $d,d-1,\dots,1$ for $k=0,\dots,d-1$ and $1$ for $k=d$, so there is no drop at $k=d-1$ and the orders are $0,1,\dots,d-2,d$: $a_n=d$, $w(P)=1$. Thus the hyperosculation points are exactly the $P$ with $dP\sim dP_0$ (a hyperplane section is a divisor in $|dP_0|$ of degree $d$, so $P$ counts $d$ times in some $H\cap X$ iff $H\cdot X=dP$), each of weight exactly $1$. Since $dP\sim dP_0$ means $d(P-P_0)\sim0$, i.e.\ $dP=0$ in the group law, we conclude that $X$ has exactly $d^2$ points $P$ with $dP=0$ (``of order $d$'', meaning order dividing $d$). \hfill$\square$

\exercise{4.7}{\emph{The Dual of a Morphism.} Let $X$ and $X'$ be elliptic curves over $k$, with base points $P_0,P_0'$.
\begin{enumerate}[label=(\alph*),leftmargin=*]
\item If $f:X\to X'$ is any morphism, use (4.11) to show that $f^*:\operatorname{Pic}X'\to\operatorname{Pic}X$ induces a homomorphism $\hat f:(X',P_0')\to(X,P_0)$. We call this the \emph{dual} of $f$.
\item If $f:X\to X'$ and $g:X'\to X''$ are two morphisms, then $(g\circ f)^\wedge=\hat f\circ\hat g$.
\item Assume $f(P_0)=P_0'$, and let $n=\deg f$. Show that if $Q\in X$ is any point and $f(Q)=Q'$, then $\hat f(Q')=n_X(Q)$. (Do the separable and purely inseparable cases separately, then combine.) Conclude that $\hat f\circ f=n_X$ and $f\circ\hat f=n_{X'}$.
\item[(d)] If $f,g:X\to X'$ are two morphisms preserving the base points $P_0,P_0'$, then $(f+g)^\wedge=\hat f+\hat g$. [\emph{Hints}: It is enough to show for any $\mathcal L\in\operatorname{Pic}X'$, that $(f+g)^*\mathcal L\cong f^*\mathcal L\otimes g^*\mathcal L$. For any $f$, let $\Gamma_f:X\to X\times X'$ be the graph morphism. Then it is enough to show (for $\mathcal L'=p_2^*\mathcal L$) that $\Gamma_{f+g}^*(\mathcal L')=\Gamma_f^*\mathcal L'\otimes\Gamma_g^*\mathcal L'$. Let $\sigma:X\to X\times X'$ be the section $x\mapsto(x,P_0')$. Define a subgroup of $\operatorname{Pic}(X\times X')$ as follows: $\operatorname{Pic}_\sigma=\{\mathcal L\in\operatorname{Pic}(X\times X')\mid\mathcal L$ has degree $0$ along each fibre of $p_1$, and $\sigma^*\mathcal L=0$ in $\operatorname{Pic}X\}$. Note that this subgroup is isomorphic to the group $\operatorname{Pic}^0(X'/X)$ used in the definition of the Jacobian variety. Hence there is a $1$-$1$ correspondence between morphisms $f:X\to X'$ and elements $\mathcal L_f\in\operatorname{Pic}_\sigma$ (this defines $\mathcal L_f$). Now compute explicitly to show that $\Gamma_g^*(\mathcal L_f)=\Gamma_f^*(\mathcal L_g)$ for any $f,g$. Use the fact that $\mathcal L_{f+g}=\mathcal L_f\otimes\mathcal L_g$, and the fact that for any $\mathcal L$ on $X'$, $p_2^*\mathcal L\in\operatorname{Pic}_\sigma$, to prove the result.] (Part (d) is starred in the book.)
\item Using (d), show that for any $n\in\mathbf Z$, $\widehat{n_X}=n_X$. Conclude that $\deg n_X=n^2$.
\item Show for any $f$ that $\deg\hat f=\deg f$.
\end{enumerate}}
\pf
(As in the text, for results on $\operatorname{End}$ we use $\operatorname{char}k\ne2$ via Prop.~4.10; for (d), whose printed hint is somewhat cryptic, I give a direct proof along the same lines, working with the universal property of Thm.~4.11 for the family $\mathcal L_f$ rather than with the graph morphisms.) Recall the group law: $P\mapsto\mathcal L(P-P_0)$ is an isomorphism $(X,P_0)\to\operatorname{Pic}^0X$; and by Thm.~4.11, $J=X$ with $\mathcal L=\mathcal L(\Delta)\otimes p_1^*\mathcal L(-P_0)$ on $X\times X$ is a Jacobian: its restriction to $X\times\{t\}$ is $\mathcal L(t-P_0)$, and for every scheme $T$ of finite type over $k$ and every $\mathcal M\in\operatorname{Pic}^0(X/T)=\{\mathcal M\in\operatorname{Pic}(X\times T)$ of degree $0$ on all fibres$\}/p_2^*\operatorname{Pic}T$ there is a unique morphism $\varphi:T\to X$ with $\varphi^*\mathcal L\cong\mathcal M$ in $\operatorname{Pic}^0(X/T)$; moreover (Remark~4.10.6 and Thm.~4.11) the group law $\mu$ of $X$ is the one for which $\mu^*\mathcal L\cong p_1^*\mathcal L\otimes p_2^*\mathcal L$ in $\operatorname{Pic}^0(X/X\times X)$.

\emph{(a)} Let $\mathcal L'$ be the universal sheaf on $X'\times X'$ and let $\mathcal M=(f\times1_{X'})^*\mathcal L'$ on $X\times X'$. Its restriction to $X\times\{t'\}$ is $f^*\mathcal L(t'-P_0')$, of degree $0$ (as $\deg f^*D=\deg f\cdot\deg D$). So $\mathcal M\in\operatorname{Pic}^0(X/X')$ and by Thm.~4.11 there is a unique morphism $\hat f:X'\to X$ with $\hat f^*\mathcal L\cong\mathcal M$ in $\operatorname{Pic}^0(X/X')$. Restricting to $X\times\{t'\}$: $\mathcal L(\hat f(t')-P_0)\cong f^*\mathcal L(t'-P_0')$, that is
\[ \hat f(t')-P_0\sim f^*(t'-P_0')\qquad(t'\in X'). \]
So $\hat f$ is the morphism corresponding to $f^*:\operatorname{Pic}^0X'\to\operatorname{Pic}^0X$ under the isomorphisms $\operatorname{Pic}^0X'\cong X'$, $\operatorname{Pic}^0X\cong X$; as $f^*$ is a homomorphism, $\hat f$ is one (it sends $P_0'\mapsto P_0$; also Lemma~4.9). It depends only on $f$ and the base points, and $\hat f(t')$ is the group-law sum $\sum n_iR_i-\sum n_i'R_i'$ of $f^*(t')=\sum n_iR_i$ and $f^*(P_0')=\sum n_i'R_i'$. (If $f$ is constant, $f^*$ of a degree-$0$ divisor is $0$ and $\hat f=0$.)

\emph{(b)} For $Q''\in X''$: $\hat g(Q'')-P_0'\sim g^*(Q''-P_0'')$ and $\hat f(\hat gQ'')-P_0\sim f^*(\hat gQ''-P_0')\sim f^*g^*(Q''-P_0'')=(gf)^*(Q''-P_0'')$. By uniqueness in (a), $(g\circ f)^\wedge(Q'')=\hat f(\hat g(Q''))$. (If $f$ or $g$ is constant both sides are $0$, since $\hat f(P_0')=P_0$ and constants have dual $0$ by (a).)

\emph{(c)} Assume $f$ nonconstant (if $f=0$ everything is $0$). Let $Q\in X$, $Q'=f(Q)$. Since $f^*(Q')-f^*(P_0')$ has degree $0$, by (a) and the group law
\[ \hat f(Q')=\sum f^*(Q')-\sum f^*(P_0')\quad(\text{sums in the group law, with multiplicities}).\tag{1}\]
\emph{$f$ separable.} $f$ is a homomorphism, so $f\tau_a=\tau_{f(a)}f$ and the ramification index $e_R$ is the same at all points; as $f$ is separable, $e=1$, so $\ker f$ has exactly $n$ points and $f^*(Q')=\sum_{\kappa\in\ker f}(Q+\kappa)$, $f^*(P_0')=\sum_{\kappa}\kappa$. Then (1) gives $\hat f(Q')=nQ+\sum\kappa-\sum\kappa=nQ$.

\emph{$f$ purely inseparable} ($n=p^e$). Then $f$ is a bijection on points, $f^{-1}(Q')=\{Q\}$, so $f^*(Q')=nQ$ and $f^*(P_0')=nP_0$, and (1) gives $\hat f(Q')=nQ-nP_0=nQ$ (here $nQ$ is the multiple in the group law, $P_0=0$).

\emph{General case.} Let $K_s$ be the separable closure of $f^*K(X')$ in $K(X)$; then $K(X)/K_s$ is purely inseparable of degree $p^e$, so $K(X)^{p^e}\subseteq K_s$ and, comparing degrees ($[K:K^{p^e}]=p^e$), $K_s=K(X)^{p^e}$. Thus $f=h\circ F$ where $F:X\to X^{(p^e)}$ is the $p^e$-th power Frobenius ($X^{(p^e)}$ is again an elliptic curve with base point $F(P_0)$) and $h:X^{(p^e)}\to X'$ is separable, $h(F(P_0))=P_0'$, $n=\deg h\cdot p^e$. By (b), $\hat f=\hat F\circ\hat h$, and
\[ \hat f(f(Q))=\hat F\bigl(\hat h(h(FQ))\bigr)=\hat F\bigl((\deg h)\cdot FQ\bigr)=\hat F\bigl(F((\deg h)Q)\bigr)=p^e\deg h\cdot Q=nQ, \]
using the separable case for $h$, that $F$ is a homomorphism, and the inseparable case for $F$. So $\hat f\circ f=n_X$.

Then $f\circ\hat f\circ f=f\circ n_X=n_{X'}\circ f$ ($f$ is a homomorphism), i.e.\ $(f\hat f-n_{X'})\circ f=0$; as $f$ is surjective, $f\circ\hat f=n_{X'}$.

\emph{(d)} Let $f,g:X\to X'$ preserve base points; $f+g=\mu'\circ(f,g)$ with $\mu'$ the group law of $X'$. For a morphism $h:X\to X'$ put $\mathcal L_h:=(1_{X'}\times h)^*\mathcal L'\in\operatorname{Pic}(X'\times X)$, where $\mathcal L'=\mathcal L(\Delta')\otimes p_1^*\mathcal L(-P_0')$. Then $\mathcal L_h\cong\mathcal O(\Gamma_h)\otimes q^*\mathcal O(-P_0')$, where $q:X'\times X\to X'$ is the projection and $\Gamma_h=\{(h(x),x)\}$ is the (transposed) graph; its restriction to $X'\times\{x\}$ is $\mathcal L(h(x)-P_0')$, of degree $0$, and its restriction to $\{Q'\}\times X$ is $\mathcal O_X(h^*Q')$ (as $\Gamma_h\cap(\{Q'\}\times X)=h^{-1}(Q')$ scheme-theoretically). By Thm.~4.11 applied to the curve $X'$ with parameter space $T=X$ we have $\mathcal L_h=(h)^*\mathcal L'$ in $\operatorname{Pic}^0(X'/X)$, and for morphisms $f,g:T\to X'=J(X')$ the group law gives
\[ (f+g)^*\mathcal L'=(f,g)^*\mu'^*\mathcal L'=(f,g)^*(p_1^*\mathcal L'\otimes p_2^*\mathcal L')=f^*\mathcal L'\otimes g^*\mathcal L'\quad\text{in }\operatorname{Pic}^0(X'/T). \]
Taking $T=X$: $\mathcal L_{f+g}\cong\mathcal L_f\otimes\mathcal L_g\otimes r^*N$ for some $N\in\operatorname{Pic}X$ ($r:X'\times X\to X$ the projection). Restrict to $\{Q'\}\times X$ (note $r^*N$ restricts to $N$) and write $E_h(Q')=\mathcal L_h|_{\{Q'\}\times X}$; then
\[ E_{f+g}(Q')\cong E_f(Q')\otimes E_g(Q')\otimes N\qquad(Q'\in X'). \]
Taking $Q'=P_0'$ and dividing the two relations eliminates $N$:
\[ E_{f+g}(Q')\otimes E_{f+g}(P_0')^{-1}\cong\bigl(E_f(Q')\otimes E_f(P_0')^{-1}\bigr)\otimes\bigl(E_g(Q')\otimes E_g(P_0')^{-1}\bigr). \]
For $h$ nonconstant, $E_h(Q')=\mathcal O_X(h^*Q')$, so $E_h(Q')\otimes E_h(P_0')^{-1}=\mathcal O_X(h^*(Q'-P_0'))$; for $h$ constant ($=P_0'$) we have $\mathcal L_h\cong\mathcal O$ (the pull-back of $\Delta'$ is $\{P_0'\}\times X$, cancelled by $q^*\mathcal O(-P_0')$), and the class is trivial, which we write $h^*(Q'-P_0')=0$. Hence in all cases
\[ (f+g)^*(Q'-P_0')\sim f^*(Q'-P_0')+g^*(Q'-P_0')\qquad\text{for all }Q'. \]
By the characterization of the dual in (a) (with the above convention for constants, whose dual is $0$) this says $(f+g)^\wedge(Q')=\hat f(Q')+\hat g(Q')$, i.e.\ $(f+g)^\wedge=\hat f+\hat g$ (equal on closed points, hence equal as morphisms). (This is the content of the hint: $\mathcal L_{f+g}=\mathcal L_f\otimes\mathcal L_g$ in $\operatorname{Pic}_\sigma$, and the term $N$ is eliminated by normalizing at $P_0'$.)

\emph{(e)} For the identity, $\hat1(Q)-P_0\sim Q-P_0$, so $\hat1_X=1_X$; for the zero map $\hat0=0$. By (d), $(n\cdot1)^\wedge=n\hat1=n_X$ for $n\ge0$ by induction, and $0=(1+(-1))^\wedge=\hat1+\widehat{(-1)}$ gives $\widehat{(-1)}=-1$; so $\widehat{n_X}=n_X$ for all $n\in\mathbf Z$. For $n\ne0$, apply (c) to $f=n_X$ (a nonconstant endomorphism, Prop.~4.10): $\widehat{n_X}\circ n_X=(\deg n_X)_X$, i.e.\ $(n^2)_X=n_X\circ n_X=(\deg n_X)_X$. Since $m\mapsto m_X$ is injective $\mathbf Z\to\operatorname{End}(X,P_0)$ (Prop.~4.10), $\deg n_X=n^2$.

\emph{(f)} If $f(P_0)=P_0'$ and $f$ is nonconstant, (c) gives $\hat f\circ f=n_X$ with $n=\deg f$; taking degrees, $\deg\hat f\cdot\deg f=\deg(n_X)=n^2=(\deg f)^2$ by (e), so $\deg\hat f=\deg f$. For a general nonconstant $f$, write $f=\tau_a\circ f_0$ with $f_0(P_0)=P_0'$; then $\deg f=\deg f_0$ and $f^*(Q'-P_0')=f_0^*\bigl((Q'-a)-(P_0'-a)\bigr)\sim f_0^*(Q'-P_0')$ because the divisor $[Q'-a]-[P_0'-a]$ and $[Q']-[P_0']$ both correspond to the point $Q'$ of $\operatorname{Pic}^0X'\cong X'$; hence $\hat f=\hat f_0$ and $\deg\hat f=\deg f$. \hfill$\square$

\exercise{4.8}{For any curve $X$, the \emph{algebraic fundamental group} $\pi_1(X)$ is defined as $\varprojlim\operatorname{Gal}(K'/K)$, where $K$ is the function field of $X$, and $K'$ runs over all Galois extensions of $K$ such that the corresponding curve $X'$ is etale over $X$ (III, Ex.~10.3). Thus, for example, $\pi_1(\mathbf P^1)=1$ (2.5.3). Show that for an elliptic curve $X$,
\[ \begin{array}{ll}\pi_1(X)=\prod_{l\ \mathrm{prime}}\mathbf Z_l\times\mathbf Z_l & \text{if }\operatorname{char}k=0;\\[2pt] \pi_1(X)=\prod_{l\ne p}\mathbf Z_l\times\mathbf Z_l & \text{if }\operatorname{char}k=p\text{ and Hasse }X=0;\\[2pt] \pi_1(X)=\mathbf Z_p\times\prod_{l\ne p}\mathbf Z_l\times\mathbf Z_l & \text{if }\operatorname{char}k=p\text{ and Hasse }X\ne0,\end{array}\]
where $\mathbf Z_l=\varprojlim\mathbf Z/l^n$ is the $l$-adic integers. [\emph{Hints}: Any Galois etale cover $X'$ of an elliptic curve is again an elliptic curve. If the degree of $X'$ over $X$ is relatively prime to $p$, then $X'$ can be dominated by the cover $n_X:X\to X$ for some $n$ with $(n,p)=1$. The Galois group of the covering $n_X$ is $\mathbf Z/n\times\mathbf Z/n$. Etale covers of degree divisible by $p$ can occur only if the Hasse invariant of $X$ is not zero.]}
\pf
Let $p=\operatorname{char}k$ ($p\ne2$ in this section; the arguments below work for $p=0$ or $p$ odd), $K=K(X)$.

\emph{Step 0: etale covers are isogenies.} Let $X'\to X$ be a connected etale cover (with $K'/K$ Galois). By Hurwitz ($\deg R=0$) $2g(X')-2=0$, so $X'$ is an elliptic curve; choose $P_0'$ over $P_0$. Then $f:X'\to X$ is a separable homomorphism of degree $N$ (the group law is the one with these base points), and $\ker f$ has $N$ points; the translations $\tau_\kappa$, $\kappa\in\ker f$, are $N$ distinct automorphisms of $X'$ over $X$, so $K'/K$ is Galois with group $\ker f$ (a finite abelian group). Conversely, a separable isogeny $f$ is etale: ramification indices are constant by translation ($f\tau_a=\tau_{f(a)}f$) and equal $1$ at a generic point. Moreover, by Ex.~4.7(c), for $f:X'\to X$ of degree $N$ there is $\hat f:X\to X'$ with $f\circ\hat f=N_X$ and $\hat f\circ f=N_{X'}$ (and $\deg \hat f=N$).

Since all groups $\ker f$ are abelian, $\pi_1(X)=\varprojlim\ker f$; we compute this limit over a cofinal family of covers, each cover in the family dominating a given one; a cover $Y$ dominating $Y'$ induces a \emph{surjection} of Galois groups.

\emph{Step 1: torsion of prime-to-$p$ order.} \textbf{Lemma A.} If $p\nmid n$ (any $n\ge1$ in char $0$), then $n_X$ is separable of degree $n^2$ and $X[n]=\ker n_X\cong\mathbf Z/n\times\mathbf Z/n$. (This is the statement in Example 4.8.1.)
\emph{Proof.} $\deg n_X=n^2$ by Ex.~4.7(e). Let $\omega\ne0$ be a regular differential on $X$ ($g=1$). Since $\Omega_X\cong\mathcal O_X$ (via $\omega$), the pull-back of $\omega$ by the group law $\mu:X\times X\to X$, as a section of the relative differentials $\Omega_{X\times X/X}=p_1^*\Omega_X\cong\mathcal O_{X\times X}$ (relative to $p_2$), is a global section of $\mathcal O_{X\times X}$, hence a constant $c$ times $p_1^*\omega$; restricting to $X\times\{P_0\}$ gives $c=1$, and restricting to $X\times\{a\}$ gives $\tau_a^*\omega=\omega$: $\omega$ is translation invariant. For $h_1,h_2:X\to X$, the differential of $h_1+h_2=\mu\circ(h_1,h_2)$ at $x$ is $d\tau_{h_2(x)}\circ dh_{1,x}+d\tau_{h_1(x)}\circ dh_{2,x}$ (as $\mu(\cdot,b)=\tau_b$, $\mu(a,\cdot)=\tau_a$), so by invariance $(h_1+h_2)^*\omega=h_1^*\omega+h_2^*\omega$. By induction $n_X^*\omega=n\omega\ne0$. If $n_X$ were inseparable it would factor through Frobenius and $n_X^*\omega=0$; so $n_X$ is separable, and a separable homomorphism has $|\ker|=\deg=n^2$. Applying this to every divisor $e$ of $n$, $|X[e]|=e^2$. A finite abelian group $A$ of order $n^2$ with $|A[e]|=e^2$ for all $e\mid n$ is $\cong(\mathbf Z/n)^2$: for a prime $l\mid n$, $l^v\|n$, the $l$-part is $\bigoplus\mathbf Z/l^{b_i}$ with $\sum b_i=2v$, $|A[l]|=l^2$ forces exactly two summands $b_1+b_2=2v$, and $|A[l^v]|=l^{\min(b_1,v)+\min(b_2,v)}=l^{2v}$ forces $b_1,b_2\ge v$, so $b_1=b_2=v$. \hfill$\square$

\emph{Step 2: characteristic $0$.} Every etale connected cover $f:X'\to X$ of degree $N$ is dominated by $N_X:X\to X$, since $N_X=f\circ\hat f$. Thus $\{N_X\}$ is cofinal and $\pi_1(X)=\varprojlim_NX[N]$, $X[N]\cong(\mathbf Z/N)^2$, the transition maps being surjections. Decomposing $N=\prod l^{v_l}$ and using the Chinese remainder theorem, $\pi_1(X)=\prod_l\varprojlim_v A_v^{(l)}$ with $A_v^{(l)}=X[l^v]\cong(\mathbf Z/l^v)^2$ and surjective transition maps. \textbf{Lemma B.} Such an inverse limit is $\cong\mathbf Z_l^2$. \emph{Proof.} A surjection $\theta:(\mathbf Z/l^{v+1})^2\to(\mathbf Z/l^v)^2$ induces a surjection, hence an isomorphism, on the quotients by $l$ (both $\cong(\mathbf Z/l)^2$); so any lifts $\tilde e,\tilde e'$ of a basis $e,e'$ of the target generate the source modulo $l$ and therefore (Nakayama) form a basis of $(\mathbf Z/l^{v+1})^2$. Thus compatible bases $(e_v,e_v')$ of all $A_v$ can be chosen inductively, giving an isomorphism of the inverse system with $(\mathbf Z/l^v)^2$ and the natural reductions, whose limit is $\mathbf Z_l^2$. \hfill$\square$
Hence $\pi_1(X)=\prod_l\mathbf Z_l\times\mathbf Z_l$ in characteristic $0$.

\emph{Step 3: characteristic $p$; Hasse invariant and $p$-torsion.} We use Ex.~4.15, applied to $Y=X^{(p)}$ (so $Y_p=X$ and $F:X\to Y$ is the $k$-linear Frobenius of the exercise): the dual $\hat F:X^{(p)}\to X$ is separable iff the Hasse invariant of $X^{(p)}$ is $1$. Now $X^{(p)}$ is $X$ with the $k$-structure twisted by the Frobenius of $k$, hence has the same underlying ringed space and the same absolute Frobenius, so $X^{(p)}$ and $X$ have the same Hasse invariant (as in the proof of 4.15). Thus $\hat F$ is separable iff the Hasse invariant of $X$ is $1$. Write $V=\hat F$; then $V\circ F=p_X$ (Ex.~4.7(c)), $F$ is bijective on points, and separable degree and kernel satisfy $|\ker g|=\deg_sg$ for a homomorphism $g$ (translation invariance). Hence $|X[p]|=\deg_sV=p$ if the Hasse invariant is $1$ and $=1$ if it is $0$.

\textbf{Lemma C.} If $f:X'\to X$ is any isogeny, then $|X'[p]|=|X[p]|$. \emph{Proof.} From $f\circ p_{X'}=p_X\circ f$ and multiplicativity of separable degrees, $\deg_sf\cdot\deg_sp_{X'}=\deg_sp_X\cdot\deg_sf$, and $\deg_sp_Y=|Y[p]|$ (as $p_Y=VF$). \hfill$\square$
In particular $X^{(p^a)}$ (isogenous to $X$ via $F^a$) has $|X^{(p^a)}[p]|=|X[p]|$: it is ordinary iff $X$ is. Also an etale cover of degree divisible by $p$ gives a point of order $p$ on $X'$ (Cauchy in $\ker f$), so $|X'[p]|\ne1$, hence by Lemma C $|X[p]|\ne1$:

\emph{If $\operatorname{Hasse}X=0$, every etale cover has degree prime to $p$.} Then $f:X'\to X$ of degree $N$, $p\nmid N$, satisfies $f\circ\hat f=N_X$, which is separable (Lemma A), so $\hat f$ is separable; $N_X$ factors through $f$, i.e.\ $\{N_X:p\nmid N\}$ is cofinal. By Lemma A, $\ker N_X\cong(\mathbf Z/N)^2$, and as in Step 2 (with $l\ne p$) $\pi_1(X)=\prod_{l\ne p}\mathbf Z_l^2$.

\emph{If $\operatorname{Hasse}X\ne0$.} Put $X_a=X^{(p^a)}$, an ordinary curve with $V_a:=(F^a)^\wedge=\hat F_{X}\hat F_{X_1}\cdots:X_a\to X$ (Ex.~4.7(b)); each factor is separable (each $X_j$ is ordinary), so $V_a$ is separable of degree $p^a$, etale, and $V_a\circ F^a=(p^a)_X$, $F^a\circ V_a=(p^a)_{X_a}$. Since $\deg_sp_{X_a}=p$ we get $|X_a[p^a]|=p^a$ and $|X_a[p]|=p$, so $X_a[p^a]$ is cyclic; $\ker V_a\subseteq X_a[p^a]$ (as $(p^a)_{X_a}=F^aV_a$) has order $p^a$, so $\ker V_a=X_a[p^a]\cong\mathbf Z/p^a$. For $a\ge0$ and $m$ prime to $p$ let
\[ c_{a,m}=V_a\circ m_{X_a}:X_a\to X, \]
etale (both factors are separable), of degree $p^am^2$, with Galois group $G_{a,m}=\ker c_{a,m}$ of order $p^am^2$. The $p$-Sylow of $G_{a,m}$ is $\ker V_a\cong\mathbf Z/p^a$ (it is contained in $G_{a,m}$ since $V_a(m x)=mV_ax$) and the prime-to-$p$ part is $\ker m_{X_a}\cong(\mathbf Z/m)^2$ (Lemma A for $X_a$): $G_{a,m}\cong\mathbf Z/p^a\times(\mathbf Z/m)^2$.

\emph{Cofinality.} Let $f:X'\to X$ be etale of degree $N=p^am$ with $p\nmid m$. Then $f\circ\hat f=N_X=m_X\circ V_a\circ F^a=V_a\circ m_{X_a}\circ F^a=c_{a,m}\circ F^a$. The inseparable degree of $N_X$ is $p^a$ (as $m_X,V_a$ are separable), $f$ is separable, so $\hat f$ has inseparable degree exactly $p^a$; therefore (as in (c) of Ex.~4.7) $\hat f=h\circ F^a$ with $h:X_a\to X'$. Then $f\circ h\circ F^a=c_{a,m}\circ F^a$, so $f\circ h=c_{a,m}$ ($F^a$ is surjective): $c_{a,m}$ dominates $f$. Thus the $c_{a,m}$ ($a\ge0$, $p\nmid m$; ordered by $a'\le a$, $m'\mid m$) form a cofinal system (directed: $V_a=V_{a'}\circ W$ with $W=\hat F_{X_{a'}}\cdots\hat F_{X_{a-1}}:X_a\to X_{a'}$, so for $a'\le a$, $m'\mid m$ one has $c_{a,m}=c_{a',m'}\circ W\circ(m/m')_{X_a}$, a morphism over $X$).

\emph{The limit.} The chain $(a_k,m_k)=(k,k!_{p'})$, $k\ge1$ ($k!_{p'}$ the prime-to-$p$ part of $k!$) is cofinal in the directed set above, so $\pi_1(X)=\varprojlim_kG_k$, $G_k=G_{k,k!_{p'}}=P_k\times Q_k$ with $P_k=\ker V_k\cong\mathbf Z/p^k$ and $Q_k=\ker(k!_{p'})_{X_k}\cong(\mathbf Z/k!_{p'})^2$. The transition maps are the surjective homomorphisms induced by the dominations (restrictions of $W\circ(\cdot)_{X_k}$ above, which are homomorphisms); they respect the decomposition into $p$-primary and prime-to-$p$ parts. Hence $\pi_1(X)=\varprojlim P_k\times\varprojlim Q_k$. The $P_k$ are cyclic $p$-groups with surjective transitions, so compatible generators can be chosen inductively and $\varprojlim P_k=\mathbf Z_p$. Decomposing $Q_k=\bigoplus_{l\ne p}Q_k^{(l)}$, $Q_k^{(l)}\cong(\mathbf Z/l^{v_l(k)})^2$, with surjective transitions, we get $\varprojlim Q_k=\prod_{l\ne p}\mathbf Z_l^2$ by Lemma B (and since $v_l(k)\to\infty$). Hence $\pi_1(X)=\mathbf Z_p\times\prod_{l\ne p}\mathbf Z_l\times\mathbf Z_l$. \hfill$\square$

\exercise{4.9}{We say two elliptic curves $X,X'$ are \emph{isogenous} if there is a finite morphism $f:X\to X'$.
\begin{enumerate}[label=(\alph*),leftmargin=*]
\item Show that isogeny is an equivalence relation.
\item For any elliptic curve $X$, show that the set of elliptic curves $X'$ isogenous to $X$, up to isomorphism, is countable. [\emph{Hint}: $X'$ is uniquely determined by $X$ and $\ker f$.]
\end{enumerate}}
\pf
A finite morphism of curves is a nonconstant morphism. Composing with a translation of $X'$ we may always assume $f(P_0)=P_0'$, so that $f$ is a homomorphism (Lemma~4.9), a \emph{finite homomorphism} with finite kernel. As above $\operatorname{char}k\ne2$ where Prop.~4.10 is used.

\emph{(a)} Reflexive: $1_X$. Transitive: the composition of nonconstant morphisms is nonconstant. Symmetric: if $f:X\to X'$ is finite with $f(P_0)=P_0'$, $n=\deg f\ge1$, then by Ex.~4.7(c) $\hat f\circ f=n_X$, which is nonconstant (Prop.~4.10); hence $\hat f:X'\to X$ is nonconstant, so finite, and $X'$ is isogenous to $X$ implies $X$ is isogenous to $X'$.

\emph{(b)} Let $p=\operatorname{char}k$ ($p=0$ allowed). Every nonconstant $f:X\to X'$ factors as $X\xrightarrow{F^a}X^{(p^a)}\xrightarrow{h}X'$ with $F^a$ the $p^a$-th power Frobenius and $h$ separable (the separable closure of $f^*K(X')$ in $K(X)$ equals $K(X)^{p^a}$, as shown in Ex.~4.7(c); $a=0$ if $p=0$); $X^{(p^a)}$ is an elliptic curve depending only on $X$ and $a$, and $h$ is a separable homomorphism (base point $F^a(P_0)\mapsto P_0'$) with kernel $\Lambda=\ker h$, a finite subgroup of $X^{(p^a)}(k)$, $|\Lambda|=\deg h$.

\emph{Claim: $X'$ is determined up to isomorphism by $(a,\Lambda)$.} Let $Y=X^{(p^a)}$ and let $h_i:Y\to X_i'$ ($i=1,2$) be separable homomorphisms with the same kernel $\Lambda$. The translations $\tau_\kappa$, $\kappa\in\Lambda$, satisfy $h_i\tau_\kappa=h_i$, so $h_i^*K(X_i')\subseteq K(Y)^\Lambda$; as $[K(Y):h_i^*K(X_i')]=\deg h_i=|\Lambda|=[K(Y):K(Y)^\Lambda]$, we get $h_i^*K(X_i')=K(Y)^\Lambda$. Hence $K(X_1')\cong K(X_2')$ over $k$ and $X_1'\cong X_2'$ (equivalence of curves and function fields).

\emph{Countability.} The set of finite subgroups of $Y(k)$ is countable: every element of a finite subgroup is a torsion point, and $Y[m]=\ker m_Y$ is finite (of order $\le\deg m_Y=m^2$, Ex.~4.7(e)) for each $m\ge1$, so the torsion subgroup $\bigcup_mY[m]$ is countable, and a countable set has countably many finite subsets. So the set of pairs $(a,\Lambda)$ ($a\ge0$ an integer, $\Lambda\subseteq X^{(p^a)}(k)$ finite subgroup) is countable. By (a), every $X'$ isogenous to $X$ receives a finite $f:X\to X'$, giving a pair $(a,\Lambda)$ from which $X'$ is determined up to isomorphism. Hence there are at most countably many isomorphism classes of elliptic curves isogenous to $X$. \hfill$\square$

\exercise{4.10}{If $X$ is an elliptic curve, show that there is an exact sequence
\[ 0\to p_1^*\operatorname{Pic}X\oplus p_2^*\operatorname{Pic}X\to\operatorname{Pic}(X\times X)\to R\to0, \]
where $R=\operatorname{End}(X,P_0)$. In particular, we see that $\operatorname{Pic}(X\times X)$ is bigger than the sum of the Picard groups of the factors. Cf.\ (III, 12.6), (V, Ex.~1.6).}
\pf
Let $p_1,p_2:X\times X\to X$ be the projections; we regard $p_2$ as the structure map of a family of curves $X\times\{t\}\cong X$ (first factor), parametrized by $t\in X$ (second factor). For $\mathcal L\in\operatorname{Pic}(X\times X)$ write $\mathcal L_t=\mathcal L|_{X\times\{t\}}\in\operatorname{Pic}X$.

\emph{Definition of the map $\Psi:\operatorname{Pic}(X\times X)\to R$.} Put $\mathcal M=\mathcal L\otimes p_1^*(\mathcal L_{P_0})^{-1}$. Then $\mathcal M_t=\mathcal L_t\otimes\mathcal L_{P_0}^{-1}$ has degree $0$ for all $t$ ($\deg\mathcal L_t$ is constant in $t$: $\mathcal L$ is flat over $X$ and the degree is locally constant, III, 9.9). So $\mathcal M\in\operatorname{Pic}^0(X/X)$, and by Thm.~4.11 (universal property of $J=X$, with $T=X$) there is a unique morphism $\Psi(\mathcal L):X\to X$ with $\mathcal L(\Psi(\mathcal L)(t)-P_0)\cong\mathcal M_t=\mathcal L_t\otimes\mathcal L_{P_0}^{-1}$ for all $t$. Since $\mathcal M_{P_0}=\mathcal O$, $\Psi(\mathcal L)(P_0)=P_0$, so $\Psi(\mathcal L)\in\operatorname{End}(X,P_0)=R$. In words: $\Psi(\mathcal L)(t)$ is the point of $X\cong\operatorname{Pic}^0X$ corresponding to $\mathcal L_t\otimes\mathcal L_{P_0}^{-1}$.

\emph{$\Psi$ is a homomorphism.} $(\mathcal L\otimes\mathcal L')_t\otimes(\mathcal L\otimes\mathcal L')_{P_0}^{-1}=\bigl(\mathcal L_t\mathcal L_{P_0}^{-1}\bigr)\otimes\bigl(\mathcal L'_t\mathcal L'^{-1}_{P_0}\bigr)$, and the identification $X\cong\operatorname{Pic}^0X$ is a group isomorphism, so $\Psi(\mathcal L\otimes\mathcal L')(t)=\Psi(\mathcal L)(t)+\Psi(\mathcal L')(t)$, i.e.\ $\Psi(\mathcal L\otimes\mathcal L')=\Psi(\mathcal L)+\Psi(\mathcal L')$ in $R$ (the sum in $R$ being pointwise).

\emph{$\Psi$ is surjective.} Let $\alpha\in R$ and $\mathcal L=(1_X\times\alpha)^*\bigl(\mathcal L(\Delta)\otimes p_1^*\mathcal L(-P_0)\bigr)$. Then $\mathcal L_t=\mathcal L(\alpha(t)-P_0)$, and $\mathcal L_{P_0}=\mathcal L(\alpha(P_0)-P_0)=\mathcal O$, so $\Psi(\mathcal L)(t)$ is the point corresponding to $\mathcal L(\alpha(t)-P_0)$, namely $\alpha(t)$: $\Psi(\mathcal L)=\alpha$.

\emph{The image of $\iota=p_1^*\oplus p_2^*:\operatorname{Pic}X\oplus\operatorname{Pic}X\to\operatorname{Pic}(X\times X)$ is $\ker\Psi$.} If $\mathcal L=p_1^*A\otimes p_2^*B$, then $\mathcal L_t\cong A$ for every $t$, so $\mathcal M_t\cong\mathcal O$ and $\Psi(\mathcal L)=0$ (as $\Psi(\mathcal L)(t)=P_0$). Conversely if $\Psi(\mathcal L)=0$, then $\mathcal M_t\cong\mathcal O_X$ for all $t\in X$, so by the criterion recalled in the text (III, Ex.~12.4: for $T$ integral of finite type, $\mathcal M\in\operatorname{Pic}(X\times T)$ is trivial on every fibre $X\times\{t\}$ iff $\mathcal M\in p_2^*\operatorname{Pic}T$) $\mathcal M\cong p_2^*B$ for some $B\in\operatorname{Pic}X$. Then $\mathcal L\cong p_1^*(\mathcal L_{P_0})\otimes p_2^*B\in\operatorname{im}\iota$.

\emph{$\iota$ is injective.} If $p_1^*A\otimes p_2^*B\cong\mathcal O$, restricting to $X\times\{t\}$ gives $A\cong\mathcal O$ (as $p_2^*B$ is trivial on it), and restricting to $\{s\}\times X$ gives $B\cong\mathcal O$.

So the sequence is exact. Since $R\supseteq\mathbf Z\cdot1_X\ne0$ (indeed $\mathbf Z\hookrightarrow R$ by Prop.~4.10), $\operatorname{Pic}(X\times X)$ is strictly larger than $p_1^*\operatorname{Pic}X\oplus p_2^*\operatorname{Pic}X$; e.g.\ $\mathcal O(\Delta)$ (with $\Psi(\mathcal O(\Delta))=1_X\ne0$; likewise $\mathcal O(\Gamma_f)$ for a nonconstant $f\in R$, where $\Psi=\hat f\ne0$) is not of the form $p_1^*A\otimes p_2^*B$. \hfill$\square$

\exercise{4.11}{Let $X$ be an elliptic curve over $\C$, defined by the elliptic functions with periods $1,\tau$. Let $R$ be the ring of endomorphisms of $X$.
\begin{enumerate}[label=(\alph*),leftmargin=*]
\item If $f\in R$ is a nonzero endomorphism corresponding to complex multiplication by $\alpha$, as in (4.18), show that $\deg f=|\alpha|^2$.
\item If $f\in R$ corresponds to $\alpha\in\C$ again, show that the dual $\widehat f$ of (Ex.~4.7) corresponds to the complex conjugate $\bar\alpha$.
\item If $\tau\in\Q(\sqrt{-d})$ happens to be integral over $\Z$, show that $R=\Z[\tau]$.
\end{enumerate}}
\pf
Let $\Lambda=\Z+\Z\tau$ (with $\tau\notin\R$), so $X\cong\C/\Lambda$ as a group and, by Prop.~4.18 and its proof, $R\to\C$, $f\mapsto\alpha_f$, is an injective ring homomorphism whose image is $\{\alpha\in\C:\alpha\Lambda\subseteq\Lambda\}$; under it $n_X\mapsto n$ (Remark 4.18.1).

\textbf{(a)} Let $f\neq0$, so $\alpha=\alpha_f\neq0$. In characteristic $0$ the morphism $f$ is finite and separable, and it is a group homomorphism, so $\deg f=\#\ker f$ (as in the proof of Cor.~4.17). Now
\[
\ker f=\{z\in\C/\Lambda:\alpha z\in\Lambda\}=\alpha^{-1}\Lambda/\Lambda,
\]
and multiplication by $\alpha$ is an isomorphism $\alpha^{-1}\Lambda/\Lambda\to\Lambda/\alpha\Lambda$, so $\deg f=[\Lambda:\alpha\Lambda]$. Write $\Lambda=\Z\omega_1\oplus\Z\omega_2$ and $\alpha\omega_i=\sum_jm_{ij}\omega_j$ with $M=(m_{ij})\in M_2(\Z)$; the sublattice $\alpha\Lambda$ has index $|\det M|$. Since $\omega_1,\omega_2$ are also an $\R$-basis of $\C$ ($\tau\notin\R$), $M$ is the matrix of the $\R$-linear map ``multiplication by $\alpha$'' of $\C=\R^2$, whose determinant is $|\alpha|^2$ (in the basis $1,i$ it is $\begin{pmatrix}a&b\\-b&a\end{pmatrix}$ for $\alpha=a+bi$, and determinants do not depend on the basis). Hence $\deg f=|\alpha|^2$.

\textbf{(b)} The dual $\widehat f$ is again in $R$ (Ex.~4.7(a), with $X'=X$, $P_0'=P_0$). If $f=0$ then $f^*$ sends every degree $0$ class to the trivial class, so $\widehat f=0$ and the claim is trivial. Let $f\ne0$, $n=\deg f=|\alpha|^2=\alpha\bar\alpha$ by (a). By Ex.~4.7(c), $\widehat f\circ f=n_X$. Applying the injective ring homomorphism $R\to\C$ and writing $\beta=\alpha_{\widehat f}$, we get $\beta\alpha=n=\alpha\bar\alpha$. Since $\alpha\ne0$, $\beta=\bar\alpha$.

\textbf{(c)} Write $\tau=r+s\sqrt{-d}$ with $r,s\in\Q$ ($s\neq0$ since $\tau\notin\R$). Since $\tau$ is integral over $\Z$, $\tau\cdot1=\tau\in\Lambda$ and $\tau\cdot\tau=\tau^2$; the minimal polynomial of $\tau$ over $\Q$ is $x^2-2rx+(r^2+ds^2)$, which lies in $\Z[x]$ because $\Z$ is integrally closed. So $\tau^2=2r\tau-(r^2+ds^2)\in\Z+\Z\tau=\Lambda$, i.e.\ $\tau\Lambda\subseteq\Lambda$, i.e.\ $\tau\in R$ (identifying $R$ with its image in $\C$). Hence $\Z[\tau]=\Z+\Z\tau=\Lambda\subseteq R$. Conversely every $\alpha\in R$ satisfies $\alpha=\alpha\cdot1\in\alpha\Lambda\subseteq\Lambda$. Therefore $R=\Lambda=\Z[\tau]$. (This is consistent with Thm.~4.19: here $2r\in\Z$ and $r^2+ds^2\in\Z$, so the conditions $2br,\,b(r^2+ds^2)\in\Z$ hold for every $b\in\Z$.) \hfill$\square$

\exercise{4.12}{Again let $X$ be an elliptic curve over $\C$ determined by the elliptic functions with periods $1,\tau$, and assume that $\tau$ lies in the region $G$ of (4.15B).
\begin{enumerate}[label=(\alph*),leftmargin=*]
\item If $X$ has any automorphisms leaving $P_0$ fixed other than $\pm1$, show that either $\tau=i$ or $\tau=\omega$, as in (4.20.1), (4.20.2). This gives another proof of the fact (4.7) that there are only two curves, up to isomorphism, having automorphisms other than $\pm1$.
\item Now show that there are exactly three values of $\tau$ for which $X$ admits an endomorphism of degree $2$. Can you match them with the three values of $j$ determined in (Ex.~4.5)? [Answers: $\tau=i$; $\tau=\sqrt{-2}$; $\tau=\frac12(-1+\sqrt{-7})$.]
\end{enumerate}}
\pf
Throughout $\Lambda=\Z+\Z\tau$, $\omega=\frac12(-1+\sqrt{-3})$ (so $\omega^3=1$), and $\operatorname{Im}\tau>0$ for $\tau\in G$ (Fig.~16). Both $i$ and $\omega$ lie in $G$ ($\operatorname{Re}\le0$ and absolute value $1$). We identify $R=\operatorname{End}(X,P_0)$ with a subring of $\C$ as in Prop.~4.18 (its image is $\{\alpha:\alpha\Lambda\subseteq\Lambda\}$), and use Thm.~4.15B as follows: if $\Lambda=\gamma(\Z+\Z\theta)$ for some $\gamma\in\C^*$ and $\theta$ with $\operatorname{Im}\theta>0$, then expressing the basis $(1,\tau)$ and $(\gamma,\gamma\theta)$ of $\Lambda$ through one another we find $\theta=(a\tau+b)/(c\tau+d)$ with $\begin{pmatrix}a&b\\c&d\end{pmatrix}\in GL_2(\Z)$, and $\operatorname{Im}\theta=(ad-bc)\operatorname{Im}\tau/|c\tau+d|^2>0$ forces $ad-bc=1$, so $J(\tau)=J(\theta)$ by (4.15B), and, if $\theta\in G$ too, the uniqueness statement of (4.15B) gives $\tau=\theta$. $(\ast)$

\textbf{A lemma on lattices.} Let $O$ be one of the rings $\Z[i]$, $\Z[\omega]$, $\Z[\sqrt{-2}]$, $\Z[\frac{1+\sqrt{-7}}2]$, and let $\Lambda\subset\C$ be a lattice with $O\Lambda\subseteq\Lambda$. Then $\Lambda=\gamma O$ for a shortest nonzero vector $\gamma$ of $\Lambda$. Indeed, every $z\in\C$ is at distance $<1$ from $O$: the lattice $O$ is tiled by a square of side $1$, an equilateral triangle of side $1$, a rectangle $1\times\sqrt2$, resp.\ acute triangles with sides $1,\sqrt2,\sqrt2$, and the largest distance from a vertex is at most $\frac{\sqrt2}2,\frac1{\sqrt3},\frac{\sqrt3}2,\frac2{\sqrt7}$ respectively (half diagonal, resp.\ circumradius of an acute triangle), all $<1$. Now let $v\in\Lambda$. Choose $\beta\in O$ with $|v/\gamma-\beta|<1$. Then $v-\beta\gamma\in\Lambda$ (as $\beta\gamma\in O\Lambda\subseteq\Lambda$) and $|v-\beta\gamma|<|\gamma|$, so $v=\beta\gamma$ by minimality of $|\gamma|$. Thus $\Lambda=O\gamma$.

\textbf{(a)} Let $\sigma\neq\pm1$ be an automorphism of $X$ with $\sigma(P_0)=P_0$. By Lemma~4.9 it is a homomorphism, so $\sigma\in R$, and $\sigma^{-1}\in R$, i.e.\ $\sigma$ is a unit of $R$. Let $\alpha\in\C$ correspond to it; then $\alpha\ne\pm1$, $\alpha^{-1}\in R$, and by Ex.~4.11(a), $|\alpha|^2=\deg\sigma=1$. As $|\alpha|=1$ and $\alpha\neq\pm1$, $\alpha\notin\R$. Since $R\cap\R=\Z$ (proof of (4.19)), $R\neq\Z$, so $X$ has complex multiplication and by Thm.~4.19, $R$ is a subring of the ring of integers $\mathcal O_K$ of an imaginary quadratic field $K$. So $\alpha$ is an algebraic integer in $K$, $t=\alpha+\bar\alpha\in\Z$ (its trace), $\alpha\bar\alpha=1$, and $|t|=2|\operatorname{Re}\alpha|\le2$. Thus $\alpha$ is a root of $x^2-tx+1$ with $t\in\{0,\pm1\}$ (for $t=\pm2$ we would get $\alpha=\pm1$), i.e.\ $\alpha\in\{\pm i,\ \omega,\omega^2,\ -\omega,-\omega^2\}$. If $\alpha=\pm i$ then $i\in R$, and if $\alpha\in\{\omega,\omega^2,-\omega,-\omega^2\}$ then $\omega\in\{\alpha,\alpha^2,\alpha^4\}\subseteq R$ (indeed $\alpha=\omega$; $\alpha=\omega^2\Rightarrow\alpha^2=\omega$; $\alpha=-\omega\Rightarrow\alpha^4=\omega^4=\omega$; $\alpha=-\omega^2\Rightarrow\alpha^2=\omega$). So $\Lambda$ is stable under $\Z[i]$ or under $\Z[\omega]$, and by the Lemma $\Lambda=\gamma\Z[i]=\gamma(\Z+\Z i)$, resp.\ $\Lambda=\gamma(\Z+\Z\omega)$. By $(\ast)$ and $i,\omega\in G$ we obtain $\tau=i$, resp.\ $\tau=\omega$.

Conversely, for $\tau=i$, $i\Lambda=\Lambda$ and for $\tau=\omega$, $\omega\Lambda=\Lambda$, so these two curves do have the automorphisms $\pm i$, resp.\ $\pm\omega,\pm\omega^2$ (Examples 4.20.1, 4.20.2). Every elliptic curve over $\C$ is isomorphic to some $X_\tau$ with $\tau\in G$ (4.14B, 4.15B), and $X_\tau\cong X_{\tau'}$ with $\tau,\tau'\in G$ forces $\tau=\tau'$ (by $J$ and 4.15B, 4.1(b)). Hence, up to isomorphism, there are exactly two elliptic curves over $\C$ with an automorphism fixing $P_0$ other than $\pm1$: $X_i$ ($j=1728$) and $X_\omega$ ($j=0$). This reproves Cor.~4.7.

\textbf{(b)} \emph{Which $\tau$.} If $f\in R$ has degree $2$, with $\alpha\in\C$ corresponding to it, then $|\alpha|^2=2$ by Ex.~4.11(a). Since $\pm\sqrt2\notin\Z=R\cap\R$, $\alpha\notin\R$, so $X$ has complex multiplication and $\alpha$ is an algebraic integer in an imaginary quadratic field (Thm.~4.19). Then $t=\alpha+\bar\alpha\in\Z$, $\alpha\bar\alpha=2$, $t^2<8$, so $t\in\{0,\pm1,\pm2\}$ and $\alpha$ is a root of $x^2-tx+2$:
\[
\alpha=\pm\sqrt{-2}\quad(t=0),\qquad \alpha=\tfrac12(\pm1\pm\sqrt{-7})\quad(t=\pm1),\qquad\alpha=\pm1\pm i\quad(t=\pm2).
\]
The subring $\Z[\alpha]$ is, respectively, $O=\Z[\sqrt{-2}]$, $O=\Z[\frac{1+\sqrt{-7}}2]$ (note $\frac{-1+\sqrt{-7}}2=\frac{1+\sqrt{-7}}2-1$), $O=\Z[i]$. As $\alpha\Lambda\subseteq\Lambda$ we have $O\Lambda\subseteq\Lambda$, so by the Lemma $\Lambda=\gamma O=\gamma(\Z+\Z\theta)$ with $\theta=\sqrt{-2}$, $\frac{1+\sqrt{-7}}2$, $i$ respectively. By $(\ast)$ (note $\Z+\Z\frac{1+\sqrt{-7}}2=\Z+\Z\frac{-1+\sqrt{-7}}2$) we get $\tau=\theta'$, where $\theta'=\sqrt{-2}$, $\frac{-1+\sqrt{-7}}2$, $i$, which all lie in $G$ ($|\theta'|^2=2,2,1$ and $\operatorname{Re}\theta'=0,-\frac12,0$).

\emph{Conversely}, for these three $\tau$ there is an endomorphism of degree $2$: for $\tau=i$, $\alpha=1+i$ satisfies $\alpha\Lambda\subseteq\Lambda$; for $\tau=\sqrt{-2}$, $\alpha=\tau$ ($\tau\cdot\tau=-2\cdot1$); for $\tau=\frac{-1+\sqrt{-7}}2$, $\alpha=\tau$ since $\tau^2+\tau+2=0$ gives $\tau\cdot\tau=-\tau-2\in\Lambda$. In each case $|\alpha|^2=2$, so by Prop.~4.18 and Ex.~4.11(a) there is an endomorphism of degree $2$. So there are exactly three values of $\tau\in G$: $i$, $\sqrt{-2}$, $\frac12(-1+\sqrt{-7})$.

\emph{Matching with $j$.} By (Ex.~4.5(d)) the curves with an endomorphism of degree $2$ have $j\in\{2^6\cdot3^3,\ 2^6\cdot5^3,\ -3^3\cdot5^3\}=\{1728,8000,-3375\}$. The three $\tau$'s give three pairwise non-isomorphic curves (4.15B, 4.1(b)), so $j$ is injective on them, and since each of the three curves has $j$ in the above set, the matching is a bijection. For $\tau=i$ we have $j=1728$ (Example 4.20.1). For the other two we use the real structure. Let $g_2,g_3$ be those of $\Lambda$ and $e_1,e_2,e_3$ the values of $\wp$ at the three nonzero half periods $\frac12,\frac\tau2,\frac{1+\tau}2$ (distinct, the roots of $4x^3-g_2x-g_3$), so $\Delta=g_2^3-27g_3^2=16\prod_{i<j}(e_i-e_j)^2$ and $j=1728g_2^3/\Delta$, $j-1728=1728\cdot27\,g_3^2/\Delta$. Since $\bar\Lambda=\Lambda$ (for $\tau=\sqrt{-2}$, $\bar\tau=-\tau$; for $\tau=\frac{-1+\sqrt{-7}}2$, $\bar\tau=-1-\tau$), we have $\wp(\bar z)=\overline{\wp(z)}$, so $g_2,g_3$ are real and complex conjugation permutes the $e_i$ according to its action on the half periods mod $\Lambda$.
\begin{itemize}
\item $\tau=\sqrt{-2}$: conjugation fixes each of $\frac12,\frac\tau2,\frac{1+\tau}2$ mod $\Lambda$ (e.g.\ $\overline{(1+\tau)/2}=(1-\tau)/2\equiv(1+\tau)/2$), so all $e_i$ are real and distinct, $\Delta>0$, hence $j-1728\ge0$.
\item $\tau=\frac{-1+\sqrt{-7}}2$: $\overline{\tau/2}=\frac{-1-\tau}2\equiv\frac{1+\tau}2$, so $e_2=\wp(\tau/2)$ and $e_3=\wp(\frac{1+\tau}2)$ are complex conjugate, and not real since $e_2\ne e_3$, while $e_1$ is real. Then $(e_1-e_2)(e_1-e_3)=|e_1-e_2|^2>0$ and $(e_2-e_3)^2<0$ ($e_2-e_3$ is purely imaginary and nonzero), so $\Delta<0$, hence $j-1728\le0$.
\end{itemize}
As $j\in\{1728,8000,-3375\}$ and $j\neq1728$ for these two curves (they are not isomorphic to $X_i$), the first gives $j=8000$ and the second $j=-3375$. (Numerical confirmation: from $j=q^{-1}+744+196884q+\cdots$, $q=e^{2\pi i\tau}$, one gets $j(\sqrt{-2})\approx8000$ and $j(\frac{-1+\sqrt{-7}}2)\approx-3375$; I also checked this with the Eisenstein series to high precision.)
Hence
\[
\tau=i\leftrightarrow j=2^6\cdot3^3=1728,\qquad \tau=\sqrt{-2}\leftrightarrow j=2^6\cdot5^3=8000,\qquad \tau=\tfrac12(-1+\sqrt{-7})\leftrightarrow j=-3^3\cdot5^3=-3375.
\]
\hfill$\square$

\exercise{4.13}{If $p=13$, there is just one value of $j$ for which the Hasse invariant of the corresponding curve is $0$. Find it. [Answer: $j=5\pmod{13}$.]}
\pf
Let $k=\overline{\mathbf F}_{13}$. Every elliptic curve over $k$ is isomorphic to some $X_\lambda:\ y^2=x(x-1)(x-\lambda)$ with $\lambda\neq0,1$ (Prop.~4.6), with $j=j(\lambda)=2^8(\lambda^2-\lambda+1)^3/(\lambda^2(\lambda-1)^2)$, and by Cor.~4.22 $X_\lambda$ has Hasse invariant $0$ iff $h(\lambda)=0$, where for $p=13$, $k'=6$,
\[
h(\lambda)=\sum_{i=0}^{6}\binom6i^2\lambda^i=1+36\lambda+225\lambda^2+400\lambda^3+225\lambda^4+36\lambda^5+\lambda^6\equiv\lambda^6+10\lambda^5+4\lambda^4+10\lambda^3+4\lambda^2+10\lambda+1\pmod{13}.
\]
Consider the polynomial $g(\lambda)=2^8(\lambda^2-\lambda+1)^3-5\lambda^2(\lambda-1)^2$, so that for $\lambda\ne0,1$, $j(\lambda)=5\iff g(\lambda)=0$. Expanding, $(\lambda^2-\lambda+1)^3=\lambda^6-3\lambda^5+6\lambda^4-7\lambda^3+6\lambda^2-3\lambda+1$ and $5\lambda^2(\lambda-1)^2=5\lambda^4-10\lambda^3+5\lambda^2$, so
\[
g=256\lambda^6-768\lambda^5+1531\lambda^4-1782\lambda^3+1531\lambda^2-768\lambda+256\equiv9\lambda^6-\lambda^5+10\lambda^4-\lambda^3+10\lambda^2-\lambda+9\pmod{13},
\]
using $256\equiv9$, $768\equiv1$, $1531\equiv10$, $1782\equiv1$. On the other hand $-4h\equiv-4\lambda^6-40\lambda^5-16\lambda^4-40\lambda^3-16\lambda^2-40\lambda-4\equiv9\lambda^6-\lambda^5+10\lambda^4-\lambda^3+10\lambda^2-\lambda+9$. Hence $g\equiv-4h\pmod{13}$, i.e.\ $h\equiv3g$. Neither polynomial vanishes at $\lambda=0$ ($h(0)=1$, $g(0)=9$) or at $\lambda=1$ ($h(1)=\binom{12}6=924\equiv1$, $g(1)=256\equiv9$). Therefore for $\lambda\in k\setminus\{0,1\}$:
\[
h(\lambda)=0\iff g(\lambda)=0\iff j(\lambda)=5.
\]
Since $h$ has degree $6$ it has roots in $k$, all different from $0,1$, so there exist curves with Hasse invariant $0$, and by the equivalence they are exactly the curves with $j=5$. So the unique value is $j=5\pmod{13}$. (Mod $13$ one finds $h=(\lambda^2+4\lambda-4)(\lambda^2-6\lambda+1)(\lambda^2-\lambda+3)$, whose six roots lie in $\mathbf F_{169}$ and form a single $\Sigma_3$-orbit, consistent with $j$ having degree $6$ in $\lambda$ and with Cor.~4.23: $[13/12]+2=3\ge1$.) \hfill$\square$

\exercise{4.14}{The Fermat curve $X:x^3+y^3=z^3$ gives a nonsingular curve in characteristic $p$ for every $p\neq3$. Determine the set $\mathfrak P=\{p\neq3\mid X_{(p)}\text{ has Hasse invariant }0\}$, and observe (modulo Dirichlet's theorem) that it is a set of primes of density $\frac12$.}
\pf
Let $f=x^3+y^3-z^3$. In characteristic $p\ne3$ the partial derivatives $3x^2,3y^2,-3z^2$ vanish simultaneously only at $x=y=z=0$, so $X_{(p)}$ is nonsingular (also for $p=2$). Proposition 4.21 applies to any nonsingular plane cubic (its proof does not use $p\ne2$): $X_{(p)}$ has Hasse invariant $0$ iff the coefficient of $(xyz)^{p-1}$ in $f^{p-1}$ is $0$ in $k$. By the multinomial theorem
\[
f^{p-1}=\sum_{a+b+c=p-1}\frac{(p-1)!}{a!\,b!\,c!}(-1)^c\,x^{3a}y^{3b}z^{3c}.
\]
The monomial $(xyz)^{p-1}$ occurs iff $3a=3b=3c=p-1$, which is possible iff $p\equiv1\pmod3$, with $a=b=c=\frac{p-1}3$. In that case the coefficient is $(-1)^{(p-1)/3}(p-1)!/((\frac{p-1}3)!)^3$, a quotient of products of integers all smaller than $p$, hence a unit mod $p$: the Hasse invariant is $1$. If $p\not\equiv1\pmod 3$ (and $p\ne3$), i.e.\ $p\equiv2\pmod3$, the monomial does not occur, the coefficient is $0$ and the Hasse invariant is $0$. (For $p=2$: $f^1$ has no $xyz$ term, so $X_{(2)}$ is supersingular.) Therefore
\[
\mathfrak P=\{p\text{ prime}:p\equiv2\pmod3\}.
\]
By Dirichlet's theorem on primes in arithmetic progressions the primes $\equiv1$ and $\equiv2\pmod3$ each have density $\frac1{\varphi(3)}=\frac12$, so $\mathfrak P$ has density $\frac12$. (This matches Remark 4.23.4 and Example 4.20.2: $X$ has $j=0$ and complex multiplication by $\Z[\omega]$, and $p\equiv2\pmod3$ iff $p$ stays prime in $\Z[\omega]$; it also matches the Note after Cor.~4.23: $j=0$ has Hasse invariant $0$ iff $p\equiv2\pmod3$.) \hfill$\square$

\exercise{4.15}{Let $X$ be an elliptic curve over a field $k$ of characteristic $p$. Let $F':X_p\to X$ be the $k$-linear Frobenius morphism (2.4.1). Use (4.10) to show that the dual morphism $\widehat{F'}:X\to X_p$ is separable if and only if the Hasse invariant of $X$ is $1$. Now use (Ex.~4.7) to show that if the Hasse invariant is $1$, then the subgroup of points of order $p$ on $X$ is isomorphic to $\Z/p$; if it is $0$, it is $0$.}
\pf
($k$ is algebraically closed.) $X_p$ is the same scheme as $X$ with the structure morphism composed with Frobenius of $k$; since $k$ is perfect it is again an elliptic curve, with base point $P_0$ (the same point), and $F'$ is the identity on underlying topological spaces and $F'^\#(s)=s^p$; it has degree $p$ (Example 2.4.3), so by Ex.~4.7(f) $\deg\widehat{F'}=p$.

\textbf{Preliminaries on tangent maps (used also in 4.16 and 4.20).} For a homomorphism $g:Y\to Z$ of elliptic curves let $dg:T_0Y\to T_0Z$ be the map on Zariski tangent spaces at the origins (both $1$-dimensional).

(T1) $d(g+h)=dg+dh$, $d(g\circ h)=dg\circ dh$, $d(n_Y)=n\cdot\mathrm{id}$. Indeed $g+h=\mu_Z\circ(g,h)$ and $d\mu_Z(u,v)=u+v$ on $T_{(0,0)}(Z\times Z)=T_0Z\oplus T_0Z$ because $\mu_Z(z,0)=z=\mu_Z(0,z)$. Also $d\rho=-\mathrm{id}$ from $\mu\circ(1,\rho)=$ const; induction gives $d(n_Y)=n$.

(T2) A nonconstant homomorphism $g:Y\to Z$ is separable iff $dg\neq0$. ($\Rightarrow$) If $g$ is separable, it is ramified at only finitely many points (Prop.~2.2(a)). Since $g\circ\tau_y=\tau_{g(y)}\circ g$ for translations $\tau_y$ (automorphisms, Prop.~4.8), the ramification index of $g$ is the same at all points; so $g$ is unramified at $0$, i.e.\ $g^*$ maps a uniformizer at $0\in Z$ to a uniformizer at $0\in Y$, i.e.\ $dg\neq0$. ($\Leftarrow$) If $g$ is inseparable, write $g=s\circ h$ with $s$ separable and $h:Y\to Y'$ purely inseparable of degree $p^r$, $r\ge1$. By Prop.~2.5, $h$ is a composition of $k$-linear Frobenius morphisms and isomorphisms; a $k$-linear Frobenius sends a local parameter $u$ to $u^p$, so $h^*(\mathfrak m_{h(0)})\subseteq\mathfrak m_0^p\subseteq\mathfrak m_0^2$, i.e.\ $dh=0$ and $dg=ds\circ dh=0$.

(T3) Let $f:Y\to Z$ be a homomorphism and $\widehat f:Z\to Y$ its dual (Ex.~4.7(a)), induced by $f^*:\operatorname{Pic}^0Z\to\operatorname{Pic}^0Y$ (under $X\cong J$, Thm.~4.11). By Remark 4.10.7 the tangent space of $J$ at $0$ is $H^1(\mathcal O)$, tangent vectors being the classes $1+\varepsilon a$, $a\in H^1(\mathcal O)$, in $\operatorname{Pic}(\,\cdot\,\times\operatorname{Spec}k[\varepsilon]/\varepsilon^2)$; $f^*$ sends $1+\varepsilon a$ to $1+\varepsilon f^*a$. Hence $d\widehat f=f^*:H^1(Z,\mathcal O_Z)\to H^1(Y,\mathcal O_Y)$.

\textbf{First assertion.} Apply (T3) to $f=F':X_p\to X$: the tangent map of $\widehat{F'}:X\to X_p$ is $F'^*:H^1(X,\mathcal O_X)\to H^1(X_p,\mathcal O_{X_p})$. Since $F'$ and the Frobenius $F$ of $X$ are the same morphism of ringed spaces, this map is, as a map of abelian groups, $F^*:H^1(X,\mathcal O_X)\to H^1(X,\mathcal O_X)$. By definition the Hasse invariant is $1$ iff $F^*\neq0$. So $d\widehat{F'}\ne0$ iff Hasse $=1$, and by (T2) $\widehat{F'}$ is separable iff Hasse $=1$.

\textbf{Points of order $p$.} Ex.~4.7(c) with $f=F'$ gives $F'\circ\widehat{F'}=p_X$. Since $F'$ is the identity on points, $F'(y)=P_0$ iff $y=P_0$. Hence for a closed point $x$: $px=0\iff F'(\widehat{F'}x)=P_0\iff\widehat{F'}(x)=P_0$, so $X[p]=\ker\widehat{F'}$ as groups of closed points. The degree of $\widehat{F'}$ is $p$. If it is separable ($\iff$ Hasse $=1$) then $\#\ker\widehat{F'}=\deg\widehat{F'}=p$ (fibres of a homomorphism are cosets of the kernel), so $X[p]$ is a group of order $p$, i.e.\ $\Z/p$. If it is inseparable ($\iff$ Hasse $=0$) then, $p$ being prime, it is purely inseparable, hence injective on points, so $X[p]=0$. \hfill$\square$

\exercise{4.16}{Again let $X$ be an elliptic curve over $k$ of characteristic $p$, and suppose $X$ is defined over the field $\mathbf F_q$ of $q=p^r$ elements, i.e., $X\subseteq\mathbf P^2$ can be defined by an equation with coefficients in $\mathbf F_q$. Assume also that $X$ has a rational point over $\mathbf F_q$. Let $F':X_q\to X$ be the $k$-linear Frobenius with respect to $q$.
\begin{enumerate}[label=(\alph*),leftmargin=*]
\item Show that $X_q\cong X$ as schemes over $k$, and that under this identification, $F':X\to X$ is the map obtained by the $q$th-power map on the coordinates of points of $X$, embedded in $\mathbf P^2$.
\item Show that $1_X-F'$ is a separable morphism and its kernel is just the set $X(\mathbf F_q)$ of points of $X$ with coordinates in $\mathbf F_q$.
\item Use (Ex.~4.7) to show that $\widehat{F'}+F'=a_X$ for some integer $a$, and that $N=q-a+1$, where $N=\#X(\mathbf F_q)$.
\item Use the fact that $\deg(m+nF')>0$ for all $m,n\in\Z$ to show that $|a|\le2\sqrt q$. This is Hasse's proof of the analogue of the Riemann hypothesis for elliptic curves (App.~C, Ex.~5.6).
\item Now assume $q=p$, and show that the Hasse invariant of $X$ is $0$ if and only if $a\equiv0\pmod p$. Conclude for $p\ge5$ that $X$ has Hasse invariant $0$ if and only if $N=p+1$.
\end{enumerate}}
\pf
Let $f\in\mathbf F_q[x,y,z]$ be an equation of $X$ and $S=k[x,y,z]/(f)$, so $X=\operatorname{Proj}S$ with base point $P_0\in X(\mathbf F_q)$. $X_q$ is the same scheme $\operatorname{Proj}S$ with $k$ acting through $a\mapsto a^q$; write $S_q$ for $S$ with this $k$-algebra structure ($a*s=a^qs$), so $X_q=\operatorname{Proj}S_q$ and $F'$ corresponds to the ring map $S\to S_q$, $s\mapsto s^q$ (it is $k$-linear: $(as)^q=a*s^q$).

\textbf{(a)} Define $\Psi:S\to S$ by $\Psi(\sum c_Ix^I)=\sum c_I^qx^I$ ($c_I\in k$). It is well defined on $S$ because $f$ has coefficients in $\mathbf F_q$ (so $\Psi(f)=f$), and it is a graded ring automorphism ($k$ is perfect). Since $\Psi(as)=a^q\Psi(s)=a*\Psi(s)$, $\Psi:S\to S_q$ is an isomorphism of graded $k$-algebras. Hence $\Theta:X\to X_q$, the morphism of schemes over $k$ corresponding to $\Psi^{-1}$, is an isomorphism. The composite $F'\circ\Theta:X\to X$ corresponds to the $k$-algebra map $S\to S_q\to S$, $s\mapsto s^q\mapsto\Psi^{-1}(s^q)$, which sends $x_i\mapsto\Psi^{-1}(x_i^q)=x_i^q$; so $F'\circ\Theta$ is the morphism $\varphi:(x:y:z)\mapsto(x^q:y^q:z^q)$. Write $\varphi=F'\circ\Theta$. As $F'$ is the identity on points, $\Theta$ is the map of points $P\mapsto\varphi(P)$; in particular $\Theta(P_0)=\varphi(P_0)=P_0$ since $P_0$ has coordinates in $\mathbf F_q$.

Below we write $\varphi\in R=\operatorname{End}(X,P_0)$ for this endomorphism (it is a homomorphism by Lemma 4.9 as $\varphi(P_0)=P_0$); "$F'$" in the statements (b)--(e) means $\varphi$ under the identification $\Theta$. Its degree is $q$ (that of $F'$; $\Theta$ is an isomorphism).

\textbf{(b)} A closed point $P=(x:y:z)$ satisfies $\varphi(P)=P$ iff $(x^q:y^q:z^q)=(x:y:z)$, i.e.\ $x_i^q=\lambda x_i$ for all $i$ and some $\lambda\ne0$; normalizing a nonzero coordinate to $1$ gives $\lambda=1$, so all $x_i^q=x_i$, i.e.\ $x_i\in\mathbf F_q$. Hence $\ker(1_X-\varphi)=X(\mathbf F_q)$ (as a set of closed points). Now $d\varphi=0$ since $\varphi^*$ sends a local parameter $u$ to $u^q\in\mathfrak m^2$ (or by (T2), $\varphi$ being purely inseparable). By (T1) $d(1_X-\varphi)=\mathrm{id}-d\varphi=\mathrm{id}\neq0$. In particular $1_X-\varphi$ is nonconstant, and by (T2) it is separable. As a separable homomorphism, $\deg(1_X-\varphi)=\#\ker(1_X-\varphi)=\#X(\mathbf F_q)=N$.

\textbf{(c)} Recall Ex.~4.7: $\widehat{\ }$ is additive (d), contravariantly multiplicative (b), $g\circ\widehat g=\widehat g\circ g=[\deg g]$ (c) for $g\in R$ (trivially also for $g=0$, whose dual is $0$), $\widehat{n_X}=n_X$ (e), $\deg\widehat g=\deg g$ (f), and $\widehat{1_X}=1_X$. Moreover integers are central in $R$, since $n_X\circ g=g\circ n_X$ for a homomorphism $g$. Dualizing the additive element $1-\varphi$: $\widehat{1-\varphi}=1-\widehat\varphi$ (from $\varphi+(-\varphi)=0$ one gets $\widehat{-\varphi}=-\widehat\varphi$). By (b) and Ex.~4.7(c),
\[
N=\deg(1-\varphi)\quad\text{and}\quad(1-\varphi)(1-\widehat\varphi)=[N].
\]
Expanding, $(1-\varphi)(1-\widehat\varphi)=1-\varphi-\widehat\varphi+\varphi\widehat\varphi=1-(\varphi+\widehat\varphi)+q$, so
\[
\varphi+\widehat\varphi=[q+1-N]=a_X\quad\text{with }a=q+1-N\in\Z,\qquad\text{i.e. }N=q-a+1,
\]
using that $\Z\to R$, $n\mapsto n_X$, is injective ($\deg n_X=n^2>0$ for $n\neq0$ by Ex.~4.7(e)).

\textbf{(d)} For $m,n\in\Z$ put $g=m+n\varphi\in R$. Then $\widehat g=m+n\widehat\varphi$ and
\[
[\deg g]=g\widehat g=m^2+mn(\varphi+\widehat\varphi)+n^2\varphi\widehat\varphi=[\,m^2+a\,mn+q\,n^2\,].
\]
So $\deg(m+n\varphi)=m^2+amn+qn^2$ (injectivity of $\Z\to R$ again). A degree is $\ge0$ (it is $>0$ unless $m+n\varphi$ is constant). Taking $(m,n)=(-a,2)$ gives $a^2-2a^2+4q=4q-a^2\ge0$, i.e.\ $|a|\le2\sqrt q$. (Remark: the printed ``$>0$ for all $m,n$'' must exclude $(0,0)$, and also fails when $m+n\varphi=0$ for some $(m,n)\neq(0,0)$, which can happen for supersingular curves over $\mathbf F_{p^{2s}}$ with $\varphi=\pm p^s$; in that case $n\varphi=-m$, $n\widehat\varphi=-m$ give $na=-2m$, $n^2q=m^2$, so $a^2=4q$ and the bound still holds, with equality. The argument with $\ge0$ covers both cases.)

\textbf{(e)} Let $q=p$, so $X_q=X_p$ and $F'\Theta=\varphi$. Then $\widehat\varphi=\widehat\Theta\circ\widehat{F'}$ by Ex.~4.7(b), and $\widehat\Theta:X_p\to X$ is an isomorphism ($\widehat\Theta\Theta=1$). Hence $\widehat\varphi$ is separable iff $\widehat{F'}$ is, iff (Ex.~4.15) the Hasse invariant is $1$. By (T1) and (c), $d\widehat\varphi=d(a_X)-d\varphi=a\cdot\mathrm{id}-0=a\cdot\mathrm{id}$, which is the multiplication by the image of $a$ in $\mathbf F_p\subseteq k$. By (T2), $\widehat\varphi$ is separable iff $a\not\equiv0\pmod p$. Hence: Hasse invariant $0\iff a\equiv0\pmod p$.

For $p\ge5$, (d) gives $|a|\le2\sqrt p<p$ (as $p>4$), so $a\equiv0\pmod p$ iff $a=0$, iff $N=p-a+1=p+1$. Thus for $p\ge5$: Hasse invariant $0\iff N=p+1$. (For $p=2,3$ this fails to be forced: $|a|=p$ is possible.) \hfill$\square$

\exercise{4.17}{Let $X$ be the curve $y^2+y=x^3-x$ of (4.23.8).
\begin{enumerate}[label=(\alph*),leftmargin=*]
\item If $Q=(a,b)$ is a point on the curve, compute the coordinates of the point $P+Q$, where $P=(0,0)$, as a function of $a,b$. Use this formula to find the coordinates of $nP$, $n=1,2,\ldots,10$. [Check: $6P=(6,14)$.]
\item This equation defines a nonsingular curve over $\mathbf F_p$ for all $p\neq37$.
\end{enumerate}}
\pf
The curve in $\mathbf P^2$ is $y^2z+yz^2=x^3-xz^2$; the line $z=0$ meets it in $x^3=0$, i.e.\ in $3P_0$ with $P_0=(0,1,0)$, which is the origin of the group law (4.23.8). So (embedding by $|3P_0|$) three points are collinear iff they sum to $0$. For a point $(x,y)$ the vertical line $x=$ const meets $X$ again at $(x,-1-y)$ (the sum of the two roots of $y^2+y-c$ is $-1$), so $-(x,y)=(x,-1-y)$.

\textbf{(a)} Let $Q=(a,b)$ with $a\neq0$ (the points with $x=0$ are only $P=(0,0)$ and $-P=(0,-1)$). The line $PQ$ is $y=mx$, $m=b/a$. Substituting: $m^2x^2+mx=x^3-x$, i.e.\ $x\,(x^2-m^2x-(m+1))=0$, whose roots are $0$, $a$ and a third value $x_3$ with $a+x_3=m^2$. So the third point of intersection is $R=(x_3,mx_3)$ with $x_3=\frac{b^2}{a^2}-a=\frac{b^2-a^3}{a^2}$. From $b^2+b=a^3-a$ we get $b^2-a^3=-a-b$, so $x_3=-\frac{a+b}{a^2}$. Then $P+Q=-R=(x_3,-1-mx_3)$, with $-1-mx_3=-1+\frac{b(a+b)}{a^3}=\frac{b^2+ab-a^3}{a^3}=\frac{ab-a-b}{a^3}$. Hence
\[
\boxed{P+Q=\Bigl(-\frac{a+b}{a^2},\ \frac{ab-a-b}{a^3}\Bigr)\qquad(a\ne0).}
\]
For $Q=P$ we compute the tangent at $P$: $(2y+1)\,dy=(3x^2-1)\,dx$ gives slope $-1$, the line $y=-x$ gives $x^2-x=x^3-x$, i.e.\ $x^2(x-1)=0$, so the third point is $(1,-1)$ and $2P=(1,0)$. Also $P+(-P)=0$.

Now $nP=P+(n-1)P$; for $n\ge3$ the formula applies as long as $x((n-1)P)\neq0$, which we see from the values themselves:
\[
\begin{array}{lll}
1P=(0,0), & 2P=(1,0), & 3P=(-1,-1),\\
4P=(2,-3), & 5P=(\frac14,-\frac58), & 6P=(6,14),\\
7P=(-\frac59,\frac8{27}), & 8P=(\frac{21}{25},-\frac{69}{125}), & 9P=(-\frac{20}{49},-\frac{435}{343}),\\
10P=(\frac{161}{16},-\frac{2065}{64}). & &
\end{array}
\]
Samples: $3P=P+2P$ with $(a,b)=(1,0)$: $(-1,\ (0-1-0)/1)=(-1,-1)$. $4P$ from $(a,b)=(-1,-1)$: $(-(-2)/1,\ (1+1+1)/(-1))=(2,-3)$. $5P$ from $(2,-3)$: $(1/4,\ (-6-2+3)/8)=(\frac14,-\frac58)$. $6P$ from $(\frac14,-\frac58)$: $a+b=-\frac38$, $a^2=\frac1{16}$, so $x=6$; $ab-a-b=-\frac5{32}+\frac38=\frac7{32}$, $a^3=\frac1{64}$, so $y=14$ ($6P=(6,14)$ as in the check; indeed $196+14=210=216-6$). $7P$ from $(6,14)$: $(-\frac{20}{36},\frac{84-20}{216})=(-\frac59,\frac8{27})$; and the remaining ones are obtained in the same way (each result lies on the curve, which is a further check; e.g.\ for $10P$, $\frac{2065^2}{64^2}-\frac{2065}{64}=\frac{4264225-132160}{4096}=\frac{4132065}{4096}=\left(\frac{161}{16}\right)^3-\frac{161}{16}$). (All of this was also verified with exact rational arithmetic.) None of $x(nP)$, $n=2,\ldots,9$, is $0$, which justifies using the formula at each step.

\textbf{(b)} Put $f=y^2+y-x^3+x$. Affine points: a singular point satisfies $f_y=2y+1=0$, $f_x=-3x^2+1=0$, $f=0$. If $p=2$, $f_y=1\ne0$ always; if $p=3$, $f_x=1\neq0$ always; so there is no singular affine point for $p=2,3$. For $p\ne2,3$: $y=-\frac12$, $x^2=\frac13$, and $f=\frac14-\frac12-x(x^2-1)=-\frac14+\frac23x=0$ gives $x=\frac38$; then $x^2=\frac9{64}=\frac13$ iff $27=64$ iff $37=0$ in $\mathbf F_p$. So for $p\neq37$ there is no singular affine point. At infinity, $z=0$ forces $x^3=0$, so the only point is $(0,1,0)$, where $\partial F/\partial z=y^2+2yz+2xz=1\neq0$ for $F=y^2z+yz^2-x^3+xz^2$: nonsingular. Hence $X_{(p)}$ is nonsingular for every $p\ne37$. (For $p=37$ the point $(5,18)$ is singular: $f(5,18)=222=6\cdot37$, $f_x=-74$, $f_y=37$.) \hfill$\square$

\exercise{4.18}{Let $X$ be the curve $y^2=x^3-7x+10$. This curve has at least $26$ points with integer coordinates. Find them (use a calculator), and verify that they are all contained in the subgroup (maybe equal to all of $X(\mathbf Q)$?) generated by $P=(1,2)$ and $Q=(2,2)$.}
\pf
Take the point at infinity $P_0=(0,1,0)$ as origin (it is a flex, $z=0$ meeting $X$ in $3P_0$), so collinear points sum to $0$ and $-(x,y)=(x,-y)$. For $A=(x_1,y_1)\ne\pm B=(x_2,y_2)$, the line through them has slope $m=\frac{y_2-y_1}{x_2-x_1}$, the third intersection point has abscissa $x_3=m^2-x_1-x_2$ and ordinate $y_1+m(x_3-x_1)$, and $A+B=(x_3,-y_1-m(x_3-x_1))$; for doubling $m=\frac{3x_1^2-7}{2y_1}$.

\textbf{The integer points.} The cubic is negative for $x<-3.2$, so one searches $x\ge-3$; a computer search over $-3\le x\le10^5$ for $x^3-7x+10$ a perfect square gives (and there are no others in this range) $x\in\{-3,-2,-1,1,2,3,5,9,13,31,41,67,302\}$, each with $y=\pm$ the following: $y=2,4,4,2,2,4,10,26,46,172,262,548,5248$. So there are $26$ integer points:
\[
(-3,\pm2),\ (-2,\pm4),\ (-1,\pm4),\ (1,\pm2),\ (2,\pm2),\ (3,\pm4),\ (5,\pm10),\ (9,\pm26),\ (13,\pm46),\ (31,\pm172),\ (41,\pm262),\ (67,\pm548),\ (302,\pm5248).
\]
(Check, e.g., $302^3-7\cdot302+10=27543608-2114+10=27541504=5248^2$.)

\textbf{They lie in $\langle P,Q\rangle$.} Since $-(x,y)=(x,-y)$ it suffices to express one point of each $\pm$ pair. With $P=(1,2)$, $Q=(2,2)$:
\begin{align*}
P+Q&:\ m=0,\ x_3=-3,\ \text{third point }(-3,2),\ P+Q=(-3,-2).\\
2P&:\ m=-1,\ x_3=-1,\ \text{third }(-1,4),\ 2P=(-1,-4).\\
P-Q&=P+(2,-2):\ m=-4,\ x_3=13,\ \text{third }(13,-46),\ P-Q=(13,46).\\
3P&=2P+P:\ m=3,\ x_3=9,\ \text{third }(9,26),\ 3P=(9,-26).\\
2P+Q&=(-1,-4)+(2,2):\ m=2,\ x_3=3,\ \text{third }(3,4),\ 2P+Q=(3,-4).\\
3P+Q&=(9,-26)+(2,2):\ m=-4,\ x_3=5,\ \text{third }(5,-10),\ 3P+Q=(5,10).\\
4P+Q&=P+(3P+Q)=(1,2)+(5,10):\ m=2,\ x_3=-2,\ \text{third }(-2,-4),\ 4P+Q=(-2,4).\\
2P+2Q&=(2P+Q)+Q=(3,-4)+(2,2):\ m=-6,\ x_3=31,\ \text{third }(31,-172),\ 2P+2Q=(31,172).\\
5P+2Q&=(4P+Q)+(P+Q)=(-2,4)+(-3,-2):\ m=6,\ x_3=41,\ \text{third }(41,262),\ 5P+2Q=(41,-262).\\
6P+Q&=(3P+Q)+3P=(5,10)+(9,-26):\ m=-9,\ x_3=67,\ \text{third }(67,-548),\ 6P+Q=(67,548).\\
4P-Q&=3P+(P-Q)=(9,-26)+(13,46):\ m=18,\ x_3=302,\ \text{third }(302,5248),\ 4P-Q=(302,-5248).
\end{align*}
Together with $P=(1,2)$, $Q=(2,2)$ these are the $13$ points $P,\ Q,\ P+Q,\ 2P,\ P-Q,\ 3P,\ 2P+Q,\ 3P+Q,\ 4P+Q,\ 2P+2Q,\ 5P+2Q,\ 6P+Q,\ 4P-Q$ of the $13$ pairs, so all $26$ integer points lie in the subgroup generated by $P$ and $Q$. (I verified these sums by exact arithmetic as well.) Whether $\langle P,Q\rangle$ is all of $X(\mathbf Q)$ would require a descent (determination of the Mordell--Weil group), which I have not carried out; the exercise only asks for the verification above. By Ex.~4.19 below the group $X(\mathbf Q)$ is torsion-free. \hfill$\square$

\exercise{4.19}{Let $X,P_0$ be an elliptic curve defined over $\mathbf Q$, represented as a curve in $\mathbf P^2$ by an equation with integer coefficients. Then $X$ can be considered as the fibre over the generic point of a scheme $\overline X$ over $\operatorname{Spec}\Z$. Let $T\subseteq\operatorname{Spec}\Z$ be the open subset consisting of all primes $p\ne2$ such that the fibre $X_{(p)}$ of $\overline X$ over $p$ is nonsingular. For any $n$, show that $n_X:X\to X$ is defined over $T$, and is a flat morphism. Show that the kernel of $n_X$ is also flat over $T$. Conclude that for any $p\in T$, the natural map $X(\mathbf Q)\to X_{(p)}(\mathbf F_p)$ induced on the groups of rational points, maps the $n$-torsion points of $X(\mathbf Q)$ injectively into the torsion subgroup of $X_{(p)}(\mathbf F_p)$, for $(n,p)=1$.

By this method one can show easily that the groups $X(\mathbf Q)$ in (Ex.~4.17) and (Ex.~4.18) are torsion-free.}
\pf
This exercise needs several results from scheme theory beyond the text; I name them where used, and indicate where the proof of a theorem of the text is used over a base other than a field.

\textbf{Setup.} Let $F$ be the integral cubic form, $\overline X=\operatorname{Proj}\Z[x,y,z]/(F)\subseteq\mathbf P^2_\Z$, $X=\overline X_{\mathbf Q}$, and let $P_0\in X(\mathbf Q)$ be given by coprime integers; it defines a section $e:\operatorname{Spec}\Z\to\overline X$ ($F(P_0)=0$ in $\Z$). $T$ is open and nonempty (the discriminant is a nonzero integer), so $T=\operatorname{Spec}\Z[S^{-1}]$ is connected, regular of dimension $1$. Put $\overline X_T=\overline X\times_{\Z}T$, $\pi:\overline X_T\to T$. Its fibres over $p\in T$ are the nonsingular plane cubics $X_{(p)}$ (nonzero mod $p$), over the generic point $X$; the Hilbert polynomial of all fibres is $3m$, so $\pi$ is flat (III, Thm.~9.9) and projective. A flat morphism of finite type with nonsingular geometric fibres is smooth (EGA IV 17.5.1; III, 10.2 is the case over a field), here of relative dimension $1$. In particular $\overline X_T$ is a regular scheme (smooth over the regular scheme $T$; alternatively Matsumura 23.7), of dimension $2$ at closed points, and all fibres are elliptic curves (over $\mathbf Q$ resp.\ $\mathbf F_p$, geometrically).

\textbf{(1) $n_X$ is defined over $T$.} We redo the proof of Thm.~4.11 with base $T$ in place of $\operatorname{Spec}k$, to show that $(\overline X_T,e)$ is a group scheme over $T$. Let $T'$ be a $T$-scheme of finite type, $X'=\overline X_T\times_TT'$ with projections $p:X'\to T'$, $q:X'\to\overline X_T$, and let $\mathcal M\in\operatorname{Pic}^0(X'/T')$ (invertible, of degree $0$ on all fibres). Put $\mathcal M'=\mathcal M\otimes q^*\mathcal L(e(T))$, where $e(T)\subseteq\overline X_T$ is an effective Cartier divisor (a section of the smooth morphism $\pi$ of relative dimension $1$ is a regular immersion of codimension $1$; likewise for the diagonal $\Delta\subseteq\overline X_T\times_T\overline X_T$, which is used at the end). 
\begin{itemize}
\item $p$ is projective and flat with geometrically integral fibres of genus $1$, and $\mathcal M'$ has degree $1$ on each fibre. By Riemann--Roch on the geometric fibre $X'_t\otimes\overline{k(t)}$ (which does not change $h^0,h^1$), $h^0(X'_t,\mathcal M'_t)=1$ and $h^1=0$ for every point $t\in T'$ (closed or not).
\item By cohomology and base change (III, 12.11) we get $R^1p_*\mathcal M'=0$ and $p_*\mathcal M'$ locally free of rank $1$, compatible with base change. (Only closed points need the explicit computation: $T'$ is of finite type over $\Z$, so every nonempty closed subset of $T'$ contains a closed point, and $\{t:h^1(\mathcal M'_t)\ge1\}$ is closed by semicontinuity, III, 12.8.)
\item Replacing $\mathcal M$ by $\mathcal M\otimes p^*(p_*\mathcal M')^{-1}$ (this does not change its class in $\operatorname{Pic}^0(X'/T')$) we may assume $p_*\mathcal M'\cong\mathcal O_{T'}$. The section $1$ gives $s\in\Gamma(X',\mathcal M')$, nonzero on every fibre (base change), hence a nonzerodivisor on each integral fibre. As $X'$ is flat over $T'$, its zero scheme $Z$ is an effective Cartier divisor flat over $T'$ (local criterion of flatness), whose fibres are single reduced points ($\deg\mathcal M'_t=1$). Thus $Z\to T'$ is finite, flat of degree $1$, i.e.\ an isomorphism, and gives a section $T'\to X'$, hence $f:T'\to\overline X_T$ (compose with $q$).
\item Since $Z$ is the graph of $f$, $Z=(f\times1)^*\Delta$ and $\mathcal M'\cong\mathcal O(Z)\cong(f\times 1)^*\mathcal L(\Delta)$ (up to a pullback from $T'$); twisting by $-e(T)$ gives $f^*\mathcal L\cong\mathcal M$ with $\mathcal L=\mathcal L(\Delta)\otimes p_2^*\mathcal L(-e(T))$ in $\operatorname{Pic}(X'/T')$. Uniqueness of $f$ is clear for the same reasons (the zero scheme of $s$ determines $f$).
\end{itemize}
Hence $(\overline X_T,\mathcal L)$ represents $T'\mapsto\operatorname{Pic}^0(X'/T')$ on $T$-schemes of finite type, and, as in Remark 4.10.6, it is a group scheme over $T$ with identity $e$, inverse $\rho$ and multiplication $\mu$ (Yoneda: the group structure of $\operatorname{Pic}^0(\cdot/T')$ is functorial in $T'$). Its fibres are the elliptic curves $X_{(p)}$ (after base change to $\overline{\mathbf F}_p$) with the group law of Prop.~4.8 relative to $e(p)$ (last statement of Thm.~4.11 over the algebraically closed field), and the generic fibre is $(X,P_0)$. Then $n_X$, the $n$-fold iterate of $\mu$ (composed with $\rho$ if $n<0$), is a $T$-morphism $\overline X_T\to\overline X_T$ restricting to $n_X$ on the generic fibre (and to $n_{X_{(p)}}$ on the fibre over $p$). Because $p\ne2$ for $p\in T$, Prop.~4.10 (which assumes $\operatorname{char}\neq2$), applied after base change to $\overline{\mathbf F}_p$, resp.\ $\overline{\mathbf Q}$, shows that for $n\ne0$ each restriction $n_{X_{(p)}}$ and $n_X$ are finite surjective morphisms of curves.

\textbf{(2) $n_X$ is flat.} Let $n\ne0$ and write $g=n_X:\overline X_T\to\overline X_T$. (For $n=0$, $0_X=e\circ\pi$ is not flat, as its image $e(T)$ is not open; so ``for any $n$'' must be read as $n\ne0$.) The morphism $g$ is projective (a $T$-morphism between projective $T$-schemes) and quasi-finite (on each fibre over $T$ it is a finite morphism of curves), hence finite (III, Ex.~11.2). Let $x\in\overline X_T$, $y=g(x)$, $A=\mathcal O_y\to B=\mathcal O_x$. Both are regular local rings, and $B/\mathfrak m_AB$ has dimension $0$ since the fibre of $g$ over $y$ is finite. Furthermore $\dim A=\dim B$: since $\pi$ is flat, $\dim\mathcal O_x=\dim\mathcal O_{T,\pi(x)}+\dim\mathcal O_{X_{\pi(x)},x}$, and the same for $y$ (note $\pi(x)=\pi(y)$); the finite surjective morphism $g$ of the fibre curve maps closed points to closed points and the generic point to the generic point, so $\dim\mathcal O_{X_t,x}=\dim\mathcal O_{X_t,y}$. By the ``miracle flatness'' criterion (Matsumura, \emph{Commutative Ring Theory}, Thm.~23.1: $A$ regular, $B$ Cohen--Macaulay, $\dim B=\dim A+\dim B/\mathfrak m_AB$ imply $B$ flat over $A$) $g$ is flat at every point. So $n_X:\overline X_T\to\overline X_T$ is finite and flat; being flat and finite it is locally free of finite rank.

\textbf{(3) The kernel.} Let $K=\overline X_T\times_{n_X,\overline X_T,e}T=\ker(n_X)$ with structure map $K\to T$. This is the base change of $n_X$ along $e:T\to\overline X_T$, hence $K\to T$ is finite and flat, i.e.\ $K=\operatorname{Spec}\mathcal A$ with $\mathcal A$ a locally free $\mathcal O_T$-algebra of finite rank; the rank is constant ($T$ connected) and equals $\deg n_X=n^2$ (generic fibre $=X[n]$ has $n^2$ points over $\overline{\mathbf Q}$, since $\deg n_X=n^2$ by Ex.~4.7(e) and $n_X$ is separable in characteristic $0$).

\textbf{(4) Injectivity on torsion.} Let $p\in T$ with $p\nmid n$ and $R_p=\Z_{(p)}$. Each $x\in X(\mathbf Q)$ extends uniquely to an $R_p$-point $\tilde x$ of $\overline X$ (valuative criterion, $\overline X$ proper), and the natural map is $x\mapsto\bar x=\tilde x\bmod p\in X_{(p)}(\mathbf F_p)$; it is a homomorphism because $\mu$ is defined over $T$ and $\tilde x+\tilde y$ is the extension of $x+y$. Let $x\in X(\mathbf Q)$ with $nx=0$. Then $x$ is a $\mathbf Q$-point of $K$, and since $K\times R_p=\operatorname{Spec}B$ with $B$ finite over $R_p$, the ring homomorphism $B\to\mathbf Q$ has image integral over $R_p$, hence contained in the integrally closed ring $R_p$: $\tilde x$ is an $R_p$-point of $K$. Next, $K\times\mathbf F_p=\ker(n_{X_{(p)}})$ is \emph{reduced}: $d(n_{X_{(p)}})=n\cdot\mathrm{id}\ne0$ in $\mathbf F_p$ (T1 in Ex.~4.15), so $n_{X_{(p)}}$ is unramified at $0$, and for $Q\in X_{(p)}[n]$, translation by $Q$ identifies neighbourhoods of $0$ and of $Q$ compatibly with $n_{X_{(p)}}$. Hence $K\times\overline{\mathbf F}_p$ is a finite reduced scheme over an algebraically closed field, i.e.\ a disjoint union of $n^2$ reduced points, so $K\times\mathbf F_p\to\operatorname{Spec}\mathbf F_p$ is etale. Since $K\times R_p\to\operatorname{Spec}R_p$ is finite flat with etale closed fibre, it is etale (the etale locus is open and contains the closed fibre, and every point of the finite scheme $K\times R_p$ specializes into the closed fibre). Now suppose $\bar x=0$. Then the two sections $\tilde x$ and $e$ of $K\times R_p\to\operatorname{Spec}R_p$ agree at the closed point. The locus where two sections of a separated etale morphism agree is open and closed, and $\operatorname{Spec}R_p$ is connected, so $\tilde x=e$, and so $x=0$. Thus the homomorphism $X(\mathbf Q)[n]\to X_{(p)}(\mathbf F_p)$ is injective, with image in the $n$-torsion subgroup $X_{(p)}(\mathbf F_p)[n]$ (it is a homomorphism). In particular, if $x\in X(\mathbf Q)$ has order $m$ with $p\nmid m$, then $\bar x$ has order $m$, so $m$ divides $\#X_{(p)}(\mathbf F_p)$.

\textbf{Application to (4.17), (4.18).} Suppose $X(\mathbf Q)$ had a point of prime order $\ell$. For every $p\in T$ with $p\ne\ell$ we would get $\ell\mid N_p:=\#X_{(p)}(\mathbf F_p)$.

\emph{(4.17)} $y^2+y=x^3-x$: $p=3,5\in T$ (nonsingular for $p\ne37$, Ex.~4.17(b)). Counting points (including the point at infinity): over $\mathbf F_3$: $x=0,1,2$ each give $2$ points, so $N_3=6+1=7$; over $\mathbf F_5$: $x^3-x=0,0,1,4,0$ for $x=0,1,2,3,4$ and $y^2+y=c$ has $1+\left(\frac{1+4c}{5}\right)$ solutions, giving $2,2,1,0,2$ points, $N_5=7+1=8$. If $\ell\ne3$ then $\ell\mid7$, so $\ell=7$, but $7\ne5$ and $7\nmid8$; if $\ell=3$ then $3\nmid N_5=8$. Contradiction: $X(\mathbf Q)$ is torsion-free.

\emph{(4.18)} $y^2=x^3-7x+10$: the discriminant is $-16(4(-7)^3+27\cdot10^2)=-2^8\cdot83$, so $3,5,7\in T$, and counting points mod $3,5,7$ gives $N_3=7$, $N_5=10$, $N_7=13$. If $\ell\neq3$, $\ell\mid7$, so $\ell=7$; but $7\ne5$ and $7\nmid10$. If $\ell=3$: $3\nmid N_5=10$. So $X(\mathbf Q)$ is torsion-free. (For $\ell=2$ one can also see directly that $x^3-7x+10$ has no rational root.) \hfill$\square$

\exercise{4.20}{Let $X$ be an elliptic curve over a field $k$ of characteristic $p>0$, and let $R=\operatorname{End}(X,P_0)$ be its ring of endomorphisms.
\begin{enumerate}[label=(\alph*),leftmargin=*]
\item Let $X_p$ be the curve over $k$ defined by changing the $k$-structure of $X$ (2.4.1). Show that $j(X_p)=j(X)^{1/p}$. Thus $X\cong X_p$ over $k$ if and only if $j\in\mathbf F_p$.
\item Show that $p_X$ in $R$ factors into a product $\pi\widehat\pi$ of two elements of degree $p$ if and only if $X\cong X_p$. In this case, the Hasse invariant of $X$ is $0$ if and only if $\pi$ and $\widehat\pi$ are associates in $R$ (i.e., differ by a unit). (Use (2.5).)
\item If $\operatorname{Hasse}(X)=0$ show in any case $j\in\mathbf F_{p^2}$.
\item For any $f\in R$, there is an induced map $f^*:H^1(\mathcal O_X)\to H^1(\mathcal O_X)$. This must be multiplication by an element $\lambda_f\in k$. So we obtain a ring homomorphism $\varphi:R\to k$ by sending $f$ to $\lambda_f$. Show that any $f\in R$ commutes with the (nonlinear) Frobenius morphism $F:X\to X$, and conclude that if $\operatorname{Hasse}(X)\ne0$, then the image of $\varphi$ is $\mathbf F_p$. Therefore, $R$ contains a prime ideal $\mathfrak p$ with $R/\mathfrak p\cong\mathbf F_p$.
\end{enumerate}}
\pf
($k$ is algebraically closed.) In (a) the $j$-invariant is the one of the text, defined for $p\ne2$ via Prop.~4.6; for $p=2$ the same proof works with the usual $j$-invariant of a Weierstrass equation, whose formula has coefficients in $\mathbf F_2$. The printed exponent is read as $1/p$ (the scan shows ``$j(X)^{1\,p}$''), which is what the proof gives; it implies the stated equivalence.

\textbf{(a)} Let $\sigma:k\to k$ be the automorphism $a\mapsto a^{1/p}$. Let $X=\operatorname{Proj}S$, $S=k[x,y,z]/(f)$, $f=\sum c_Ix^I$, and let $X^\sigma$ be the curve defined by $f^\sigma=\sum\sigma(c_I)x^I$. As in Ex.~4.16(a), $X_p=\operatorname{Proj}S_p$ with $S_p$ the $k$-algebra $S$ with $a*s=a^ps$. The ring map $\Psi:k[x,y,z]/(f^\sigma)\to S_p$, $\sum d_Ix^I\mapsto\sum d_I^px^I$, is well defined ($\Psi(f^\sigma)=\sum\sigma(c_I)^px^I=f$), graded, bijective, and $k$-linear for the structure $*$ ($\Psi(as)=a^p\Psi(s)=a*\Psi(s)$). Hence $X_p\cong X^\sigma$ over $k$. For $p\neq2$ take $X:y^2=x(x-1)(x-\lambda)$ (Prop.~4.6); then $X_p\cong y^2=x(x-1)(x-\sigma(\lambda))$, so $\lambda(X_p)=\sigma(\lambda)$ up to $\Sigma_3$, and since $j(\lambda)=2^8(\lambda^2-\lambda+1)^3/(\lambda^2(\lambda-1)^2)$ has coefficients in $\mathbf F_p$,
\[
j(X_p)=j(\sigma(\lambda))=\sigma(j(\lambda))=j(X)^{1/p}.
\]
By Thm.~4.1(b), $X\cong X_p$ iff $j^{1/p}=j$ iff $j^p=j$ iff $j\in\mathbf F_p$.

\textbf{(b)} Recall (Ex.~4.7(c),(f)): every $\pi\in R$ of degree $m$ satisfies $\pi\widehat\pi=\widehat\pi\pi=m_X$ and $\deg\widehat\pi=m$; so $p_X=\pi\widehat\pi$ with both of degree $p$ holds for \emph{every} $\pi\in R$ of degree $p$, and the first claim says: $R$ contains an element of degree $p$ iff $X\cong X_p$. Also $p_X=F'\widehat{F'}$ (proof of Ex.~4.15).

($\Leftarrow$) Let $\theta:X_p\to X$ be an isomorphism; composing with a translation we may assume $\theta(P_0)=P_0$. Then $\pi=F'\circ\theta^{-1}\in R$ has degree $\deg F'=p$.

($\Rightarrow$) Let $\pi\in R$ have degree $p$. By (T1), $0=p\cdot\mathrm{id}=d(p_X)=d\pi\circ d\widehat\pi$ in $k$, so $d\pi=0$ or $d\widehat\pi=0$, i.e.\ one of $\pi,\widehat\pi$ is inseparable by (T2), hence (degree $p$ prime) purely inseparable of degree $p$. By Prop.~2.5 (with $r=1$, $Y=X$) such a morphism is, up to an isomorphism, the $k$-linear Frobenius $X_p\to X$, and in particular $X\cong X_p$ over $k$. (Concretely: if $L=\rho^*K(X)\subseteq K(X)$ for $\rho$ purely inseparable of degree $p$, then $K(X)^p\subseteq L$ and $[K(X):K(X)^p]=p$, so $L=K(X)^p$.)

\emph{Hasse invariant.} Assume $X\cong X_p$ and let $\pi\in R$ be any element of degree $p$. The statement is symmetric in $\pi,\widehat\pi$ (as $\widehat{\widehat\pi}=\pi$), and by the argument of ($\Rightarrow$) at least one of $\pi,\widehat\pi$ is purely inseparable; replacing $\pi$ by $\widehat\pi$ if necessary we may assume $\pi$ is purely inseparable of degree $p$. By Prop.~2.5 it is then $\pi=F'\theta^{-1}$ for an isomorphism $\theta:X_p\to X$, and $\theta(P_0)=P_0$ because $\pi(P_0)=P_0$ and $F'$ is the identity on points. Since $\theta^{-1}$ is an isomorphism, $\widehat{\theta^{-1}}=\theta$ ($\widehat g g=1$), so $\widehat\pi=\theta\widehat{F'}$ (Ex.~4.7(b)). Now $\pi$ is purely inseparable of degree $p$, and $\widehat\pi$ is separable iff $\widehat{F'}$ is, iff Hasse $=1$ (Ex.~4.15). If $\operatorname{Hasse}(X)=1$, $\widehat\pi$ is separable and $\pi$ is not, so they are not associates (composing with a unit preserves (in)separability). If $\operatorname{Hasse}(X)=0$ both $\pi$ and $\widehat\pi$ are purely inseparable of degree $p$, so $\pi^*K(X)=K(X)^p=\widehat\pi^*K(X)$, and then $v=(\pi^*)^{-1}\circ\widehat\pi^*$ is an automorphism of the field $K(X)$ with $\widehat\pi^*=\pi^*\circ v$. By the equivalence of categories (I, 6.12) $v=u^*$ for an automorphism $u$ of $X$, and $\widehat\pi^*=\pi^*u^*=(u\circ\pi)^*$, i.e.\ $\widehat\pi=u\circ\pi$. Moreover $u(P_0)=u(\pi(P_0))=\widehat\pi(P_0)=P_0$. So $u$ is a unit of $R$ and $\widehat\pi=u\pi$: $\pi,\widehat\pi$ are associates. This proves: Hasse $=0$ iff $\pi,\widehat\pi$ are associates.

\textbf{(c)} If Hasse $=0$, $\widehat{F'}:X\to X_p$ is purely inseparable of degree $p$ (not separable, Ex.~4.15). By Prop.~2.5 applied to $\widehat{F'}$, $X\cong(X_p)_p=X_{p^2}$ over $k$. By (a) applied twice, $j(X_{p^2})=j(X)^{1/p^2}$, so $j=j^{1/p^2}$, i.e.\ $j^{p^2}=j$, i.e.\ $j\in\mathbf F_{p^2}$.

\textbf{(d)} For $f\in R$, $f$ is a $k$-morphism so $f^*:H^1(X,\mathcal O_X)\to H^1(X,\mathcal O_X)$ is $k$-linear, and $H^1$ is $1$-dimensional, so $f^*=\lambda_f\cdot\mathrm{id}$. $\varphi(1)=1$; $(fg)^*=g^*f^*$ gives $\varphi(fg)=\varphi(f)\varphi(g)$; additivity: by (T3) (Ex.~4.15) the tangent map of $\widehat f$ is $f^*$, and $(f+g)^\wedge=\widehat f+\widehat g$ (Ex.~4.7(d)), so by (T1) $(f+g)^*=d(\widehat f+\widehat g)=f^*+g^*$. So $\varphi:R\to k$ is a ring homomorphism.

The absolute Frobenius commutes with every morphism of schemes: $f\circ F=F\circ f$ (both send $s\mapsto f^\#(s)^p=f^\#(s^p)$). Hence $F^*f^*=f^*F^*$ on $H^1(X,\mathcal O_X)$. Take $a\in H^1$ and recall $F^*$ is $p$-linear: $F^*(\lambda_fa)=\lambda_f^pF^*a$, while $f^*(F^*a)=\lambda_fF^*a$. So $(\lambda_f^p-\lambda_f)F^*(a)=0$ for all $a$. If $\operatorname{Hasse}(X)\ne0$, $F^*$ is bijective and we can choose $a$ with $F^*a\neq0$, so $\lambda_f^p=\lambda_f$, i.e.\ $\lambda_f\in\mathbf F_p$. Thus $\varphi:R\to\mathbf F_p$, and it is onto since $\varphi(1)=1$. Its kernel $\mathfrak p$ is a two-sided ideal with $R/\mathfrak p\cong\mathbf F_p$ a field, hence a (maximal and) prime ideal. (It contains $p_X$.) \hfill$\square$

\exercise{4.21}{Let $O$ be the ring of integers in a quadratic number field $\Q(\sqrt{-d})$. Show that any subring $R\subseteq O$, $R\ne\Z$, is of the form $R=\Z+f\cdot O$, for a uniquely determined integer $f\ge1$. This integer $f$ is called the conductor of the ring $R$.}
\pf
Here $d>0$ is squarefree (the field is imaginary quadratic; this is the case of interest, cf.\ Thm.~4.19) and subrings contain $1$. Then $O=\Z\oplus\Z\eta$ with $\eta=\sqrt{-d}$ if $d\equiv1,2\pmod4$ and $\eta=\frac{1+\sqrt{-d}}2$ if $d\equiv3\pmod4$. Let $\pi:O\to\Z$, $a+b\eta\mapsto b$, a surjective homomorphism of groups with kernel $\Z\cdot1$.

Let $R\subseteq O$ be a subring, $R\neq\Z$. Then $\Z\subseteq R$ and $\pi(R)=f\Z$ for a unique integer $f\ge0$. If $f=0$ then $R\subseteq\ker\pi=\Z$, so $R=\Z$, excluded; hence $f\ge1$. Since $\Z\subseteq R$ and $\pi(R)=f\Z$, we have $R=\pi^{-1}(f\Z)=\Z\oplus\Z f\eta$ (every element $a+b\eta$ of $R$ has $f\mid b$; conversely there is $r_0\in R$ with $\pi(r_0)=f$, $r_0=a_0+f\eta$, and $a_0\in\Z\subseteq R$, so $f\eta\in R$ and then $a+bf\eta\in R$ for all $a,b\in\Z$). On the other hand
\[
\Z+fO=\Z+\Z f+\Z f\eta=\Z+\Z f\eta .
\]
So $R=\Z+fO$. (Conversely $\Z+fO$ is a subring for every $f\ge1$: $(fu)(fv)=f(fuv)\in fO$.)

\emph{Uniqueness.} $[O:\Z+fO]=[\Z\eta:\Z f\eta]=f$, because $O/R\cong\Z/\pi(R)=\Z/f$ (via $\pi$, as $R=\pi^{-1}(f\Z)$). So $f=[O:R]$ is determined by $R$. \hfill$\square$

\exercise{4.22}{If $X\to\mathbf A^1_{\mathbf C}$ is a family of elliptic curves having a section, show that the family is trivial. [Hints: Use the section to fix the group structure on the fibres. Show that the points of order $2$ on the fibres form an etale cover of $\mathbf A^1_{\mathbf C}$, which must be trivial, since $\mathbf A^1_{\mathbf C}$ is simply connected. This implies that $\lambda$ can be defined on the family, so it gives a map $\mathbf A^1_{\mathbf C}\to\mathbf A^1_{\mathbf C}-\{0,1\}$. Any such map is constant, so $\lambda$ is constant, so the family is trivial.] (Starred in the book.)}
\pf
\emph{Interpretation.} The scan prints the base as $\mathbf A^1_{\mathbf C}$ (affine line over $\C$, as the hint about simple connectedness requires). A family of elliptic curves is a smooth projective morphism $\pi:X\to\mathbf A^1=\operatorname{Spec}A$, $A=\C[t]$, with all fibres of genus $1$, and a section $e$; ``trivial'' means $X\cong E\times\mathbf A^1$ over $\mathbf A^1$ for an elliptic curve $E$ (compatibly with the section). I give a direct proof in which the three hinted steps are replaced by a Weierstrass equation over $A$: the role of ``the $2$-torsion is a trivial cover'' is played by the fact that $A$ is a PID, and ``$\lambda$ is constant'' by the fact that a family with a Weierstrass equation over $A$ has constant discriminant, forcing a map from $\mathbf A^1$ to a curve of genus $1$ to be constant.

\textbf{Step 1: a global Weierstrass equation.} Since $\pi$ is smooth and $e$ is a section, $D=e(\mathbf A^1)$ is an effective Cartier divisor on $X$, flat over $A$. Put $\mathcal L_n=\mathcal O_X(nD)$ and $M_n=H^0(X,\mathcal L_n)$. On each fibre $X_t$ (a genus $1$ curve with origin $e(t)$), $\mathcal L_n|_{X_t}=\mathcal O(n\,e(t))$ has degree $n\ge1$, so by Riemann--Roch $h^0=n$ and $h^1=0$ for all $t$ (also for non-closed points, by semicontinuity). By cohomology and base change (III, 12.11, or Grauert 12.9), $\pi_*\mathcal L_n$ is locally free of rank $n$ and $M_n\otimes k(t)=H^0(X_t,\mathcal L_n|_{X_t})$ for every $t$. $A$ is a PID, so $M_n$ is free of rank $n$. The inclusions $M_n\subseteq M_{n+1}$ stay injective on all fibres $H^0(X_t,\mathcal O(ne))\subseteq H^0(X_t,\mathcal O((n+1)e))$, so $M_{n+1}/M_n$ is locally free of rank $1$, hence free. We have $M_1=A\cdot1$. Choose $x\in M_2$ with $M_2=A\oplus Ax$ and $y\in M_3$ with $M_3=M_2\oplus Ay$. On every fibre $x_t$ has a pole of order exactly $2$ at $e(t)$ and $y_t$ of order exactly $3$ (they are not in the previous spaces), so $x^2,xy,x^3,y^2$ have exact pole orders $4,5,6,6$ on every fibre, and by Nakayama $x^2,xy$ generate $M_4/M_3$, $M_5/M_4$, and $x^3$ generates $M_6/M_5$. Therefore $M_6$ has $A$-basis $1,x,y,x^2,xy,x^3$ and $y^2=u\,x^3+(\text{combination of }1,x,y,x^2,xy)$ with $u\in A$ nowhere vanishing (the pole order of $y_t^2$ is exactly $6$), i.e.\ $u\in A^*=\C^*$. Replacing $(x,y)$ by $(ux,uy)$ gives
\[
y^2+a_1xy+a_3y=x^3+a_2x^2+a_4x+a_6,\qquad a_i\in A.
\]
As $\operatorname{char}=0$, replacing $y$ by $y+\frac12(a_1x+a_3)$ and $x$ by $x-\frac13a_2$ (changes inside $M_3$, $M_2$) we get $y^2=x^3+ax+b$ with $a,b\in A=\C[t]$.

The sections $1,x,y$ of $\mathcal L_3$ generate it on every fibre and give on each fibre $X_t$ the closed embedding $X_t\hookrightarrow\mathbf P^2$ by $|3e(t)|$ ((3.3.3), as in the proof of Prop.~4.6). Hence $\iota:X\to\mathbf P^2_A$ is a proper morphism which is a closed immersion on every fibre, hence a closed immersion (proper + injective on points + unramified; II, 7.3 / EGA IV 18.12.6), with image inside $W=\{y^2z=x^3+axz^2+bz^3\}\subseteq\mathbf P^2_A$. Both $X$ and $W$ are flat over $A$ (the cubic form is nonzero on every fibre) and $\iota$ is an isomorphism of $X_t$ onto $W_t$ for every closed $t$ ($X_t$ is an integral curve contained in the integral plane cubic $W_t$ --- $y^2z-(\text{cubic in }x,z)$ is irreducible --- so it has the same underlying space, and a closed subscheme of the reduced scheme $W_t$ with the same underlying space is $W_t$). If $\mathcal I$ is the ideal of $X$ in $W$, then $\mathcal I\otimes k(t)=0$ for all closed $t$ (flatness of $\mathcal O_X$ gives injectivity of $\mathcal I\otimes k(t)\to\mathcal O_{W_t}$ with image the ideal of $X_t$ in $W_t$, which is $0$), so $\mathcal I=0$ by Nakayama and since the support of $\mathcal I$ is closed and every nonempty closed subset of $W$ contains a closed point lying over a closed point of $\mathbf A^1$. Thus $X\cong W$ over $A$.

\textbf{Step 2: the discriminant is constant.} Each fibre $X_t$ is nonsingular, so the cubic $x^3+a(t)x+b(t)$ has distinct roots: $\delta(t)=4a(t)^3+27b(t)^2\ne0$ for all $t\in\C$. So $\delta\in A=\C[t]$ has no zeros, hence $\delta\equiv\delta_0\in\C^*$ is a nonzero constant.

\textbf{Step 3: $a,b$ are constant.} The affine curve $C:4a^3+27b^2=\delta_0$ in $\mathbf A^2$ is nonsingular (the partials $12a^2,54b$ vanish only at $(0,0)\notin C$), and its closure $\overline C:4a^3+27b^2c=\delta_0c^3$ in $\mathbf P^2$ (coordinates $a,b,c$) is a nonsingular plane cubic (the only point at infinity, $c=0$, is $(0:1:0)$ since then $a^3=0$, and there $\partial/\partial c=27b^2-3\delta_0c^2=27\ne0$), so $\overline C$ is irreducible (two components of a plane curve would meet in a singular point) of genus $1$ ((I, Ex.~7.2) / (II, 8.20.3)). The polynomial map $t\mapsto(a(t),b(t))$ is a morphism $\mathbf A^1\to C\subseteq\overline C$, which extends to a morphism $\mathbf P^1\to\overline C$ ($\mathbf P^1$ is a nonsingular curve and $\overline C$ is projective). If it were nonconstant, it would be a finite morphism and Example 2.5.4 would give $g(\mathbf P^1)=0\ge g(\overline C)=1$, which is absurd. Hence $a(t)\equiv a_0$, $b(t)\equiv b_0$ are constants.

\textbf{Conclusion.} $X\cong\{y^2z=x^3+a_0xz^2+b_0z^3\}\times\mathbf A^1=E\times\mathbf A^1$ over $\mathbf A^1$, where $E$ is the elliptic curve $y^2=x^3+a_0x+b_0$ (nonsingular since $\delta_0\neq0$), and the section $e$ corresponds to the point at infinity $(0:1:0)\times\mathbf A^1$ (the pole of $x,y$). So the family is trivial. (The $\lambda$-invariant of the family is then constant, as in the hint: in the notation of Step 1 the roots of $x^3+ax+b$ are the abscissas of the points of order $2$, and they are constants.) \hfill$\square$

\section*{5. The canonical embedding}

Conventions. $k$ is algebraically closed; curves are nonsingular projective curves; $K$ is a canonical divisor, $h^0(D)=\dim H^0(X,\mathcal O_X(D))$ and $\dim|D|=h^0(D)-1$. Results of the text are quoted by number (Prop.~5.2, Thm.~5.4, Ex.~4.6, \dots). Several exercises (5.2, 5.3, 5.5, 5.7) use that $\operatorname{Aut}X$ is finite, Bertini's theorem, Hurwitz' theorem and $|\operatorname{Aut}X|\le 84(g-1)$ (Ex.~2.5); there I assume $\operatorname{char}k=0$, as the book's Ex.~2.5 does. Two standard facts beyond the text are used and named where needed: the Plücker formula for the total weight of the Weierstrass points (this is the weighted form of Ex.~4.6(b)), and the elementary Castelnuovo--Severi inequality, which is proved below as Lemma~A.

\medskip
\noindent\textbf{Lemma A (Castelnuovo--Severi, simple form).} Let $f,f'$ be base-point-free pencils on a curve $X$ of genus $g$, of degrees $a,b$, i.e.\ morphisms $f,f':X\to\mathbf P^1$. Then either $g\le(a-1)(b-1)$, or there is a morphism $h:X\to Z$ of degree $m>1$ with $m\mid a$, $m\mid b$, through which both $f$ and $f'$ factor. In particular: (i) if $\gcd(a,b)=1$ then $g\le (a-1)(b-1)$; (ii) if $a=b=3$ and $g\ge5$ then $f=\alpha\circ h$, $f'=\beta\circ h$ with $h:X\to\mathbf P^1$ of degree $3$ and $\alpha,\beta\in\operatorname{Aut}\mathbf P^1$, so $f$ and $f'$ define the same linear system.

\emph{Proof.} Let $\varphi=(f,f'):X\to\mathbf P^1\times\mathbf P^1$, $Y=\varphi(X)$, $m=\deg(X\to Y)$, $\nu:\tilde Y\to Y$ the normalization, so $\varphi=\nu\circ h$ with $h:X\to\tilde Y$ of degree $m$. The projections $\mathrm{pr}_1\circ\nu$, $\mathrm{pr}_2\circ\nu:\tilde Y\to\mathbf P^1$ are finite of degrees $a/m$, $b/m$, so $m\mid a,b$. If $m=1$, $X=\tilde Y$ is the normalization of the integral curve $Y\subset\mathbf P^1\times\mathbf P^1$. $Y$ is an effective divisor of some type $(a',b')$; intersecting with the two rulings shows $\{a',b'\}=\{a,b\}$, and by (III, Ex.~5.6) $p_a(Y)=(a-1)(b-1)$. By (Ex.~1.8) $g\le p_a(Y)$. If $m>1$ we are in the second alternative. For (ii): $m=3$, so $\tilde Y\to\mathbf P^1$ has degree $1$ and is finite onto $\mathbf P^1$ with $\tilde Y$ normal, hence an isomorphism; so $f=\alpha\circ h$ with $\alpha:\tilde Y\cong\mathbf P^1$, and likewise $f'=\beta\circ h$ with $\beta:\tilde Y\cong\mathbf P^1$ (here $h:X\to\tilde Y$). \hfill$\square$

\exercise{5.1}{Show that a hyperelliptic curve can never be a complete intersection in any projective space. Cf.\ (Ex.~3.3).}
\pf
Let $X$ be hyperelliptic (so $g\ge2$) and suppose $X\subseteq\mathbf P^n$ is a (nonsingular) complete intersection of hypersurfaces of degrees $d_1,\dots,d_{n-1}$. By (II, Ex.~8.4(c)), $\omega_X\cong\mathcal O_X(a)$ with $a=\sum d_i-n-1$. Since $\deg\omega_X=2g-2>0$ and $\deg\mathcal O_X(1)>0$ we get $a\ge1$.

We show that $|K|=|aH|$ ($H$ a hyperplane section) is very ample, contradicting Prop.~5.2. By (3.1) we must show that for all $P,Q\in X$ (possibly equal) $h^0(K-P-Q)=h^0(K)-2$, i.e.\ that the evaluation map $H^0(X,\mathcal O_X(a))\to H^0(\mathcal O_{P+Q}\otimes\mathcal O_X(a))$ onto the $2$-dimensional space of $\mathcal O_X(a)$ restricted to the subscheme $P+Q$ is surjective. Choose a linear form $\ell_0$ not vanishing at $P$ or $Q$ and let $V\subseteq H^0(\mathcal O_X(a))$ be the image of $\{\ell_0^{a-1}\ell:\ell\text{ a linear form}\}$. Since $\ell_0^{a-1}$ is a unit near $P,Q$, the image of $V$ in $H^0(\mathcal O_{P+Q}\otimes\mathcal O_X(a))$ is the same as that of the linear forms $\ell$ in $H^0(\mathcal O_{P+Q}\otimes\mathcal O_X(1))$ up to a unit. But linear forms on $X\subseteq\mathbf P^n$ separate points and tangent vectors ($\mathcal O_X(1)$ is very ample), i.e.\ they surject onto $H^0(\mathcal O_{P+Q}\otimes\mathcal O_X(1))$. Hence $H^0(\mathcal O_X(a))\to H^0(\mathcal O_{P+Q}\otimes\mathcal O_X(a))$ is surjective, so $K$ is very ample, contradicting Prop.~5.2 since $X$ is hyperelliptic. \hfill$\square$

\exercise{5.2}{If $X$ is a curve of genus $\ge2$ over a field of characteristic $0$, show that the group $\operatorname{Aut}X$ of automorphisms of $X$ is finite. [\emph{Hint}: If $X$ is hyperelliptic, use the unique $g^1_2$ to show that $\operatorname{Aut}X$ permutes the ramification points of the $2$-fold covering $X\to\mathbf P^1$. If $X$ is not hyperelliptic, show that $\operatorname{Aut}X$ permutes the hyperosculation points (Ex.~4.6) of the canonical embedding.] Cf.\ (Ex.~2.5).}
\pf
\textbf{A fixed point bound.} \emph{If $\sigma\ne\mathrm{id}$ is an automorphism of $X$ (genus $g\ge2$), then $\sigma$ has at most $2g+2$ fixed points.} The fixed set is closed and proper (an equalizer of two morphisms to a separated curve), hence finite, so there is $Q$ with $\sigma Q\ne Q$. By Riemann--Roch $h^0((g+1)Q)\ge2$, so there is a nonconstant $f$ with poles only at $Q$, of order $m\le g+1$ (a nonconstant function has a pole). Then $f\circ\sigma$ has its only pole at $\sigma^{-1}Q\ne Q$, so $h=f-f\circ\sigma\ne0$ has poles of total degree $2m\le2g+2$, hence at most $2g+2$ zeros. If $P\ne Q$ is fixed by $\sigma$ then $P\ne\sigma^{-1}Q$ (otherwise $Q=\sigma P=P$), $f$ is regular at $P$ and $h(P)=f(P)-f(\sigma P)=0$. So $\sigma$ has at most $2g+2$ fixed points besides possibly $Q$ (which is not fixed). \hfill$(\ast)$

\textbf{Hyperelliptic case.} Let $f:X\to\mathbf P^1$ be of degree 2 and $L=f^*\mathcal O(1)$. Any $g^1_2$ is complete (Clifford: a $g^r_2$ has $r\le1$, so $|L|$ is the complete system), and by Prop.~5.3 the $g^1_2$ is unique. If $\sigma\in\operatorname{Aut}X$ then $\sigma^*|L|$ is a $g^1_2$, so $\sigma^*L\cong L$ and $\sigma^*$ acts on the $2$-dimensional space $H^0(L)$ (after choosing an isomorphism) by some $A\in GL_2$ up to scalar; thus $f\circ\sigma=\tau_\sigma\circ f$ with $\tau_\sigma\in PGL_2$ uniquely determined (because $f$ is surjective), and $\sigma\mapsto\tau_\sigma$ is a homomorphism $\operatorname{Aut}X\to\operatorname{Aut}\mathbf P^1$. Let $R\subseteq X$ be the ramification points of $f$ and $B=f(R)$. By Hurwitz (II, 2.4 / the formula $2g-2=2(-2)+\deg R$, valid in characteristic $0$) $|B|=2g+2\ge6$. $f$ is ramified at $P$ iff $f\circ\sigma=\tau_\sigma f$ is ramified at $\sigma^{-1}P$, so $\sigma$ permutes $R$ and $\tau_\sigma$ permutes $B$. An element of $\operatorname{Aut}\mathbf P^1$ fixing three points is the identity, hence the image of $\sigma\mapsto\tau_\sigma$ embeds into the finite group $\operatorname{Sym}(B)$. The kernel consists of $\sigma$ with $f\circ\sigma=f$, i.e.\ of elements of the Galois group of the degree-$2$ extension $K(X)/K(\mathbf P^1)$: at most $2$ elements. So $|\operatorname{Aut}X|\le2\,(2g+2)!$.

\textbf{Nonhyperelliptic case} ($g\ge3$; for $g=2$ the curve is hyperelliptic). Let $W=\{P\in X: h^0(K-gP)>0\}$, i.e.\ the points through which some canonical divisor has multiplicity $\ge g$. In the canonical embedding $X\subset\mathbf P^{g-1}$ (Prop.~5.2) this says that some hyperplane meets $X$ at $P$ with multiplicity $\ge g=n+1$ ($n=g-1$): the hyperosculation points of Ex.~4.6(b). If $\sigma\in\operatorname{Aut}X$ then $\sigma^*K\sim K$, so $\sigma$ maps canonical divisors to canonical divisors, and therefore $\sigma(W)=W$. Hence there is a homomorphism $\operatorname{Aut}X\to\operatorname{Sym}(W)$ once we know $W$ is finite, and its kernel consists of automorphisms with $|W|$ fixed points. By $(\ast)$ the kernel is trivial as soon as $|W|\ge2g+3$. So it suffices to show: $W$ is finite and $|W|\ge2g+3$.

Let $P\in X$ and $0=a_0<a_1<\dots<a_{g-1}$ the vanishing orders at $P$ of the $g$ sections of $H^0(K)$; $h^0(K-aP)=\#\{i:a_i\ge a\}$. Define the weight $w(P)=\sum_i(a_i-i)\ge0$; $P\in W$ iff $a_{g-1}\ge g$ iff $w(P)>0$. By Ex.~4.6(b) (here $n=g-1$, $d=2g-2$, so $n(n+1)(g-1)+(n+1)d=g(g-1)(g+1)$ hyperosculation points counted with weights), $\sum_Pw(P)=g(g-1)(g+1)$. For completeness, this is the following Wronskian argument: if $s_1,\dots,s_g$ is a basis of $H^0(\omega_X)$ and $s_j=f_j\,dz$ in a local parameter $z$, the Wronskian $\det(d^if_j/dz^i)_{0\le i<g}$ is a global section of $\omega_X^{\otimes g(g+1)/2}$ (it is multiplied by the right power of the Jacobian under a change of parameter), of degree $\frac12g(g+1)(2g-2)=g(g-1)(g+1)$; in characteristic $0$ it is not identically $0$ because the $f_j$ are linearly independent over $k$, and its order of vanishing at $P$ is $\sum_i(a_i-i)=w(P)$. In particular $W$ is finite and nonempty.

\emph{Weight bound.} For each $j\ge0$ with $a:=a_{g-1-j}$, the divisor $D=K-aP$ is effective and special (as $K-D=aP$ is effective) with $\dim|D|=j$ and $\deg D=2g-2-a$. By Clifford (Thm.~5.4) $j\le\frac12(2g-2-a)$, i.e.\ $a_{g-1-j}\le2g-2-2j$, i.e.\ $a_i\le2i$ for all $i$. So $w(P)\le\sum i=g(g-1)/2$. If $a_i=2i$ for some $1\le i\le g-2$, then for $D=K-a_iP$ we have equality in Clifford, with $D\ne0$ ($\deg D=2g-2-2i>0$) and $D\ne K$; by Thm.~5.4, $X$ would be hyperelliptic, which it is not. So $a_i\le2i-1$ for every $1\le i\le g-2$ (and $g\ge3$ so this range is nonempty), whence $w(P)\le g(g-1)/2-1$ for every $P$ (in particular for $P\in W$). (For $g=3$ this is $2$.) Therefore
\[
|W|\ \ge\ \frac{g(g-1)(g+1)}{g(g-1)/2-1}\ >\ \frac{g(g-1)(g+1)}{g(g-1)/2}=2g+2,
\]
and $|W|\ge2g+3$. (For $g=3$ this reads $|W|\ge12$; indeed the flexes of a smooth quartic are $24$ counted with weight.) This proves finiteness of $\operatorname{Aut}X$. \hfill$\square$

\exercise{5.3}{\emph{Moduli of Curves of Genus 4.} The hyperelliptic curves of genus $4$ form an irreducible family of dimension $7$. The nonhyperelliptic ones form an irreducible family of dimension $9$. The subset of those having only one $g^1_3$ is an irreducible family of dimension $8$. [\emph{Hint}: Use (5.2.2) to count how many complete intersections $Q\cap F_3$ there are.]}
\pf
Throughout $\operatorname{char}k=0$ (Ex.~5.2: finiteness of $\operatorname{Aut}$, Bertini), and ``dimension of a family'' is the dimension of its image in the (coarse) moduli space $\mathfrak M_4$, as in Examples 5.5.5, 5.5.6.

\textbf{Hyperelliptic.} By Example 5.5.5 a hyperelliptic curve of genus $g$ is the double cover of $\mathbf P^1$ branched in $2g+2$ distinct points, and by Prop.~5.3 the $g^1_2$ and hence the branch set $B$ is determined by $X$ up to $\operatorname{Aut}\mathbf P^1$ (and conversely $B$ determines $X$: $y^2=\prod(x-b_i)$). The space of branch sets is an open subset of $\operatorname{Sym}^{10}\mathbf P^1=\mathbf P^{10}$ (irreducible, dim $10$) for $g=4$, and $PGL_2$ (dim 3) acts with finite stabilizers (a set of $\ge3$ points). So the family has dimension $10-3=7$ and is irreducible (image of an irreducible variety).

\textbf{Nonhyperelliptic.} By Example 5.2.2, the canonical models are exactly the nonsingular complete intersections $X=Q\cap F$ of a quadric $Q$ and a cubic $F$ in $\mathbf P^3$ (conversely such $X$ has $\omega_X=\mathcal O_X(2+3-4)=\mathcal O_X(1)$, degree $6=2g-2$, $g=4$, and $H^0(\mathcal O_{\mathbf P^3}(1))\to H^0(\omega_X)$ is injective of dimension $4=g$, so it is the canonical embedding). $X$ determines $Q$ (unique quadric through $X$, 5.2.2). By the Koszul resolution $0\to\mathcal O(-5)\to\mathcal O(-2)\oplus\mathcal O(-3)\to\mathcal I_X\to0$ and $H^1(\mathcal O_{\mathbf P^3}(-2))=0$ we get $H^0(\mathcal I_X(3))=Q\cdot H^0(\mathcal O(1))\oplus kF$ of dimension $5$; so $F$ is determined by $X$ modulo $Q\cdot(\text{linear forms})$ and scalars. Hence the set of such curves is parametrised by an open subset $\mathcal U$ (nonsingularity is open, nonempty by Bertini) of the $\mathbf P^{15}$-bundle over $\mathbf P^9$ whose fibre over $Q$ is $\mathbf P(S_3/Q\cdot S_1)$, $\dim S_3=20$. So $\mathcal U$ is irreducible of dimension $9+15=24$. Two such curves are isomorphic iff they differ by an element of $PGL_4$ (canonical embedding is unique up to $PGL_4$), and stabilizers are the finite groups $\operatorname{Aut}X$ (Ex.~5.2). So $\mathcal U\to\mathfrak M_4$ has fibres of dimension $\dim PGL_4=15$, and the image is irreducible of dimension $24-15=9$.

\textbf{Exactly one $g^1_3$.} By Example 5.5.2, $X=Q\cap F$ has exactly two $g^1_3$ if $Q$ is nonsingular and exactly one if $Q$ is a quadric cone. Quadric cones form the (irreducible) determinant hypersurface of rank $\le3$ quadrics in $\mathbf P^9$, of dimension $8$. Such a $Q$ with vertex $V$: for $F$ general, nonsingular $X=Q\cap F$ exist: $|\mathcal O_Q(3)|$ is base point free so a general $F$ misses $V$, and by Bertini (II 8.18) a general $F\cap Q$ is nonsingular on $Q\setminus V$. (Note: if $V\in X$ then $X$ is singular at $V$, since $Q$ is singular there; so for nonsingular $X$ the vertex is not on $X$.) The same count gives an irreducible family of dimension $8+15=23$ in the parameter space, and $23-15=8$ in $\mathfrak M_4$. \hfill$\square$

\exercise{5.4}{Another way of distinguishing curves of genus $g$ is to ask, what is the least degree of a birational plane model with only nodes as singularities (3.11)? Let $X$ be nonhyperelliptic of genus $4$. Then:
\begin{enumerate}[label=(\alph*),leftmargin=*]
\item if $X$ has two $g^1_3$'s, it can be represented as a plane quintic with two nodes, and conversely;
\item if $X$ has only one $g^1_3$, then it can be represented as a plane quintic with a tacnode (I, Ex.~5.14d), but the least degree of a plane representation with only nodes is $6$.
\end{enumerate}}
\pf
Here ``represented as a plane curve $X'$'' means a birational morphism $\nu:X\to X'\subset\mathbf P^2$ (so $X$ is the normalization). We use: for an integral plane curve $X'$ of degree $e$ with normalization of genus $g$, $g=\frac12(e-1)(e-2)-\sum\delta_x$, $\delta_x=\operatorname{length}(\tilde{\mathcal O}_x/\mathcal O_x)$ (Ex.~1.8); a singular point consisting of two smooth branches has $\delta$ equal to their intersection multiplicity, so $\delta=1$ iff it is a node. Since $X$ is nonhyperelliptic of genus $4$ every $g^1_3$ is complete and base-point free (a base point would give a $g^1_2$; Clifford gives $\dim\le1$).

Write $X\subset Q\cap F\subset\mathbf P^3$ for the canonical model (5.2.2), $K=H$ the hyperplane class.

\textbf{Preliminary: projection from a point.} For $P\in X$, $|K-P|$ has dimension $2$ and is base point free: $h^0(K-P-Q')=2$ for all $Q'$, since $h^0(P+Q')=1$ ($X$ is nonhyperelliptic) and Riemann--Roch give $h^0(K-P-Q')=h^0(P+Q')-(2+1-4)=2$. It is the projection $\psi_P$ from $P$ of the canonical curve to $\mathbf P^2$; it is birational onto a quintic $X_P$ (the image has degree $m$ with $m\mid5$ and cannot be a line, so $m=5$). $p_a(X_P)=6$, so $\sum\delta=2$. Two points $R,R'\ne P$ have the same image iff $P,R,R'$ are collinear, and $\psi_P(P)=\psi_P(R)$ iff $R$ lies on the tangent line $T_P$; $\psi_P$ fails to be an immersion at $R\ne P$ iff $T_R\ni P$, at $P$ iff $T_P$ meets $X$ with multiplicity $\ge3$ at $P$. A line meeting $X=Q\cap F$ in $\ge3$ points (counted with multiplicity) meets $Q$ in $\ge3$ points, hence lies in $Q$.

\textbf{(a)} Suppose $X$ has two $g^1_3$. By Example 5.5.2, $Q$ is nonsingular, $Q\cong\mathbf P^1\times\mathbf P^1$, and $X$ has type $(3,3)$ on $Q$ (II 6.6.1); the two $g^1_3$ are cut out by the two rulings. Choose $P\in X$ such that $P$ is not a ramification point of either projection $p_i:X\to\mathbf P^1$ (finitely many exceptions, char $0$) and such that the two ruling lines $L_1,L_2$ through $P$ meet $X$ in three distinct points each: $L_i\cap X=\{P,Q_i,R_i\}$ (the fibre of $p_i$ through $P$ is reduced for $P$ outside the finitely many branch fibres). Since $L_1\cap L_2=\{P\}$, the points $Q_1,R_1,Q_2,R_2$ are distinct and different from $P$. Also $T_P\ne L_i$ (transversality), $T_P\subset T_PQ$, and $T_P\cap Q=\{P\}$ with multiplicity $2$, so $T_P$ meets $X$ only at $P$ (a double point there). Lines through $P$ meeting $X$ in two more points or tangent elsewhere lie in $Q$, so are $L_1,L_2$, which meet $X$ transversally. Consequently $\psi_P$ is an immersion everywhere (at $P$: if $T_P$ met $X$ with multiplicity $3$ at $P$, then, as $Q(\gamma(t))=0$ along a local parametrization with the line $T_P$ contained in $T_PQ$, $Q|_{T_P}$ vanishes to order $\ge3$ and being quadratic it vanishes identically, so $T_P\subset Q$, contradiction), injective except $\psi_P(Q_i)=\psi_P(R_i)$, $i=1,2$. So $X_P$ has exactly two singular points, each with two smooth branches, hence $\delta\ge1$ at each; as $\sum\delta=2$, both are nodes. So $X$ is a plane quintic with two nodes $N_1,N_2$.

Conversely, let $\nu:X\to X'$ be a quintic with two nodes $N_1,N_2$ and $H$ the pullback of the line class. If $\nu^{-1}(N_i)=\{N_i',N_i''\}$, the linear forms vanishing at $N_i$ give a $2$-dimensional subspace of $H^0(H-N_i'-N_i'')$ (degree $3$), i.e.\ a $g^1_3$ $D_i$. If $D_1\sim D_2$ then $N_1'+N_1''\sim N_2'+N_2''$ is a divisor of degree $2$ with $h^0\ge2$ unless the two divisors coincide ($h^0=1$ since $X$ is nonhyperelliptic), which is impossible as $N_1\ne N_2$. So $X$ has two distinct $g^1_3$.

\textbf{(b)} Suppose $X$ has only one $g^1_3$. By Example 5.5.2, $Q$ is a quadric cone with vertex $V$, $V\notin X$ (see 5.3), and the lines of the cone cut out the $g^1_3$ ($X$ meets each line of the cone in three points). Take $P\in X$ general so that the line $L$ of the cone through $P$ and $V$ meets $X$ in three distinct points $P,Q',R'$ transversally and $T_P\ne L$. The only lines in $Q$ through $P$ is $L$. As in (a), $\psi_P$ is an immersion everywhere, injective except $\psi_P(Q')=\psi_P(R')$. So $X_P$ has a unique singular point $c$, with two smooth branches (images of neighbourhoods of $Q',R'$), and $\delta_c=2$, i.e.\ the two smooth branches have intersection multiplicity $2$, hence are tangent, and the singularity is analytically $y(y-x^2)=0$: this is a tacnode (I, Ex.~5.14d). So $X$ is a plane quintic with a tacnode.

Least degree with only nodes. If $\nu:X\to X'$ is a birational plane model of degree $e$ with only nodes, then $4=\frac12(e-1)(e-2)-\#\text{nodes}$, so $\frac12(e-1)(e-2)\ge4$ forces $e\ge5$, and $e=5$ forces exactly two nodes, so by the converse in (a) $X$ would have two $g^1_3$, contrary to hypothesis. Hence $e\ge6$. On the other hand projecting the canonical sextic in $\mathbf P^3$ from a general point of $\mathbf P^3$ (Thm.~3.10 and the proof of 3.6) gives a birational morphism onto a plane sextic with only nodes. So the least degree is $6$. \hfill$\square$

\exercise{5.5}{\emph{Curves of Genus 5.} Assume $X$ is not hyperelliptic.
\begin{enumerate}[label=(\alph*),leftmargin=*]
\item The curves of genus $5$ whose canonical model in $\mathbf P^4$ is a complete intersection $F_2\cdot F_2\cdot F_2$ form a family of dimension $12$.
\item $X$ has a $g^1_3$ if and only if it can be represented as a plane quintic with one node. These form an irreducible family of dimension $11$. [\emph{Hint}: If $D\in g^1_3$, use $K-D$ to map $X\to\mathbf P^2$.]
\item[*(c)] In that case, the conics through the node cut out the canonical system (not counting the fixed points at the node). Mapping $\mathbf P^2\to\mathbf P^4$ by this linear system of conics, show that the canonical curve $X$ is contained in a cubic surface $V\subseteq\mathbf P^4$, with $V$ isomorphic to $\mathbf P^2$ with one point blown up (II, Ex.~7.7). Furthermore, $V$ is the union of all the trisecants of $X$ corresponding to the $g^1_3$ (5.5.3), so $V$ is contained in the intersection of all the quadric hypersurfaces containing $X$. Thus $V$ and the $g^1_3$ are unique.
\end{enumerate}
\emph{Note.} Conversely, if $X$ does not have a $g^1_3$, then its canonical embedding is a complete intersection, as in (a). More generally, a classical theorem of Enriques and Petri shows that for any nonhyperelliptic curve of genus $g\ge3$ the canonical model is projectively normal and is an intersection of quadric hypersurfaces unless $X$ has a $g^1_3$ or $g=6$ and $X$ has a $g^2_5$ (see Saint-Donat [1]). (Starred in the book: part (c).)}
\pf
Again $\operatorname{char}k=0$. (Only the dimension counts need this.)

\textbf{(a)} A complete intersection $X=Q_1\cap Q_2\cap Q_3\subset\mathbf P^4$ of three quadrics has $\omega_X=\mathcal O_X(2+2+2-5)=\mathcal O_X(1)$ (II, Ex.~8.4), so $\deg\omega_X=8=2g-2$, $g=5$. The map $H^0(\mathcal O_{\mathbf P^4}(1))\to H^0(\omega_X)$ is injective ($X$ nondegenerate, as $I_X(1)=0$ by the Koszul complex) between spaces of dimension $5=g$, so $X\subset\mathbf P^4$ is the canonical embedding (and $X$ is nonhyperelliptic since $K$ is very ample). By the Koszul complex $H^0(\mathcal I_X(2))=\langle Q_1,Q_2,Q_3\rangle$ is $3$-dimensional, so $X$ determines a point of the Grassmannian $G(3,15)$ of $3$-dimensional subspaces of the $15$-dimensional space of quadratic forms, $\dim G(3,15)=3\cdot12=36$. Nets with nonsingular intersection form a nonempty open subset $\mathcal U$ (Bertini applied to the Veronese embedding $v_2(\mathbf P^4)\subset\mathbf P^{14}$ three times), irreducible of dimension $36$. Two such curves are isomorphic iff their nets are $PGL_5$-conjugate (uniqueness of the canonical embedding up to projective equivalence), and stabilizers are finite (Ex.~5.2). So the image in $\mathfrak M_5$ is irreducible of dimension $36-24=12$ (which is $\dim\mathfrak M_5=3g-3$, i.e.\ these curves are the general curves of genus $5$).

\textbf{(b)} \emph{($\Leftarrow$)} Suppose $\nu:X\to X'\subset\mathbf P^2$ is a plane quintic with a node $N$, $\nu^{-1}(N)=\{N',N''\}$, $H$ the line class. The pencil of lines through $N$ gives a $2$-dimensional subspace of $H^0(H-N'-N'')$, of degree $3$: so $X$ has a $g^1_3$.

\emph{($\Rightarrow$)} Let $D$ be a $g^1_3$; it is complete with $h^0(D)=2$ (Clifford, $\dim|D|\le3/2$). $\deg(K-D)=5$ and Riemann--Roch gives $h^0(D)-h^0(K-D)=3+1-5$, so $h^0(K-D)=3$: $|K-D|$ is a $g^2_5$. It has no base points: if $P$ were one, $h^0(K-D-P)=3$ and Riemann--Roch gives $h^0(D+P)=h^0(K-D-P)+4+1-5=3$, so $D+P$ is a $g^2_4$, contradicting Clifford (equality $\dim|D+P|=\frac12\deg$ forces $D+P=0$, $K$ or hyperelliptic). So $\psi=\varphi_{|K-D|}:X\to\mathbf P^2$ is a morphism, nondegenerate (the system is complete with $h^0=3$), of image $X'$ of degree $m$ with $m\cdot\deg\psi=5$ and $m\ge2$: so $m=5$, $\psi$ is birational onto a plane quintic $X'$. $p_a(X')=6$, so $\sum\delta=1$: $X'$ has exactly one singular point, a double point with $\delta=1$, i.e.\ a node or an ordinary cusp.

\emph{Remark (honest correction).} The cusp can occur, so the literal ``only if'' needs ``node or ordinary cusp''. Let $X'$ be a plane quintic with an ordinary cusp $c$ and no other singularity (such quintics exist: a general member of the linear system of quintics with an ordinary cusp at a given point, by Bertini), and let $X$ be its normalization, of genus $5$; $c$ has one preimage $P_0$. The lines through $c$ cut a base-point-free $g^1_3$ $D=H-2P_0$, so by Lemma~A(i) with $(a,b)=(2,3)$ $X$ is not hyperelliptic. For a plane quintic whose only singularity is a double point of $\delta=1$ one has $K=2H-C$ with $C$ the conductor divisor ($N'+N''$ for a node, $2P_0$ for a cusp; adjunction on the plane), so in both cases $D=H-C$ and $K-D=H$: the model is $\varphi_{|K-D|}$ for the pencil $D$ given by the lines through the double point. Now if this $X$ also had a nodal quintic model, its pencil of lines through the node would be a second base-point-free $g^1_3$ $D'$, equal to $D$ by Lemma~A(ii) ($g=5$); then both models are $\varphi_{|K-D|}$, hence projectively equivalent, which is absurd (one has a node, the other a cusp). So such $X$ has a $g^1_3$ but no nodal quintic model: the nodal case is the general one (the cuspidal ones form a family of dimension $16+2-8=10$: quintics with an ordinary cusp at a given point form a $\mathbf P^{16}$, the cusp point varies in $\mathbf P^2$, and $\dim PGL_3=8$). In (c) below the arguments are identical for the cusp (conics through the cusp), and I assume the singular point is a node as in the statement.

\emph{Dimension $11$.} The plane quintics with a node at a point $N$ form a linear system of dimension $21-3=18$, i.e.\ $\mathbf P^{17}$; as $N$ varies we get an irreducible variety (a $\mathbf P^{17}$-bundle over $\mathbf P^2$) of dimension $19$, in which the curves with exactly one node, an ordinary one, and nonsingular elsewhere form a nonempty open subset (nonempty by Bertini on the blow-up of $N$ with the very ample system $|5L-2E|$). Isomorphism classes of the normalizations correspond to $PGL_3$-orbits: if $X'_1,X'_2$ have isomorphic normalizations $X$, the $g^1_3$ is unique (Lemma~A(ii), $g=5$), so the model is $\varphi_{|K-D|}$ for the same $D$, hence unique up to $PGL_3$; stabilizers are finite (Ex.~5.2). So the family is irreducible of dimension $19-8=11$.

\textbf{(c)} Let $\nu:X\to X'$ be a nodal quintic as in (b), $H=K-D$, $\nu^{-1}(N)=\{N',N''\}$. The pencil of lines through $N$ gives a $g^1_3$ $D'\subseteq|H-N'-N''|$; by Lemma~A(ii) $D'=D$ (unique $g^1_3$, $g=5$), so $K=H+D=2H-N'-N''$. Conics through $N$ form a $5$-dimensional space which injects into $H^0(\mathcal O_X(2H-N'-N''))=H^0(K)$ ($X'$ is not on a conic as it is an irreducible quintic), which has dimension $5$; hence the conics through $N$ cut out exactly $|K|$ outside $N',N''$.

Put $\mathbf P^2$ with coordinates $x,y,z$ and $N=(0:0:1)$. The conics through $N$ are spanned by $x^2,xy,y^2,xz,yz$, and $\theta=(x^2:xy:y^2:xz:yz):\mathbf P^2\dashrightarrow\mathbf P^4$ is the canonical map composed with $\nu$. Let $\pi:\tilde P\to\mathbf P^2$ be the blow-up at $N$, with exceptional curve $E$, line class $L$; then $\theta$ extends to $\tilde\theta:\tilde P\to\mathbf P^4$ given by $|2L-E|$ (the blow-up resolves the base point). As $\nu$ separates the two branches at the node, $\nu$ lifts to $\tilde\nu:X\to\tilde P$ onto the proper transform $\tilde X'$, of class $5L-2E$, and $\tilde\theta\circ\tilde\nu$ is the canonical map. Put $V=\tilde\theta(\tilde P)\supseteq X$.

\emph{$V$ is the cubic scroll, $\tilde\theta$ is an isomorphism onto $V$.} Write $(a_0,\dots,a_4)=(x^2,xy,y^2,xz,yz)$. The three quadrics $a_0a_2-a_1^2$, $a_0a_4-a_1a_3$, $a_1a_4-a_2a_3$ (the $2\times2$ minors of $\begin{pmatrix}a_0&a_1&a_3\\a_1&a_2&a_4\end{pmatrix}$) vanish on $\theta$, hence on $V$. Conversely let $a$ be a point where all minors vanish. If $(a_3,a_4)=(0,0)$ then $a_0a_2=a_1^2$, $(a_0:a_1:a_2)=(x^2:xy:y^2)$ and $a=\theta(x:y:0)$. If $(a_3,a_4)=(u,v)\ne0$, the minors give $va_0=ua_1$, $va_1=ua_2$, so $(a_0,a_1,a_2)=t(u^2,uv,v^2)$; for $t\ne0$, $a=\theta(u:v:1/t)$, and for $t=0$, $a=(0:0:0:u:v)$ is the limit point of $\theta$ along the exceptional direction. So $V=V(\text{minors})$ as a set, and $V$ is the union of $\theta(\mathbf P^2\setminus N)$ and the line $\{a_0=a_1=a_2=0\}=\tilde\theta(E)$. Now $\tilde\theta$ is injective on points by this description (a point $a$ with $(a_3,a_4)=(u,v)\ne0$ determines $(u:v)$ and $t$, hence the point $(u:v:1/t)$ of $\mathbf P^2\setminus N$, or, if $t=0$, the direction $(u:v)$ at $N$; a point with $(a_3,a_4)=0$ determines $(x:y)$). It is an immersion: on $\mathbf P^2\setminus N$ use the charts $x=1$ and $y=1$ (which do not contain $N$): in the chart $x=1$ the map is $(y,z)\mapsto(1:y:y^2:z:yz)$, and $a_1/a_0=y$, $a_3/a_0=z$ are local coordinates on the target, so it is an immersion (likewise for $y=1$); the two charts $x=1$, $y=1$ cover $\mathbf P^2\setminus N$. Along the exceptional curve use the blow-up chart $y=xs$ (coordinates $(x,s)$, $E=\{x=0\}$), where the map reads $(x,s)\mapsto(x:xs:xs^2:1:s)$, i.e.\ $(x,xs,xs^2,s)$ in the affine chart $a_3=1$, with Jacobian columns $(1,s,s^2,0)$ and $(0,x,2xs,1)$ always independent; the chart $x=ys'$ is symmetric. A proper morphism that is injective and immersive is a closed immersion, so $\tilde\theta:\tilde P\to V$ is an isomorphism and $V\cong\mathbf P^2$ blown up at a point; $\deg V=(2L-E)^2=3$.

\emph{$V$ is the union of the trisecants.} For a line $\ell\ni N$ in $\mathbf P^2$ let $\tilde\ell$ be its proper transform, of class $L-E$. Then $(2L-E)\cdot(L-E)=1$, so $\tilde\theta(\tilde\ell)$ is a line $\Lambda_\ell$ in $\mathbf P^4$, and $\tilde X'\cdot\tilde\ell=(5L-2E)(L-E)=3$, so $\Lambda_\ell$ meets $X$ in a divisor of degree $3$: it is the trisecant spanned by the members of the divisor of the $g^1_3$ cut by the pencil of lines through $N$ (the divisor $D'$ above; compare 5.5.3). The lines $\Lambda_\ell$, together with the line $\tilde\theta(E)$ (the image of the directrix, limit of $\Lambda_\ell$), sweep out $V$: $V=\bigcup_\ell\Lambda_\ell$. If $Q'$ is any quadric containing $X$, then $\tilde\theta^*Q'$ is a section of $\mathcal O_{\tilde P}(2(2L-E))$ vanishing on $\tilde X'$; its restriction to $\tilde\ell\cong\mathbf P^1$ is a section of $\mathcal O(2)$ vanishing on the degree-$3$ divisor $\tilde X'\cap\tilde\ell$, hence is zero. So $\Lambda_\ell\subseteq Q'$ for every $\ell$, hence $V\subseteq Q'$. Thus $V\subseteq\bigcap_{Q'\supseteq X}Q'$. Since the three minors above are quadrics containing $V\supseteq X$ and cut out $V$ set-theoretically, the intersection of all quadrics containing $X$ \emph{equals} $V$ (as a set). So $V$ depends only on $X\subset\mathbf P^4$, and the $g^1_3$ is unique by Lemma~A(ii). \hfill$\square$

\exercise{5.6}{Show that a nonsingular plane curve of degree $5$ has no $g^1_3$. Show that there are nonhyperelliptic curves of genus $6$ which cannot be represented as a nonsingular plane quintic curve.}
\pf
Let $X\subset\mathbf P^2$ be a nonsingular quintic, $H$ the line class. Then $g=6$ and $\omega_X=\mathcal O_X(2)$ (II, 8.20.3 / Ex.~8.4). The restriction $H^0(\mathbf P^2,\mathcal O(2))\to H^0(\mathcal O_X(2))=H^0(K)$ is injective ($X$ not in a conic) between spaces of dimension $6$, hence bijective: \emph{sections of $K$ are the conics}.

Suppose $D$ is a $g^1_3$; by Clifford $D$ is complete with $h^0(D)=2$ (and we may replace $D$ by an effective member). Riemann--Roch: $h^0(K-D)=h^0(D)-(3+1-6)=4$. By the above, $h^0(K-D)$ is the dimension of the space of conics containing the scheme $D\subset X\subset\mathbf P^2$ (i.e.\ whose restriction to $X$ vanishes on $D$). We show that conics through a degree-$3$ divisor $D$ on $X$ form a space of dimension exactly $3$, a contradiction.

\emph{Claim.} If $D$ is an effective divisor of degree $3$ on a nonsingular plane curve $X$, the conics containing $D$ form a $3$-dimensional space. Let $\ell_D$ be a line with $D\le\ell_D\cdot X$ if one exists. (i) If $D\le\ell\cdot X$ for a line $\ell$: a conic containing $D$ restricts to $\ell$ as a binary quadratic form vanishing on the degree-$3$ subscheme $D\subset\ell$, so it vanishes on $\ell$ and equals $\ell\cdot\ell'$; these form a $3$-dimensional space. (ii) Otherwise pick $P\in\operatorname{Supp}D$ and $D'=D-P$ (degree $2$); $D'$ lies on a unique line $\ell'$ and imposes $2$ independent conditions on conics (conics surject onto binary quadratics on $\ell'$, and then onto the length-$2$ subscheme), so conics containing $D'$ form a $4$-dimensional space. Since $D\not\le\ell'\cdot X$ (else we are in (i)), choose a line $\ell''$ missing $\operatorname{Supp}D$: then $\ell'\ell''$ contains $D'$ but not $D$ (its divisor on $X$ is $\ell'X+\ell''X$, which is $\ge D$ iff $\ell'X\ge D$). So conics containing $D$ form a hyperplane of the $4$-dimensional space: dimension $3$. This proves the claim and hence that $X$ has no $g^1_3$.

\emph{Second part.} A nonsingular curve of type $(3,4)$ on $Q=\mathbf P^1\times\mathbf P^1$ exists (the system $|\mathcal O_Q(3,4)|$ is very ample, so a general member is nonsingular by Bertini, and it is connected, hence irreducible, being an ample divisor on a nonsingular surface, III 7.9), and has genus $(3-1)(4-1)=6$ (III, Ex.~5.6). Its projections to the two factors are base-point-free pencils of degrees $3$ and $4$ (intersection with the rulings), in particular $X$ has a $g^1_3$. $X$ is not hyperelliptic: a base-point-free $g^1_2$ together with the base-point-free $g^1_3$ would give by Lemma~A(i) $g\le(2-1)(3-1)=2$, but $g=6$. So $X$ is nonhyperelliptic of genus $6$ with a $g^1_3$, hence by the first part it cannot be represented as a nonsingular plane quintic. \hfill$\square$

\exercise{5.7}{\begin{enumerate}[label=(\alph*),leftmargin=*]
\item Any automorphism of a curve of genus $3$ is induced by an automorphism of $\mathbf P^2$ via the canonical embedding.
\item[*(b)] Assume $\operatorname{char}k\ne3$. If $X$ is the curve given by
\[
x^3y+y^3z+z^3x=0,
\]
the group $\operatorname{Aut}X$ is the simple group of order $168$, whose order is the maximum $84(g-1)$ allowed by (Ex.~2.5). See Burnside [1, Section 232] or Klein [1].
\item[*(c)] Most curves of genus $3$ have no automorphisms except the identity. [\emph{Hint}: For each $n$, count the dimension of the family of curves with an automorphism $T$ of order $n$. For example, if $n=2$, then for suitable choice of coordinates, $T$ can be written $x\to-x$, $y\to y$, $z\to z$. Then there is an $8$-dimensional family of curves fixed by $T$; changing coordinates there is a $4$-dimensional family of such $T$, so the curves having an automorphism of degree $2$ form a family of dimension $12$ inside the $14$-dimensional family of all plane curves of degree $4$.]
\end{enumerate}
\emph{Note.} More generally it is true (at least over $\mathbf C$) that for any $g\ge3$, a ``sufficiently general'' curve of genus $g$ has no automorphisms except the identity (Baily [1]). (Starred in the book: parts (b), (c).)}
\pf
\textbf{(a)} Let $X$ have genus $g\ge2$ (in particular $g=3$), $V=H^0(X,\omega_X)$, $\varphi:X\to\mathbf P(V^\vee)$ the canonical morphism: $\varphi(P)$ is the hyperplane $V_P=\{s:s(P)=0\}\subset V$ (a base-point-free system, Lemma~5.1). An automorphism $\sigma$ acts on differentials by pull-back $\rho=\sigma^*:V\to V$, a linear isomorphism, and for $s\in V$: $(\sigma^*s)(P)=0\iff s(\sigma P)=0$. Hence $\rho(V_{\sigma P})=V_P$. Let $A$ be the projective linear map of $\mathbf P(V^\vee)$ (the space of hyperplanes of $V$) induced by $\rho^{-1}$ acting on subspaces of $V$: it sends the hyperplane $V_P$ to $\rho^{-1}(V_P)=V_{\sigma P}$. That is, $\varphi\circ\sigma=A\circ\varphi$. For $g=3$ nonhyperelliptic, $\varphi$ is an embedding onto a plane quartic (Example 5.2.1), and $\sigma$ is the restriction of $A\in PGL_3=\operatorname{Aut}\mathbf P^2$. (For hyperelliptic $X$ the same equation $\varphi\sigma=A\varphi$ holds with $\varphi$ the canonical $2:1$ map onto a conic.) \hfill$\square$

\textbf{(b)} \emph{Hypotheses.} The hypothesis printed in the book is $\operatorname{char}k\ne3$. I prove the statement in characteristic $0$ (hence over $\mathbf C$), which is what Ex.~2.5 and Ex.~5.2 (finiteness of $\operatorname{Aut}X$, the bound $84(g-1)$, Hurwitz) require. Two remarks on the printed hypothesis: the curve is singular in characteristic $7$ (the point $(1:2:4)$ is a singular point), so $\operatorname{char}k\ne7$ at least is needed for $X$ to be a genus-$3$ curve; and I have not tried to prove the statement in positive characteristic, where $|\operatorname{Aut}X|\le84(g-1)$ is not available in general.

Let $F=x^3y+y^3z+z^3x$. \emph{$X$ is nonsingular:} $F_x=3x^2y+z^3$, $F_y=x^3+3y^2z$, $F_z=y^3+3z^2x$. At a singular point $3x^2y=-z^3$, $3y^2z=-x^3$, $3z^2x=-y^3$; multiplying, $27x^3y^3z^3=-x^3y^3z^3$, so $28x^3y^3z^3=0$, $xyz=0$, say $x=0$, then $z=0$ (first equation), then $y=0$ (third); similarly in the other cases. So $X$ is a nonsingular quartic, $g=3$, nonhyperelliptic.

By (a) $G=\operatorname{Aut}X=\{A\in PGL_3:F\circ A=cF\text{ for a scalar }c\}$ (as $F$ is irreducible), a finite group with $|G|\le168$ (Ex.~5.2, 2.5).

\emph{Elements.} Let $\zeta$ be a primitive $7$-th root of unity. $S:(x,y,z)\mapsto(y,z,x)$ and $T=\operatorname{diag}(\zeta,\zeta^4,\zeta^2)$ lie in $G$ (check on monomials: $x^3y\mapsto\zeta^{3+4}=1$, $y^3z\mapsto\zeta^{12+2}=1$, $z^3x\mapsto\zeta^{6+1}=1$). $STS^{-1}=T^4$, so $H=\langle S,T\rangle$ has order $21$. Put $s_a=\zeta^a-\zeta^{-a}$ and $p=s_1$, $q=s_2$, $r=s_4$, and $M=\begin{pmatrix}p&q&r\\q&r&p\\r&p&q\end{pmatrix}$ (so $M_{ij}=s_{a_ia_j}$ with $(a_1,a_2,a_3)=(1,2,4)$ and indices mod $7$, $s_{-a}=-s_a$). We show $F(Mv)=49F(v)$. Let $\delta=p+q+r=\eta-\bar\eta$, $\eta=\zeta+\zeta^2+\zeta^4$; then $\delta^2=-7$ ($\eta+\bar\eta=-1$, $\eta\bar\eta=2$). Direct computation in $\mathbf Z[\zeta]$ using $\zeta^7=1$, $\sum\zeta^i=0$: $p^2+q^2+r^2=(\zeta^2+\zeta^5)+(\zeta^4+\zeta^3)+(\zeta+\zeta^6)-6=-7$, hence $pq+qr+rp=0$; $pq=\zeta^3+\zeta^4-\zeta-\zeta^6$, $qr=\zeta^6+\zeta-\zeta^5-\zeta^2$, $rp=\zeta^5+\zeta^2-\zeta^3-\zeta^4$, whence $pqr=\delta$. So $e_1=\delta$, $e_2=0$, $e_3=\delta$ and $t^3=\delta t^2+\delta$ for $t=p,q,r$. Moreover $u:=p^2q+q^2r+r^2p=-\delta$ (expand with $p^2=\zeta^2+\zeta^5-2$, $q^2=\zeta^3+\zeta^4-2$, $r^2=\zeta+\zeta^6-2$: $u=-\zeta-\zeta^2-\zeta^4+\zeta^3+\zeta^5+\zeta^6=-\delta$) and $v:=pq^2+qr^2+rp^2=-3\delta-u=-2\delta$.
The linear forms $X'=px+qy+rz$, $Y'=qx+ry+pz$, $Z'=rx+py+qz$ are permuted cyclically by $(x,y,z)\mapsto(y,z,x)$, so $F(X',Y',Z')$ is invariant under this cyclic permutation. Expanding, the coefficients are (cyclic classes): $x^4$: $p^3q+q^3r+r^3p$; $x^3y$: $p^3r+q^3p+r^3q+3(p^2q^2+q^2r^2+r^2p^2)$; $x^3z$: $p^4+q^4+r^4+3pqr(p+q+r)$; $x^2y^2$: $3(p^3r+q^3p+r^3q+pqr(p+q+r))$; $x^2yz$: $3(p+q+r)(u+pqr)$. Using $t^3=\delta(t^2+1)$: $p^3q+q^3r+r^3p=\delta(u+\delta)=0$; $p^3r+q^3p+r^3q=\delta(v+\delta)=-\delta^2=7$; $p^2q^2+q^2r^2+r^2p^2=e_2^2-2e_1e_3=-2\delta^2=14$; $p^4+q^4+r^4=-3\delta^2=21$ (Newton: $P_1=\delta$, $P_2=-7$, $P_3=-4\delta$, $P_4=21$). So the coefficients are $0$, $7+42=49$, $21+3\delta^2=0$, $3(7+\delta^2)=0$, $3\delta(-\delta+\delta)=0$. Hence $F(Mv)=49(x^3y+y^3z+z^3x)$. Also $M^2=-7I$: $(M^2)_{ac}=\sum_{b\in\{1,2,4\}}s_{ab}s_{bc}=\sum_{b\in(\mathbf Z/7)^*}\zeta^{b(a+c)}-\sum_{b\in(\mathbf Z/7)^*}\zeta^{b(a-c)}=7[a\equiv-c]-7[a\equiv c]=-7[a=c]$ (squares mod $7$ contain no $a,c$ with $a\equiv-c$). So $R=M/\delta\in G$ has $F\circ R=F$, $R^2=I$, and $R$ is not monomial (all entries nonzero).

\emph{$|G|=168$.} Let $P=\langle T\rangle$. The eigenvalues $\zeta,\zeta^4,\zeta^2$ of $T$ are distinct, so any $A$ with $ATA^{-1}\in\langle T\rangle$ (up to scalar) maps coordinate axes to coordinate axes: $N_G(P)$ consists of monomial matrices. A monomial $A=DP_\pi$ ($D$ diagonal, $P_\pi$ permutation) satisfies $F\circ A=cF$ iff $\pi$ preserves the set of monomials $\{x^3y,y^3z,z^3x\}$, i.e.\ $\pi\in\langle S\rangle$, and $D=\operatorname{diag}(a,b,c')$ satisfies $a^3b=b^3c'=c'^3a$; normalizing $a=1$: $b=\lambda$, $c'=\lambda^{-2}$, $\lambda^{-6}=\lambda$, i.e.\ $\lambda^7=1$: 7 solutions. So the monomial subgroup of $G$ is $H$, of order $21$, and $N_G(P)=H$ (it contains $H$ and is monomial). Since $R\in G$ is not monomial, $R\notin N_G(P)$, so $P$ is not normal and $n_7=|G:N_G(P)|=|G|/21\ge2$. By Sylow $n_7\equiv1\pmod7$ and $n_7$ divides $|G|/7\le24$: so $n_7=8$ and $|G|=168$.

\emph{Simplicity.} Facts: $N_G(P)=H$ of order $21$; $C_G(P)=P$ (in $H$, $S$ does not commute with $T$); $n_7=8$, so $G$ has $48$ elements of order $7$; every element of $G$ commuting with an element of order $7$ lies in the corresponding Sylow $7$-subgroup, so an element whose order is divisible by $7$ has order exactly $7$. Also $PGL_3$ has no subgroup $(\mathbf Z/2)^3$ (lift generators $A_i$ with $A_i^2=1$; commutators are scalars $\omega$ with $\omega^3=1$ (determinants) and $\omega^2=1$ (conjugating $A_j^2$), so the $A_i$ commute, are simultaneously diagonal with entries $\pm1$, and the image of $\{\pm1\}^3$ in $PGL_3$ has order $4$). Let $1\ne N\lhd G$. \emph{Case $7\nmid|N|$.} Then $P\cap N=1$, and $P$ acts on $N$ by conjugation; if $P$ centralized $N$ then $N\subseteq C_G(P)=P$, so $N=1$. So $P\to\operatorname{Aut}N$ is injective and $7\mid|\operatorname{Aut}N|$. The group $NP$ has $n_7(NP)=|N:N\cap H|\in\{|N|,|N|/3\}$ dividing $24$ and $\equiv1\pmod 7$ so equal to $1$ or $8$. If $1$, $P$ is normal in $NP$ and $[N,P]\subseteq N\cap P=1$, so $P$ centralizes $N$, contradiction. If $8$: $|N|=24$, or $|N|=8$. If $|N|=24$ then $G/N\cong\mathbf Z/7$, so all $144$ elements outside $N$ have order divisible by $7$, hence order $7$; but $G$ has only $48$ elements of order $7$: contradiction. If $|N|=8$, $7\mid|\operatorname{Aut}N|$ forces $N\cong(\mathbf Z/2)^3$ (the automorphism groups of the other groups of order $8$ have orders $4,8,8,24$), impossible by the above. \emph{Case $7\mid|N|$.} $N$ contains all $8$ Sylow $7$-subgroups, so $|N|\ge49$ and $|N|\in\{56,84,168\}$. $|N|=84$: $n_7(N)=8$ must divide $12$, no. $|N|=56$: the $48$ elements of order $7$ leave $8$ elements, forming the unique Sylow $2$-subgroup of $N$, characteristic in $N$ hence normal in $G$ of order $8$, contradicting the first case. So $N=G$: $G$ is simple of order $168$ (the unique such group, $PSL_2(\mathbf F_7)\cong GL_3(\mathbf F_2)$), and $168=84(3-1)$. \hfill$\square$

\textbf{(c)} Let $X$ be nonhyperelliptic of genus $3$ (the hyperelliptic curves form a $5$-dimensional family in the $6$-dimensional $\mathfrak M_3$, Example 5.5.6, so they are not general). By (a) and Example 5.5.6, $X$ is a nonsingular quartic and an automorphism is a linear map $A\in PGL_3$ with $F\circ A=cF$. We show that the quartics with a nontrivial automorphism lie in a finite union of constructible sets of dimension $\le13$ in $\mathbf P^{14}$, so the general quartic has trivial automorphism group. (Then in $\mathfrak M_3$ the locus has dimension $\le13-8=5<6$.)

By Ex.~5.2 and 2.5 ($|\operatorname{Aut}X|\le168$) a nontrivial automorphism group contains an element $T$ of prime order $p\le168$, so finitely many $p$. In characteristic $0$ such $T\in PGL_3$ is diagonalizable: $T=\operatorname{diag}(1,\varepsilon^a,\varepsilon^b)$, $\varepsilon$ a primitive $p$-th root of unity; this gives finitely many conjugacy classes of $T$ (types). For a given type let $C\subset PGL_3$ be its conjugacy class and $\mathrm{Cent}$ the centralizer; $\dim C=8-\dim\mathrm{Cent}=9-\sum m_i^2$ where $m_i$ are the multiplicities of the distinct eigenvalues ($\dim\mathrm{Cent}=\sum m_i^2-1$): $\dim C=6$ if the eigenvalues are distinct and $4$ if two coincide. The quartics with $F\circ T=cF$ are those whose form lies in an eigenspace $E_\chi$ of $T$ on the $15$-dimensional space of quartic forms: the monomial $x^iy^jz^k$ has eigenvalue $\varepsilon^{aj+bk}$. The union over conjugates $gTg^{-1}$ of the curves invariant under $T$ is the image of $PGL_3\times\mathbf P(E_\chi)\to\mathbf P^{14}$, of dimension $\le\dim C+\dim E_\chi-1$. We check this is $\le13$, i.e.\ $\dim E_\chi\le14-\dim C$:

\begin{itemize}
\item Two equal eigenvalues, $T\sim\operatorname{diag}(1,1,\varepsilon^c)$, need $\dim E_\chi\le10$: the eigenvalue is $\varepsilon^{ck}$, $k$ the $z$-degree. For $p=2$: $k$ even gives $5+3+1=9$ monomials, $k$ odd gives $4+2=6$ (so $\dim E_\chi\le9$: the $8$-dimensional family of the hint, $8+4=12$). For $p=3$: $k\equiv0$: $5+2=7$; $k\equiv1$: $4+1=5$; $k\equiv2$: $3$. For $p\ge5$ the $k=0,\dots,4$ are distinct mod $p$ so $\dim E_\chi\le5$.
\item Distinct eigenvalues ($p\ge3$), need $\dim E_\chi\le8$: for $p=3$, $T=\operatorname{diag}(1,\omega,\omega^2)$ and the eigenvalue of $x^iy^jz^k$ is $\omega^{j-k}$; the $15$ pairs $(j,k)$, $j+k\le4$, give $5,5,5$ for the three classes of $j-k\bmod3$. For $p\ge5$ the eigenvalue is $\varepsilon^{aj+bk}$ with $b\ne0$, and $j$ determines $k\in\{0,\dots,4\}$ uniquely mod $p$ for a given value, so $\dim E_\chi\le5$.
\end{itemize}
In every case the dimension is $\le13<14$. Taking the finite union over types, over $p$, and over the characters $\chi$, the quartics with a nontrivial automorphism lie in a constructible subset of dimension $\le13$ of $\mathbf P^{14}$. Hence the general plane quartic, i.e.\ the general curve of genus $3$, has no automorphism except the identity. \hfill$\square$

\section*{6. Classification of curves in projective 3-space}

Conventions. $k$ is algebraically closed; ``curve in $\mathbf P^3$'' means an irreducible nonsingular projective curve embedded in $\mathbf P^3$, $d$ its degree, $g$ its genus, $D$ (or $H$) its hyperplane section, $\mathcal O_X(n)=\mathcal O_X(nD)$. A curve is \emph{nondegenerate} if it lies in no plane. A \emph{complete intersection} of surfaces of degrees $a,b$ has $\omega_X=\mathcal O_X(a+b-4)$, degree $ab$ and genus $\frac12ab(a+b-4)+1$ (II, Ex.~8.4); a nonsingular curve on $\mathbf P^1\times\mathbf P^1\cong$ a nonsingular quadric has type $(a,b)$, degree $a+b$ and genus $(a-1)(b-1)$ (III, Ex.~5.6). Results of the text are quoted by number. The following lemma on quadric cones is stated in the text as Remark~6.4.1(d) (proved in V, Ex.~2.9); I include a short proof.

\medskip
\noindent\textbf{Lemma C (curves on a quadric cone).} Let $X$ be a nonsingular irreducible curve of degree $d$ and genus $g$ on a quadric cone $Q\subset\mathbf P^3$ with vertex $V$. If $d=2a$ then $V\notin X$ and $g=(a-1)^2$; if $d=2a+1$ then $V\in X$ and $g=a^2-a$.

\emph{Proof.} The blow-up of $Q$ at $V$ is the ruled surface $\tilde Q=\mathbf F_2$ (V, Section~2), with $\operatorname{Pic}=\mathbf ZC_0\oplus\mathbf Zf$, $C_0^2=-2$, $C_0f=1$, $f^2=0$, $K=-2C_0-4f$, hyperplane class $H=C_0+2f$, exceptional curve $C_0$. Since $X$ is nonsingular, its proper transform $\tilde X\cong X$, and $\tilde X\sim\alpha C_0+\beta f$ for integers $\alpha,\beta$. Then $d=\tilde X\cdot H=\beta$, and $\tilde X\cdot C_0=\beta-2\alpha=\operatorname{mult}_VX\in\{0,1\}$ ($X$ nonsingular), so $\beta$ is even iff $V\notin X$. By adjunction $2g-2=\tilde X(\tilde X+K)=-2\alpha^2+2\alpha\beta-2\beta$, i.e.\ $g=1-\alpha^2+\alpha\beta-\beta$. If $\beta=2a$, $\alpha=a$, $g=1-a^2+2a^2-2a=(a-1)^2$; if $\beta=2a+1$, $\alpha=a$, $g=1-a^2+a(2a+1)-2a-1=a^2-a$. \hfill$\square$

\exercise{6.1}{A rational curve of degree $4$ in $\mathbf P^3$ is contained in a unique quadric surface $Q$, and $Q$ is necessarily nonsingular.}
\pf
Let $X\cong\mathbf P^1$ be nonsingular of degree $4$. $X$ is nondegenerate (a plane curve of degree $4$ has genus $3$). $\mathcal O_X(2)=\mathcal O_{\mathbf P^1}(8)$ has $h^0=9<10=\dim H^0(\mathcal O_{\mathbf P^3}(2))$, so some quadric $Q$ contains $X$. Every quadric containing $X$ is irreducible: otherwise it is a union of two planes (or a double plane) and the irreducible $X$ would lie in a plane. If $Q\ne Q'$ both contain $X$, they have no common component, so $Q\cap Q'$ is a complete intersection curve of degree $4=\deg X$ containing $X$; hence $Q\cap Q'$ has the same support as $X$ with multiplicity $1$, and, having no embedded points, it equals $X$ scheme-theoretically (a complete intersection is Cohen--Macaulay and generically reduced, so reduced). Then $X$ is a complete intersection of two quadrics and has genus $1\ne0$ (see conventions), a contradiction. So $Q$ is unique.

$Q$ is irreducible, so it is nonsingular or a cone. If $Q$ were a cone, Lemma~C with $d=4=2a$, $a=2$ gives $g=(2-1)^2=1\ne0$, a contradiction. So $Q$ is nonsingular. (Then $X$ has type $(a,b)$ with $a+b=4$, $(a-1)(b-1)=0$: type $(1,3)$.) \hfill$\square$

\exercise{6.2}{A rational curve of degree $5$ in $\mathbf P^3$ is always contained in a cubic surface, but there are such curves which are not contained in any quadric surface.}
\pf
Let $X\cong\mathbf P^1$ have degree $5$ in $\mathbf P^3$. Then $\mathcal O_X(3)=\mathcal O_{\mathbf P^1}(15)$ has $h^0=16<20=\dim H^0(\mathcal O_{\mathbf P^3}(3))$, so the restriction map has a kernel of dimension $\ge4$: $X$ lies on a cubic surface (indeed on a $\ge4$-dimensional space of cubic forms).

\emph{An example.} Let $\varphi:\mathbf P^1\to\mathbf P^3$, $(s:t)\mapsto(s^5:s^4t:s^3t^2+st^4:t^5)$. It is a closed immersion: it has no base point ($s=0$ gives $t^5\ne0$, $t=0$ gives $s^5\ne0$); on the chart $s=1$ it is $t\mapsto(1:t:t^2+t^4:t^5)$, which is injective with $x_1/x_0=t$ an immersion, and the point $t=\infty$ maps to $(0:0:0:1)$, not in the image of the chart $s=1$ (where $x_0=1$); near $t=\infty$ (chart $t=1$, coordinate $s$): $(s^5:s^4:s^3+s:1)$ and $x_2/x_3=s^3+s$ has derivative $1\ne0$ at $s=0$. So $X=\varphi(\mathbf P^1)$ is a nonsingular rational quintic. A quadric $\sum c_{ij}x_ix_j$ contains $X$ iff $\sum c_{ij}\varphi_i\varphi_j=0$ identically; setting $s=1$ the ten products are $1,\ t,\ t^2+t^4,\ t^5,\ t^2,\ t^3+t^5,\ t^6,\ t^4+2t^6+t^8,\ t^7+t^9,\ t^{10}$ (for $x_0^2,x_0x_1,x_0x_2,x_0x_3,x_1^2,x_1x_2,x_1x_3,x_2^2,x_2x_3,x_3^2$). These are linearly independent: the coefficients of $1,t,t^{10}$ force $c_{00}=c_{01}=c_{33}=0$; $t^9$ and $t^8$ force $c_{23}=c_{22}=0$; $t^6$ forces $c_{13}=0$; $t^3$ forces $c_{12}=0$; $t^5$ forces $c_{03}=0$; and $t^2+t^4$, $t^2$ are independent, so $c_{02}=c_{11}=0$. Hence $X$ lies on no quadric. \hfill$\square$

\exercise{6.3}{A curve of degree $5$ and genus $2$ in $\mathbf P^3$ is contained in a unique irreducible quadric surface $Q$. Show that for any abstract curve $X$ of genus $2$, there exist embeddings of degree $5$ in $\mathbf P^3$ for which $Q$ is nonsingular, and there exist other embeddings of degree $5$ for which $Q$ is singular.}
\pf
\textbf{Uniqueness.} Let $X\subset\mathbf P^3$ be of degree $5$, genus $2$. It is nondegenerate (a plane quintic has genus $6$). $\mathcal O_X(2)$ has degree $10>2g-2$, so $h^0=10+1-2=9<10$ and $X$ lies on some quadric; any such quadric is irreducible (else $X$ is in a plane). Two distinct quadrics containing $X$ would meet in a curve of degree $4<5=\deg X$ containing $X$ (no common component, both irreducible), impossible. So $Q$ is unique and irreducible.

\textbf{Embeddings with $Q$ nonsingular.} Let $X$ have genus $2$ and let $K$ be canonical; $|K|$ is the (unique, Prop.~5.3) $g^1_2$, $h^0(K)=2$. Choose distinct points $P_1,P_2,P_3$ no two of which are conjugate under the hyperelliptic involution, and let $D_3=P_1+P_2+P_3$. By Riemann--Roch $h^0(D_3)=2$. The pencil $|D_3|$ has no base point: a base point must lie in $D_3$, say $P_1$, and then $h^0(D_3-P_1)=2$, so $P_2+P_3\sim K$ ($h^0=2$ for a degree $2$ divisor forces $\sim K$), i.e.\ $P_2,P_3$ are conjugate (or equal Weierstrass points $2P\sim K$, excluded by $P_2\ne P_3$); contradiction. Also $D_3\not\sim K+P$ for all $P$ (that would give $h^0(D_3-P)=h^0(K)=2$, so $P$ is a base point). Now consider multiplication $\mu:H^0(K)\otimes H^0(D_3)\to H^0(K+D_3)$. It is injective: if $sa+tb=0$ with $s,t$ a basis of $H^0(K)$ (without common zeros) and $a,b\in H^0(D_3)$, then $sa=-tb$ and $\operatorname{div}(t)\le\operatorname{div}(a)+\operatorname{div}(s)$, with $\operatorname{div}(t),\operatorname{div}(s)$ disjoint, so $a$ vanishes on $\operatorname{div}(t)$, i.e.\ $a\in H^0(D_3-K)$ which is $0$ since $D_3-K$ has degree $1$ and $D_3\not\sim K+P$; so $a=0$, then $b=0$. Both sides have dimension $4$ ($h^0(K+D_3)=5+1-2=4$), so $\mu$ is an isomorphism. The divisor $D=K+D_3$ has degree $5=2g+1$, hence is very ample (3.2(b)) and $|D|$ embeds $X$ in $\mathbf P^3$ with coordinates $x_{ij}=s_ia_j$ ($s_0,s_1$ a basis of $H^0(K)$, $a_0,a_1$ a basis of $H^0(D_3)$). Then $x_{00}x_{11}=x_{01}x_{10}$ on $X$: $X$ lies on the nonsingular quadric $Q:x_{00}x_{11}=x_{01}x_{10}$, which by the uniqueness above is \emph{the} quadric of $X$.

\textbf{Embeddings with $Q$ singular.} Let $P\in X$ and $D=2K+P$, degree $5$, very ample (3.2(b)), $h^0(D)=4$. Let $\sigma_P\in H^0(\mathcal O(P))$ be the section vanishing at $P$ and let $s,t$ be a basis of $H^0(K)$; then $\sigma_Ps^2,\sigma_Pst,\sigma_Pt^2$ are independent sections of $D$ (they are $\sigma_P\cdot$ independent elements of $H^0(2K)$, as $s^2,st,t^2$ are independent: $t/s$ is a nonconstant function), vanishing at $P$. Since $|D|$ is base point free, there is a section $u$ with $u(P)\ne0$; so $\{\sigma_Ps^2,\sigma_Pst,\sigma_Pt^2,u\}$ is a basis of $H^0(D)$. In the embedding $(x_0:x_1:x_2:x_3)=(\sigma_Ps^2:\sigma_Pst:\sigma_Pt^2:u)$ the curve lies on the quadric cone $x_0x_2=x_1^2$ with vertex $(0:0:0:1)=$ the image of $P$. By uniqueness this is the quadric $Q$ of $X$, and it is singular. \hfill$\square$

\exercise{6.4}{There is no curve of degree $9$ and genus $11$ in $\mathbf P^3$. [\emph{Hint}: Show that it would have to lie on a quadric surface, then use (6.4.1).]}
\pf
Let $X\subset\mathbf P^3$ be such a curve. It is nondegenerate (a plane curve of degree $9$ has genus $28$). $\mathcal O_X(3)$ has degree $27>2g-2=20$, so $h^0(\mathcal O_X(3))=27+1-11=17<20$ and there are at least $3$ linearly independent cubic forms $F_1,F_2,F_3$ vanishing on $X$.

\emph{$X$ lies on a quadric.} Suppose not. $X$ is irreducible and lies in no plane, so for each $i$ it lies in an irreducible component of $F_i=0$ that is a surface of degree $\ge2$; and not of degree $2$ (it would be an irreducible quadric containing $X$). So that component has degree $3$, i.e.\ $F_i$ is irreducible. Then $F_1,F_2$ are two non-proportional irreducible cubics without common component, so $F_1\cap F_2$ is a complete intersection curve of degree $9=\deg X$ containing $X$; as before it equals $X$, so $X$ is a complete intersection of type $(3,3)$, of genus $\frac12\cdot9\cdot2+1=10\ne11$. Contradiction. Hence $X$ lies on an irreducible quadric $Q$.

\emph{Using 6.4.1.} If $Q$ is nonsingular, $X$ has type $(a,b)$, $a+b=9$, $(a-1)(b-1)=11$, i.e.\ $ab=19$ with $a+b=9$: $t^2-9t+19$ has discriminant $5$, not a square: impossible. If $Q$ is a cone, Lemma~C (Remark 6.4.1(d)) with $d=9=2a+1$, $a=4$ gives $g=a^2-a=12\ne11$. So no such curve exists. \hfill$\square$

\exercise{6.5}{If $X$ is a complete intersection of surfaces of degrees $a,b$ in $\mathbf P^3$, then $X$ does not lie on any surface of degree $<\min(a,b)$.}
\pf
Let $X=F_a\cap F_b$ scheme-theoretically (the ideal sheaf $\mathcal I_X$ is generated locally by $F_a,F_b$, which form a regular sequence since $X$ has codimension $2$). The Koszul complex gives an exact sequence $0\to\mathcal O(-a-b)\to\mathcal O(-a)\oplus\mathcal O(-b)\to\mathcal I_X\to0$. Twisting by $c$ and taking cohomology, using $H^1(\mathcal O_{\mathbf P^3}(c-a-b))=0$ (III, 5.1), we get that $H^0(\mathcal I_X(c))$ is the image of $H^0(\mathcal O(c-a))\oplus H^0(\mathcal O(c-b))$, which is $0$ for $c<\min(a,b)$. So no nonzero form of degree $c<\min(a,b)$ vanishes on $X$. \hfill$\square$

\exercise{6.6}{Let $X$ be a projectively normal curve in $\mathbf P^3$, not contained in any plane. If $d=6$, then $g=3$ or $4$. If $d=7$, then $g=5$ or $6$. Cf.\ (II, Ex.~8.4) and (III, Ex.~5.6).}
\pf
By (II, Ex.~5.14(d)) ($X$ nonsingular, hence normal) projective normality means that $H^0(\mathbf P^3,\mathcal O(n))\to H^0(X,\mathcal O_X(n))$ is surjective for all $n\ge0$. For $n=1$ it is injective ($X$ nondegenerate), so $h^0(\mathcal O_X(1))=4$. Riemann--Roch: $4=d+1-g+h^1(D)$, so $h^1(D)=g+3-d\ge0$, i.e.\ $g\ge d-3$. By Castelnuovo (6.4): $g\le\frac14d^2-d+1$ for $d$ even and $g\le\frac14(d^2-1)-d+1$ for $d$ odd. For $d=6$: $3\le g\le4$. For $d=7$: $4\le g\le6$. It remains to exclude $g=4$ for $d=7$. For $n=2$: $\deg\mathcal O_X(2)=14>2g-2=6$ so $h^0(\mathcal O_X(2))=14+1-4=11>10=h^0(\mathcal O_{\mathbf P^3}(2))$, so the restriction map is not surjective: $X$ is not projectively normal. Hence $g=5$ or $6$ when $d=7$. \hfill$\square$

\exercise{6.7}{The line, the conic, the twisted cubic curve, and the elliptic quartic curve in $\mathbf P^3$ have no multisecants. Every other curve in $\mathbf P^3$ has infinitely many multisecants. [\emph{Hint}: Consider a projection from a point of the curve to $\mathbf P^2$.]}
\pf
A \emph{multisecant} is a line meeting $X$ in at least three distinct points (text before 3.7); for $X$ a line one excludes the line $X$ itself.

\textbf{No multisecants.} \emph{Line:} any other line meets it in at most one point. \emph{Conic:} a line meeting the conic in three points lies in its plane and meets it in at most $2$ points (Bezout). \emph{Twisted cubic and elliptic quartic:} let $X$ be nondegenerate of degree $d$ and $L$ a line meeting $X$ in distinct points $P_1,P_2,P_3$. The linear forms vanishing on $L$ form a $2$-dimensional space, which injects into $H^0(\mathcal O_X(1))$ ($X$ nondegenerate) and lands in $H^0(\mathcal O_X(D-P_1-P_2-P_3))$, so $h^0(D-P_1-P_2-P_3)\ge2$, where $\deg(D-\sum P_i)=d-3$. For $d=3$ this is a divisor of degree $0$, whose $h^0\le1$; for $d=4$, $g=1$ it is a divisor of degree $1$ on an elliptic curve, with $h^0=1$ by Riemann--Roch. Contradiction in both cases.

\textbf{Every other curve has infinitely many.} Let $X$ be a curve in $\mathbf P^3$ other than these four. Throughout assume $\operatorname{char}k=0$ for this part (separability of the pencils used).

\emph{Plane curves.} If $X$ is a plane curve of degree $d\ge3$, a general line of its plane meets $X$ transversally in $d\ge3$ distinct points (Bertini, II 8.18), and there are infinitely many such lines.

\emph{Curves on a quadric.} Let $X$ be nondegenerate and lie on an irreducible quadric $Q$. If $Q$ is nonsingular of type $(a,b)$, $a\le b$, then $d=a+b$, $g=(a-1)(b-1)$; the pair $(a,b)$ is not $(1,1),(1,2),(2,2)$ (conic is planar, twisted cubic, elliptic quartic), so $b\ge3$ (note $(2,2)$ has $g=1$, $d=4$, and $(1,2)$ has $g=0$, $d=3$; the pairs with $a=0$ are lines). A general line of the ruling meeting $X$ in $b$ points meets it in $b\ge3$ distinct points (general fibre of the degree-$b$ morphism $X\to\mathbf P^1$, separable), giving infinitely many multisecants. If $Q$ is a cone, by Lemma~C: for $d=2a$, the lines of the cone meet $X$ in $a$ points each, and $a\ge3$ unless $d=4$ ($a=2$, $g=1$: the elliptic quartic, excluded) or $d=2$; for $d=2a+1$, $V\in X$ and each line of the cone through $V$ meets $X$ in $V$ and $a$ further points, $a\ge2$ for $d\ge5$ (and $d=3$ is the twisted cubic). Again infinitely many multisecants.

\emph{General case: reduction via the hint.} Let $X$ be nondegenerate, not one of the four types, and suppose, for a contradiction, that $X$ has only finitely many multisecants. They contain only finitely many points of $X$, so we may take $R\in X$ on none of them (``general'' below means: outside a finite set). Project from $R$: $\psi_R:X\to\mathbf P^2$, given by the system of planes through $R$ with the base point $R$ removed (a $g^2_{d-1}$ inside $|D-R|$; it has no base point, since a plane through $R$ transversal to $T_R$ meets $X$ simply at $R$). If $\psi_R(Q')=\psi_R(Q'')$ with $Q',Q''\ne R$ distinct, then $R,Q',Q''$ are collinear, a multisecant through $R$: excluded. So $\psi_R$ is injective on $X\setminus R$; in characteristic $0$ it is separable, so a map of degree $>1$ would have fibres with $\ge2$ points off $R$, and $\psi_R$ is birational onto its image $Y_R$, a plane curve of degree $d-1$. Now $\psi_R$ is defined by $|D-R|$, so by (3.1(b)) it is a closed immersion iff for all $P,Q$ (possibly equal) planes through $R$ separate them: this means

\begin{itemize}
\item[(1)] no line through $R$ meets $X$ in two further distinct points (true by the choice of $R$);
\item[(2)] $T_R$ meets $X$ only at $R$, and with multiplicity exactly $2$ (so $R$ is not a flex);
\item[(3)] no tangent line $T_Q$, $Q\ne R$, passes through $R$.
\end{itemize}

If $\psi_R$ were a closed immersion, $X\cong Y_R$ would be a nonsingular plane curve of degree $d-1$ and $g=\frac12(d-2)(d-3)$. By Castelnuovo (6.4), $g\le\lfloor d^2/4\rfloor-d+1$ and $\frac12(d-2)(d-3)-(\frac14d^2-d+1)=\frac14(d-2)(d-4)>0$ for $d\ge5$, so $d\ge5$ is impossible; $d=4$ would give $g=1$ (elliptic quartic, excluded), $d=3$ would give $g=0$ (twisted cubic, excluded), and $d\le2$ is a line or planar. Hence, for $X$ not among the four types, \emph{(2) or (3) fails for every $R$ off a finite set.}

\emph{The remaining step (not proved).} \emph{Flexes.} If every point of $X$ were a flex, i.e.\ $T_R\cdot X\ge3$ at $R$ for all $R$, then in a local parameter $t$ one has $\gamma''\parallel\gamma'$ identically; writing $\gamma''=f\gamma'$ gives $\gamma'=e^{\int f}v$ with $v$ constant (characteristic $0$), so $X$ is a line. So for general $R$ the failure of (2),(3) is: \emph{$T_R$ meets $X$ at a second point, or some $T_Q$ ($Q\ne R$) passes through $R$.} Both say that the correspondence $\mathcal C=\{(Q,R):Q\ne R,\ R\in T_Q\cap X\}\subset X\times X$ has a component of dimension $1$ whose two projections are dominant (if the projection to the first factor were finite-to-one onto a finite set of $Q$'s, the $R$'s would lie in finitely many lines and could not fill $X$). That is: \emph{the tangent line at a general point of $X$ meets $X$ again.} So, in characteristic $0$, what remains to prove the second assertion in general is

\medskip
\noindent\textbf{Lemma T (not proved here).} \emph{If $X\subset\mathbf P^3$ is a nondegenerate curve over a field of characteristic $0$, then $T_P\cap X=\{P\}$ for $P$ in a dense open subset of $X$.}
\medskip

Lemma T says that the tangent developable has multiplicity exactly $2$ along $X$ at a general point of $X$. I believe it is true (it is the classical fact that the tangent line at a general point of a non-planar space curve is not a ``tangential secant''), but \emph{I could not give a complete proof of it}. What I verified: it fails exactly for plane curves (so nondegeneracy is essential, and Lemma T is where it enters); it holds for nondegenerate curves on a quadric $Q$ (a tangent line $T_P\not\subset Q$ satisfies $T_P\cdot Q=2P$, so it meets $X\subset Q$ only at $P$, and $T_P\subset Q$ for only finitely many $P$); and it holds when $X$ is the image of $t\mapsto(1:t^a:t^b:t^c)$ for all small exponents (such a tangent line meets $X$ again iff there is $\tau\ne1$ with $\tau^a=1+a(\tau-1)$, and likewise for $b,c$; a computer search over $|a|,|b|,|c|\le9$ found no solution, but this is only evidence). The argument above is therefore complete for: (i) the four exceptional curves having no multisecants; (ii) plane curves and curves on a quadric (nonsingular or a cone) having infinitely many multisecants; and for (iii) all other nondegenerate curves it is complete \emph{only conditionally on Lemma T}, which is used exactly once, to pass from ``(2) or (3) fails for general $R$'' to a contradiction. I flag this as the one gap in Section 6. \hfill$\square$

\exercise{6.8}{A curve $X$ of genus $g$ has a nonspecial divisor $D$ of degree $d$ such that $|D|$ has no base points if and only if $d\ge g+1$.}
\pf
\emph{Convention.} $|D|$ is understood to be nonempty (and we assume $g\ge1$, or $d\ge1$ when $g=0$: for $g=0$ and $d=0$ the trivial divisor $D=0$ is nonspecial with no base points, an evident degenerate exception).

\textbf{($\Rightarrow$)} Let $D$ be nonspecial of degree $d$ with $|D|$ base point free. If $d\le g$ then $h^0(D)=d+1-g\le1$, so $|D|$ consists of at most one divisor $E\ge0$ of degree $d$ and has base points $\operatorname{Supp}E$, unless $E=0$, i.e.\ $d=0$, $D\sim0$; but $D\sim0$ nonspecial means $h^0(K)=0$, $g=0$, the excluded case. So $d\ge g+1$.

\textbf{($\Leftarrow$)} Let $d\ge g+1$. \emph{Reformulation.} If $D$ is nonspecial then $|D|$ is base point free iff $h^0(D-P)=h^0(D)-1$ for all $P$ (3.1); Riemann--Roch gives $h^0(D-P)=d-g+h^0(K-D+P)$ and $h^0(D)=d+1-g$, so $|D|$ is base point free iff $h^0(K-D+P)=0$ for all $P$ (i.e.\ all $D-P$ are nonspecial).

\emph{Case $g=0$:} take $D=dP_0$, $d\ge1$; $h^0=d+1$, base point free.

\emph{Case $g\ge1$, $d=g+1$.} Choose points $P_1,\dots,P_g$ inductively: if $h^0(K-P_1-\dots-P_{i-1})\ge1$ choose $P_i$ outside the (finite) base locus of this complete system, so the $h^0$ drops by one. Then $h^0(K-D_g)=0$ for $D_g=\sum_{i\le g}P_i$ ($h^0(K)=g$). By Riemann--Roch $h^0(D_g)=1$, so $|D_g|=\{D_g\}$. For a point $P$: $h^0(K-D_g+P)=h^0(D_g-P)-(g-1+1-g)=h^0(D_g-P)$ (degree $g-1$), which is $\le1$ and equals $1$ iff $D_g-P$ is linearly equivalent to an effective divisor, i.e.\ iff $P\in\operatorname{Supp}D_g$ (since $|D_g|$ is a point). For $P=P_i$ the system $|K-D_g+P_i|$ has a unique member $F_i$, effective of degree $g-1$. Choose $P_{g+1}\notin\bigcup_iF_i$, and put $D=D_g+P_{g+1}$ (degree $g+1$). Then for all $P$: if $P\notin\operatorname{Supp}D_g$, $h^0(K-D+P)\le h^0(K-D_g+P)=0$; if $P=P_i$, $h^0(K-D+P_i)=h^0(F_i-P_{g+1})=0$ since $F_i$ is the unique member of its system and $P_{g+1}\notin F_i$. Also $h^0(K-D)\le h^0(K-D_g)=0$, so $D$ is nonspecial; so $D$ is nonspecial and base point free of degree $g+1$.

\emph{Case $d>g+1$.} Let $D_0$ be as above of degree $g+1$ and $E\ge0$ any effective divisor of degree $d-g-1$, $D=D_0+E$. Then $h^0(K-D)=0$ and $h^0(K-D+P)\le h^0(K-D_0+P)=0$ for all $P$, so $D$ is nonspecial and base point free. \hfill$\square$

\exercise{6.9}{Let $X$ be an irreducible nonsingular curve in $\mathbf P^3$. Then for each $m\gg0$, there is a nonsingular surface $F$ of degree $m$ containing $X$. [\emph{Hint}: Let $\pi:\tilde P\to\mathbf P^3$ be the blowing-up of $X$ and let $Y=\pi^{-1}(X)$. Apply Bertini's theorem to the projective embedding of $\tilde P$ corresponding to $\mathcal I_Y\otimes\pi^*\mathcal O_{\mathbf P^3}(m)$.] (Starred in the book.)}
\pf
(The ambient space is $\mathbf P^3$, as for a ``surface $F$''; the superscript is blurred in the scan. The same proof gives a nonsingular hypersurface in $\mathbf P^r$, $r\ge3$.)

Let $\pi:\tilde P\to\mathbf P^3$ be the blow-up along $X$. It is nonsingular (II, 8.24, the center $X$ being nonsingular), $Y=\pi^{-1}(X)$ is an effective Cartier divisor, $\mathcal I_Y=\mathcal O_{\tilde P}(-Y)$ is invertible and $\pi$ is an isomorphism off $Y$. By (II, 7.10(b)) for $m\gg0$ the invertible sheaf $\mathcal L_m=\mathcal I_Y\otimes\pi^*\mathcal O(m)$ is very ample on $\tilde P$ (I used $\mathcal I_Y=\mathcal I_X\mathcal O_{\tilde P}$). Embed $\tilde P\subset\mathbf P^N$ by the complete system $|\mathcal L_m|$. By Bertini (II, 8.18) a general hyperplane section $G$ of $\tilde P$ is nonsingular (and irreducible, $\dim\tilde P=3$), and, $Y$ being irreducible and $\mathcal L_m$ very ample, $G\not\supseteq Y$ for general $G$.

\emph{Identification of sections.} A section of $\mathcal L_m$ is a section of $\pi^*\mathcal O(m)$ vanishing on $Y$; it restricts on $\tilde P\setminus Y\cong\mathbf P^3\setminus X$ to a section of $\mathcal O(m)$ and extends over $X$ (codimension $2$, $\mathbf P^3$ normal), so it comes from a form $f$ of degree $m$, which vanishes on $X$. So $H^0(\tilde P,\mathcal L_m)=H^0(\mathbf P^3,\mathcal I_X(m))=:W_m$, and the hyperplane section $G$ corresponding to $f$ satisfies $\pi^*(F)=G+Y$, where $F=V(f)\supseteq X$ is a surface of degree $m$; $G$ is the proper transform of $F$ (as $G\not\supseteq Y$). For $f$ general, $G$ is nonsingular, so $F\setminus X\cong G\setminus Y$ is nonsingular.

\emph{Nonsingularity of $F$ along $X$.} Let $m$ be large enough that $\mathcal I_X(m)$ is globally generated (II, 5.17), so $W_m\to\mathcal I_{X,P}/\mathfrak m_P\mathcal I_{X,P}$ is surjective for each $P\in X$ (a $2$-dimensional space, as $X$ is a nonsingular curve in a nonsingular threefold: $\mathcal I_{X,P}=(u,v)$ with $u,v$ part of a regular system of parameters). Now $f\in\mathfrak m_P^2$ iff the class of $f$ in $\mathcal I_P/\mathfrak m_P\mathcal I_P$ is $0$ (because $u,v$ are linearly independent in $\mathfrak m_P/\mathfrak m_P^2$): two independent linear conditions. So $B_P=\{f\in W_m:f\in\mathfrak m_P^2\}$ has codimension $2$, and $B=\{(f,P):f\in B_P\}\subseteq W_m\times X$ is closed of dimension $\dim W_m-1$. Its image in $W_m$ (a closed set, $X$ projective) is a proper closed subset: for $f$ outside it, $F=V(f)$ is nonsingular at every point of $X$ ($F$ is singular at $P\in X$ iff $f\in\mathfrak m_P^2$).

The set of $f\in W_m$ for which $G_f$ is nonsingular and $\not\supseteq Y$ contains a nonempty open subset of $\mathbf P(W_m)$ (Bertini), and the set of $f$ with $F$ nonsingular along $X$ is a nonempty open subset. $\mathbf P(W_m)$ being irreducible, the two meet: for such $f$, $F$ is nonsingular off $X$ and along $X$, so $F$ is a nonsingular surface of degree $m$ containing $X$. \hfill$\square$

\end{document}
